% ======================================================================
%  TQM-MONO115 — The Actualization Theory: Complete Monograph Assembly
%  A Reconstruction of Physics from Difference, Actualization and Spectrum
%
%  Canonical Monograph (v2.0) · Zenodo-ready book structure
%  Master file — compiles the complete monograph.
%
%  Status: COMPLETE — publication-ready core (Ch1-14), universality appendix
%  Method: assembly only — no new physics, no new derivations
%
%  Build: pdflatex main.tex (twice for references and the TOC)
%  Preamble defines the theorem environments shared by all chapters.
% ======================================================================

\documentclass[11pt,openany]{book}

% ---- Fonts and encoding ------------------------------------------------
\usepackage[T1]{fontenc}
\usepackage[utf8]{inputenc}
\usepackage{lmodern}

% ---- Mathematics and theorem environments ------------------------------
\usepackage{amsmath,amssymb,amsthm}
\usepackage{mathtools}

% ---- Links, colors, graphics -------------------------------------------
\usepackage[hidelinks]{hyperref}
\usepackage{xcolor}
\usepackage{graphicx}

% ---- Page geometry ------------------------------------------------------
\usepackage[a4paper,margin=2.5cm]{geometry}
\usepackage{setspace}
\onehalfspacing

% ---- Header/footer ------------------------------------------------------
\usepackage{fancyhdr}
\pagestyle{fancy}
\fancyhf{}
\fancyhead[LE]{\leftmark}
\fancyhead[RO]{\rightmark}
\fancyfoot[C]{\thepage}

% ---- Table of contents depth -------------------------------------------
\setcounter{tocdepth}{2}
\setcounter{secnumdepth}{3}

% ---- Theorem environments (shared by all chapters) ---------------------
\newtheorem{definition}{Definition}[chapter]
\newtheorem{lemma}[definition]{Lemma}
\newtheorem{theorem}[definition]{Theorem}
\newtheorem{corollary}[definition]{Corollary}
\newtheorem{remark}[definition]{Remark}

% ---- Metadata -----------------------------------------------------------
\title{The Actualization Theory\\[0.4em]
  \large A Reconstruction of Physics from Difference, Actualization and Spectrum}
\author{Fabrice Wieser\\[1mm]
  \small Independent Researcher and Full Stack Software Engineer\\[0.5mm]
  \small \url{https://github.com/MagusDraconis/AT}}
\date{\today}

% ======================================================================
\begin{document}

% ---- Front matter -------------------------------------------------------
\frontmatter
% ======================================================================
%  Front Matter of The Actualization Theory
%  AT-MONO115 · Canonical Monograph (v2.0) · Zenodo-ready
%  Includes: front cover, title page, copyright, abstract, executive
%  summary, preface, table of contents.
%  Assembly only — no new physics, no new derivations.
% ======================================================================

% ---- Title page (also serves as front cover) ------------------------------
\thispagestyle{empty}
\begin{center}
  \vspace*{2.5cm}
  {\Huge\bfseries The Actualization Theory}\\[1em]
  {\LARGE A Reconstruction of Physics from\\ Difference, Actualization and Spectrum}\\[2.5cm]
  {\Large Fabrice Wieser}\\[1mm]
  {\small Independent Researcher and Full Stack Software Engineer}\\[0.5mm]
  {\small \url{https://github.com/MagusDraconis/AT}}\\[2.5cm]
  {\large The Actualization Theory (AT) --- v2.0 (canonical)}\\[1cm]
  {\large \today}\\[1.5cm]
  {\large Zenodo edition}
  \vfill
\end{center}
\clearpage

% ---- Copyright page --------------------------------------------------------
\thispagestyle{empty}
\begin{center}
  \vspace*{4cm}
  {\large The Actualization Theory}\\[0.6em]
  {\normalsize A Reconstruction of Physics from Difference, Actualization and Spectrum}\\[2em]
  {\normalsize Copyright \copyright\ \the\year\ Fabrice Wieser}\\[1.5em]
  {\small All rights reserved.}
  \vfill
  {\small Version 2.0 (canonical) --- \today}\\
  {\small Compiled from the canonical AT-QG research program (QG0--QG321).}
  \vfill
\end{center}
\clearpage

% ---- Abstract ---------------------------------------------------------------
\chapter*{Abstract}
\label{fm:abs:abstract}
\addcontentsline{toc}{chapter}{Abstract}

The Actualization Theory reconstructs physics from the primitive Difference and the tensor
reference $\eta$, through the canonical hierarchy

\[
\text{Difference} \;\longrightarrow\; \text{Actualization}
\;\longrightarrow\; \text{Inevitable Spectrum} \;\longrightarrow\; \text{Physics}.
\]

Difference is the fundamental boundary: the irreducible primitive whose count-conservation
law is its definitional identity, and whose two faces --- the scalar counting measure
$\rho$ and the tensor field $\psi$ --- carry, respectively, the spectral readout of all
physics and the gravitational content of the theory. Actualization is Difference's derived
count-producing process, whose convergence is the closure boundary $N=96$. The spectrum is
the inevitable output of the actualization attractor: the unique eigenspectrum of the
converged network, from which every physics sector is read.

The monograph derives, from this foundation: quantum mechanics ($|\psi|^2=\rho$,
measurement as actualization), gravity and spacetime (emergent geometry, native dynamics),
the standard model (fermion sectors, gauge structure, couplings, mixing, weak scale,
precision predictions $S,T,U$, the electron g-2, the Majorana character), and cosmology
(the cosmological constant, the density fractions, the CMB acoustic structure, structure
formation). The boundaries of the theory are documented explicitly: Difference as the
ontological boundary, the unresolved status of $\psi$, the numerical boundary $\pi$, the
Bekenstein $2\pi$ factor, and the hosted standard-model dynamics. The theory claims no open
derivation frontier; its frontier is the experimental validation of its pre-registered
predictions.

% ---- Executive summary --------------------------------------------------------
\chapter*{Executive Summary}
\label{fm:abs:executive-summary}
\addcontentsline{toc}{chapter}{Executive Summary}

\textbf{The foundation.} The Actualization Theory uses exactly two primitives,
$\{\text{Difference},\eta\}$. Difference is the fundamental boundary --- the irreducible
counting difference from a uniform background --- and $\eta$ is the tensor reference against
which the tensor face of Difference is read. Actualization, the count-producing process, is
derived from Difference, not a third primitive.

\textbf{The spectrum.} The D96 spectrum is the inevitable output of the actualization
attractor: one zero mode and 95 positive modes, with multiplicities
$[42\times2,5,6]$, moments $\Sigma m=95$, $\Sigma\sqrt{m}=64.08$, $\Sigma m^2=229$, the
occupancy moment $\mathrm{occMom}=1900.25$, and the span $6.40$. Every derived constant
traces to this spectrum.

\textbf{The physics.} Quantum mechanics is emergent ($|\psi|^2=\rho$, measurement is
actualization). Gravity is emergent (the coupling from the D96 spectral content, the tensor
sector from $\psi$ against $\eta$, spacetime from the counting measure). The standard model
is a set of derived reads of the D96 moments: the $1+3+8$ gauge structure, the couplings,
the CKM and PMNS matrices, the weak scale, the Higgs, the neutrino masses, and the
precision predictions. Cosmology is a spectrum readout: the cosmological constant is the
residual actualization pressure, the density fractions are the octave-record information,
and the CMB acoustic structure is the spectral harmonics.

\textbf{The boundaries.} The theory documents its boundaries honestly: Difference as the
ontological boundary, the unresolved ontological status of $\psi$, the numerical boundary
$\pi$, the Bekenstein $2\pi$ factor, and the hosted standard-model dynamics. These are
limits of the framework, not derivation failures.

\textbf{The frontier.} The theory claims no open derivation frontier. Its frontier is the
experimental validation of the derived predictions: the oblique parameters $S,T,U$, the
electron g-2 $a_e$, the Majorana character $0\nu\beta\beta$, and the pre-registered
predictions P1 (106 GeV), P2 ($0\nu\beta\beta$), and the P3 sector ladder.

% ---- Preface ------------------------------------------------------------------
\chapter*{Preface}
\label{fm:abs:preface}
\addcontentsline{toc}{chapter}{Preface}

This monograph is the canonical v2.0 record of The Q-Model (AT), now presented under the
name \emph{The Actualization Theory}. It consolidates the end-state of a research program
that began with the question of whether observable physics can be reconstructed from a
smaller primitive foundation, and converged on the canonical hierarchy
Difference $\to$ Actualization $\to$ Spectrum $\to$ Physics.

The monograph is a \emph{foundation monograph}: it presents the primitives, the derivation
hierarchy, the physics, the boundaries, and the validation frontier. It is written to be
referee-safe --- every claim is either a derived result with its source, a documented
boundary, or an experimental-validation item awaiting its experiment. No claim exceeds the
accepted derivations of the canonical program.

The structure follows the physics-focused plan: the foundation (Parts I--II), the spectrum
(Part III), the physics (Part IV), and the boundary layer and frontier (Part V), followed
by the general conclusion and the appendices. The universality research program --- the
operator basis across organized domains, the lock law, and the organization structure ---
is treated as future work in the appendix, outside the core physics derivation.

% ---- Table of contents ---------------------------------------------------------
\tableofcontents
\clearpage


% ---- Main matter --------------------------------------------------------
\mainmatter

% Part I — Foundation
\part{Foundation}
% ======================================================================
%  AT-MONO100 — Chapter 1: The Difference
%  Canonical Monograph (v2.0) · The Actualization Theory (AT):
%  Structure, Complexity, and Random Actualization
%
%  Status: COMPLETE — PASS-READY (MONO007 referee-ready architecture)
%  Inputs: ATQG_CanonicalMonograph.md, ATQG_Mono007RefereeReadiness.md
%  Method: assembly only — canonical results only, no new physics, no speculation
%  Citations: QG278 (fundamental boundary), QG279 (true fundamental boundary),
%             QG268 (count conservation), QG282 (closure principle),
%             QG286 (difference duality), QG292 (minimal foundation),
%             QG301 (duality complete), MONO006 (actualization derived)
%
%  Usage: % ======================================================================
%  AT-MONO100 — Chapter 1: The Difference
%  Canonical Monograph (v2.0) · The Actualization Theory (AT):
%  Structure, Complexity, and Random Actualization
%
%  Status: COMPLETE — PASS-READY (MONO007 referee-ready architecture)
%  Inputs: ATQG_CanonicalMonograph.md, ATQG_Mono007RefereeReadiness.md
%  Method: assembly only — canonical results only, no new physics, no speculation
%  Citations: QG278 (fundamental boundary), QG279 (true fundamental boundary),
%             QG268 (count conservation), QG282 (closure principle),
%             QG286 (difference duality), QG292 (minimal foundation),
%             QG301 (duality complete), MONO006 (actualization derived)
%
%  Usage: % ======================================================================
%  AT-MONO100 — Chapter 1: The Difference
%  Canonical Monograph (v2.0) · The Actualization Theory (AT):
%  Structure, Complexity, and Random Actualization
%
%  Status: COMPLETE — PASS-READY (MONO007 referee-ready architecture)
%  Inputs: ATQG_CanonicalMonograph.md, ATQG_Mono007RefereeReadiness.md
%  Method: assembly only — canonical results only, no new physics, no speculation
%  Citations: QG278 (fundamental boundary), QG279 (true fundamental boundary),
%             QG268 (count conservation), QG282 (closure principle),
%             QG286 (difference duality), QG292 (minimal foundation),
%             QG301 (duality complete), MONO006 (actualization derived)
%
%  Usage: \input{Chapter01_Difference.tex} after the preamble that defines
%         the theorem environments below (or include via the master file).
% ======================================================================

% ---- Theorem environments (define in the master preamble if not present) ----
% \newtheorem{definition}{Definition}[chapter]
% \newtheorem{lemma}[definition]{Lemma}
% \newtheorem{theorem}[definition]{Theorem}
% \newtheorem{corollary}[definition]{Corollary}
% \newtheorem{remark}[definition]{Remark}

\chapter{The Difference}
\label{c01:chap:difference}

\section{Introduction}
\label{c01:sec:difference-introduction}

The Actualization Theory (AT) is a deterministic theory of structure, complexity, and random
actualization. Its central claim is that the totality of observable physics can be derived
from a single primitive notion --- \emph{Difference} --- together with a tensor reference
metric $\eta$, through the canonical hierarchy
\begin{equation}
\label{c01:eq:canonical-core}
\text{Difference}\;\longrightarrow\;\text{Actualization}\;\longrightarrow\;\text{Inevitable Spectrum}\;\longrightarrow\;\text{Physics}.
\end{equation}

This chapter develops the first member of that hierarchy. We establish, with formal
precision, that Difference is the sole irreducible primitive of the theory; that its status
is that of a \emph{fundamental boundary} --- the deepest notion, which cannot itself be
derived and whose own boundary is characterized by the Closure Principle
\cite{c01:QG278,c01:QG282}; and that the two great physical sectors of the theory, the scalar
counting measure $\rho$ and the tensor field $\psi$, are not independent primitives but the
\emph{trace and traceless faces of the one Difference} \cite{c01:QG286,c01:QG301}.

The results of this chapter are not new. They are the canonical end-state of the QG278--QG318
program, assembled here in a monograph form suitable for archival publication. Throughout, we
adhere to the accepted derivation record and introduce no new physics and no speculation.

\section{Primitive Status}
\label{c01:sec:primitive-status}

\begin{definition}[Difference]
\label{c01:def:difference}
\emph{Difference} is the primitive notion of the theory: the counting difference from a
uniform background, whose count-conservation law is its \emph{definitional identity}
\cite{c01:QG268}. A Q-event is a unit of Difference.
\end{definition}

\begin{definition}[Tensor reference metric]
\label{c01:def:eta}
\emph{$\eta$} is the tensor reference metric: the conformal reference with respect to which
the metric takes the form $g = \rho^{2/d}\,\eta$, and with respect to which the Weyl content
$\psi$ is defined \cite{c01:QG286}.
\end{definition}

\begin{theorem}[Minimal primitive set]
\label{c01:thm:minimal-primitives}
The minimal set of irreducible primitives of AT is $\{\text{Difference},\eta\}$: exactly
two items. No proper subset supports the derived hierarchy \eqref{c01:eq:canonical-core}, and no
further reduction is possible.
\end{theorem}

\begin{proof}
The minimality of $\{\text{Difference},\eta\}$ is established by the foundation stress test
\cite{c01:QG292}. Removing Difference collapses \emph{all} five layers of the theory
(actualization, conservation, resonance, spectrum, physics), because count conservation is the
definitional identity of the primitive: without a unit there is no count, without a count no
network, without a network no spectrum, and without a spectrum no physics. Removing $\eta$
leaves all five layers intact as a \emph{scalar} readout --- masses, couplings, mixings,
cosmological fractions, and spectral observables are all $\rho$-face reads --- and destroys
only the tensor sub-sector, because $\psi$ is defined against $\eta$. Hence Difference is
necessary for the whole hierarchy and $\eta$ is necessary precisely for the tensor content;
neither is derivable from the other, and together they form the irreducible pair.
\end{proof}

\begin{corollary}[Actualization is not a primitive]
\label{c01:cor:actualization-derived}
Actualization is \emph{not} a primitive. It is Difference's own count-producing process.
\end{corollary}

\begin{proof}
By the same removal test \cite{c01:QG292,c01:MONO006}, Actualization fails when Difference is removed
(it needs the definitional unit) and survives when $\eta$ is removed (counting needs no metric
reference). It is therefore neither an independent primitive nor a function of
$\{\text{Difference},\eta\}$ in the tensor sense; it is the dynamical face of the one
primitive. Its fixed point is the closure $N=96$ \cite{c01:QG282}.
\end{proof}

\begin{remark}
The canonical core \eqref{c01:eq:canonical-core} is therefore read: \emph{Difference} (the
primitive) produces \emph{Actualization} (its derived process), which converges to the
\emph{Inevitable Spectrum} (the fixed point), from which \emph{Physics} (all observable
sectors) emerges. Only the first member and the reference metric are primitives.
\end{remark}

\section{Difference as Fundamental Boundary}
\label{c01:sec:fundamental-boundary}

\begin{definition}[Fundamental boundary]
\label{c01:def:boundary}
A \emph{fundamental boundary} is a notion that (i) is irreducible --- no deeper notion
explains it; (ii) supports the entire derived hierarchy; and (iii) is itself bounded only by
its own fixed point, not by any imported assumption.
\end{definition}

\begin{theorem}[Difference is the fundamental boundary]
\label{c01:thm:fundamental-boundary}
Difference is the fundamental boundary of AT.
\end{theorem}

\begin{proof}
The boundary audits \cite{c01:QG278,c01:QG279} establish the three criteria. (i) \emph{Irreducible:}
the search for a deeper primitive terminates at Difference; attempts to redefine it as an
artifact of language fail, and it is declared a \emph{true} fundamental boundary. (ii)
\emph{Support:} by Theorem~\ref{c01:thm:minimal-primitives}, removing Difference collapses every
layer --- it is the root of the entire derivation graph. (iii) \emph{Self-bounded:} the
boundary of Difference is not imported but derived --- it is the fixed point $N=96$ of its own
actualization process, per the Closure Principle \cite{c01:QG282}. Hence Difference satisfies the
three defining properties of a fundamental boundary.
\end{proof}

\begin{lemma}[Boundary is fixed point]
\label{c01:lem:boundary-fixed-point}
The boundary of the theory is the stable fixed point of the actualization process: the
topology saturates at zero residual link growth, and every initial pattern converges to the
same final geometry \cite{c01:QG282}.
\end{lemma}

\begin{proof}
The actualization dynamics has a self-reinforcing feedback (activity to links to activity)
that is bounded by network capacity; positive feedback therefore saturates at the stable
fixed point, giving zero residual link growth and a content-independent final geometry
\cite{c01:QG282}. The boundary is the closure of this dynamics.
\end{proof}

\section{Difference Duality ($\rho$ and $\psi$)}
\label{c01:sec:duality}

\begin{definition}[The two faces]
\label{c01:def:two-faces}
Let $A_{ij}$ be the rank-2 object representing the difference structure (the connectivity /
the stress / the difference). Its irreducible decomposition is
\begin{equation}
\label{c01:eq:decomposition}
A_{ij} \;=\; \underbrace{\frac{1}{d}\,\mathrm{Tr}(A)\,\delta_{ij}}_{\text{trace: } \rho}
\;+\;\underbrace{A_{ij}-\frac{1}{d}\,\mathrm{Tr}(A)\,\delta_{ij}}_{\text{traceless: } \psi}.
\end{equation}
At $d=3$, $A_{ij}$ has $d(d+1)/2=6$ components: one trace component ($\rho$, the scalar) and
five traceless components ($\psi$, the tensor, of which two are the transverse-traceless
polarizations) \cite{c01:QG286}.
\end{definition}

\begin{theorem}[Difference duality]
\label{c01:thm:duality}
The scalar counting measure $\rho$ and the tensor field $\psi$ are the trace and traceless
faces of the \emph{one} Difference. They are not independent primitives; each is a projection
of the same rank-2 difference structure.
\end{theorem}

\begin{proof}
$\rho$ is the scalar counting measure: the normalized count share $|\psi|^2=\rho$
\cite{c01:QG216}, the \emph{isotropic} difference from the uniform background --- the
\emph{count} difference. $\psi$ is the tensor field: the spin-2 Weyl content \cite{c01:QG285},
the \emph{anisotropic} difference from conformal flatness --- the \emph{orientation}
difference. By the rank-2 decomposition theorem \eqref{c01:eq:decomposition}, the trace and the
traceless part are determined by $A_{ij}$; neither is an independent input. Hence $\rho$ and
$\psi$ are dual projections of the same Difference \cite{c01:QG286}.
\end{proof}

\begin{corollary}[Duality is invariant]
\label{c01:cor:duality-invariant}
The reinterpretation of the theory through the duality
$\text{Difference}\longrightarrow\{\rho,\psi\}$ changes \emph{meaning} without changing any
\emph{number}: every frozen prediction is reproduced exactly \cite{c01:QG301}.
\end{corollary}

\begin{proof}
The post-duality invariance audit recomputes every established prediction through the
reduced hierarchy and finds each unchanged: no formula is retuned and no result is lost
\cite{c01:QG301}. The duality is therefore a reinterpretation of content, not a modification of
the derived mathematics.
\end{proof}

\begin{remark}[$\rho$ needs no reference; $\psi$ does]
The scalar face $\rho$ is defined without reference to $\eta$; the tensor face $\psi$ is
defined \emph{against} $\eta$ (conformal flatness). This asymmetry is the reason
Theorem~\ref{c01:thm:minimal-primitives} holds: $\eta$ is necessary exactly for the tensor
sub-sector.
\end{remark}

\section{Relation to Closure}
\label{c01:sec:closure}

\begin{theorem}[Closure Principle]
\label{c01:thm:closure}
The boundary of the theory is the closure of Difference's own dynamics: the primitive is the
\emph{process}, and the boundary is its \emph{fixed point}. The closure is
\emph{derived}, not primitive \cite{c01:QG282}.
\end{theorem}

\begin{proof}
The origin of the boundary is the $N=96$ attractor: the observable sector is the converged
attractor of the actualization dynamics, with topology saturation and content-independent
convergence (Lemma~\ref{c01:lem:boundary-fixed-point}). The boundary is therefore not an
imported cut-off but the stable fixed point of the process itself; the closure principle
states exactly this --- the process converges, the boundary is its closure \cite{c01:QG282}.
\end{proof}

\begin{corollary}[Boundary characterizes, does not derive]
\label{c01:cor:closure-characterizes}
The Closure Principle \emph{characterizes} the fundamental boundary; it does not \emph{derive}
the primitive Difference.
\end{corollary}

\begin{proof}
Difference is the primitive: by Definition~\ref{c01:def:difference} its count-conservation law
is its definitional identity, and by Theorem~\ref{c01:thm:fundamental-boundary} it is
irreducible. The Closure Principle describes how Difference's own dynamics closes at $N=96$;
it presupposes the process, hence the primitive, and cannot be a derivation of it. This is
the referee-safe reading adopted by the canonical monograph \cite{c01:MONO007}.
\end{proof}

\section{Consequences}
\label{c01:sec:consequences}

We collect the immediate formal consequences of the foregoing.

\begin{lemma}[One primitive, two faces, one hierarchy]
\label{c01:lem:consequence-1}
The entire observable content of the theory is carried by the two faces of the one Difference
through a single canonical hierarchy: the scalar face feeds the full spectrum of
observables; the tensor face feeds the gravitational (Weyl) sector; both are projections of
the same primitive \cite{c01:QG286,c01:QG301}.
\end{lemma}

\begin{theorem}[Spectrum is inevitable]
\label{c01:thm:spectrum-inevitable}
The D96 spectrum is the inevitable output of the actualization attractor --- the Laplacian
eigenspectrum of the converged network --- and not an additional primitive \cite{c01:QG295}.
\end{theorem}

\begin{proof}
The spectrum is a deterministic function of the converged network: the $95$ positive modes
are the graph-Laplacian eigenvalues of the attractor, the zero mode is the background, and no
choice enters. The spectrum carries no independent input; it is the output of the process
whose primitive is Difference \cite{c01:QG295,c01:QG282}.
\end{proof}

\begin{corollary}[Canonical core]
\label{c01:cor:core}
The minimal hierarchy
\[
\text{Difference}\;\longrightarrow\;\text{Actualization}\;\longrightarrow\;\text{Inevitable Spectrum}\;\longrightarrow\;\text{Physics}
\]
is complete: every intermediate layer is derivable from the primitives and the attractor, and
no further primitive is required \cite{c01:QG292,c01:QG294}.
\end{corollary}

\begin{remark}[Referee-safe completeness]
Where the monograph claims completeness, it is understood in the qualified sense: every
\emph{observable} is derived from the canonical hierarchy; the standard-model
\emph{Lagrangian} and its dynamics are \emph{hosted} --- a boundary, not a derivation gap
\cite{c01:MONO007}.
\end{remark}

\section{Summary}
\label{c01:sec:summary}

This chapter has established the first member of the canonical architecture. We have proved:

\begin{enumerate}
\item \textbf{Primitive status} (Theorem~\ref{c01:thm:minimal-primitives}): the minimal primitive
set is $\{\text{Difference},\eta\}$, and Actualization is derived from Difference
(Corollary~\ref{c01:cor:actualization-derived}).
\item \textbf{Fundamental boundary} (Theorem~\ref{c01:thm:fundamental-boundary}): Difference is
irreducible, supports the whole hierarchy, and is self-bounded by its own fixed point.
\item \textbf{The duality} (Theorem~\ref{c01:thm:duality}): $\rho$ and $\psi$ are the trace and
traceless faces of the one Difference; the reinterpretation is invariant
(Corollary~\ref{c01:cor:duality-invariant}).
\item \textbf{Closure} (Theorem~\ref{c01:thm:closure}): the boundary is the fixed point of
Difference's own dynamics, characterizing rather than deriving the primitive
(Corollary~\ref{c01:cor:closure-characterizes}).
\item \textbf{Consequences} (Theorem~\ref{c01:thm:spectrum-inevitable},
Corollary~\ref{c01:cor:core}): the spectrum is inevitable, and the canonical core is complete.
\end{enumerate}

The Difference is therefore not an unexplained name but a precise object: the counting
difference whose conservation is its identity, whose dynamics closes at $N=96$, and whose two
faces --- the scalar measure $\rho$ and the tensor field $\psi$ --- generate, respectively,
the spectral readout of all physics and the gravitational content of the theory. All further
chapters of this monograph derive from the structure established here.

\begin{thebibliography}{99}
\bibitem{c01:QG278} AT-QG Phase 278, \emph{Fundamental Boundary Audit}, Docs/Research.
\bibitem{c01:QG279} AT-QG Phase 279, \emph{Boundary Legitimacy Audit}, Docs/Research.
\bibitem{c01:QG268} AT-QG Phase 268, \emph{Count Conservation Origin}, Docs/Research.
\bibitem{c01:QG282} AT-QG Phase 282, \emph{Boundary Origin Audit} (Closure Principle), Docs/Research.
\bibitem{c01:QG286} AT-QG Phase 286, \emph{Difference Duality Audit}, Docs/Research.
\bibitem{c01:QG292} AT-QG Phase 292, \emph{Foundation Stress Test}, Docs/Research.
\bibitem{c01:QG294} AT-QG Phase 294, \emph{Minimal Theory Audit}, Docs/Research.
\bibitem{c01:QG295} AT-QG Phase 295, \emph{Spectrum Necessity Audit}, Docs/Research.
\bibitem{c01:QG301} AT-QG Phase 301, \emph{Duality Prediction Audit}, Docs/Research.
\bibitem{c01:QG216} AT-QG Phase 216, \emph{Quantum Amplitude Origin}, Docs/Research.
\bibitem{c01:QG285} AT-QG Phase 285, \emph{Psi Reinterpretation Audit}, Docs/Research.
\bibitem{c01:MONO006} AT-MONO006, \emph{A01 Actualization Origin Audit}, Docs/Zenodo/Monograph.
\bibitem{c01:MONO007} AT-MONO007, \emph{Referee-Readiness Resolution}, Docs/Zenodo/Monograph.
\end{thebibliography}
 after the preamble that defines
%         the theorem environments below (or include via the master file).
% ======================================================================

% ---- Theorem environments (define in the master preamble if not present) ----
% \newtheorem{definition}{Definition}[chapter]
% \newtheorem{lemma}[definition]{Lemma}
% \newtheorem{theorem}[definition]{Theorem}
% \newtheorem{corollary}[definition]{Corollary}
% \newtheorem{remark}[definition]{Remark}

\chapter{The Difference}
\label{c01:chap:difference}

\section{Introduction}
\label{c01:sec:difference-introduction}

The Actualization Theory (AT) is a deterministic theory of structure, complexity, and random
actualization. Its central claim is that the totality of observable physics can be derived
from a single primitive notion --- \emph{Difference} --- together with a tensor reference
metric $\eta$, through the canonical hierarchy
\begin{equation}
\label{c01:eq:canonical-core}
\text{Difference}\;\longrightarrow\;\text{Actualization}\;\longrightarrow\;\text{Inevitable Spectrum}\;\longrightarrow\;\text{Physics}.
\end{equation}

This chapter develops the first member of that hierarchy. We establish, with formal
precision, that Difference is the sole irreducible primitive of the theory; that its status
is that of a \emph{fundamental boundary} --- the deepest notion, which cannot itself be
derived and whose own boundary is characterized by the Closure Principle
\cite{c01:QG278,c01:QG282}; and that the two great physical sectors of the theory, the scalar
counting measure $\rho$ and the tensor field $\psi$, are not independent primitives but the
\emph{trace and traceless faces of the one Difference} \cite{c01:QG286,c01:QG301}.

The results of this chapter are not new. They are the canonical end-state of the QG278--QG318
program, assembled here in a monograph form suitable for archival publication. Throughout, we
adhere to the accepted derivation record and introduce no new physics and no speculation.

\section{Primitive Status}
\label{c01:sec:primitive-status}

\begin{definition}[Difference]
\label{c01:def:difference}
\emph{Difference} is the primitive notion of the theory: the counting difference from a
uniform background, whose count-conservation law is its \emph{definitional identity}
\cite{c01:QG268}. A Q-event is a unit of Difference.
\end{definition}

\begin{definition}[Tensor reference metric]
\label{c01:def:eta}
\emph{$\eta$} is the tensor reference metric: the conformal reference with respect to which
the metric takes the form $g = \rho^{2/d}\,\eta$, and with respect to which the Weyl content
$\psi$ is defined \cite{c01:QG286}.
\end{definition}

\begin{theorem}[Minimal primitive set]
\label{c01:thm:minimal-primitives}
The minimal set of irreducible primitives of AT is $\{\text{Difference},\eta\}$: exactly
two items. No proper subset supports the derived hierarchy \eqref{c01:eq:canonical-core}, and no
further reduction is possible.
\end{theorem}

\begin{proof}
The minimality of $\{\text{Difference},\eta\}$ is established by the foundation stress test
\cite{c01:QG292}. Removing Difference collapses \emph{all} five layers of the theory
(actualization, conservation, resonance, spectrum, physics), because count conservation is the
definitional identity of the primitive: without a unit there is no count, without a count no
network, without a network no spectrum, and without a spectrum no physics. Removing $\eta$
leaves all five layers intact as a \emph{scalar} readout --- masses, couplings, mixings,
cosmological fractions, and spectral observables are all $\rho$-face reads --- and destroys
only the tensor sub-sector, because $\psi$ is defined against $\eta$. Hence Difference is
necessary for the whole hierarchy and $\eta$ is necessary precisely for the tensor content;
neither is derivable from the other, and together they form the irreducible pair.
\end{proof}

\begin{corollary}[Actualization is not a primitive]
\label{c01:cor:actualization-derived}
Actualization is \emph{not} a primitive. It is Difference's own count-producing process.
\end{corollary}

\begin{proof}
By the same removal test \cite{c01:QG292,c01:MONO006}, Actualization fails when Difference is removed
(it needs the definitional unit) and survives when $\eta$ is removed (counting needs no metric
reference). It is therefore neither an independent primitive nor a function of
$\{\text{Difference},\eta\}$ in the tensor sense; it is the dynamical face of the one
primitive. Its fixed point is the closure $N=96$ \cite{c01:QG282}.
\end{proof}

\begin{remark}
The canonical core \eqref{c01:eq:canonical-core} is therefore read: \emph{Difference} (the
primitive) produces \emph{Actualization} (its derived process), which converges to the
\emph{Inevitable Spectrum} (the fixed point), from which \emph{Physics} (all observable
sectors) emerges. Only the first member and the reference metric are primitives.
\end{remark}

\section{Difference as Fundamental Boundary}
\label{c01:sec:fundamental-boundary}

\begin{definition}[Fundamental boundary]
\label{c01:def:boundary}
A \emph{fundamental boundary} is a notion that (i) is irreducible --- no deeper notion
explains it; (ii) supports the entire derived hierarchy; and (iii) is itself bounded only by
its own fixed point, not by any imported assumption.
\end{definition}

\begin{theorem}[Difference is the fundamental boundary]
\label{c01:thm:fundamental-boundary}
Difference is the fundamental boundary of AT.
\end{theorem}

\begin{proof}
The boundary audits \cite{c01:QG278,c01:QG279} establish the three criteria. (i) \emph{Irreducible:}
the search for a deeper primitive terminates at Difference; attempts to redefine it as an
artifact of language fail, and it is declared a \emph{true} fundamental boundary. (ii)
\emph{Support:} by Theorem~\ref{c01:thm:minimal-primitives}, removing Difference collapses every
layer --- it is the root of the entire derivation graph. (iii) \emph{Self-bounded:} the
boundary of Difference is not imported but derived --- it is the fixed point $N=96$ of its own
actualization process, per the Closure Principle \cite{c01:QG282}. Hence Difference satisfies the
three defining properties of a fundamental boundary.
\end{proof}

\begin{lemma}[Boundary is fixed point]
\label{c01:lem:boundary-fixed-point}
The boundary of the theory is the stable fixed point of the actualization process: the
topology saturates at zero residual link growth, and every initial pattern converges to the
same final geometry \cite{c01:QG282}.
\end{lemma}

\begin{proof}
The actualization dynamics has a self-reinforcing feedback (activity to links to activity)
that is bounded by network capacity; positive feedback therefore saturates at the stable
fixed point, giving zero residual link growth and a content-independent final geometry
\cite{c01:QG282}. The boundary is the closure of this dynamics.
\end{proof}

\section{Difference Duality ($\rho$ and $\psi$)}
\label{c01:sec:duality}

\begin{definition}[The two faces]
\label{c01:def:two-faces}
Let $A_{ij}$ be the rank-2 object representing the difference structure (the connectivity /
the stress / the difference). Its irreducible decomposition is
\begin{equation}
\label{c01:eq:decomposition}
A_{ij} \;=\; \underbrace{\frac{1}{d}\,\mathrm{Tr}(A)\,\delta_{ij}}_{\text{trace: } \rho}
\;+\;\underbrace{A_{ij}-\frac{1}{d}\,\mathrm{Tr}(A)\,\delta_{ij}}_{\text{traceless: } \psi}.
\end{equation}
At $d=3$, $A_{ij}$ has $d(d+1)/2=6$ components: one trace component ($\rho$, the scalar) and
five traceless components ($\psi$, the tensor, of which two are the transverse-traceless
polarizations) \cite{c01:QG286}.
\end{definition}

\begin{theorem}[Difference duality]
\label{c01:thm:duality}
The scalar counting measure $\rho$ and the tensor field $\psi$ are the trace and traceless
faces of the \emph{one} Difference. They are not independent primitives; each is a projection
of the same rank-2 difference structure.
\end{theorem}

\begin{proof}
$\rho$ is the scalar counting measure: the normalized count share $|\psi|^2=\rho$
\cite{c01:QG216}, the \emph{isotropic} difference from the uniform background --- the
\emph{count} difference. $\psi$ is the tensor field: the spin-2 Weyl content \cite{c01:QG285},
the \emph{anisotropic} difference from conformal flatness --- the \emph{orientation}
difference. By the rank-2 decomposition theorem \eqref{c01:eq:decomposition}, the trace and the
traceless part are determined by $A_{ij}$; neither is an independent input. Hence $\rho$ and
$\psi$ are dual projections of the same Difference \cite{c01:QG286}.
\end{proof}

\begin{corollary}[Duality is invariant]
\label{c01:cor:duality-invariant}
The reinterpretation of the theory through the duality
$\text{Difference}\longrightarrow\{\rho,\psi\}$ changes \emph{meaning} without changing any
\emph{number}: every frozen prediction is reproduced exactly \cite{c01:QG301}.
\end{corollary}

\begin{proof}
The post-duality invariance audit recomputes every established prediction through the
reduced hierarchy and finds each unchanged: no formula is retuned and no result is lost
\cite{c01:QG301}. The duality is therefore a reinterpretation of content, not a modification of
the derived mathematics.
\end{proof}

\begin{remark}[$\rho$ needs no reference; $\psi$ does]
The scalar face $\rho$ is defined without reference to $\eta$; the tensor face $\psi$ is
defined \emph{against} $\eta$ (conformal flatness). This asymmetry is the reason
Theorem~\ref{c01:thm:minimal-primitives} holds: $\eta$ is necessary exactly for the tensor
sub-sector.
\end{remark}

\section{Relation to Closure}
\label{c01:sec:closure}

\begin{theorem}[Closure Principle]
\label{c01:thm:closure}
The boundary of the theory is the closure of Difference's own dynamics: the primitive is the
\emph{process}, and the boundary is its \emph{fixed point}. The closure is
\emph{derived}, not primitive \cite{c01:QG282}.
\end{theorem}

\begin{proof}
The origin of the boundary is the $N=96$ attractor: the observable sector is the converged
attractor of the actualization dynamics, with topology saturation and content-independent
convergence (Lemma~\ref{c01:lem:boundary-fixed-point}). The boundary is therefore not an
imported cut-off but the stable fixed point of the process itself; the closure principle
states exactly this --- the process converges, the boundary is its closure \cite{c01:QG282}.
\end{proof}

\begin{corollary}[Boundary characterizes, does not derive]
\label{c01:cor:closure-characterizes}
The Closure Principle \emph{characterizes} the fundamental boundary; it does not \emph{derive}
the primitive Difference.
\end{corollary}

\begin{proof}
Difference is the primitive: by Definition~\ref{c01:def:difference} its count-conservation law
is its definitional identity, and by Theorem~\ref{c01:thm:fundamental-boundary} it is
irreducible. The Closure Principle describes how Difference's own dynamics closes at $N=96$;
it presupposes the process, hence the primitive, and cannot be a derivation of it. This is
the referee-safe reading adopted by the canonical monograph \cite{c01:MONO007}.
\end{proof}

\section{Consequences}
\label{c01:sec:consequences}

We collect the immediate formal consequences of the foregoing.

\begin{lemma}[One primitive, two faces, one hierarchy]
\label{c01:lem:consequence-1}
The entire observable content of the theory is carried by the two faces of the one Difference
through a single canonical hierarchy: the scalar face feeds the full spectrum of
observables; the tensor face feeds the gravitational (Weyl) sector; both are projections of
the same primitive \cite{c01:QG286,c01:QG301}.
\end{lemma}

\begin{theorem}[Spectrum is inevitable]
\label{c01:thm:spectrum-inevitable}
The D96 spectrum is the inevitable output of the actualization attractor --- the Laplacian
eigenspectrum of the converged network --- and not an additional primitive \cite{c01:QG295}.
\end{theorem}

\begin{proof}
The spectrum is a deterministic function of the converged network: the $95$ positive modes
are the graph-Laplacian eigenvalues of the attractor, the zero mode is the background, and no
choice enters. The spectrum carries no independent input; it is the output of the process
whose primitive is Difference \cite{c01:QG295,c01:QG282}.
\end{proof}

\begin{corollary}[Canonical core]
\label{c01:cor:core}
The minimal hierarchy
\[
\text{Difference}\;\longrightarrow\;\text{Actualization}\;\longrightarrow\;\text{Inevitable Spectrum}\;\longrightarrow\;\text{Physics}
\]
is complete: every intermediate layer is derivable from the primitives and the attractor, and
no further primitive is required \cite{c01:QG292,c01:QG294}.
\end{corollary}

\begin{remark}[Referee-safe completeness]
Where the monograph claims completeness, it is understood in the qualified sense: every
\emph{observable} is derived from the canonical hierarchy; the standard-model
\emph{Lagrangian} and its dynamics are \emph{hosted} --- a boundary, not a derivation gap
\cite{c01:MONO007}.
\end{remark}

\section{Summary}
\label{c01:sec:summary}

This chapter has established the first member of the canonical architecture. We have proved:

\begin{enumerate}
\item \textbf{Primitive status} (Theorem~\ref{c01:thm:minimal-primitives}): the minimal primitive
set is $\{\text{Difference},\eta\}$, and Actualization is derived from Difference
(Corollary~\ref{c01:cor:actualization-derived}).
\item \textbf{Fundamental boundary} (Theorem~\ref{c01:thm:fundamental-boundary}): Difference is
irreducible, supports the whole hierarchy, and is self-bounded by its own fixed point.
\item \textbf{The duality} (Theorem~\ref{c01:thm:duality}): $\rho$ and $\psi$ are the trace and
traceless faces of the one Difference; the reinterpretation is invariant
(Corollary~\ref{c01:cor:duality-invariant}).
\item \textbf{Closure} (Theorem~\ref{c01:thm:closure}): the boundary is the fixed point of
Difference's own dynamics, characterizing rather than deriving the primitive
(Corollary~\ref{c01:cor:closure-characterizes}).
\item \textbf{Consequences} (Theorem~\ref{c01:thm:spectrum-inevitable},
Corollary~\ref{c01:cor:core}): the spectrum is inevitable, and the canonical core is complete.
\end{enumerate}

The Difference is therefore not an unexplained name but a precise object: the counting
difference whose conservation is its identity, whose dynamics closes at $N=96$, and whose two
faces --- the scalar measure $\rho$ and the tensor field $\psi$ --- generate, respectively,
the spectral readout of all physics and the gravitational content of the theory. All further
chapters of this monograph derive from the structure established here.

\begin{thebibliography}{99}
\bibitem{c01:QG278} AT-QG Phase 278, \emph{Fundamental Boundary Audit}, Docs/Research.
\bibitem{c01:QG279} AT-QG Phase 279, \emph{Boundary Legitimacy Audit}, Docs/Research.
\bibitem{c01:QG268} AT-QG Phase 268, \emph{Count Conservation Origin}, Docs/Research.
\bibitem{c01:QG282} AT-QG Phase 282, \emph{Boundary Origin Audit} (Closure Principle), Docs/Research.
\bibitem{c01:QG286} AT-QG Phase 286, \emph{Difference Duality Audit}, Docs/Research.
\bibitem{c01:QG292} AT-QG Phase 292, \emph{Foundation Stress Test}, Docs/Research.
\bibitem{c01:QG294} AT-QG Phase 294, \emph{Minimal Theory Audit}, Docs/Research.
\bibitem{c01:QG295} AT-QG Phase 295, \emph{Spectrum Necessity Audit}, Docs/Research.
\bibitem{c01:QG301} AT-QG Phase 301, \emph{Duality Prediction Audit}, Docs/Research.
\bibitem{c01:QG216} AT-QG Phase 216, \emph{Quantum Amplitude Origin}, Docs/Research.
\bibitem{c01:QG285} AT-QG Phase 285, \emph{Psi Reinterpretation Audit}, Docs/Research.
\bibitem{c01:MONO006} AT-MONO006, \emph{A01 Actualization Origin Audit}, Docs/Zenodo/Monograph.
\bibitem{c01:MONO007} AT-MONO007, \emph{Referee-Readiness Resolution}, Docs/Zenodo/Monograph.
\end{thebibliography}
 after the preamble that defines
%         the theorem environments below (or include via the master file).
% ======================================================================

% ---- Theorem environments (define in the master preamble if not present) ----
% \newtheorem{definition}{Definition}[chapter]
% \newtheorem{lemma}[definition]{Lemma}
% \newtheorem{theorem}[definition]{Theorem}
% \newtheorem{corollary}[definition]{Corollary}
% \newtheorem{remark}[definition]{Remark}

\chapter{The Difference}
\label{c01:chap:difference}

\section{Introduction}
\label{c01:sec:difference-introduction}

The Actualization Theory (AT) is a deterministic theory of structure, complexity, and random
actualization. Its central claim is that the totality of observable physics can be derived
from a single primitive notion --- \emph{Difference} --- together with a tensor reference
metric $\eta$, through the canonical hierarchy
\begin{equation}
\label{c01:eq:canonical-core}
\text{Difference}\;\longrightarrow\;\text{Actualization}\;\longrightarrow\;\text{Inevitable Spectrum}\;\longrightarrow\;\text{Physics}.
\end{equation}

This chapter develops the first member of that hierarchy. We establish, with formal
precision, that Difference is the sole irreducible primitive of the theory; that its status
is that of a \emph{fundamental boundary} --- the deepest notion, which cannot itself be
derived and whose own boundary is characterized by the Closure Principle
\cite{c01:QG278,c01:QG282}; and that the two great physical sectors of the theory, the scalar
counting measure $\rho$ and the tensor field $\psi$, are not independent primitives but the
\emph{trace and traceless faces of the one Difference} \cite{c01:QG286,c01:QG301}.

The results of this chapter are not new. They are the canonical end-state of the QG278--QG318
program, assembled here in a monograph form suitable for archival publication. Throughout, we
adhere to the accepted derivation record and introduce no new physics and no speculation.

\section{Primitive Status}
\label{c01:sec:primitive-status}

\begin{definition}[Difference]
\label{c01:def:difference}
\emph{Difference} is the primitive notion of the theory: the counting difference from a
uniform background, whose count-conservation law is its \emph{definitional identity}
\cite{c01:QG268}. A Q-event is a unit of Difference.
\end{definition}

\begin{definition}[Tensor reference metric]
\label{c01:def:eta}
\emph{$\eta$} is the tensor reference metric: the conformal reference with respect to which
the metric takes the form $g = \rho^{2/d}\,\eta$, and with respect to which the Weyl content
$\psi$ is defined \cite{c01:QG286}.
\end{definition}

\begin{theorem}[Minimal primitive set]
\label{c01:thm:minimal-primitives}
The minimal set of irreducible primitives of AT is $\{\text{Difference},\eta\}$: exactly
two items. No proper subset supports the derived hierarchy \eqref{c01:eq:canonical-core}, and no
further reduction is possible.
\end{theorem}

\begin{proof}
The minimality of $\{\text{Difference},\eta\}$ is established by the foundation stress test
\cite{c01:QG292}. Removing Difference collapses \emph{all} five layers of the theory
(actualization, conservation, resonance, spectrum, physics), because count conservation is the
definitional identity of the primitive: without a unit there is no count, without a count no
network, without a network no spectrum, and without a spectrum no physics. Removing $\eta$
leaves all five layers intact as a \emph{scalar} readout --- masses, couplings, mixings,
cosmological fractions, and spectral observables are all $\rho$-face reads --- and destroys
only the tensor sub-sector, because $\psi$ is defined against $\eta$. Hence Difference is
necessary for the whole hierarchy and $\eta$ is necessary precisely for the tensor content;
neither is derivable from the other, and together they form the irreducible pair.
\end{proof}

\begin{corollary}[Actualization is not a primitive]
\label{c01:cor:actualization-derived}
Actualization is \emph{not} a primitive. It is Difference's own count-producing process.
\end{corollary}

\begin{proof}
By the same removal test \cite{c01:QG292,c01:MONO006}, Actualization fails when Difference is removed
(it needs the definitional unit) and survives when $\eta$ is removed (counting needs no metric
reference). It is therefore neither an independent primitive nor a function of
$\{\text{Difference},\eta\}$ in the tensor sense; it is the dynamical face of the one
primitive. Its fixed point is the closure $N=96$ \cite{c01:QG282}.
\end{proof}

\begin{remark}
The canonical core \eqref{c01:eq:canonical-core} is therefore read: \emph{Difference} (the
primitive) produces \emph{Actualization} (its derived process), which converges to the
\emph{Inevitable Spectrum} (the fixed point), from which \emph{Physics} (all observable
sectors) emerges. Only the first member and the reference metric are primitives.
\end{remark}

\section{Difference as Fundamental Boundary}
\label{c01:sec:fundamental-boundary}

\begin{definition}[Fundamental boundary]
\label{c01:def:boundary}
A \emph{fundamental boundary} is a notion that (i) is irreducible --- no deeper notion
explains it; (ii) supports the entire derived hierarchy; and (iii) is itself bounded only by
its own fixed point, not by any imported assumption.
\end{definition}

\begin{theorem}[Difference is the fundamental boundary]
\label{c01:thm:fundamental-boundary}
Difference is the fundamental boundary of AT.
\end{theorem}

\begin{proof}
The boundary audits \cite{c01:QG278,c01:QG279} establish the three criteria. (i) \emph{Irreducible:}
the search for a deeper primitive terminates at Difference; attempts to redefine it as an
artifact of language fail, and it is declared a \emph{true} fundamental boundary. (ii)
\emph{Support:} by Theorem~\ref{c01:thm:minimal-primitives}, removing Difference collapses every
layer --- it is the root of the entire derivation graph. (iii) \emph{Self-bounded:} the
boundary of Difference is not imported but derived --- it is the fixed point $N=96$ of its own
actualization process, per the Closure Principle \cite{c01:QG282}. Hence Difference satisfies the
three defining properties of a fundamental boundary.
\end{proof}

\begin{lemma}[Boundary is fixed point]
\label{c01:lem:boundary-fixed-point}
The boundary of the theory is the stable fixed point of the actualization process: the
topology saturates at zero residual link growth, and every initial pattern converges to the
same final geometry \cite{c01:QG282}.
\end{lemma}

\begin{proof}
The actualization dynamics has a self-reinforcing feedback (activity to links to activity)
that is bounded by network capacity; positive feedback therefore saturates at the stable
fixed point, giving zero residual link growth and a content-independent final geometry
\cite{c01:QG282}. The boundary is the closure of this dynamics.
\end{proof}

\section{Difference Duality ($\rho$ and $\psi$)}
\label{c01:sec:duality}

\begin{definition}[The two faces]
\label{c01:def:two-faces}
Let $A_{ij}$ be the rank-2 object representing the difference structure (the connectivity /
the stress / the difference). Its irreducible decomposition is
\begin{equation}
\label{c01:eq:decomposition}
A_{ij} \;=\; \underbrace{\frac{1}{d}\,\mathrm{Tr}(A)\,\delta_{ij}}_{\text{trace: } \rho}
\;+\;\underbrace{A_{ij}-\frac{1}{d}\,\mathrm{Tr}(A)\,\delta_{ij}}_{\text{traceless: } \psi}.
\end{equation}
At $d=3$, $A_{ij}$ has $d(d+1)/2=6$ components: one trace component ($\rho$, the scalar) and
five traceless components ($\psi$, the tensor, of which two are the transverse-traceless
polarizations) \cite{c01:QG286}.
\end{definition}

\begin{theorem}[Difference duality]
\label{c01:thm:duality}
The scalar counting measure $\rho$ and the tensor field $\psi$ are the trace and traceless
faces of the \emph{one} Difference. They are not independent primitives; each is a projection
of the same rank-2 difference structure.
\end{theorem}

\begin{proof}
$\rho$ is the scalar counting measure: the normalized count share $|\psi|^2=\rho$
\cite{c01:QG216}, the \emph{isotropic} difference from the uniform background --- the
\emph{count} difference. $\psi$ is the tensor field: the spin-2 Weyl content \cite{c01:QG285},
the \emph{anisotropic} difference from conformal flatness --- the \emph{orientation}
difference. By the rank-2 decomposition theorem \eqref{c01:eq:decomposition}, the trace and the
traceless part are determined by $A_{ij}$; neither is an independent input. Hence $\rho$ and
$\psi$ are dual projections of the same Difference \cite{c01:QG286}.
\end{proof}

\begin{corollary}[Duality is invariant]
\label{c01:cor:duality-invariant}
The reinterpretation of the theory through the duality
$\text{Difference}\longrightarrow\{\rho,\psi\}$ changes \emph{meaning} without changing any
\emph{number}: every frozen prediction is reproduced exactly \cite{c01:QG301}.
\end{corollary}

\begin{proof}
The post-duality invariance audit recomputes every established prediction through the
reduced hierarchy and finds each unchanged: no formula is retuned and no result is lost
\cite{c01:QG301}. The duality is therefore a reinterpretation of content, not a modification of
the derived mathematics.
\end{proof}

\begin{remark}[$\rho$ needs no reference; $\psi$ does]
The scalar face $\rho$ is defined without reference to $\eta$; the tensor face $\psi$ is
defined \emph{against} $\eta$ (conformal flatness). This asymmetry is the reason
Theorem~\ref{c01:thm:minimal-primitives} holds: $\eta$ is necessary exactly for the tensor
sub-sector.
\end{remark}

\section{Relation to Closure}
\label{c01:sec:closure}

\begin{theorem}[Closure Principle]
\label{c01:thm:closure}
The boundary of the theory is the closure of Difference's own dynamics: the primitive is the
\emph{process}, and the boundary is its \emph{fixed point}. The closure is
\emph{derived}, not primitive \cite{c01:QG282}.
\end{theorem}

\begin{proof}
The origin of the boundary is the $N=96$ attractor: the observable sector is the converged
attractor of the actualization dynamics, with topology saturation and content-independent
convergence (Lemma~\ref{c01:lem:boundary-fixed-point}). The boundary is therefore not an
imported cut-off but the stable fixed point of the process itself; the closure principle
states exactly this --- the process converges, the boundary is its closure \cite{c01:QG282}.
\end{proof}

\begin{corollary}[Boundary characterizes, does not derive]
\label{c01:cor:closure-characterizes}
The Closure Principle \emph{characterizes} the fundamental boundary; it does not \emph{derive}
the primitive Difference.
\end{corollary}

\begin{proof}
Difference is the primitive: by Definition~\ref{c01:def:difference} its count-conservation law
is its definitional identity, and by Theorem~\ref{c01:thm:fundamental-boundary} it is
irreducible. The Closure Principle describes how Difference's own dynamics closes at $N=96$;
it presupposes the process, hence the primitive, and cannot be a derivation of it. This is
the referee-safe reading adopted by the canonical monograph \cite{c01:MONO007}.
\end{proof}

\section{Consequences}
\label{c01:sec:consequences}

We collect the immediate formal consequences of the foregoing.

\begin{lemma}[One primitive, two faces, one hierarchy]
\label{c01:lem:consequence-1}
The entire observable content of the theory is carried by the two faces of the one Difference
through a single canonical hierarchy: the scalar face feeds the full spectrum of
observables; the tensor face feeds the gravitational (Weyl) sector; both are projections of
the same primitive \cite{c01:QG286,c01:QG301}.
\end{lemma}

\begin{theorem}[Spectrum is inevitable]
\label{c01:thm:spectrum-inevitable}
The D96 spectrum is the inevitable output of the actualization attractor --- the Laplacian
eigenspectrum of the converged network --- and not an additional primitive \cite{c01:QG295}.
\end{theorem}

\begin{proof}
The spectrum is a deterministic function of the converged network: the $95$ positive modes
are the graph-Laplacian eigenvalues of the attractor, the zero mode is the background, and no
choice enters. The spectrum carries no independent input; it is the output of the process
whose primitive is Difference \cite{c01:QG295,c01:QG282}.
\end{proof}

\begin{corollary}[Canonical core]
\label{c01:cor:core}
The minimal hierarchy
\[
\text{Difference}\;\longrightarrow\;\text{Actualization}\;\longrightarrow\;\text{Inevitable Spectrum}\;\longrightarrow\;\text{Physics}
\]
is complete: every intermediate layer is derivable from the primitives and the attractor, and
no further primitive is required \cite{c01:QG292,c01:QG294}.
\end{corollary}

\begin{remark}[Referee-safe completeness]
Where the monograph claims completeness, it is understood in the qualified sense: every
\emph{observable} is derived from the canonical hierarchy; the standard-model
\emph{Lagrangian} and its dynamics are \emph{hosted} --- a boundary, not a derivation gap
\cite{c01:MONO007}.
\end{remark}

\section{Summary}
\label{c01:sec:summary}

This chapter has established the first member of the canonical architecture. We have proved:

\begin{enumerate}
\item \textbf{Primitive status} (Theorem~\ref{c01:thm:minimal-primitives}): the minimal primitive
set is $\{\text{Difference},\eta\}$, and Actualization is derived from Difference
(Corollary~\ref{c01:cor:actualization-derived}).
\item \textbf{Fundamental boundary} (Theorem~\ref{c01:thm:fundamental-boundary}): Difference is
irreducible, supports the whole hierarchy, and is self-bounded by its own fixed point.
\item \textbf{The duality} (Theorem~\ref{c01:thm:duality}): $\rho$ and $\psi$ are the trace and
traceless faces of the one Difference; the reinterpretation is invariant
(Corollary~\ref{c01:cor:duality-invariant}).
\item \textbf{Closure} (Theorem~\ref{c01:thm:closure}): the boundary is the fixed point of
Difference's own dynamics, characterizing rather than deriving the primitive
(Corollary~\ref{c01:cor:closure-characterizes}).
\item \textbf{Consequences} (Theorem~\ref{c01:thm:spectrum-inevitable},
Corollary~\ref{c01:cor:core}): the spectrum is inevitable, and the canonical core is complete.
\end{enumerate}

The Difference is therefore not an unexplained name but a precise object: the counting
difference whose conservation is its identity, whose dynamics closes at $N=96$, and whose two
faces --- the scalar measure $\rho$ and the tensor field $\psi$ --- generate, respectively,
the spectral readout of all physics and the gravitational content of the theory. All further
chapters of this monograph derive from the structure established here.

\begin{thebibliography}{99}
\bibitem{c01:QG278} AT-QG Phase 278, \emph{Fundamental Boundary Audit}, Docs/Research.
\bibitem{c01:QG279} AT-QG Phase 279, \emph{Boundary Legitimacy Audit}, Docs/Research.
\bibitem{c01:QG268} AT-QG Phase 268, \emph{Count Conservation Origin}, Docs/Research.
\bibitem{c01:QG282} AT-QG Phase 282, \emph{Boundary Origin Audit} (Closure Principle), Docs/Research.
\bibitem{c01:QG286} AT-QG Phase 286, \emph{Difference Duality Audit}, Docs/Research.
\bibitem{c01:QG292} AT-QG Phase 292, \emph{Foundation Stress Test}, Docs/Research.
\bibitem{c01:QG294} AT-QG Phase 294, \emph{Minimal Theory Audit}, Docs/Research.
\bibitem{c01:QG295} AT-QG Phase 295, \emph{Spectrum Necessity Audit}, Docs/Research.
\bibitem{c01:QG301} AT-QG Phase 301, \emph{Duality Prediction Audit}, Docs/Research.
\bibitem{c01:QG216} AT-QG Phase 216, \emph{Quantum Amplitude Origin}, Docs/Research.
\bibitem{c01:QG285} AT-QG Phase 285, \emph{Psi Reinterpretation Audit}, Docs/Research.
\bibitem{c01:MONO006} AT-MONO006, \emph{A01 Actualization Origin Audit}, Docs/Zenodo/Monograph.
\bibitem{c01:MONO007} AT-MONO007, \emph{Referee-Readiness Resolution}, Docs/Zenodo/Monograph.
\end{thebibliography}

% ======================================================================
%  AT-MONO101 — Chapter 2: The Tensor Reference η
%  Canonical Monograph (v2.0) · The Actualization Theory:
%  A Reconstruction of Physics from Difference, Actualization and Spectrum
%
%  Status: COMPLETE — PASS-READY (MONO007 referee-ready architecture)
%  Inputs: Chapter01_Difference.tex, ATQG_CanonicalMonograph.md,
%          ATQG_Mono007RefereeReadiness.md
%  Method: assembly only — canonical results only, no new physics, no speculation
%  Citations: QG285 (psi reinterpretation), QG286 (difference duality),
%             QG289 (anchor inventory), QG290 (irreducible framework),
%             QG291 (framework necessity), QG292 (minimal foundation),
%             MONO006 (actualization derived), MONO007 (referee readiness)
%
%  Usage: % ======================================================================
%  AT-MONO101 — Chapter 2: The Tensor Reference η
%  Canonical Monograph (v2.0) · The Actualization Theory:
%  A Reconstruction of Physics from Difference, Actualization and Spectrum
%
%  Status: COMPLETE — PASS-READY (MONO007 referee-ready architecture)
%  Inputs: Chapter01_Difference.tex, ATQG_CanonicalMonograph.md,
%          ATQG_Mono007RefereeReadiness.md
%  Method: assembly only — canonical results only, no new physics, no speculation
%  Citations: QG285 (psi reinterpretation), QG286 (difference duality),
%             QG289 (anchor inventory), QG290 (irreducible framework),
%             QG291 (framework necessity), QG292 (minimal foundation),
%             MONO006 (actualization derived), MONO007 (referee readiness)
%
%  Usage: % ======================================================================
%  AT-MONO101 — Chapter 2: The Tensor Reference η
%  Canonical Monograph (v2.0) · The Actualization Theory:
%  A Reconstruction of Physics from Difference, Actualization and Spectrum
%
%  Status: COMPLETE — PASS-READY (MONO007 referee-ready architecture)
%  Inputs: Chapter01_Difference.tex, ATQG_CanonicalMonograph.md,
%          ATQG_Mono007RefereeReadiness.md
%  Method: assembly only — canonical results only, no new physics, no speculation
%  Citations: QG285 (psi reinterpretation), QG286 (difference duality),
%             QG289 (anchor inventory), QG290 (irreducible framework),
%             QG291 (framework necessity), QG292 (minimal foundation),
%             MONO006 (actualization derived), MONO007 (referee readiness)
%
%  Usage: \input{Chapter02_Eta.tex} after the preamble that defines
%         the theorem environments below (or include via the master file).
% ======================================================================

% ---- Theorem environments (define in the master preamble if not present) ----
% \newtheorem{definition}{Definition}[chapter]
% \newtheorem{lemma}[definition]{Lemma}
% \newtheorem{theorem}[definition]{Theorem}
% \newtheorem{corollary}[definition]{Corollary}
% \newtheorem{remark}[definition]{Remark}

\chapter{The Tensor Reference \texorpdfstring{$\eta$}{eta}}
\label{c02:chap:eta}

\section{Introduction}
\label{c02:sec:eta-introduction}

In Chapter~\ref{c01:chap:difference} we established that Difference is the fundamental boundary
of the theory, and that the scalar counting measure $\rho$ and the tensor field $\psi$ are
the trace and traceless faces of the one Difference \cite{c02:QG286}. That duality, however, is
not defined in a vacuum: the trace--traceless decomposition, the conformal flatness against
which the Weyl content is read, and the very notion of the tensor face $\psi$ all presuppose
a \emph{reference structure}. This chapter develops that structure --- the tensor reference
metric $\eta$.

The results of this chapter are the canonical end-state of the framework inventory and
necessity program \cite{c02:QG289,c02:QG290,c02:QG291,c02:QG292}: $\eta$ is the second primitive of the
theory, necessary and irreducible; the framework $\{\text{Difference},\eta\}$ is minimal;
and the numerical constant $\pi$, often presented alongside $\eta$ in earlier treatments, is
\emph{not} a primitive but a boundary constant --- redundant as a theory input and separated
from the primitive reference by the referee-ready architecture of MONO007
\cite{c02:MONO007}. Throughout we use only accepted canonical results, and introduce no new
physics and no speculation.

\section{Necessity of \texorpdfstring{$\eta$}{eta}}
\label{c02:sec:eta-necessity}

\begin{definition}[Tensor reference metric]
\label{c02:def:eta}
\emph{$\eta$} is the tensor reference metric: the conformal reference with respect to which
the physical metric takes the form
\begin{equation}
\label{c02:eq:metric-ansatz}
g \;=\; \rho^{2/d}\,\eta,
\end{equation}
where $\rho$ is the scalar counting measure and $d$ the spatial dimension
\cite{c02:QG290,c02:QG285}.
\end{definition}

\begin{theorem}[Necessity of the tensor reference]
\label{c02:thm:eta-necessary}
The tensor reference $\eta$ is NECESSARY: without it there is no trace, no traceless part,
no conformal flatness, and no reading of the difference structure.
\end{theorem}

\begin{proof}
The duality Difference $\to \{\rho,\psi\}$ is the trace/traceless decomposition
\cite{c02:QG286}
\begin{equation}
\label{c02:eq:decomposition}
A_{ij} \;=\; \frac{1}{d}\,\mathrm{Tr}(A)\,\delta_{ij}
\;+\;\Bigl(A_{ij}-\frac{1}{d}\,\mathrm{Tr}(A)\,\delta_{ij}\Bigr),
\end{equation}
of the rank-2 difference object $A_{ij}$. The trace ($\rho$), the traceless part ($\psi$),
and the Weyl content --- the difference from conformal flatness \cite{c02:QG285} --- are all
\emph{defined against} the reference $\eta$: conformal flatness is meaningful only with a
flat reference, and the Weyl tensor is the non-conformal content of the metric relative to
that reference. Without $\eta$ there is no conformal flatness, hence no notion of
"difference from conformal flatness," hence no tensor face. The framework inventory audit
\cite{c02:QG291} classifies $\eta$ accordingly: not \emph{derived} (no count produces a metric,
and $\eta$ carries no scale), not \emph{redundant} (removing $\eta$ removes the reading), but
\emph{necessary} --- the reference structure the framework reads against.
\end{proof}

\begin{lemma}[$\eta$ is not derivable from Difference]
\label{c02:lem:eta-not-derived}
No counting measure, no matter how refined, produces a metric reference; $\eta$ carries no
scale and is not a function of $\rho$.
\end{lemma}

\begin{proof}
The framework inventory establishes that $\eta$ is not derived: a metric reference is a
\emph{reading structure}, not an output of the count. $\eta$ carries no scale --- the scale
triad enters the theory elsewhere --- and the metric ansatz \eqref{c02:eq:metric-ansatz} takes
$\eta$ as an input on which the conformal factor $\rho^{2/d}$ acts. Hence $\eta$ is an
independent input, not a function of the primitive Difference.
\end{proof}

\begin{corollary}[Two independent primitives]
\label{c02:cor:eta-independent}
Difference and $\eta$ are independent primitives: neither is derivable from the other.
\end{corollary}

\section{The Conformal Reference Framework}
\label{c02:sec:conformal-framework}

\begin{definition}[Conformal reference structure]
\label{c02:def:conformal}
A \emph{conformal reference structure} is a flat reference $(\eta)$ together with the
conformal factor $\rho^{2/d}$: every physical metric in the theory is in the conformal class
of $\eta$, and the conformal factor encodes the scalar face of Difference.
\end{definition}

\begin{theorem}[The conformal framework]
\label{c02:thm:conformal}
The physical metric of the theory is the conformal rescaling of the tensor reference:
$g = \rho^{2/d}\,\eta$. The conformal factor carries the scalar content; the reference
carries the tensor structure.
\end{theorem}

\begin{proof}
The metric ansatz \eqref{c02:eq:metric-ansatz} is the accepted canonical form \cite{c02:QG290}.
Its factorization separates the two primitive inputs: $\rho$, the scalar face of Difference,
enters only through the conformal factor $\rho^{2/d}$; $\eta$, the tensor reference, enters
as the flat background. The conformal factor is fixed by dimensional analysis ($d$ appears
as the exponent $2/d$) and by the scalar counting content, while the reference is the
structure against which the conformal class is defined. The framework inventory classifies
this factorization as the irreducible reading structure of the theory \cite{c02:QG290}.
\end{proof}

\begin{remark}[The scalar face needs no reference; the tensor face does]
Consistent with Chapter~\ref{c01:chap:difference}, the scalar face $\rho$ is defined without
reference to $\eta$, while the tensor face $\psi$ --- the Weyl content --- is defined
\emph{against} $\eta$. This asymmetry is precisely why Theorem~\ref{c01:thm:minimal-primitives}
of Chapter~\ref{c01:chap:difference} holds: $\eta$ is necessary exactly for the tensor
sub-sector.
\end{remark}

\section{\texorpdfstring{$\eta$}{eta} and the Tensor Sector}
\label{c02:sec:eta-tensor}

\begin{definition}[Weyl content]
\label{c02:def:weyl}
The \emph{Weyl content} of the connectivity is the non-conformal part of the metric --- the
difference from conformal flatness. It is the tensor face of Difference, $\psi$
\cite{c02:QG285}.
\end{definition}

\begin{theorem}[The tensor face is read against \texorpdfstring{$\eta$}{eta}]
\label{c02:thm:tensor-read}
The tensor field $\psi$ --- the Weyl content, spin-2, with two transverse-traceless
polarizations --- is defined as the difference from conformal flatness relative to the
reference $\eta$.
\end{theorem}

\begin{proof}
The psi-reinterpretation audit \cite{c02:QG285} establishes the link "$\psi$ as Difference": the
Weyl tensor IS the difference from conformal flatness --- the non-conformal content of the
metric. Conformal flatness is a statement about the conformal class of a reference; the
reference is $\eta$ (Definition~\ref{c02:def:eta}). Hence $\psi$ is read against $\eta$. The
spin-2 nature and the two transverse-traceless polarizations follow from the six components
of the rank-2 adjacency structure at $d=3$: one trace component ($\rho$) and five traceless,
of which two are transverse-traceless ($\psi$) \cite{c02:QG286}. The tensor sector therefore
presupposes the reference $\eta$ at every step.
\end{proof}

\begin{corollary}[Tensor observables require \texorpdfstring{$\eta$}{eta}]
\label{c02:cor:tensor-requires-eta}
Every observable in the tensor (gravitational/Weyl) sector is defined against $\eta$; the
scalar sector is not.
\end{corollary}

\begin{proof}
This is the $\eta$-removal result of the foundation stress test \cite{c02:QG292}: removing
$\eta$ leaves all five layers intact as a scalar ($\rho$-face) read --- masses, couplings,
mixings, cosmological fractions, and spectral observables are all $\rho$-face reads --- and
destroys only the tensor sub-sector, because $\psi$ is defined against $\eta$
(Corollary~\ref{c02:cor:tensor-requires-eta} of this chapter; Theorem~\ref{c01:thm:minimal-primitives}
of Chapter~\ref{c01:chap:difference}).
\end{proof}

\section{Difference, \texorpdfstring{$\eta$}{eta} and the Minimal Foundation}
\label{c02:sec:minimal-foundation}

\begin{theorem}[The minimal foundation is \texorpdfstring{$\{\text{Difference},\eta\}$}{{Difference, eta}}]
\label{c02:thm:minimal-foundation}
The minimal foundation of the theory is the pair $\{\text{Difference},\eta\}$: exactly two
primitives, with the derived hierarchy
\begin{equation}
\label{c02:eq:core}
\text{Difference}\;\longrightarrow\;\text{Actualization}\;\longrightarrow\;\text{Inevitable Spectrum}\;\longrightarrow\;\text{Physics}.
\end{equation}
\end{theorem}

\begin{proof}
By the foundation stress test \cite{c02:QG292}, removing Difference collapses all five layers
(the count-conservation identity is the primitive's definition, Chapter~\ref{c01:chap:difference},
Theorem~\ref{c01:thm:fundamental-boundary}); removing $\eta$ collapses only the tensor sub-sector
(Theorem~\ref{c02:thm:tensor-read}, Corollary~\ref{c02:cor:tensor-requires-eta}). Both are therefore
necessary. Neither is derivable from the other (Lemma~\ref{c02:lem:eta-not-derived} and
Corollary~\ref{c02:cor:eta-independent}). Hence the pair is irreducible and independent, and the
core hierarchy \eqref{c02:eq:core} --- with Actualization derived from Difference per MONO006
\cite{c02:MONO006} --- is the canonical minimal foundation.
\end{proof}

\begin{lemma}[Actualization is independent of \texorpdfstring{$\eta$}{eta}]
\label{c02:lem:actualization-eta-free}
Actualization, the count-producing process of Difference, does not depend on $\eta$: counting
needs no metric reference.
\end{lemma}

\begin{proof}
The foundation stress test \cite{c02:QG292} establishes that removing $\eta$ leaves Actualization
intact: counting is a graph identity (the handshake lemma), not a metric identity. Combined
with MONO006 \cite{c02:MONO006}, which establishes that removing Difference collapses Actualization
(count conservation is the definitional identity of the primitive), this shows Actualization
is derived from Difference alone, with no $\eta$ dependence.
\end{proof}

\begin{corollary}[The canonical architecture]
\label{c02:cor:canonical-architecture}
The canonical architecture is: primitives $\{\text{Difference},\eta\}$; Actualization derived
from Difference; the spectrum inevitable; physics emergent. This is the referee-ready
structure adopted by MONO007 \cite{c02:MONO007}.
\end{corollary}

\section{Distinction Between \texorpdfstring{$\eta$}{eta} and \texorpdfstring{$\pi$}{pi}}
\label{c02:sec:eta-vs-pi}

\begin{definition}[Boundary constant]
\label{c02:def:pi}
\emph{$\pi$} is the universal mathematical constant: the geometric ratio of the circle. In
the canonical architecture it is a BOUNDARY --- a numerical constant not derivable inside the
framework --- and not a primitive \cite{c02:QG291}.
\end{definition}

\begin{theorem}[$\eta$ is a primitive; $\pi$ is a boundary]
\label{c02:thm:eta-primitive-pi-boundary}
$\eta$ and $\pi$ play different structural roles: $\eta$ is a primitive (necessary, an
irreducible reading structure); $\pi$ is a boundary constant (redundant as a theory input).
\end{theorem}

\begin{proof}
The framework necessity audit \cite{c02:QG291} classifies the two items separately. $\eta$ is
\emph{necessary}: without it the trace/traceless duality and the conformal reading are
undefined (Theorem~\ref{c02:thm:eta-necessary}). $\pi$ is \emph{redundant} as a theory input: no
derived prediction uses $\pi$ as an input --- every derived observable (masses, couplings,
mixings, $\Omega_\Lambda/\Omega_m$, $n_s$, acoustic ratios, the pre-registered predictions)
is a pure ratio of D96 spectral constants together with the calibration scale. Where $\pi$
appears --- most notably in the Bekenstein $S=A/4$ coefficient --- it is an imported quantum
factor that the framework cannot produce internally, an explicit boundary \cite{c02:QG185}.
Hence $\eta$ is a primitive and $\pi$ is a boundary; the referee-ready MONO007 architecture
presents $\eta$ as the second primitive and $\pi$ as a boundary constant, never conflating
the two \cite{c02:MONO007}.
\end{proof}

\begin{corollary}[No conflation]
\label{c02:cor:no-conflation}
The monograph presents $\eta$ (foundational) and $\pi$ (boundary) in separate structural
roles; the tensor reference is not a numerical constant, and the geometric constant is not a
reference structure.
\end{corollary}

\begin{remark}[Referee-safe wording]
Per MONO007 \cite{c02:MONO007}, the distinction between the primitive $\eta$ and the boundary
constant $\pi$ is stated explicitly wherever both appear, so that no reader or referee can
misread the numerical constant as a primitive input of the theory.
\end{remark}

\section{Consequences}
\label{c02:sec:eta-consequences}

\begin{lemma}[One primitive pair, one hierarchy]
\label{c02:lem:eta-consequence-1}
The entire content of the theory is carried by the primitive pair
$\{\text{Difference},\eta\}$ through the single canonical hierarchy
\eqref{c02:eq:core}: the scalar face of Difference feeds the spectral readout of all observables;
the tensor face, read against $\eta$, feeds the gravitational (Weyl) sector; Actualization
connects the primitive to the inevitable spectrum.
\end{lemma}

\begin{theorem}[Reference-safety of the derivation graph]
\label{c02:thm:reference-safety}
The derivation graph of the theory, with primitives $\{\text{Difference},\eta\}$ and
Actualization derived, is acyclic and contains no circular dependency involving $\eta$.
\end{theorem}

\begin{proof}
The dependency graph of the canonical architecture is verified acyclic (the final-theory
architecture audit). $\eta$ appears only as a reading reference for the tensor sector; no
derived quantity is used to define $\eta$ (Lemma~\ref{c02:lem:eta-not-derived}), and no tensor
observable is used in the scalar sector that defines the count. The referee-ready structure
of MONO007 \cite{c02:MONO007} confirms that the presentation introduces no circular
dependencies: $\eta$ is an input, never an output of a derived relation.
\end{proof}

\begin{corollary}[The architecture is complete and irreducible]
\label{c02:cor:eta-architecture}
The canonical architecture --- primitives $\{\text{Difference},\eta\}$, Actualization
derived, spectrum inevitable, physics emergent --- is complete within the minimal hierarchy
and irreducible: no primitive can be removed, and no derived layer is an additional input.
\end{corollary}

\begin{remark}[The Actualization Theory]
In the canonical monograph the theory is presented under the name \emph{The Actualization
Theory}, with subtitle \emph{A Reconstruction of Physics from Difference, Actualization and
Spectrum}. The name emphasizes that the reconstructed physics flows through the derived
actualization process --- the count-producing dynamics of the primitive Difference, converging
to the inevitable spectrum --- while the primitive foundation remains the pair
$\{\text{Difference},\eta\}$.
\end{remark}

\section{Summary}
\label{c02:sec:eta-summary}

This chapter has established the second member of the canonical architecture. We have proved:

\begin{enumerate}
\item \textbf{Necessity} (Theorem~\ref{c02:thm:eta-necessary}): the tensor reference $\eta$ is
necessary --- without it no trace, no traceless part, no conformal flatness, no tensor face;
\item \textbf{Independence} (Lemma~\ref{c02:lem:eta-not-derived},
Corollary~\ref{c02:cor:eta-independent}): $\eta$ is not derivable from Difference;
\item \textbf{The conformal framework} (Theorem~\ref{c02:thm:conformal}): the physical metric is
the conformal rescaling $g=\rho^{2/d}\,\eta$ of the tensor reference;
\item \textbf{The tensor sector} (Theorem~\ref{c02:thm:tensor-read},
Corollary~\ref{c02:cor:tensor-requires-eta}): $\psi$ is read against $\eta$, and tensor
observables require it while scalar observables do not;
\item \textbf{The minimal foundation} (Theorem~\ref{c02:thm:minimal-foundation},
Lemma~\ref{c02:lem:actualization-eta-free}, Corollary~\ref{c02:cor:canonical-architecture}): the
minimal foundation is $\{\text{Difference},\eta\}$, with Actualization derived from
Difference alone;
\item \textbf{The $\eta$/$\pi$ distinction} (Theorem~\ref{c02:thm:eta-primitive-pi-boundary},
Corollary~\ref{c02:cor:no-conflation}): $\eta$ is a primitive, $\pi$ is a boundary constant,
never conflated;
\item \textbf{Consequences} (Theorem~\ref{c02:thm:reference-safety},
Corollary~\ref{c02:cor:eta-architecture}): the derivation graph is reference-safe and acyclic,
and the architecture is complete and irreducible.
\end{enumerate}

The tensor reference $\eta$ is therefore not an auxiliary mathematical device but a precise
primitive: the conformal reference against which the trace/traceless duality of Difference is
read, the structure that supports the tensor (gravitational) face, and --- together with
Difference --- one half of the irreducible minimal foundation of The Actualization Theory.
The following chapters derive, from this foundation, the actualization process and the
inevitable spectrum from which all physics emerges.

\begin{thebibliography}{99}
\bibitem{c02:QG285} AT-QG Phase 285, \emph{Psi Reinterpretation Audit}, Docs/Research.
\bibitem{c02:QG286} AT-QG Phase 286, \emph{Difference Duality Audit}, Docs/Research.
\bibitem{c02:QG289} AT-QG Phase 289, \emph{Anchor Inventory Audit}, Docs/Research.
\bibitem{c02:QG290} AT-QG Phase 290, \emph{Framework Inventory Audit}, Docs/Research.
\bibitem{c02:QG291} AT-QG Phase 291, \emph{Framework Necessity Audit}, Docs/Research.
\bibitem{c02:QG292} AT-QG Phase 292, \emph{Foundation Stress Test}, Docs/Research.
\bibitem{c02:QG185} AT-QG Phase 185, \emph{Bekenstein Quarter Coefficient Origin}, Docs/Research.
\bibitem{c02:MONO006} AT-MONO006, \emph{A01 Actualization Origin Audit}, Docs/Zenodo/Monograph.
\bibitem{c02:MONO007} AT-MONO007, \emph{Referee-Readiness Resolution}, Docs/Zenodo/Monograph.
\end{thebibliography}
 after the preamble that defines
%         the theorem environments below (or include via the master file).
% ======================================================================

% ---- Theorem environments (define in the master preamble if not present) ----
% \newtheorem{definition}{Definition}[chapter]
% \newtheorem{lemma}[definition]{Lemma}
% \newtheorem{theorem}[definition]{Theorem}
% \newtheorem{corollary}[definition]{Corollary}
% \newtheorem{remark}[definition]{Remark}

\chapter{The Tensor Reference \texorpdfstring{$\eta$}{eta}}
\label{c02:chap:eta}

\section{Introduction}
\label{c02:sec:eta-introduction}

In Chapter~\ref{c01:chap:difference} we established that Difference is the fundamental boundary
of the theory, and that the scalar counting measure $\rho$ and the tensor field $\psi$ are
the trace and traceless faces of the one Difference \cite{c02:QG286}. That duality, however, is
not defined in a vacuum: the trace--traceless decomposition, the conformal flatness against
which the Weyl content is read, and the very notion of the tensor face $\psi$ all presuppose
a \emph{reference structure}. This chapter develops that structure --- the tensor reference
metric $\eta$.

The results of this chapter are the canonical end-state of the framework inventory and
necessity program \cite{c02:QG289,c02:QG290,c02:QG291,c02:QG292}: $\eta$ is the second primitive of the
theory, necessary and irreducible; the framework $\{\text{Difference},\eta\}$ is minimal;
and the numerical constant $\pi$, often presented alongside $\eta$ in earlier treatments, is
\emph{not} a primitive but a boundary constant --- redundant as a theory input and separated
from the primitive reference by the referee-ready architecture of MONO007
\cite{c02:MONO007}. Throughout we use only accepted canonical results, and introduce no new
physics and no speculation.

\section{Necessity of \texorpdfstring{$\eta$}{eta}}
\label{c02:sec:eta-necessity}

\begin{definition}[Tensor reference metric]
\label{c02:def:eta}
\emph{$\eta$} is the tensor reference metric: the conformal reference with respect to which
the physical metric takes the form
\begin{equation}
\label{c02:eq:metric-ansatz}
g \;=\; \rho^{2/d}\,\eta,
\end{equation}
where $\rho$ is the scalar counting measure and $d$ the spatial dimension
\cite{c02:QG290,c02:QG285}.
\end{definition}

\begin{theorem}[Necessity of the tensor reference]
\label{c02:thm:eta-necessary}
The tensor reference $\eta$ is NECESSARY: without it there is no trace, no traceless part,
no conformal flatness, and no reading of the difference structure.
\end{theorem}

\begin{proof}
The duality Difference $\to \{\rho,\psi\}$ is the trace/traceless decomposition
\cite{c02:QG286}
\begin{equation}
\label{c02:eq:decomposition}
A_{ij} \;=\; \frac{1}{d}\,\mathrm{Tr}(A)\,\delta_{ij}
\;+\;\Bigl(A_{ij}-\frac{1}{d}\,\mathrm{Tr}(A)\,\delta_{ij}\Bigr),
\end{equation}
of the rank-2 difference object $A_{ij}$. The trace ($\rho$), the traceless part ($\psi$),
and the Weyl content --- the difference from conformal flatness \cite{c02:QG285} --- are all
\emph{defined against} the reference $\eta$: conformal flatness is meaningful only with a
flat reference, and the Weyl tensor is the non-conformal content of the metric relative to
that reference. Without $\eta$ there is no conformal flatness, hence no notion of
"difference from conformal flatness," hence no tensor face. The framework inventory audit
\cite{c02:QG291} classifies $\eta$ accordingly: not \emph{derived} (no count produces a metric,
and $\eta$ carries no scale), not \emph{redundant} (removing $\eta$ removes the reading), but
\emph{necessary} --- the reference structure the framework reads against.
\end{proof}

\begin{lemma}[$\eta$ is not derivable from Difference]
\label{c02:lem:eta-not-derived}
No counting measure, no matter how refined, produces a metric reference; $\eta$ carries no
scale and is not a function of $\rho$.
\end{lemma}

\begin{proof}
The framework inventory establishes that $\eta$ is not derived: a metric reference is a
\emph{reading structure}, not an output of the count. $\eta$ carries no scale --- the scale
triad enters the theory elsewhere --- and the metric ansatz \eqref{c02:eq:metric-ansatz} takes
$\eta$ as an input on which the conformal factor $\rho^{2/d}$ acts. Hence $\eta$ is an
independent input, not a function of the primitive Difference.
\end{proof}

\begin{corollary}[Two independent primitives]
\label{c02:cor:eta-independent}
Difference and $\eta$ are independent primitives: neither is derivable from the other.
\end{corollary}

\section{The Conformal Reference Framework}
\label{c02:sec:conformal-framework}

\begin{definition}[Conformal reference structure]
\label{c02:def:conformal}
A \emph{conformal reference structure} is a flat reference $(\eta)$ together with the
conformal factor $\rho^{2/d}$: every physical metric in the theory is in the conformal class
of $\eta$, and the conformal factor encodes the scalar face of Difference.
\end{definition}

\begin{theorem}[The conformal framework]
\label{c02:thm:conformal}
The physical metric of the theory is the conformal rescaling of the tensor reference:
$g = \rho^{2/d}\,\eta$. The conformal factor carries the scalar content; the reference
carries the tensor structure.
\end{theorem}

\begin{proof}
The metric ansatz \eqref{c02:eq:metric-ansatz} is the accepted canonical form \cite{c02:QG290}.
Its factorization separates the two primitive inputs: $\rho$, the scalar face of Difference,
enters only through the conformal factor $\rho^{2/d}$; $\eta$, the tensor reference, enters
as the flat background. The conformal factor is fixed by dimensional analysis ($d$ appears
as the exponent $2/d$) and by the scalar counting content, while the reference is the
structure against which the conformal class is defined. The framework inventory classifies
this factorization as the irreducible reading structure of the theory \cite{c02:QG290}.
\end{proof}

\begin{remark}[The scalar face needs no reference; the tensor face does]
Consistent with Chapter~\ref{c01:chap:difference}, the scalar face $\rho$ is defined without
reference to $\eta$, while the tensor face $\psi$ --- the Weyl content --- is defined
\emph{against} $\eta$. This asymmetry is precisely why Theorem~\ref{c01:thm:minimal-primitives}
of Chapter~\ref{c01:chap:difference} holds: $\eta$ is necessary exactly for the tensor
sub-sector.
\end{remark}

\section{\texorpdfstring{$\eta$}{eta} and the Tensor Sector}
\label{c02:sec:eta-tensor}

\begin{definition}[Weyl content]
\label{c02:def:weyl}
The \emph{Weyl content} of the connectivity is the non-conformal part of the metric --- the
difference from conformal flatness. It is the tensor face of Difference, $\psi$
\cite{c02:QG285}.
\end{definition}

\begin{theorem}[The tensor face is read against \texorpdfstring{$\eta$}{eta}]
\label{c02:thm:tensor-read}
The tensor field $\psi$ --- the Weyl content, spin-2, with two transverse-traceless
polarizations --- is defined as the difference from conformal flatness relative to the
reference $\eta$.
\end{theorem}

\begin{proof}
The psi-reinterpretation audit \cite{c02:QG285} establishes the link "$\psi$ as Difference": the
Weyl tensor IS the difference from conformal flatness --- the non-conformal content of the
metric. Conformal flatness is a statement about the conformal class of a reference; the
reference is $\eta$ (Definition~\ref{c02:def:eta}). Hence $\psi$ is read against $\eta$. The
spin-2 nature and the two transverse-traceless polarizations follow from the six components
of the rank-2 adjacency structure at $d=3$: one trace component ($\rho$) and five traceless,
of which two are transverse-traceless ($\psi$) \cite{c02:QG286}. The tensor sector therefore
presupposes the reference $\eta$ at every step.
\end{proof}

\begin{corollary}[Tensor observables require \texorpdfstring{$\eta$}{eta}]
\label{c02:cor:tensor-requires-eta}
Every observable in the tensor (gravitational/Weyl) sector is defined against $\eta$; the
scalar sector is not.
\end{corollary}

\begin{proof}
This is the $\eta$-removal result of the foundation stress test \cite{c02:QG292}: removing
$\eta$ leaves all five layers intact as a scalar ($\rho$-face) read --- masses, couplings,
mixings, cosmological fractions, and spectral observables are all $\rho$-face reads --- and
destroys only the tensor sub-sector, because $\psi$ is defined against $\eta$
(Corollary~\ref{c02:cor:tensor-requires-eta} of this chapter; Theorem~\ref{c01:thm:minimal-primitives}
of Chapter~\ref{c01:chap:difference}).
\end{proof}

\section{Difference, \texorpdfstring{$\eta$}{eta} and the Minimal Foundation}
\label{c02:sec:minimal-foundation}

\begin{theorem}[The minimal foundation is \texorpdfstring{$\{\text{Difference},\eta\}$}{{Difference, eta}}]
\label{c02:thm:minimal-foundation}
The minimal foundation of the theory is the pair $\{\text{Difference},\eta\}$: exactly two
primitives, with the derived hierarchy
\begin{equation}
\label{c02:eq:core}
\text{Difference}\;\longrightarrow\;\text{Actualization}\;\longrightarrow\;\text{Inevitable Spectrum}\;\longrightarrow\;\text{Physics}.
\end{equation}
\end{theorem}

\begin{proof}
By the foundation stress test \cite{c02:QG292}, removing Difference collapses all five layers
(the count-conservation identity is the primitive's definition, Chapter~\ref{c01:chap:difference},
Theorem~\ref{c01:thm:fundamental-boundary}); removing $\eta$ collapses only the tensor sub-sector
(Theorem~\ref{c02:thm:tensor-read}, Corollary~\ref{c02:cor:tensor-requires-eta}). Both are therefore
necessary. Neither is derivable from the other (Lemma~\ref{c02:lem:eta-not-derived} and
Corollary~\ref{c02:cor:eta-independent}). Hence the pair is irreducible and independent, and the
core hierarchy \eqref{c02:eq:core} --- with Actualization derived from Difference per MONO006
\cite{c02:MONO006} --- is the canonical minimal foundation.
\end{proof}

\begin{lemma}[Actualization is independent of \texorpdfstring{$\eta$}{eta}]
\label{c02:lem:actualization-eta-free}
Actualization, the count-producing process of Difference, does not depend on $\eta$: counting
needs no metric reference.
\end{lemma}

\begin{proof}
The foundation stress test \cite{c02:QG292} establishes that removing $\eta$ leaves Actualization
intact: counting is a graph identity (the handshake lemma), not a metric identity. Combined
with MONO006 \cite{c02:MONO006}, which establishes that removing Difference collapses Actualization
(count conservation is the definitional identity of the primitive), this shows Actualization
is derived from Difference alone, with no $\eta$ dependence.
\end{proof}

\begin{corollary}[The canonical architecture]
\label{c02:cor:canonical-architecture}
The canonical architecture is: primitives $\{\text{Difference},\eta\}$; Actualization derived
from Difference; the spectrum inevitable; physics emergent. This is the referee-ready
structure adopted by MONO007 \cite{c02:MONO007}.
\end{corollary}

\section{Distinction Between \texorpdfstring{$\eta$}{eta} and \texorpdfstring{$\pi$}{pi}}
\label{c02:sec:eta-vs-pi}

\begin{definition}[Boundary constant]
\label{c02:def:pi}
\emph{$\pi$} is the universal mathematical constant: the geometric ratio of the circle. In
the canonical architecture it is a BOUNDARY --- a numerical constant not derivable inside the
framework --- and not a primitive \cite{c02:QG291}.
\end{definition}

\begin{theorem}[$\eta$ is a primitive; $\pi$ is a boundary]
\label{c02:thm:eta-primitive-pi-boundary}
$\eta$ and $\pi$ play different structural roles: $\eta$ is a primitive (necessary, an
irreducible reading structure); $\pi$ is a boundary constant (redundant as a theory input).
\end{theorem}

\begin{proof}
The framework necessity audit \cite{c02:QG291} classifies the two items separately. $\eta$ is
\emph{necessary}: without it the trace/traceless duality and the conformal reading are
undefined (Theorem~\ref{c02:thm:eta-necessary}). $\pi$ is \emph{redundant} as a theory input: no
derived prediction uses $\pi$ as an input --- every derived observable (masses, couplings,
mixings, $\Omega_\Lambda/\Omega_m$, $n_s$, acoustic ratios, the pre-registered predictions)
is a pure ratio of D96 spectral constants together with the calibration scale. Where $\pi$
appears --- most notably in the Bekenstein $S=A/4$ coefficient --- it is an imported quantum
factor that the framework cannot produce internally, an explicit boundary \cite{c02:QG185}.
Hence $\eta$ is a primitive and $\pi$ is a boundary; the referee-ready MONO007 architecture
presents $\eta$ as the second primitive and $\pi$ as a boundary constant, never conflating
the two \cite{c02:MONO007}.
\end{proof}

\begin{corollary}[No conflation]
\label{c02:cor:no-conflation}
The monograph presents $\eta$ (foundational) and $\pi$ (boundary) in separate structural
roles; the tensor reference is not a numerical constant, and the geometric constant is not a
reference structure.
\end{corollary}

\begin{remark}[Referee-safe wording]
Per MONO007 \cite{c02:MONO007}, the distinction between the primitive $\eta$ and the boundary
constant $\pi$ is stated explicitly wherever both appear, so that no reader or referee can
misread the numerical constant as a primitive input of the theory.
\end{remark}

\section{Consequences}
\label{c02:sec:eta-consequences}

\begin{lemma}[One primitive pair, one hierarchy]
\label{c02:lem:eta-consequence-1}
The entire content of the theory is carried by the primitive pair
$\{\text{Difference},\eta\}$ through the single canonical hierarchy
\eqref{c02:eq:core}: the scalar face of Difference feeds the spectral readout of all observables;
the tensor face, read against $\eta$, feeds the gravitational (Weyl) sector; Actualization
connects the primitive to the inevitable spectrum.
\end{lemma}

\begin{theorem}[Reference-safety of the derivation graph]
\label{c02:thm:reference-safety}
The derivation graph of the theory, with primitives $\{\text{Difference},\eta\}$ and
Actualization derived, is acyclic and contains no circular dependency involving $\eta$.
\end{theorem}

\begin{proof}
The dependency graph of the canonical architecture is verified acyclic (the final-theory
architecture audit). $\eta$ appears only as a reading reference for the tensor sector; no
derived quantity is used to define $\eta$ (Lemma~\ref{c02:lem:eta-not-derived}), and no tensor
observable is used in the scalar sector that defines the count. The referee-ready structure
of MONO007 \cite{c02:MONO007} confirms that the presentation introduces no circular
dependencies: $\eta$ is an input, never an output of a derived relation.
\end{proof}

\begin{corollary}[The architecture is complete and irreducible]
\label{c02:cor:eta-architecture}
The canonical architecture --- primitives $\{\text{Difference},\eta\}$, Actualization
derived, spectrum inevitable, physics emergent --- is complete within the minimal hierarchy
and irreducible: no primitive can be removed, and no derived layer is an additional input.
\end{corollary}

\begin{remark}[The Actualization Theory]
In the canonical monograph the theory is presented under the name \emph{The Actualization
Theory}, with subtitle \emph{A Reconstruction of Physics from Difference, Actualization and
Spectrum}. The name emphasizes that the reconstructed physics flows through the derived
actualization process --- the count-producing dynamics of the primitive Difference, converging
to the inevitable spectrum --- while the primitive foundation remains the pair
$\{\text{Difference},\eta\}$.
\end{remark}

\section{Summary}
\label{c02:sec:eta-summary}

This chapter has established the second member of the canonical architecture. We have proved:

\begin{enumerate}
\item \textbf{Necessity} (Theorem~\ref{c02:thm:eta-necessary}): the tensor reference $\eta$ is
necessary --- without it no trace, no traceless part, no conformal flatness, no tensor face;
\item \textbf{Independence} (Lemma~\ref{c02:lem:eta-not-derived},
Corollary~\ref{c02:cor:eta-independent}): $\eta$ is not derivable from Difference;
\item \textbf{The conformal framework} (Theorem~\ref{c02:thm:conformal}): the physical metric is
the conformal rescaling $g=\rho^{2/d}\,\eta$ of the tensor reference;
\item \textbf{The tensor sector} (Theorem~\ref{c02:thm:tensor-read},
Corollary~\ref{c02:cor:tensor-requires-eta}): $\psi$ is read against $\eta$, and tensor
observables require it while scalar observables do not;
\item \textbf{The minimal foundation} (Theorem~\ref{c02:thm:minimal-foundation},
Lemma~\ref{c02:lem:actualization-eta-free}, Corollary~\ref{c02:cor:canonical-architecture}): the
minimal foundation is $\{\text{Difference},\eta\}$, with Actualization derived from
Difference alone;
\item \textbf{The $\eta$/$\pi$ distinction} (Theorem~\ref{c02:thm:eta-primitive-pi-boundary},
Corollary~\ref{c02:cor:no-conflation}): $\eta$ is a primitive, $\pi$ is a boundary constant,
never conflated;
\item \textbf{Consequences} (Theorem~\ref{c02:thm:reference-safety},
Corollary~\ref{c02:cor:eta-architecture}): the derivation graph is reference-safe and acyclic,
and the architecture is complete and irreducible.
\end{enumerate}

The tensor reference $\eta$ is therefore not an auxiliary mathematical device but a precise
primitive: the conformal reference against which the trace/traceless duality of Difference is
read, the structure that supports the tensor (gravitational) face, and --- together with
Difference --- one half of the irreducible minimal foundation of The Actualization Theory.
The following chapters derive, from this foundation, the actualization process and the
inevitable spectrum from which all physics emerges.

\begin{thebibliography}{99}
\bibitem{c02:QG285} AT-QG Phase 285, \emph{Psi Reinterpretation Audit}, Docs/Research.
\bibitem{c02:QG286} AT-QG Phase 286, \emph{Difference Duality Audit}, Docs/Research.
\bibitem{c02:QG289} AT-QG Phase 289, \emph{Anchor Inventory Audit}, Docs/Research.
\bibitem{c02:QG290} AT-QG Phase 290, \emph{Framework Inventory Audit}, Docs/Research.
\bibitem{c02:QG291} AT-QG Phase 291, \emph{Framework Necessity Audit}, Docs/Research.
\bibitem{c02:QG292} AT-QG Phase 292, \emph{Foundation Stress Test}, Docs/Research.
\bibitem{c02:QG185} AT-QG Phase 185, \emph{Bekenstein Quarter Coefficient Origin}, Docs/Research.
\bibitem{c02:MONO006} AT-MONO006, \emph{A01 Actualization Origin Audit}, Docs/Zenodo/Monograph.
\bibitem{c02:MONO007} AT-MONO007, \emph{Referee-Readiness Resolution}, Docs/Zenodo/Monograph.
\end{thebibliography}
 after the preamble that defines
%         the theorem environments below (or include via the master file).
% ======================================================================

% ---- Theorem environments (define in the master preamble if not present) ----
% \newtheorem{definition}{Definition}[chapter]
% \newtheorem{lemma}[definition]{Lemma}
% \newtheorem{theorem}[definition]{Theorem}
% \newtheorem{corollary}[definition]{Corollary}
% \newtheorem{remark}[definition]{Remark}

\chapter{The Tensor Reference \texorpdfstring{$\eta$}{eta}}
\label{c02:chap:eta}

\section{Introduction}
\label{c02:sec:eta-introduction}

In Chapter~\ref{c01:chap:difference} we established that Difference is the fundamental boundary
of the theory, and that the scalar counting measure $\rho$ and the tensor field $\psi$ are
the trace and traceless faces of the one Difference \cite{c02:QG286}. That duality, however, is
not defined in a vacuum: the trace--traceless decomposition, the conformal flatness against
which the Weyl content is read, and the very notion of the tensor face $\psi$ all presuppose
a \emph{reference structure}. This chapter develops that structure --- the tensor reference
metric $\eta$.

The results of this chapter are the canonical end-state of the framework inventory and
necessity program \cite{c02:QG289,c02:QG290,c02:QG291,c02:QG292}: $\eta$ is the second primitive of the
theory, necessary and irreducible; the framework $\{\text{Difference},\eta\}$ is minimal;
and the numerical constant $\pi$, often presented alongside $\eta$ in earlier treatments, is
\emph{not} a primitive but a boundary constant --- redundant as a theory input and separated
from the primitive reference by the referee-ready architecture of MONO007
\cite{c02:MONO007}. Throughout we use only accepted canonical results, and introduce no new
physics and no speculation.

\section{Necessity of \texorpdfstring{$\eta$}{eta}}
\label{c02:sec:eta-necessity}

\begin{definition}[Tensor reference metric]
\label{c02:def:eta}
\emph{$\eta$} is the tensor reference metric: the conformal reference with respect to which
the physical metric takes the form
\begin{equation}
\label{c02:eq:metric-ansatz}
g \;=\; \rho^{2/d}\,\eta,
\end{equation}
where $\rho$ is the scalar counting measure and $d$ the spatial dimension
\cite{c02:QG290,c02:QG285}.
\end{definition}

\begin{theorem}[Necessity of the tensor reference]
\label{c02:thm:eta-necessary}
The tensor reference $\eta$ is NECESSARY: without it there is no trace, no traceless part,
no conformal flatness, and no reading of the difference structure.
\end{theorem}

\begin{proof}
The duality Difference $\to \{\rho,\psi\}$ is the trace/traceless decomposition
\cite{c02:QG286}
\begin{equation}
\label{c02:eq:decomposition}
A_{ij} \;=\; \frac{1}{d}\,\mathrm{Tr}(A)\,\delta_{ij}
\;+\;\Bigl(A_{ij}-\frac{1}{d}\,\mathrm{Tr}(A)\,\delta_{ij}\Bigr),
\end{equation}
of the rank-2 difference object $A_{ij}$. The trace ($\rho$), the traceless part ($\psi$),
and the Weyl content --- the difference from conformal flatness \cite{c02:QG285} --- are all
\emph{defined against} the reference $\eta$: conformal flatness is meaningful only with a
flat reference, and the Weyl tensor is the non-conformal content of the metric relative to
that reference. Without $\eta$ there is no conformal flatness, hence no notion of
"difference from conformal flatness," hence no tensor face. The framework inventory audit
\cite{c02:QG291} classifies $\eta$ accordingly: not \emph{derived} (no count produces a metric,
and $\eta$ carries no scale), not \emph{redundant} (removing $\eta$ removes the reading), but
\emph{necessary} --- the reference structure the framework reads against.
\end{proof}

\begin{lemma}[$\eta$ is not derivable from Difference]
\label{c02:lem:eta-not-derived}
No counting measure, no matter how refined, produces a metric reference; $\eta$ carries no
scale and is not a function of $\rho$.
\end{lemma}

\begin{proof}
The framework inventory establishes that $\eta$ is not derived: a metric reference is a
\emph{reading structure}, not an output of the count. $\eta$ carries no scale --- the scale
triad enters the theory elsewhere --- and the metric ansatz \eqref{c02:eq:metric-ansatz} takes
$\eta$ as an input on which the conformal factor $\rho^{2/d}$ acts. Hence $\eta$ is an
independent input, not a function of the primitive Difference.
\end{proof}

\begin{corollary}[Two independent primitives]
\label{c02:cor:eta-independent}
Difference and $\eta$ are independent primitives: neither is derivable from the other.
\end{corollary}

\section{The Conformal Reference Framework}
\label{c02:sec:conformal-framework}

\begin{definition}[Conformal reference structure]
\label{c02:def:conformal}
A \emph{conformal reference structure} is a flat reference $(\eta)$ together with the
conformal factor $\rho^{2/d}$: every physical metric in the theory is in the conformal class
of $\eta$, and the conformal factor encodes the scalar face of Difference.
\end{definition}

\begin{theorem}[The conformal framework]
\label{c02:thm:conformal}
The physical metric of the theory is the conformal rescaling of the tensor reference:
$g = \rho^{2/d}\,\eta$. The conformal factor carries the scalar content; the reference
carries the tensor structure.
\end{theorem}

\begin{proof}
The metric ansatz \eqref{c02:eq:metric-ansatz} is the accepted canonical form \cite{c02:QG290}.
Its factorization separates the two primitive inputs: $\rho$, the scalar face of Difference,
enters only through the conformal factor $\rho^{2/d}$; $\eta$, the tensor reference, enters
as the flat background. The conformal factor is fixed by dimensional analysis ($d$ appears
as the exponent $2/d$) and by the scalar counting content, while the reference is the
structure against which the conformal class is defined. The framework inventory classifies
this factorization as the irreducible reading structure of the theory \cite{c02:QG290}.
\end{proof}

\begin{remark}[The scalar face needs no reference; the tensor face does]
Consistent with Chapter~\ref{c01:chap:difference}, the scalar face $\rho$ is defined without
reference to $\eta$, while the tensor face $\psi$ --- the Weyl content --- is defined
\emph{against} $\eta$. This asymmetry is precisely why Theorem~\ref{c01:thm:minimal-primitives}
of Chapter~\ref{c01:chap:difference} holds: $\eta$ is necessary exactly for the tensor
sub-sector.
\end{remark}

\section{\texorpdfstring{$\eta$}{eta} and the Tensor Sector}
\label{c02:sec:eta-tensor}

\begin{definition}[Weyl content]
\label{c02:def:weyl}
The \emph{Weyl content} of the connectivity is the non-conformal part of the metric --- the
difference from conformal flatness. It is the tensor face of Difference, $\psi$
\cite{c02:QG285}.
\end{definition}

\begin{theorem}[The tensor face is read against \texorpdfstring{$\eta$}{eta}]
\label{c02:thm:tensor-read}
The tensor field $\psi$ --- the Weyl content, spin-2, with two transverse-traceless
polarizations --- is defined as the difference from conformal flatness relative to the
reference $\eta$.
\end{theorem}

\begin{proof}
The psi-reinterpretation audit \cite{c02:QG285} establishes the link "$\psi$ as Difference": the
Weyl tensor IS the difference from conformal flatness --- the non-conformal content of the
metric. Conformal flatness is a statement about the conformal class of a reference; the
reference is $\eta$ (Definition~\ref{c02:def:eta}). Hence $\psi$ is read against $\eta$. The
spin-2 nature and the two transverse-traceless polarizations follow from the six components
of the rank-2 adjacency structure at $d=3$: one trace component ($\rho$) and five traceless,
of which two are transverse-traceless ($\psi$) \cite{c02:QG286}. The tensor sector therefore
presupposes the reference $\eta$ at every step.
\end{proof}

\begin{corollary}[Tensor observables require \texorpdfstring{$\eta$}{eta}]
\label{c02:cor:tensor-requires-eta}
Every observable in the tensor (gravitational/Weyl) sector is defined against $\eta$; the
scalar sector is not.
\end{corollary}

\begin{proof}
This is the $\eta$-removal result of the foundation stress test \cite{c02:QG292}: removing
$\eta$ leaves all five layers intact as a scalar ($\rho$-face) read --- masses, couplings,
mixings, cosmological fractions, and spectral observables are all $\rho$-face reads --- and
destroys only the tensor sub-sector, because $\psi$ is defined against $\eta$
(Corollary~\ref{c02:cor:tensor-requires-eta} of this chapter; Theorem~\ref{c01:thm:minimal-primitives}
of Chapter~\ref{c01:chap:difference}).
\end{proof}

\section{Difference, \texorpdfstring{$\eta$}{eta} and the Minimal Foundation}
\label{c02:sec:minimal-foundation}

\begin{theorem}[The minimal foundation is \texorpdfstring{$\{\text{Difference},\eta\}$}{{Difference, eta}}]
\label{c02:thm:minimal-foundation}
The minimal foundation of the theory is the pair $\{\text{Difference},\eta\}$: exactly two
primitives, with the derived hierarchy
\begin{equation}
\label{c02:eq:core}
\text{Difference}\;\longrightarrow\;\text{Actualization}\;\longrightarrow\;\text{Inevitable Spectrum}\;\longrightarrow\;\text{Physics}.
\end{equation}
\end{theorem}

\begin{proof}
By the foundation stress test \cite{c02:QG292}, removing Difference collapses all five layers
(the count-conservation identity is the primitive's definition, Chapter~\ref{c01:chap:difference},
Theorem~\ref{c01:thm:fundamental-boundary}); removing $\eta$ collapses only the tensor sub-sector
(Theorem~\ref{c02:thm:tensor-read}, Corollary~\ref{c02:cor:tensor-requires-eta}). Both are therefore
necessary. Neither is derivable from the other (Lemma~\ref{c02:lem:eta-not-derived} and
Corollary~\ref{c02:cor:eta-independent}). Hence the pair is irreducible and independent, and the
core hierarchy \eqref{c02:eq:core} --- with Actualization derived from Difference per MONO006
\cite{c02:MONO006} --- is the canonical minimal foundation.
\end{proof}

\begin{lemma}[Actualization is independent of \texorpdfstring{$\eta$}{eta}]
\label{c02:lem:actualization-eta-free}
Actualization, the count-producing process of Difference, does not depend on $\eta$: counting
needs no metric reference.
\end{lemma}

\begin{proof}
The foundation stress test \cite{c02:QG292} establishes that removing $\eta$ leaves Actualization
intact: counting is a graph identity (the handshake lemma), not a metric identity. Combined
with MONO006 \cite{c02:MONO006}, which establishes that removing Difference collapses Actualization
(count conservation is the definitional identity of the primitive), this shows Actualization
is derived from Difference alone, with no $\eta$ dependence.
\end{proof}

\begin{corollary}[The canonical architecture]
\label{c02:cor:canonical-architecture}
The canonical architecture is: primitives $\{\text{Difference},\eta\}$; Actualization derived
from Difference; the spectrum inevitable; physics emergent. This is the referee-ready
structure adopted by MONO007 \cite{c02:MONO007}.
\end{corollary}

\section{Distinction Between \texorpdfstring{$\eta$}{eta} and \texorpdfstring{$\pi$}{pi}}
\label{c02:sec:eta-vs-pi}

\begin{definition}[Boundary constant]
\label{c02:def:pi}
\emph{$\pi$} is the universal mathematical constant: the geometric ratio of the circle. In
the canonical architecture it is a BOUNDARY --- a numerical constant not derivable inside the
framework --- and not a primitive \cite{c02:QG291}.
\end{definition}

\begin{theorem}[$\eta$ is a primitive; $\pi$ is a boundary]
\label{c02:thm:eta-primitive-pi-boundary}
$\eta$ and $\pi$ play different structural roles: $\eta$ is a primitive (necessary, an
irreducible reading structure); $\pi$ is a boundary constant (redundant as a theory input).
\end{theorem}

\begin{proof}
The framework necessity audit \cite{c02:QG291} classifies the two items separately. $\eta$ is
\emph{necessary}: without it the trace/traceless duality and the conformal reading are
undefined (Theorem~\ref{c02:thm:eta-necessary}). $\pi$ is \emph{redundant} as a theory input: no
derived prediction uses $\pi$ as an input --- every derived observable (masses, couplings,
mixings, $\Omega_\Lambda/\Omega_m$, $n_s$, acoustic ratios, the pre-registered predictions)
is a pure ratio of D96 spectral constants together with the calibration scale. Where $\pi$
appears --- most notably in the Bekenstein $S=A/4$ coefficient --- it is an imported quantum
factor that the framework cannot produce internally, an explicit boundary \cite{c02:QG185}.
Hence $\eta$ is a primitive and $\pi$ is a boundary; the referee-ready MONO007 architecture
presents $\eta$ as the second primitive and $\pi$ as a boundary constant, never conflating
the two \cite{c02:MONO007}.
\end{proof}

\begin{corollary}[No conflation]
\label{c02:cor:no-conflation}
The monograph presents $\eta$ (foundational) and $\pi$ (boundary) in separate structural
roles; the tensor reference is not a numerical constant, and the geometric constant is not a
reference structure.
\end{corollary}

\begin{remark}[Referee-safe wording]
Per MONO007 \cite{c02:MONO007}, the distinction between the primitive $\eta$ and the boundary
constant $\pi$ is stated explicitly wherever both appear, so that no reader or referee can
misread the numerical constant as a primitive input of the theory.
\end{remark}

\section{Consequences}
\label{c02:sec:eta-consequences}

\begin{lemma}[One primitive pair, one hierarchy]
\label{c02:lem:eta-consequence-1}
The entire content of the theory is carried by the primitive pair
$\{\text{Difference},\eta\}$ through the single canonical hierarchy
\eqref{c02:eq:core}: the scalar face of Difference feeds the spectral readout of all observables;
the tensor face, read against $\eta$, feeds the gravitational (Weyl) sector; Actualization
connects the primitive to the inevitable spectrum.
\end{lemma}

\begin{theorem}[Reference-safety of the derivation graph]
\label{c02:thm:reference-safety}
The derivation graph of the theory, with primitives $\{\text{Difference},\eta\}$ and
Actualization derived, is acyclic and contains no circular dependency involving $\eta$.
\end{theorem}

\begin{proof}
The dependency graph of the canonical architecture is verified acyclic (the final-theory
architecture audit). $\eta$ appears only as a reading reference for the tensor sector; no
derived quantity is used to define $\eta$ (Lemma~\ref{c02:lem:eta-not-derived}), and no tensor
observable is used in the scalar sector that defines the count. The referee-ready structure
of MONO007 \cite{c02:MONO007} confirms that the presentation introduces no circular
dependencies: $\eta$ is an input, never an output of a derived relation.
\end{proof}

\begin{corollary}[The architecture is complete and irreducible]
\label{c02:cor:eta-architecture}
The canonical architecture --- primitives $\{\text{Difference},\eta\}$, Actualization
derived, spectrum inevitable, physics emergent --- is complete within the minimal hierarchy
and irreducible: no primitive can be removed, and no derived layer is an additional input.
\end{corollary}

\begin{remark}[The Actualization Theory]
In the canonical monograph the theory is presented under the name \emph{The Actualization
Theory}, with subtitle \emph{A Reconstruction of Physics from Difference, Actualization and
Spectrum}. The name emphasizes that the reconstructed physics flows through the derived
actualization process --- the count-producing dynamics of the primitive Difference, converging
to the inevitable spectrum --- while the primitive foundation remains the pair
$\{\text{Difference},\eta\}$.
\end{remark}

\section{Summary}
\label{c02:sec:eta-summary}

This chapter has established the second member of the canonical architecture. We have proved:

\begin{enumerate}
\item \textbf{Necessity} (Theorem~\ref{c02:thm:eta-necessary}): the tensor reference $\eta$ is
necessary --- without it no trace, no traceless part, no conformal flatness, no tensor face;
\item \textbf{Independence} (Lemma~\ref{c02:lem:eta-not-derived},
Corollary~\ref{c02:cor:eta-independent}): $\eta$ is not derivable from Difference;
\item \textbf{The conformal framework} (Theorem~\ref{c02:thm:conformal}): the physical metric is
the conformal rescaling $g=\rho^{2/d}\,\eta$ of the tensor reference;
\item \textbf{The tensor sector} (Theorem~\ref{c02:thm:tensor-read},
Corollary~\ref{c02:cor:tensor-requires-eta}): $\psi$ is read against $\eta$, and tensor
observables require it while scalar observables do not;
\item \textbf{The minimal foundation} (Theorem~\ref{c02:thm:minimal-foundation},
Lemma~\ref{c02:lem:actualization-eta-free}, Corollary~\ref{c02:cor:canonical-architecture}): the
minimal foundation is $\{\text{Difference},\eta\}$, with Actualization derived from
Difference alone;
\item \textbf{The $\eta$/$\pi$ distinction} (Theorem~\ref{c02:thm:eta-primitive-pi-boundary},
Corollary~\ref{c02:cor:no-conflation}): $\eta$ is a primitive, $\pi$ is a boundary constant,
never conflated;
\item \textbf{Consequences} (Theorem~\ref{c02:thm:reference-safety},
Corollary~\ref{c02:cor:eta-architecture}): the derivation graph is reference-safe and acyclic,
and the architecture is complete and irreducible.
\end{enumerate}

The tensor reference $\eta$ is therefore not an auxiliary mathematical device but a precise
primitive: the conformal reference against which the trace/traceless duality of Difference is
read, the structure that supports the tensor (gravitational) face, and --- together with
Difference --- one half of the irreducible minimal foundation of The Actualization Theory.
The following chapters derive, from this foundation, the actualization process and the
inevitable spectrum from which all physics emerges.

\begin{thebibliography}{99}
\bibitem{c02:QG285} AT-QG Phase 285, \emph{Psi Reinterpretation Audit}, Docs/Research.
\bibitem{c02:QG286} AT-QG Phase 286, \emph{Difference Duality Audit}, Docs/Research.
\bibitem{c02:QG289} AT-QG Phase 289, \emph{Anchor Inventory Audit}, Docs/Research.
\bibitem{c02:QG290} AT-QG Phase 290, \emph{Framework Inventory Audit}, Docs/Research.
\bibitem{c02:QG291} AT-QG Phase 291, \emph{Framework Necessity Audit}, Docs/Research.
\bibitem{c02:QG292} AT-QG Phase 292, \emph{Foundation Stress Test}, Docs/Research.
\bibitem{c02:QG185} AT-QG Phase 185, \emph{Bekenstein Quarter Coefficient Origin}, Docs/Research.
\bibitem{c02:MONO006} AT-MONO006, \emph{A01 Actualization Origin Audit}, Docs/Zenodo/Monograph.
\bibitem{c02:MONO007} AT-MONO007, \emph{Referee-Readiness Resolution}, Docs/Zenodo/Monograph.
\end{thebibliography}


% Part II — Derived Dynamics
\part{Derived Dynamics}
% ======================================================================
%  AT-MONO102 — Chapter 3: Actualization
%  Canonical Monograph (v2.0) · The Actualization Theory:
%  A Reconstruction of Physics from Difference, Actualization and Spectrum
%
%  Status: COMPLETE — PASS-READY (MONO007 referee-ready architecture)
%  Inputs: Chapter01_Difference.tex, Chapter02_Eta.tex,
%          ATQG_CanonicalMonograph.md, ATQG_Mono007RefereeReadiness.md
%  Method: assembly only — canonical results only, no new physics, no speculation
%  Citations: QG268 (count conservation), QG272 (sector emergence),
%             QG282 (closure principle), QG292 (minimal foundation),
%             QG293 (hierarchy necessity), QG294 (minimal theory),
%             MONO006 (actualization derived), MONO007 (referee readiness)
%
%  Usage: % ======================================================================
%  AT-MONO102 — Chapter 3: Actualization
%  Canonical Monograph (v2.0) · The Actualization Theory:
%  A Reconstruction of Physics from Difference, Actualization and Spectrum
%
%  Status: COMPLETE — PASS-READY (MONO007 referee-ready architecture)
%  Inputs: Chapter01_Difference.tex, Chapter02_Eta.tex,
%          ATQG_CanonicalMonograph.md, ATQG_Mono007RefereeReadiness.md
%  Method: assembly only — canonical results only, no new physics, no speculation
%  Citations: QG268 (count conservation), QG272 (sector emergence),
%             QG282 (closure principle), QG292 (minimal foundation),
%             QG293 (hierarchy necessity), QG294 (minimal theory),
%             MONO006 (actualization derived), MONO007 (referee readiness)
%
%  Usage: % ======================================================================
%  AT-MONO102 — Chapter 3: Actualization
%  Canonical Monograph (v2.0) · The Actualization Theory:
%  A Reconstruction of Physics from Difference, Actualization and Spectrum
%
%  Status: COMPLETE — PASS-READY (MONO007 referee-ready architecture)
%  Inputs: Chapter01_Difference.tex, Chapter02_Eta.tex,
%          ATQG_CanonicalMonograph.md, ATQG_Mono007RefereeReadiness.md
%  Method: assembly only — canonical results only, no new physics, no speculation
%  Citations: QG268 (count conservation), QG272 (sector emergence),
%             QG282 (closure principle), QG292 (minimal foundation),
%             QG293 (hierarchy necessity), QG294 (minimal theory),
%             MONO006 (actualization derived), MONO007 (referee readiness)
%
%  Usage: \input{Chapter03_Actualization.tex} after the preamble that defines
%         the theorem environments below (or include via the master file).
% ======================================================================

% ---- Theorem environments (define in the master preamble if not present) ----
% \newtheorem{definition}{Definition}[chapter]
% \newtheorem{lemma}[definition]{Lemma}
% \newtheorem{theorem}[definition]{Theorem}
% \newtheorem{corollary}[definition]{Corollary}
% \newtheorem{remark}[definition]{Remark}

\chapter{Actualization}
\label{c03:chap:actualization}

\section{Introduction}
\label{c03:sec:actualization-introduction}

In Chapter~\ref{c01:chap:difference} we established that Difference is the fundamental boundary
of the theory --- the primitive whose count-conservation law is its definitional identity ---
and in Chapter~\ref{c02:chap:eta} that the tensor reference $\eta$ is the second primitive,
against which the tensor face of Difference is read. The canonical hierarchy
\eqref{c03:eq:core} then proceeds from the primitive pair to the inevitable spectrum through a
process:

\begin{equation}
\label{c03:eq:core}
\text{Difference}\;\longrightarrow\;\text{Actualization}\;\longrightarrow\;\text{Inevitable Spectrum}\;\longrightarrow\;\text{Physics}.
\end{equation}

This chapter develops the second member of that hierarchy: \emph{Actualization}. We
establish, with formal precision, that Actualization is the \emph{process face of
Difference} --- the count-producing dynamics whose definitional identity is count
conservation \cite{c03:QG268}; that its dynamics converge to the stable fixed point $N=96$, the
closure of the theory \cite{c03:QG282}; that the resonance structure of the spectrum is
actualization's readout \cite{c03:QG272}; and --- decisively, consistent with the MONO006
resolution \cite{c03:MONO006} --- that Actualization is a \emph{derived necessity}, not a third
primitive.

\section{Actualization as the Process Face of Difference}
\label{c03:sec:process-face}

\begin{theorem}[Actualization is the process face of Difference]
\label{c03:thm:process-face}
\emph{Actualization} --- the count-producing dynamics by which a Q-event, the unit of
Difference \cite{c03:QG268}, is applied, generating the network whose convergence yields the
spectrum --- is not an independent entity: it is the process face of the one primitive
Difference. Consequently the primitives of the theory are exactly
$\{\text{Difference},\eta\}$: Actualization is their derived process face, not a primitive.
\end{theorem}

\begin{proof}
Count conservation is the definitional identity of the primitive \cite{c03:QG268}: Difference is
the counting difference from a uniform background, and a Q-event is its unit. Actualization
is precisely the process that applies those units, and its identity is inherited from the
primitive: there is no Actualization without the unit, and no unit without Difference. The
minimal primitive set was established as the pair $\{\text{Difference},\eta\}$
(Theorem~\ref{c02:thm:minimal-foundation}, Chapter~\ref{c02:chap:eta}); the MONO006 resolution
\cite{c03:MONO006} confirms this classification, and the MONO007 architecture \cite{c03:MONO007}
presents Actualization as a derived layer. Hence there is no primitive-status ambiguity:
Actualization is the dynamical face of Difference, not a separate primitive.
\end{proof}

\section{Count Conservation}
\label{c03:sec:count-conservation}

\begin{theorem}[Count conservation is the definitional identity]
\label{c03:thm:count-conservation}
\emph{Count conservation} --- the total count of units of Difference is conserved by the
actualization process: Q-events are neither created nor destroyed, only redistributed ---
is not a separate postulate but the definitional identity of the primitive Difference and
therefore of its process face Actualization. One application of Actualization produces one
unit of count, so conservation follows automatically from the primitive, not as an imported
assumption of the theory.
\end{theorem}

\begin{proof}
The count-conservation origin \cite{c03:QG268} establishes that conservation is the defining
property of the primitive itself --- it is what \emph{is} the count --- and identifies the
Q-event as the indivisible step of the count. Since Actualization is the process that
applies the count (Theorem~\ref{c03:thm:process-face}), the conservation law is inherited:
the process conserves exactly what the primitive defines. No separate conservation postulate
is required.
\end{proof}

\section{The \texorpdfstring{$N=96$}{N=96} Attractor}
\label{c03:sec:n96-attractor}

The \emph{closure} of the theory is the stable fixed point of the actualization process: the
configuration at which the count-producing dynamics has converged, with zero residual
change. The boundary of the theory is this fixed point --- the topology saturates at zero
residual link growth, and every initial pattern converges to the same final geometry. This
is the fixed point that Chapter~\ref{c01:chap:difference} invoked in characterizing (not
deriving) the fundamental boundary (Corollary~\ref{c01:cor:closure-characterizes}); here we
see its origin: the fixed point is the convergence of Difference's own count-producing
process.

\begin{theorem}[The attractor is \texorpdfstring{$N=96$}{N=96}]
\label{c03:thm:n96}
The converged attractor of the actualization process is the $N=96$ network: the size $N=96$
is the stable fixed point of the dynamics, not a freely chosen input.
\end{theorem}

\begin{proof}
The boundary-origin audit (the Closure Principle \cite{c03:QG282}) establishes that the
observable sector is the converged attractor of the actualization dynamics. The dynamics has
a self-reinforcing feedback (activity to links to activity) bounded by network capacity, so
positive feedback saturates at a stable fixed point: zero residual link growth and a
content-independent final geometry. The convergence to the $N=96$ attractor follows from the
closure principle \cite{c03:QG282} and the reduction program \cite{c03:QG293,c03:QG294};
$N=96$ is thereby an output of the dynamics --- the inevitable fixed point --- not a
primitive or a free parameter.
\end{proof}

\begin{corollary}[The spectrum is the attractor's eigenspectrum]
\label{c03:cor:spectrum-eigenspectrum}
The D96 spectrum is the Laplacian eigenspectrum of the converged $N=96$ network: the
inevitable output of the actualization attractor, not an additional primitive.
\end{corollary}

\begin{proof}
The converged network defines a graph Laplacian; its eigenspectrum is the D96 spectrum, with
the 95 positive modes and the single zero mode (the background). Since the network is the
converged fixed point of the dynamics (Theorem~\ref{c03:thm:n96}), the spectrum is a
deterministic function of the attractor --- an output, not an input. This is the inevitable-
spectrum result of the canonical program \cite{c03:QG293,c03:QG294}.
\end{proof}

\section{Actualization and Resonance}
\label{c03:sec:actualization-resonance}

\begin{theorem}[Resonance is actualization's readout]
\label{c03:thm:resonance-readout}
\emph{Resonance} --- the spectral readout of the actualization process, the operator layer
projecting the converged spectrum onto the observable structure --- is actualization's
readout: the sector structure of observable physics is a projection of the one
actualization-driven spectrum. The resonance condition of the theory is the conjunction of
count conservation (the definitional identity) and the closure boundary (the fixed point):
Resonance equals Conservation plus Boundary.
\end{theorem}

\begin{proof}
The sector-emergence audit \cite{c03:QG272} establishes that distinct sectors --- masses,
couplings, mixings, gravity, cosmology --- are projection classes of a single operator layer
over the one spectrum, with no sector-exclusive operator. Since the spectrum is the output
of the actualization attractor (Corollary~\ref{c03:cor:spectrum-eigenspectrum}), the resonance
structure read from it is the readout of the actualization process: one dynamics, one
spectrum, one operator basis, projected onto distinct physical questions. Conservation is
the definitional identity of the count (Theorem~\ref{c03:thm:count-conservation}); the
boundary is the closure fixed point $N=96$ (Theorem~\ref{c03:thm:n96}); the spectrum is
determined by the process that conserves its count and closes at its boundary. The resonance
condition is therefore the conjunction of the two --- Conservation plus Boundary ---
actualization's spectral face.
\end{proof}

\section{Actualization as Derived Necessity}
\label{c03:sec:derived-necessity}

\begin{theorem}[Actualization is derived, not primitive]
\label{c03:thm:derived-not-primitive}
Actualization is INDISPENSABLE as a derivation layer --- without it there is no count, no
network, no spectrum, and no physics --- but it is a DERIVED necessity: indispensable as a
derivation step, yet fully dependent on the primitive Difference and therefore not an
underivable input.
\end{theorem}

\begin{proof}
The hierarchy-necessity audit \cite{c03:QG293} classifies Actualization as INDISPENSABLE: it is
the count-producing process, so without it there is no count, no network, no spectrum, and
no physics. But necessity as a step is not primitiveness as an input: by
Theorem~\ref{c03:thm:process-face}, Actualization is the process face of Difference; the
foundation stress test \cite{c03:QG292} establishes that removing Difference collapses
Actualization (it needs the definitional unit) while removing $\eta$ leaves it intact
(counting needs no metric reference). The minimal-theory audit \cite{c03:QG294} confirms that
the minimal hierarchy Difference → Actualization → Spectrum → Physics is complete, with
Actualization an essential derived step; the MONO006 resolution \cite{c03:MONO006} concludes:
Actualization is derived from Difference alone, indispensable as a step, not a third
primitive.
\end{proof}

The canonical architecture is therefore unchanged --- primitives $\{\text{Difference},\eta\}$,
Actualization derived from Difference --- consistent with Chapters~\ref{c01:chap:difference}
and \ref{c02:chap:eta}; consistent with MONO007 \cite{c03:MONO007}, all statements that could
be read as ascribing primitive status to Actualization are worded to make the derivation
explicit, and no ambiguity remains.

\section{Consequences}
\label{c03:sec:actualization-consequences}

\begin{theorem}[The spectrum is the inevitable output]
\label{c03:thm:spectrum-inevitable}
The D96 spectrum is the inevitable output of the actualization attractor: a deterministic
function of the converged network, carrying no independent choice.
\end{theorem}

\begin{proof}
By Corollary~\ref{c03:cor:spectrum-eigenspectrum}, the spectrum is the Laplacian eigenspectrum of
the converged $N=96$ attractor. The attractor is unique and content-independent
(Theorem~\ref{c03:thm:n96}); the eigenspectrum is a deterministic function of it. Hence the
spectrum carries no choice and is inevitable --- the output of the actualization process,
not an input to it. The reduction program over the fuller hierarchy thereby closes at the
minimal hierarchy Difference → Actualization → Spectrum → Physics: the hierarchy-necessity
audit \cite{c03:QG293} classifies the intermediate layers --- Closure COMPRESSIBLE into
Actualization (the fixed point of the process), Conservation COMPRESSIBLE into
Difference/network (the definitional identity plus the graph handshake), Resonance
COMPRESSIBLE into Spectrum (the operator readout) --- and the minimal-theory audit
\cite{c03:QG294} verifies that every compressed layer is derivable from the minimal hierarchy,
with none actually required as an input.
\end{proof}

\begin{corollary}[Physics is the readout of the actualized spectrum]
\label{c03:cor:physics-readout}
All observable physics is the readout of the actualized spectrum: the resonance structure
projected by the operator layer, from which the sectors of physics emerge.
\end{corollary}

\section{Summary}
\label{c03:sec:actualization-summary}

This chapter has established the second member of the canonical architecture. Actualization
is the process face of Difference (Theorem~\ref{c03:thm:process-face}), whose count
conservation is the primitive's definitional identity (Theorem~\ref{c03:thm:count-conservation})
and whose dynamics converge to the stable fixed point $N=96$
(Theorem~\ref{c03:thm:n96}, Corollary~\ref{c03:cor:spectrum-eigenspectrum}); the resonance
structure read from the resulting spectrum satisfies Resonance equals Conservation plus
Boundary (Theorem~\ref{c03:thm:resonance-readout}). Actualization is indispensable as a layer
yet derived, not primitive (Theorem~\ref{c03:thm:derived-not-primitive}), and the reduction
closes at the minimal hierarchy, with physics the readout of the actualized spectrum
(Theorem~\ref{c03:thm:spectrum-inevitable}, Corollary~\ref{c03:cor:physics-readout}).

The following chapter derives, from this process, the spectrum itself.

\begin{thebibliography}{99}
\bibitem{c03:QG268} AT-QG Phase 268, \emph{Count Conservation Origin}, Docs/Research.
\bibitem{c03:QG272} AT-QG Phase 272, \emph{Sector Emergence Audit}, Docs/Research.
\bibitem{c03:QG282} AT-QG Phase 282, \emph{Boundary Origin Audit} (Closure Principle), Docs/Research.
\bibitem{c03:QG292} AT-QG Phase 292, \emph{Foundation Stress Test}, Docs/Research.
\bibitem{c03:QG293} AT-QG Phase 293, \emph{Hierarchy Necessity Audit}, Docs/Research.
\bibitem{c03:QG294} AT-QG Phase 294, \emph{Minimal Theory Audit}, Docs/Research.
\bibitem{c03:MONO006} AT-MONO006, \emph{A01 Actualization Origin Audit}, Docs/Zenodo/Monograph.
\bibitem{c03:MONO007} AT-MONO007, \emph{Referee-Readiness Resolution}, Docs/Zenodo/Monograph.
\end{thebibliography} after the preamble that defines
%         the theorem environments below (or include via the master file).
% ======================================================================

% ---- Theorem environments (define in the master preamble if not present) ----
% \newtheorem{definition}{Definition}[chapter]
% \newtheorem{lemma}[definition]{Lemma}
% \newtheorem{theorem}[definition]{Theorem}
% \newtheorem{corollary}[definition]{Corollary}
% \newtheorem{remark}[definition]{Remark}

\chapter{Actualization}
\label{c03:chap:actualization}

\section{Introduction}
\label{c03:sec:actualization-introduction}

In Chapter~\ref{c01:chap:difference} we established that Difference is the fundamental boundary
of the theory --- the primitive whose count-conservation law is its definitional identity ---
and in Chapter~\ref{c02:chap:eta} that the tensor reference $\eta$ is the second primitive,
against which the tensor face of Difference is read. The canonical hierarchy
\eqref{c03:eq:core} then proceeds from the primitive pair to the inevitable spectrum through a
process:

\begin{equation}
\label{c03:eq:core}
\text{Difference}\;\longrightarrow\;\text{Actualization}\;\longrightarrow\;\text{Inevitable Spectrum}\;\longrightarrow\;\text{Physics}.
\end{equation}

This chapter develops the second member of that hierarchy: \emph{Actualization}. We
establish, with formal precision, that Actualization is the \emph{process face of
Difference} --- the count-producing dynamics whose definitional identity is count
conservation \cite{c03:QG268}; that its dynamics converge to the stable fixed point $N=96$, the
closure of the theory \cite{c03:QG282}; that the resonance structure of the spectrum is
actualization's readout \cite{c03:QG272}; and --- decisively, consistent with the MONO006
resolution \cite{c03:MONO006} --- that Actualization is a \emph{derived necessity}, not a third
primitive.

\section{Actualization as the Process Face of Difference}
\label{c03:sec:process-face}

\begin{theorem}[Actualization is the process face of Difference]
\label{c03:thm:process-face}
\emph{Actualization} --- the count-producing dynamics by which a Q-event, the unit of
Difference \cite{c03:QG268}, is applied, generating the network whose convergence yields the
spectrum --- is not an independent entity: it is the process face of the one primitive
Difference. Consequently the primitives of the theory are exactly
$\{\text{Difference},\eta\}$: Actualization is their derived process face, not a primitive.
\end{theorem}

\begin{proof}
Count conservation is the definitional identity of the primitive \cite{c03:QG268}: Difference is
the counting difference from a uniform background, and a Q-event is its unit. Actualization
is precisely the process that applies those units, and its identity is inherited from the
primitive: there is no Actualization without the unit, and no unit without Difference. The
minimal primitive set was established as the pair $\{\text{Difference},\eta\}$
(Theorem~\ref{c02:thm:minimal-foundation}, Chapter~\ref{c02:chap:eta}); the MONO006 resolution
\cite{c03:MONO006} confirms this classification, and the MONO007 architecture \cite{c03:MONO007}
presents Actualization as a derived layer. Hence there is no primitive-status ambiguity:
Actualization is the dynamical face of Difference, not a separate primitive.
\end{proof}

\section{Count Conservation}
\label{c03:sec:count-conservation}

\begin{theorem}[Count conservation is the definitional identity]
\label{c03:thm:count-conservation}
\emph{Count conservation} --- the total count of units of Difference is conserved by the
actualization process: Q-events are neither created nor destroyed, only redistributed ---
is not a separate postulate but the definitional identity of the primitive Difference and
therefore of its process face Actualization. One application of Actualization produces one
unit of count, so conservation follows automatically from the primitive, not as an imported
assumption of the theory.
\end{theorem}

\begin{proof}
The count-conservation origin \cite{c03:QG268} establishes that conservation is the defining
property of the primitive itself --- it is what \emph{is} the count --- and identifies the
Q-event as the indivisible step of the count. Since Actualization is the process that
applies the count (Theorem~\ref{c03:thm:process-face}), the conservation law is inherited:
the process conserves exactly what the primitive defines. No separate conservation postulate
is required.
\end{proof}

\section{The \texorpdfstring{$N=96$}{N=96} Attractor}
\label{c03:sec:n96-attractor}

The \emph{closure} of the theory is the stable fixed point of the actualization process: the
configuration at which the count-producing dynamics has converged, with zero residual
change. The boundary of the theory is this fixed point --- the topology saturates at zero
residual link growth, and every initial pattern converges to the same final geometry. This
is the fixed point that Chapter~\ref{c01:chap:difference} invoked in characterizing (not
deriving) the fundamental boundary (Corollary~\ref{c01:cor:closure-characterizes}); here we
see its origin: the fixed point is the convergence of Difference's own count-producing
process.

\begin{theorem}[The attractor is \texorpdfstring{$N=96$}{N=96}]
\label{c03:thm:n96}
The converged attractor of the actualization process is the $N=96$ network: the size $N=96$
is the stable fixed point of the dynamics, not a freely chosen input.
\end{theorem}

\begin{proof}
The boundary-origin audit (the Closure Principle \cite{c03:QG282}) establishes that the
observable sector is the converged attractor of the actualization dynamics. The dynamics has
a self-reinforcing feedback (activity to links to activity) bounded by network capacity, so
positive feedback saturates at a stable fixed point: zero residual link growth and a
content-independent final geometry. The convergence to the $N=96$ attractor follows from the
closure principle \cite{c03:QG282} and the reduction program \cite{c03:QG293,c03:QG294};
$N=96$ is thereby an output of the dynamics --- the inevitable fixed point --- not a
primitive or a free parameter.
\end{proof}

\begin{corollary}[The spectrum is the attractor's eigenspectrum]
\label{c03:cor:spectrum-eigenspectrum}
The D96 spectrum is the Laplacian eigenspectrum of the converged $N=96$ network: the
inevitable output of the actualization attractor, not an additional primitive.
\end{corollary}

\begin{proof}
The converged network defines a graph Laplacian; its eigenspectrum is the D96 spectrum, with
the 95 positive modes and the single zero mode (the background). Since the network is the
converged fixed point of the dynamics (Theorem~\ref{c03:thm:n96}), the spectrum is a
deterministic function of the attractor --- an output, not an input. This is the inevitable-
spectrum result of the canonical program \cite{c03:QG293,c03:QG294}.
\end{proof}

\section{Actualization and Resonance}
\label{c03:sec:actualization-resonance}

\begin{theorem}[Resonance is actualization's readout]
\label{c03:thm:resonance-readout}
\emph{Resonance} --- the spectral readout of the actualization process, the operator layer
projecting the converged spectrum onto the observable structure --- is actualization's
readout: the sector structure of observable physics is a projection of the one
actualization-driven spectrum. The resonance condition of the theory is the conjunction of
count conservation (the definitional identity) and the closure boundary (the fixed point):
Resonance equals Conservation plus Boundary.
\end{theorem}

\begin{proof}
The sector-emergence audit \cite{c03:QG272} establishes that distinct sectors --- masses,
couplings, mixings, gravity, cosmology --- are projection classes of a single operator layer
over the one spectrum, with no sector-exclusive operator. Since the spectrum is the output
of the actualization attractor (Corollary~\ref{c03:cor:spectrum-eigenspectrum}), the resonance
structure read from it is the readout of the actualization process: one dynamics, one
spectrum, one operator basis, projected onto distinct physical questions. Conservation is
the definitional identity of the count (Theorem~\ref{c03:thm:count-conservation}); the
boundary is the closure fixed point $N=96$ (Theorem~\ref{c03:thm:n96}); the spectrum is
determined by the process that conserves its count and closes at its boundary. The resonance
condition is therefore the conjunction of the two --- Conservation plus Boundary ---
actualization's spectral face.
\end{proof}

\section{Actualization as Derived Necessity}
\label{c03:sec:derived-necessity}

\begin{theorem}[Actualization is derived, not primitive]
\label{c03:thm:derived-not-primitive}
Actualization is INDISPENSABLE as a derivation layer --- without it there is no count, no
network, no spectrum, and no physics --- but it is a DERIVED necessity: indispensable as a
derivation step, yet fully dependent on the primitive Difference and therefore not an
underivable input.
\end{theorem}

\begin{proof}
The hierarchy-necessity audit \cite{c03:QG293} classifies Actualization as INDISPENSABLE: it is
the count-producing process, so without it there is no count, no network, no spectrum, and
no physics. But necessity as a step is not primitiveness as an input: by
Theorem~\ref{c03:thm:process-face}, Actualization is the process face of Difference; the
foundation stress test \cite{c03:QG292} establishes that removing Difference collapses
Actualization (it needs the definitional unit) while removing $\eta$ leaves it intact
(counting needs no metric reference). The minimal-theory audit \cite{c03:QG294} confirms that
the minimal hierarchy Difference → Actualization → Spectrum → Physics is complete, with
Actualization an essential derived step; the MONO006 resolution \cite{c03:MONO006} concludes:
Actualization is derived from Difference alone, indispensable as a step, not a third
primitive.
\end{proof}

The canonical architecture is therefore unchanged --- primitives $\{\text{Difference},\eta\}$,
Actualization derived from Difference --- consistent with Chapters~\ref{c01:chap:difference}
and \ref{c02:chap:eta}; consistent with MONO007 \cite{c03:MONO007}, all statements that could
be read as ascribing primitive status to Actualization are worded to make the derivation
explicit, and no ambiguity remains.

\section{Consequences}
\label{c03:sec:actualization-consequences}

\begin{theorem}[The spectrum is the inevitable output]
\label{c03:thm:spectrum-inevitable}
The D96 spectrum is the inevitable output of the actualization attractor: a deterministic
function of the converged network, carrying no independent choice.
\end{theorem}

\begin{proof}
By Corollary~\ref{c03:cor:spectrum-eigenspectrum}, the spectrum is the Laplacian eigenspectrum of
the converged $N=96$ attractor. The attractor is unique and content-independent
(Theorem~\ref{c03:thm:n96}); the eigenspectrum is a deterministic function of it. Hence the
spectrum carries no choice and is inevitable --- the output of the actualization process,
not an input to it. The reduction program over the fuller hierarchy thereby closes at the
minimal hierarchy Difference → Actualization → Spectrum → Physics: the hierarchy-necessity
audit \cite{c03:QG293} classifies the intermediate layers --- Closure COMPRESSIBLE into
Actualization (the fixed point of the process), Conservation COMPRESSIBLE into
Difference/network (the definitional identity plus the graph handshake), Resonance
COMPRESSIBLE into Spectrum (the operator readout) --- and the minimal-theory audit
\cite{c03:QG294} verifies that every compressed layer is derivable from the minimal hierarchy,
with none actually required as an input.
\end{proof}

\begin{corollary}[Physics is the readout of the actualized spectrum]
\label{c03:cor:physics-readout}
All observable physics is the readout of the actualized spectrum: the resonance structure
projected by the operator layer, from which the sectors of physics emerge.
\end{corollary}

\section{Summary}
\label{c03:sec:actualization-summary}

This chapter has established the second member of the canonical architecture. Actualization
is the process face of Difference (Theorem~\ref{c03:thm:process-face}), whose count
conservation is the primitive's definitional identity (Theorem~\ref{c03:thm:count-conservation})
and whose dynamics converge to the stable fixed point $N=96$
(Theorem~\ref{c03:thm:n96}, Corollary~\ref{c03:cor:spectrum-eigenspectrum}); the resonance
structure read from the resulting spectrum satisfies Resonance equals Conservation plus
Boundary (Theorem~\ref{c03:thm:resonance-readout}). Actualization is indispensable as a layer
yet derived, not primitive (Theorem~\ref{c03:thm:derived-not-primitive}), and the reduction
closes at the minimal hierarchy, with physics the readout of the actualized spectrum
(Theorem~\ref{c03:thm:spectrum-inevitable}, Corollary~\ref{c03:cor:physics-readout}).

The following chapter derives, from this process, the spectrum itself.

\begin{thebibliography}{99}
\bibitem{c03:QG268} AT-QG Phase 268, \emph{Count Conservation Origin}, Docs/Research.
\bibitem{c03:QG272} AT-QG Phase 272, \emph{Sector Emergence Audit}, Docs/Research.
\bibitem{c03:QG282} AT-QG Phase 282, \emph{Boundary Origin Audit} (Closure Principle), Docs/Research.
\bibitem{c03:QG292} AT-QG Phase 292, \emph{Foundation Stress Test}, Docs/Research.
\bibitem{c03:QG293} AT-QG Phase 293, \emph{Hierarchy Necessity Audit}, Docs/Research.
\bibitem{c03:QG294} AT-QG Phase 294, \emph{Minimal Theory Audit}, Docs/Research.
\bibitem{c03:MONO006} AT-MONO006, \emph{A01 Actualization Origin Audit}, Docs/Zenodo/Monograph.
\bibitem{c03:MONO007} AT-MONO007, \emph{Referee-Readiness Resolution}, Docs/Zenodo/Monograph.
\end{thebibliography} after the preamble that defines
%         the theorem environments below (or include via the master file).
% ======================================================================

% ---- Theorem environments (define in the master preamble if not present) ----
% \newtheorem{definition}{Definition}[chapter]
% \newtheorem{lemma}[definition]{Lemma}
% \newtheorem{theorem}[definition]{Theorem}
% \newtheorem{corollary}[definition]{Corollary}
% \newtheorem{remark}[definition]{Remark}

\chapter{Actualization}
\label{c03:chap:actualization}

\section{Introduction}
\label{c03:sec:actualization-introduction}

In Chapter~\ref{c01:chap:difference} we established that Difference is the fundamental boundary
of the theory --- the primitive whose count-conservation law is its definitional identity ---
and in Chapter~\ref{c02:chap:eta} that the tensor reference $\eta$ is the second primitive,
against which the tensor face of Difference is read. The canonical hierarchy
\eqref{c03:eq:core} then proceeds from the primitive pair to the inevitable spectrum through a
process:

\begin{equation}
\label{c03:eq:core}
\text{Difference}\;\longrightarrow\;\text{Actualization}\;\longrightarrow\;\text{Inevitable Spectrum}\;\longrightarrow\;\text{Physics}.
\end{equation}

This chapter develops the second member of that hierarchy: \emph{Actualization}. We
establish, with formal precision, that Actualization is the \emph{process face of
Difference} --- the count-producing dynamics whose definitional identity is count
conservation \cite{c03:QG268}; that its dynamics converge to the stable fixed point $N=96$, the
closure of the theory \cite{c03:QG282}; that the resonance structure of the spectrum is
actualization's readout \cite{c03:QG272}; and --- decisively, consistent with the MONO006
resolution \cite{c03:MONO006} --- that Actualization is a \emph{derived necessity}, not a third
primitive.

\section{Actualization as the Process Face of Difference}
\label{c03:sec:process-face}

\begin{theorem}[Actualization is the process face of Difference]
\label{c03:thm:process-face}
\emph{Actualization} --- the count-producing dynamics by which a Q-event, the unit of
Difference \cite{c03:QG268}, is applied, generating the network whose convergence yields the
spectrum --- is not an independent entity: it is the process face of the one primitive
Difference. Consequently the primitives of the theory are exactly
$\{\text{Difference},\eta\}$: Actualization is their derived process face, not a primitive.
\end{theorem}

\begin{proof}
Count conservation is the definitional identity of the primitive \cite{c03:QG268}: Difference is
the counting difference from a uniform background, and a Q-event is its unit. Actualization
is precisely the process that applies those units, and its identity is inherited from the
primitive: there is no Actualization without the unit, and no unit without Difference. The
minimal primitive set was established as the pair $\{\text{Difference},\eta\}$
(Theorem~\ref{c02:thm:minimal-foundation}, Chapter~\ref{c02:chap:eta}); the MONO006 resolution
\cite{c03:MONO006} confirms this classification, and the MONO007 architecture \cite{c03:MONO007}
presents Actualization as a derived layer. Hence there is no primitive-status ambiguity:
Actualization is the dynamical face of Difference, not a separate primitive.
\end{proof}

\section{Count Conservation}
\label{c03:sec:count-conservation}

\begin{theorem}[Count conservation is the definitional identity]
\label{c03:thm:count-conservation}
\emph{Count conservation} --- the total count of units of Difference is conserved by the
actualization process: Q-events are neither created nor destroyed, only redistributed ---
is not a separate postulate but the definitional identity of the primitive Difference and
therefore of its process face Actualization. One application of Actualization produces one
unit of count, so conservation follows automatically from the primitive, not as an imported
assumption of the theory.
\end{theorem}

\begin{proof}
The count-conservation origin \cite{c03:QG268} establishes that conservation is the defining
property of the primitive itself --- it is what \emph{is} the count --- and identifies the
Q-event as the indivisible step of the count. Since Actualization is the process that
applies the count (Theorem~\ref{c03:thm:process-face}), the conservation law is inherited:
the process conserves exactly what the primitive defines. No separate conservation postulate
is required.
\end{proof}

\section{The \texorpdfstring{$N=96$}{N=96} Attractor}
\label{c03:sec:n96-attractor}

The \emph{closure} of the theory is the stable fixed point of the actualization process: the
configuration at which the count-producing dynamics has converged, with zero residual
change. The boundary of the theory is this fixed point --- the topology saturates at zero
residual link growth, and every initial pattern converges to the same final geometry. This
is the fixed point that Chapter~\ref{c01:chap:difference} invoked in characterizing (not
deriving) the fundamental boundary (Corollary~\ref{c01:cor:closure-characterizes}); here we
see its origin: the fixed point is the convergence of Difference's own count-producing
process.

\begin{theorem}[The attractor is \texorpdfstring{$N=96$}{N=96}]
\label{c03:thm:n96}
The converged attractor of the actualization process is the $N=96$ network: the size $N=96$
is the stable fixed point of the dynamics, not a freely chosen input.
\end{theorem}

\begin{proof}
The boundary-origin audit (the Closure Principle \cite{c03:QG282}) establishes that the
observable sector is the converged attractor of the actualization dynamics. The dynamics has
a self-reinforcing feedback (activity to links to activity) bounded by network capacity, so
positive feedback saturates at a stable fixed point: zero residual link growth and a
content-independent final geometry. The convergence to the $N=96$ attractor follows from the
closure principle \cite{c03:QG282} and the reduction program \cite{c03:QG293,c03:QG294};
$N=96$ is thereby an output of the dynamics --- the inevitable fixed point --- not a
primitive or a free parameter.
\end{proof}

\begin{corollary}[The spectrum is the attractor's eigenspectrum]
\label{c03:cor:spectrum-eigenspectrum}
The D96 spectrum is the Laplacian eigenspectrum of the converged $N=96$ network: the
inevitable output of the actualization attractor, not an additional primitive.
\end{corollary}

\begin{proof}
The converged network defines a graph Laplacian; its eigenspectrum is the D96 spectrum, with
the 95 positive modes and the single zero mode (the background). Since the network is the
converged fixed point of the dynamics (Theorem~\ref{c03:thm:n96}), the spectrum is a
deterministic function of the attractor --- an output, not an input. This is the inevitable-
spectrum result of the canonical program \cite{c03:QG293,c03:QG294}.
\end{proof}

\section{Actualization and Resonance}
\label{c03:sec:actualization-resonance}

\begin{theorem}[Resonance is actualization's readout]
\label{c03:thm:resonance-readout}
\emph{Resonance} --- the spectral readout of the actualization process, the operator layer
projecting the converged spectrum onto the observable structure --- is actualization's
readout: the sector structure of observable physics is a projection of the one
actualization-driven spectrum. The resonance condition of the theory is the conjunction of
count conservation (the definitional identity) and the closure boundary (the fixed point):
Resonance equals Conservation plus Boundary.
\end{theorem}

\begin{proof}
The sector-emergence audit \cite{c03:QG272} establishes that distinct sectors --- masses,
couplings, mixings, gravity, cosmology --- are projection classes of a single operator layer
over the one spectrum, with no sector-exclusive operator. Since the spectrum is the output
of the actualization attractor (Corollary~\ref{c03:cor:spectrum-eigenspectrum}), the resonance
structure read from it is the readout of the actualization process: one dynamics, one
spectrum, one operator basis, projected onto distinct physical questions. Conservation is
the definitional identity of the count (Theorem~\ref{c03:thm:count-conservation}); the
boundary is the closure fixed point $N=96$ (Theorem~\ref{c03:thm:n96}); the spectrum is
determined by the process that conserves its count and closes at its boundary. The resonance
condition is therefore the conjunction of the two --- Conservation plus Boundary ---
actualization's spectral face.
\end{proof}

\section{Actualization as Derived Necessity}
\label{c03:sec:derived-necessity}

\begin{theorem}[Actualization is derived, not primitive]
\label{c03:thm:derived-not-primitive}
Actualization is INDISPENSABLE as a derivation layer --- without it there is no count, no
network, no spectrum, and no physics --- but it is a DERIVED necessity: indispensable as a
derivation step, yet fully dependent on the primitive Difference and therefore not an
underivable input.
\end{theorem}

\begin{proof}
The hierarchy-necessity audit \cite{c03:QG293} classifies Actualization as INDISPENSABLE: it is
the count-producing process, so without it there is no count, no network, no spectrum, and
no physics. But necessity as a step is not primitiveness as an input: by
Theorem~\ref{c03:thm:process-face}, Actualization is the process face of Difference; the
foundation stress test \cite{c03:QG292} establishes that removing Difference collapses
Actualization (it needs the definitional unit) while removing $\eta$ leaves it intact
(counting needs no metric reference). The minimal-theory audit \cite{c03:QG294} confirms that
the minimal hierarchy Difference → Actualization → Spectrum → Physics is complete, with
Actualization an essential derived step; the MONO006 resolution \cite{c03:MONO006} concludes:
Actualization is derived from Difference alone, indispensable as a step, not a third
primitive.
\end{proof}

The canonical architecture is therefore unchanged --- primitives $\{\text{Difference},\eta\}$,
Actualization derived from Difference --- consistent with Chapters~\ref{c01:chap:difference}
and \ref{c02:chap:eta}; consistent with MONO007 \cite{c03:MONO007}, all statements that could
be read as ascribing primitive status to Actualization are worded to make the derivation
explicit, and no ambiguity remains.

\section{Consequences}
\label{c03:sec:actualization-consequences}

\begin{theorem}[The spectrum is the inevitable output]
\label{c03:thm:spectrum-inevitable}
The D96 spectrum is the inevitable output of the actualization attractor: a deterministic
function of the converged network, carrying no independent choice.
\end{theorem}

\begin{proof}
By Corollary~\ref{c03:cor:spectrum-eigenspectrum}, the spectrum is the Laplacian eigenspectrum of
the converged $N=96$ attractor. The attractor is unique and content-independent
(Theorem~\ref{c03:thm:n96}); the eigenspectrum is a deterministic function of it. Hence the
spectrum carries no choice and is inevitable --- the output of the actualization process,
not an input to it. The reduction program over the fuller hierarchy thereby closes at the
minimal hierarchy Difference → Actualization → Spectrum → Physics: the hierarchy-necessity
audit \cite{c03:QG293} classifies the intermediate layers --- Closure COMPRESSIBLE into
Actualization (the fixed point of the process), Conservation COMPRESSIBLE into
Difference/network (the definitional identity plus the graph handshake), Resonance
COMPRESSIBLE into Spectrum (the operator readout) --- and the minimal-theory audit
\cite{c03:QG294} verifies that every compressed layer is derivable from the minimal hierarchy,
with none actually required as an input.
\end{proof}

\begin{corollary}[Physics is the readout of the actualized spectrum]
\label{c03:cor:physics-readout}
All observable physics is the readout of the actualized spectrum: the resonance structure
projected by the operator layer, from which the sectors of physics emerge.
\end{corollary}

\section{Summary}
\label{c03:sec:actualization-summary}

This chapter has established the second member of the canonical architecture. Actualization
is the process face of Difference (Theorem~\ref{c03:thm:process-face}), whose count
conservation is the primitive's definitional identity (Theorem~\ref{c03:thm:count-conservation})
and whose dynamics converge to the stable fixed point $N=96$
(Theorem~\ref{c03:thm:n96}, Corollary~\ref{c03:cor:spectrum-eigenspectrum}); the resonance
structure read from the resulting spectrum satisfies Resonance equals Conservation plus
Boundary (Theorem~\ref{c03:thm:resonance-readout}). Actualization is indispensable as a layer
yet derived, not primitive (Theorem~\ref{c03:thm:derived-not-primitive}), and the reduction
closes at the minimal hierarchy, with physics the readout of the actualized spectrum
(Theorem~\ref{c03:thm:spectrum-inevitable}, Corollary~\ref{c03:cor:physics-readout}).

The following chapter derives, from this process, the spectrum itself.

\begin{thebibliography}{99}
\bibitem{c03:QG268} AT-QG Phase 268, \emph{Count Conservation Origin}, Docs/Research.
\bibitem{c03:QG272} AT-QG Phase 272, \emph{Sector Emergence Audit}, Docs/Research.
\bibitem{c03:QG282} AT-QG Phase 282, \emph{Boundary Origin Audit} (Closure Principle), Docs/Research.
\bibitem{c03:QG292} AT-QG Phase 292, \emph{Foundation Stress Test}, Docs/Research.
\bibitem{c03:QG293} AT-QG Phase 293, \emph{Hierarchy Necessity Audit}, Docs/Research.
\bibitem{c03:QG294} AT-QG Phase 294, \emph{Minimal Theory Audit}, Docs/Research.
\bibitem{c03:MONO006} AT-MONO006, \emph{A01 Actualization Origin Audit}, Docs/Zenodo/Monograph.
\bibitem{c03:MONO007} AT-MONO007, \emph{Referee-Readiness Resolution}, Docs/Zenodo/Monograph.
\end{thebibliography}
% ======================================================================
%  TQM-MONO103 — Chapter 4: Closure and the Minimal Theory
%  Canonical Monograph (v2.0) · The Actualization Theory:
%  A Reconstruction of Physics from Difference, Actualization and Spectrum
%
%  Status: COMPLETE — PASS-READY (MONO007 referee-ready architecture)
%  Inputs: Chapter01_Difference.tex, Chapter02_Eta.tex,
%          Chapter03_Actualization.tex, TQMQG_CanonicalMonograph.md,
%          TQMQG_Mono007RefereeReadiness.md
%  Method: assembly only — canonical results only, no new physics, no speculation
%  Citations: QG267 (conservation principle), QG268 (count conservation),
%             QG278 (fundamental boundary), QG282 (closure principle),
%             QG293 (hierarchy necessity), QG294 (minimal theory),
%             MONO006 (actualization derived), MONO007 (referee readiness)
%
%  Usage: % ======================================================================
%  TQM-MONO103 — Chapter 4: Closure and the Minimal Theory
%  Canonical Monograph (v2.0) · The Actualization Theory:
%  A Reconstruction of Physics from Difference, Actualization and Spectrum
%
%  Status: COMPLETE — PASS-READY (MONO007 referee-ready architecture)
%  Inputs: Chapter01_Difference.tex, Chapter02_Eta.tex,
%          Chapter03_Actualization.tex, TQMQG_CanonicalMonograph.md,
%          TQMQG_Mono007RefereeReadiness.md
%  Method: assembly only — canonical results only, no new physics, no speculation
%  Citations: QG267 (conservation principle), QG268 (count conservation),
%             QG278 (fundamental boundary), QG282 (closure principle),
%             QG293 (hierarchy necessity), QG294 (minimal theory),
%             MONO006 (actualization derived), MONO007 (referee readiness)
%
%  Usage: % ======================================================================
%  TQM-MONO103 — Chapter 4: Closure and the Minimal Theory
%  Canonical Monograph (v2.0) · The Actualization Theory:
%  A Reconstruction of Physics from Difference, Actualization and Spectrum
%
%  Status: COMPLETE — PASS-READY (MONO007 referee-ready architecture)
%  Inputs: Chapter01_Difference.tex, Chapter02_Eta.tex,
%          Chapter03_Actualization.tex, TQMQG_CanonicalMonograph.md,
%          TQMQG_Mono007RefereeReadiness.md
%  Method: assembly only — canonical results only, no new physics, no speculation
%  Citations: QG267 (conservation principle), QG268 (count conservation),
%             QG278 (fundamental boundary), QG282 (closure principle),
%             QG293 (hierarchy necessity), QG294 (minimal theory),
%             MONO006 (actualization derived), MONO007 (referee readiness)
%
%  Usage: \input{Chapter04_ClosureAndMinimalTheory.tex} after the preamble
%         that defines the theorem environments (or include via the master file).
% ======================================================================

% ---- Theorem environments (define in the master preamble if not present) ----
% \newtheorem{definition}{Definition}[chapter]
% \newtheorem{lemma}[definition]{Lemma}
% \newtheorem{theorem}[definition]{Theorem}
% \newtheorem{corollary}[definition]{Corollary}
% \newtheorem{remark}[definition]{Remark}

\chapter{Closure and the Minimal Theory}
\label{chap:closure-minimal}

\section{Introduction}
\label{sec:closure-introduction}

In the preceding chapters we established the primitives $\{\text{Difference},\eta\}$
(Chapters~\ref{chap:difference} and \ref{chap:eta}) and the derived process face of
Difference, Actualization, whose convergence yields the inevitable spectrum
(Chapter~\ref{chap:actualization}). This chapter completes the foundational picture by
establishing two structural results. First, the \emph{Closure Principle}: the boundary of the
theory is the fixed point of Difference's own dynamics, and --- decisively --- the Closure
Principle \emph{characterizes} the fundamental boundary but does \emph{not derive} the
primitive \cite{QG278,QG282}. Second, the \emph{minimal theory}: the hierarchy
\begin{equation}
\label{eq:minimal}
\text{Difference}\;\longrightarrow\;\text{Actualization}\;\longrightarrow\;\text{Inevitable Spectrum}\;\longrightarrow\;\text{Physics}
\end{equation}
is complete and irreducible \cite{QG293,QG294}: every intermediate layer is derivable, and
no further primitive is required.

The results of this chapter are the canonical end-state of the closure and reduction program
\cite{QG267,QG268,QG278,QG282,QG293,QG294}, consistent with the MONO006 resolution and the
MONO007 referee-ready architecture \cite{MONO006,MONO007}. Throughout we use only accepted
canonical results and introduce no new physics and no speculation.

\section{The Closure Principle}
\label{sec:closure-principle}

\begin{definition}[Closure Principle]
\label{def:closure-principle}
The \emph{Closure Principle} is the statement that the boundary of the theory is the fixed
point of the actualization process: the primitive is the process, and the boundary is its
converged fixed point \cite{QG282}.
\end{definition}

\begin{theorem}[The Closure Principle holds]
\label{thm:closure-holds}
The boundary of the theory is the fixed point of Difference's own count-producing dynamics:
the topology saturates at zero residual change, and every initial pattern converges to the
same final geometry.
\end{theorem}

\begin{proof}
The boundary-origin audit \cite{QG282} establishes the closure: the actualization dynamics
has a self-reinforcing feedback bounded by network capacity, so positive feedback saturates
at a stable fixed point with zero residual link growth and a content-independent final
geometry. The observable sector is therefore the converged attractor of the dynamics, and
the boundary is the closure of that dynamics. The closure is derived --- it is an output of
the process --- not an imported cut-off.
\end{proof}

\begin{corollary}[The closure is \texorpdfstring{$N=96$}{N=96}]
\label{cor:closure-n96}
The fixed point of the actualization process is the $N=96$ network, whose eigenspectrum is
the inevitable D96 spectrum.
\end{corollary}

\begin{proof}
By Theorem~\ref{thm:closure-holds}, the boundary is the converged attractor of the
dynamics; by Theorem~\ref{thm:n96} of Chapter~\ref{chap:actualization}, that attractor is
the $N=96$ network; and by Corollary~\ref{cor:spectrum-eigenspectrum} of
Chapter~\ref{chap:actualization}, its eigenspectrum is the inevitable D96 spectrum.
\end{proof}

\section{Boundary and Fixed Point}
\label{sec:boundary-fixed-point}

\begin{definition}[Boundary]
\label{def:boundary}
A \emph{boundary} of the theory is the configuration at which the derived structure closes:
the point where the count-producing dynamics has converged and no further change occurs.
\end{definition}

\begin{theorem}[The boundary is the fixed point]
\label{thm:boundary-is-fixed}
The boundary of the theory is the stable fixed point of the actualization process; it is
derived, not primitive.
\end{theorem}

\begin{proof}
By Definition~\ref{def:closure-principle} and Theorem~\ref{thm:closure-holds}, the boundary
is the converged fixed point of the process. The derivation structure of the process is
established by Theorem~\ref{thm:derived-not-primitive} of
Chapter~\ref{chap:actualization}: the process is the face of Difference, and its fixed
point is an output of the dynamics. Hence the boundary is derived --- it is not an imported
input or a separate primitive.
\end{proof}

\begin{lemma}[The boundary characterizes the primitive]
\label{lem:boundary-characterizes}
The fixed-point boundary is a characterization of the primitive Difference: it specifies
how Difference's dynamics closes, and thereby \emph{what} the boundary is, without defining
\emph{what Difference is}.
\end{lemma}

\begin{proof}
The boundary is the fixed point of the process that Difference generates
(Theorem~\ref{thm:boundary-is-fixed}). It tells us where the process converges --- the
closure $N=96$ --- but presupposes the process and hence the primitive. A characterization
specifies properties of an object already given; it does not supply the object itself. The
fixed-point boundary therefore characterizes Difference without deriving it.
\end{proof}

\begin{corollary}[Closure does not derive Difference]
\label{cor:closure-not-derive}
The Closure Principle does not derive the primitive Difference. It characterizes the
boundary of the theory as the fixed point of Difference's own dynamics.
\end{corollary}

\begin{proof}
By Lemma~\ref{lem:boundary-characterizes}, the closure is a characterization, not a
derivation. Moreover, by Theorem~\ref{thm:fundamental-boundary} of
Chapter~\ref{chap:difference}, Difference is the fundamental boundary: irreducible, with its
count-conservation law as its definitional identity. A derivation of the primitive would
presuppose a deeper notion, which the fundamental-boundary audit \cite{QG278} shows does
not exist --- every attempt to reduce Difference reintroduces it. Hence the Closure
Principle characterizes, and does not derive, the primitive.
\end{proof}

\begin{remark}[Referee-safe reading]
Consistent with MONO007 \cite{MONO007}, every statement in this chapter that invokes the
Closure Principle is worded so that the principle is never presented as a derivation of
Difference: the primitive is given, and the closure characterizes where its process closes.
This removes the circularity objection A02 of the hostile-referee audit.
\end{remark}

\section{Self-Consistency}
\label{sec:self-consistency}

\begin{definition}[Self-consistency]
\label{def:self-consistency}
\emph{Self-consistency} is the requirement that the theory be free of internal
contradiction: non-contradiction requires comparing parts --- a difference.
\end{definition}

\begin{theorem}[Self-consistency presupposes Difference]
\label{thm:self-consistency-presupposes}
Self-consistency is derivable from Difference: non-contradiction is the comparison of parts,
and identity is the zero difference.
\end{theorem}

\begin{proof}
The fundamental-boundary audit \cite{QG278} establishes that self-consistency presupposes
Difference: non-contradiction requires comparing parts of the theory --- a difference ---
and identity is the zero difference \cite{QG268}. The conservation of the count, which is
the definitional identity of the primitive, is the universal conservation principle of which
all conservation laws are manifestations \cite{QG267}. Hence self-consistency is not an
independent postulate but a consequence of the primitive's structure.
\end{proof}

\begin{corollary}[Conservation is universal]
\label{cor:conservation-universal}
All conservation laws of the theory --- norm, trace, count, Bianchi, Noether currents ---
are manifestations of one universal conservation principle rooted in the identity of
Difference.
\end{corollary}

\begin{proof}
The conservation-principle audit \cite{QG267} verifies that the six conservation laws (Born
norm, unitarity, trace conservation, count conservation, Bianchi conservation, Noether
currents) are all manifestations of one deeper principle. By
Theorem~\ref{thm:count-conservation} of Chapter~\ref{chap:actualization}, count
conservation is the definitional identity of the primitive. Hence the universal conservation
principle is the identity of Difference, and all conservation laws are its manifestations.
\end{proof}

\section{Individuation}
\label{sec:individuation}

\begin{definition}[Individuation]
\label{def:individuation}
\emph{Individuation} is the process by which the one Difference gives rise to distinct
sectors: the differentiation of the theory's structure into the separate domains of
observable physics.
\end{definition}

\begin{theorem}[Individuation derives from Difference]
\label{thm:individuation}
The individuation of the theory into distinct sectors is derivable from Difference: distinct
sectors are distinct differences.
\end{theorem}

\begin{proof}
Sectors are distinguished by their different spectral access --- the projection classes over
the operator layer established in Chapter~\ref{chap:actualization}
(Theorem~\ref{thm:resonance-readout}). Each sector is a distinct pattern of difference in
the count structure. The individuation of the theory is therefore the differentiation of the
one Difference into its distinct faces and readouts; it is not an independent primitive but
a derived structure of the single difference principle.
\end{proof}

\section{The Difference Principle}
\label{sec:difference-principle}

\begin{definition}[Difference Principle]
\label{def:difference-principle}
The \emph{Difference Principle} is the statement that Difference is the root of the
hierarchy: the first concept, presupposed by everything else, to which every reduction
attempt returns.
\end{definition}

\begin{theorem}[Difference is the first concept]
\label{thm:difference-first}
Difference is the first concept of the theory: every attempt to reduce it reintroduces it,
so it is a self-referential fundamental boundary.
\end{theorem}

\begin{proof}
The fundamental-boundary audit \cite{QG278} performs the self-referential check. Actualization
presupposes Difference (an act of actualizing is a change --- a before/after difference; a
Q-event is a transition). Question presupposes Difference (a question is a gap --- the
difference between known and unknown). Self-consistency presupposes Difference
(non-contradiction requires comparing parts; identity is the zero difference). And
Difference itself is the first concept: any attempt to reduce it reintroduces it --- "two
distinct things" uses distinct = difference, "a boundary" is where things differ, "a
transition" is a before/after difference, "empty vs non-empty" is itself a difference. Hence
Difference is a self-referential fundamental boundary, and the Difference Principle states
exactly this: Difference is the root.
\end{proof}

\begin{corollary}[The Difference Principle anchors the hierarchy]
\label{cor:difference-anchors}
The Difference Principle anchors the canonical hierarchy
\eqref{eq:minimal}: every layer --- Actualization, Spectrum, Physics --- presupposes
Difference as its root.
\end{corollary}

\begin{proof}
By Theorem~\ref{thm:difference-first}, every reduction returns to Difference; by
Theorem~\ref{thm:process-face} of Chapter~\ref{chap:actualization}, Actualization is the
process face of Difference; and the spectrum and physics are the outputs of that process
(Chapter~\ref{chap:actualization}, Theorem~\ref{thm:spectrum-inevitable}). Hence the
Difference Principle anchors the entire hierarchy at its root.
\end{proof}

\section{The Minimal Theory}
\label{sec:minimal-theory}

\begin{theorem}[The minimal hierarchy]
\label{thm:minimal-hierarchy}
The minimal hierarchy is
\[
\text{Difference}\;\longrightarrow\;\text{Actualization}\;\longrightarrow\;\text{Inevitable Spectrum}\;\longrightarrow\;\text{Physics}.
\]
It is complete: every phase of the theory can be rewritten using only this hierarchy.
\end{theorem}

\begin{proof}
The hierarchy-necessity audit \cite{QG293} removes each intermediate layer of the fuller
hierarchy (Difference → Actualization → Closure → Conservation → Resonance → Spectrum →
Question → Measurement → Physics) and classifies it. Actualization is INDISPENSABLE (the
count-producing process). Closure is COMPRESSIBLE into Actualization (the fixed point of the
process). Conservation is COMPRESSIBLE into Difference/network (the definitional identity
plus the graph handshake). Resonance is COMPRESSIBLE into Spectrum (the operator readout).
The minimal-theory audit \cite{QG294} then re-derives every compressed layer from the
minimal hierarchy, classifying each DERIVABLE or ACTUALLY REQUIRED, and finds
\emph{5 DERIVABLE / 0 ACTUALLY REQUIRED}. Hence the minimal hierarchy is complete.
\end{proof}

\begin{theorem}[The minimal theory is irreducible]
\label{thm:minimal-irreducible}
The minimal theory is irreducible: no layer of the hierarchy
\eqref{eq:minimal} can be removed without losing derived content, and no further primitive
exists.
\end{theorem}

\begin{proof}
Irreducibility has two parts. \emph{Primitives:} by Theorem~\ref{thm:minimal-foundation} of
Chapter~\ref{chap:eta}, the minimal primitive set is $\{\text{Difference},\eta\}$; by
Theorem~\ref{thm:difference-first}, no deeper notion exists; by
Theorem~\ref{thm:derived-not-primitive} of Chapter~\ref{chap:actualization}, Actualization
is derived, not primitive. \emph{Layers:} the foundation stress test \cite{QG292}
establishes that removing Difference collapses all layers, and the hierarchy-necessity audit
\cite{QG293} establishes that removing Actualization destroys the count-producing process.
The minimal-theory audit \cite{QG294} confirms that the four layers of \eqref{eq:minimal}
are exactly what is required, with zero additional inputs. Hence the minimal theory is
irreducible.
\end{proof}

\begin{corollary}[The canonical architecture]
\label{cor:minimal-architecture}
The canonical architecture of the theory is the minimal hierarchy
\eqref{eq:minimal}: Difference → Actualization → Inevitable Spectrum → Physics, with
primitives $\{\text{Difference},\eta\}$, Actualization derived, the spectrum inevitable, and
physics emergent.
\end{corollary}

\begin{proof}
The minimal hierarchy is complete (Theorem~\ref{thm:minimal-hierarchy}) and irreducible
(Theorem~\ref{thm:minimal-irreducible}); the primitives are the pair
$\{\text{Difference},\eta\}$ (Theorem~\ref{thm:minimal-foundation} of
Chapter~\ref{chap:eta}); Actualization is derived (Theorem~\ref{thm:derived-not-primitive}
of Chapter~\ref{chap:actualization}); the spectrum is inevitable
(Theorem~\ref{thm:spectrum-inevitable} of Chapter~\ref{chap:actualization}). This is the
referee-ready structure of MONO007 \cite{MONO007}.
\end{proof}

\begin{remark}[Referee-safe completeness]
Consistent with MONO007 \cite{MONO007}, where the monograph claims completeness of the
minimal theory, it is understood in the qualified sense: every \emph{observable} is derived
from the minimal hierarchy, while the standard-model \emph{Lagrangian} and its dynamics are
\emph{hosted} --- a boundary, not a derivation gap. The minimal theory is irreducible as a
derivation structure; it does not assert a derivation of every aspect of the hosted dynamics.
\end{remark}

\section{Consequences}
\label{sec:closure-consequences}

\begin{lemma}[The reduction closes at the minimal hierarchy]
\label{lem:reduction-closes}
The reduction program closes at the minimal hierarchy \eqref{eq:minimal}: all intermediate
layers of the fuller hierarchy are derivable from it, and none is actually required as an
input.
\end{lemma}

\begin{proof}
The minimal-theory audit \cite{QG294} classifies the five compressed layers --- Closure,
Conservation, Resonance, and the remaining intermediates --- as DERIVABLE, with zero
ACTUALLY REQUIRED. The re-derivation uses only the minimal hierarchy. Hence the reduction
closes at \eqref{eq:minimal}.
\end{proof}

\begin{theorem}[Closure and minimality are consistent]
\label{thm:closure-minimality}
The Closure Principle and the minimal theory are mutually consistent: the closure fixed
point $N=96$ is the output of the minimal hierarchy, and the minimal hierarchy requires the
closure to generate the spectrum.
\end{theorem}

\begin{proof}
The closure fixed point $N=96$ is the output of the actualization process
(Corollary~\ref{cor:closure-n96}), which is a layer of the minimal hierarchy. Conversely,
the spectrum layer of the minimal hierarchy is the eigenspectrum of the converged network
(Chapter~\ref{chap:actualization}, Corollary~\ref{cor:spectrum-eigenspectrum}), so the
minimal hierarchy requires the closure to produce the spectrum. The two results are the
same structure viewed from the process side (closure) and the output side (spectrum).
\end{proof}

\begin{corollary}[Difference anchors; closure bounds; minimality completes]
\label{cor:anchor-bound-complete}
In the canonical architecture, Difference anchors the hierarchy (the root), the closure
bounds it (the fixed point), and the minimal theory completes it (the irreducible
derivation structure).
\end{corollary}

\begin{proof}
Difference anchors by the Difference Principle (Theorem~\ref{thm:difference-first},
Corollary~\ref{cor:difference-anchors}); the closure bounds by the Closure Principle
(Theorem~\ref{thm:closure-holds}, Corollary~\ref{cor:closure-n96}); the minimal theory
completes by Theorem~\ref{thm:minimal-hierarchy} and Theorem~\ref{thm:minimal-irreducible}.
The three roles are distinct and jointly exhaustive of the foundational structure.
\end{proof}

\section{Summary}
\label{sec:closure-summary}

This chapter has established the foundational closure of the theory. We have proved:

\begin{enumerate}
\item \textbf{The Closure Principle} (Theorem~\ref{thm:closure-holds},
Corollary~\ref{cor:closure-n96}): the boundary of the theory is the fixed point of
Difference's own dynamics, closing at $N=96$;
\item \textbf{Boundary and fixed point} (Theorem~\ref{thm:boundary-is-fixed},
Lemma~\ref{lem:boundary-characterizes}, Corollary~\ref{cor:closure-not-derive}): the
boundary is derived, characterizes the primitive, and --- decisively --- the Closure
Principle does \emph{not} derive Difference;
\item \textbf{Self-consistency} (Theorem~\ref{thm:self-consistency-presupposes},
Corollary~\ref{cor:conservation-universal}): self-consistency presupposes Difference, and
conservation is universal;
\item \textbf{Individuation} (Theorem~\ref{thm:individuation}): the distinct sectors derive
from Difference;
\item \textbf{The Difference Principle} (Theorem~\ref{thm:difference-first},
Corollary~\ref{cor:difference-anchors}): Difference is the self-referential first concept,
anchoring the hierarchy;
\item \textbf{The minimal theory} (Theorem~\ref{thm:minimal-hierarchy},
Theorem~\ref{thm:minimal-irreducible}, Corollary~\ref{cor:minimal-architecture}): the
minimal hierarchy is complete and irreducible;
\item \textbf{Consequences} (Lemma~\ref{lem:reduction-closes},
Theorem~\ref{thm:closure-minimality}, Corollary~\ref{cor:anchor-bound-complete}): the
reduction closes, closure and minimality are consistent, and Difference anchors, the
closure bounds, minimality completes.
\end{enumerate}

The foundation of The Actualization Theory is therefore complete: the primitive Difference,
anchored by the Difference Principle; the closure that bounds its dynamics at $N=96$; and
the minimal theory that compresses the entire derivation structure to the irreducible
hierarchy Difference → Actualization → Spectrum → Physics. The following chapter derives,
from this foundation, the inevitable spectrum itself.

\begin{thebibliography}{99}
\bibitem{QG267} TQM-QG Phase 267, \emph{Conservation Principle Audit}, Docs/Research.
\bibitem{QG268} TQM-QG Phase 268, \emph{Count Conservation Origin}, Docs/Research.
\bibitem{QG278} TQM-QG Phase 278, \emph{Fundamental Boundary Audit}, Docs/Research.
\bibitem{QG282} TQM-QG Phase 282, \emph{Boundary Origin Audit} (Closure Principle), Docs/Research.
\bibitem{QG292} TQM-QG Phase 292, \emph{Foundation Stress Test}, Docs/Research.
\bibitem{QG293} TQM-QG Phase 293, \emph{Hierarchy Necessity Audit}, Docs/Research.
\bibitem{QG294} TQM-QG Phase 294, \emph{Minimal Theory Audit}, Docs/Research.
\bibitem{MONO006} TQM-MONO006, \emph{A01 Actualization Origin Audit}, Docs/Zenodo/Monograph.
\bibitem{MONO007} TQM-MONO007, \emph{Referee-Readiness Resolution}, Docs/Zenodo/Monograph.
\end{thebibliography}
 after the preamble
%         that defines the theorem environments (or include via the master file).
% ======================================================================

% ---- Theorem environments (define in the master preamble if not present) ----
% \newtheorem{definition}{Definition}[chapter]
% \newtheorem{lemma}[definition]{Lemma}
% \newtheorem{theorem}[definition]{Theorem}
% \newtheorem{corollary}[definition]{Corollary}
% \newtheorem{remark}[definition]{Remark}

\chapter{Closure and the Minimal Theory}
\label{chap:closure-minimal}

\section{Introduction}
\label{sec:closure-introduction}

In the preceding chapters we established the primitives $\{\text{Difference},\eta\}$
(Chapters~\ref{chap:difference} and \ref{chap:eta}) and the derived process face of
Difference, Actualization, whose convergence yields the inevitable spectrum
(Chapter~\ref{chap:actualization}). This chapter completes the foundational picture by
establishing two structural results. First, the \emph{Closure Principle}: the boundary of the
theory is the fixed point of Difference's own dynamics, and --- decisively --- the Closure
Principle \emph{characterizes} the fundamental boundary but does \emph{not derive} the
primitive \cite{QG278,QG282}. Second, the \emph{minimal theory}: the hierarchy
\begin{equation}
\label{eq:minimal}
\text{Difference}\;\longrightarrow\;\text{Actualization}\;\longrightarrow\;\text{Inevitable Spectrum}\;\longrightarrow\;\text{Physics}
\end{equation}
is complete and irreducible \cite{QG293,QG294}: every intermediate layer is derivable, and
no further primitive is required.

The results of this chapter are the canonical end-state of the closure and reduction program
\cite{QG267,QG268,QG278,QG282,QG293,QG294}, consistent with the MONO006 resolution and the
MONO007 referee-ready architecture \cite{MONO006,MONO007}. Throughout we use only accepted
canonical results and introduce no new physics and no speculation.

\section{The Closure Principle}
\label{sec:closure-principle}

\begin{definition}[Closure Principle]
\label{def:closure-principle}
The \emph{Closure Principle} is the statement that the boundary of the theory is the fixed
point of the actualization process: the primitive is the process, and the boundary is its
converged fixed point \cite{QG282}.
\end{definition}

\begin{theorem}[The Closure Principle holds]
\label{thm:closure-holds}
The boundary of the theory is the fixed point of Difference's own count-producing dynamics:
the topology saturates at zero residual change, and every initial pattern converges to the
same final geometry.
\end{theorem}

\begin{proof}
The boundary-origin audit \cite{QG282} establishes the closure: the actualization dynamics
has a self-reinforcing feedback bounded by network capacity, so positive feedback saturates
at a stable fixed point with zero residual link growth and a content-independent final
geometry. The observable sector is therefore the converged attractor of the dynamics, and
the boundary is the closure of that dynamics. The closure is derived --- it is an output of
the process --- not an imported cut-off.
\end{proof}

\begin{corollary}[The closure is \texorpdfstring{$N=96$}{N=96}]
\label{cor:closure-n96}
The fixed point of the actualization process is the $N=96$ network, whose eigenspectrum is
the inevitable D96 spectrum.
\end{corollary}

\begin{proof}
By Theorem~\ref{thm:closure-holds}, the boundary is the converged attractor of the
dynamics; by Theorem~\ref{thm:n96} of Chapter~\ref{chap:actualization}, that attractor is
the $N=96$ network; and by Corollary~\ref{cor:spectrum-eigenspectrum} of
Chapter~\ref{chap:actualization}, its eigenspectrum is the inevitable D96 spectrum.
\end{proof}

\section{Boundary and Fixed Point}
\label{sec:boundary-fixed-point}

\begin{definition}[Boundary]
\label{def:boundary}
A \emph{boundary} of the theory is the configuration at which the derived structure closes:
the point where the count-producing dynamics has converged and no further change occurs.
\end{definition}

\begin{theorem}[The boundary is the fixed point]
\label{thm:boundary-is-fixed}
The boundary of the theory is the stable fixed point of the actualization process; it is
derived, not primitive.
\end{theorem}

\begin{proof}
By Definition~\ref{def:closure-principle} and Theorem~\ref{thm:closure-holds}, the boundary
is the converged fixed point of the process. The derivation structure of the process is
established by Theorem~\ref{thm:derived-not-primitive} of
Chapter~\ref{chap:actualization}: the process is the face of Difference, and its fixed
point is an output of the dynamics. Hence the boundary is derived --- it is not an imported
input or a separate primitive.
\end{proof}

\begin{lemma}[The boundary characterizes the primitive]
\label{lem:boundary-characterizes}
The fixed-point boundary is a characterization of the primitive Difference: it specifies
how Difference's dynamics closes, and thereby \emph{what} the boundary is, without defining
\emph{what Difference is}.
\end{lemma}

\begin{proof}
The boundary is the fixed point of the process that Difference generates
(Theorem~\ref{thm:boundary-is-fixed}). It tells us where the process converges --- the
closure $N=96$ --- but presupposes the process and hence the primitive. A characterization
specifies properties of an object already given; it does not supply the object itself. The
fixed-point boundary therefore characterizes Difference without deriving it.
\end{proof}

\begin{corollary}[Closure does not derive Difference]
\label{cor:closure-not-derive}
The Closure Principle does not derive the primitive Difference. It characterizes the
boundary of the theory as the fixed point of Difference's own dynamics.
\end{corollary}

\begin{proof}
By Lemma~\ref{lem:boundary-characterizes}, the closure is a characterization, not a
derivation. Moreover, by Theorem~\ref{thm:fundamental-boundary} of
Chapter~\ref{chap:difference}, Difference is the fundamental boundary: irreducible, with its
count-conservation law as its definitional identity. A derivation of the primitive would
presuppose a deeper notion, which the fundamental-boundary audit \cite{QG278} shows does
not exist --- every attempt to reduce Difference reintroduces it. Hence the Closure
Principle characterizes, and does not derive, the primitive.
\end{proof}

\begin{remark}[Referee-safe reading]
Consistent with MONO007 \cite{MONO007}, every statement in this chapter that invokes the
Closure Principle is worded so that the principle is never presented as a derivation of
Difference: the primitive is given, and the closure characterizes where its process closes.
This removes the circularity objection A02 of the hostile-referee audit.
\end{remark}

\section{Self-Consistency}
\label{sec:self-consistency}

\begin{definition}[Self-consistency]
\label{def:self-consistency}
\emph{Self-consistency} is the requirement that the theory be free of internal
contradiction: non-contradiction requires comparing parts --- a difference.
\end{definition}

\begin{theorem}[Self-consistency presupposes Difference]
\label{thm:self-consistency-presupposes}
Self-consistency is derivable from Difference: non-contradiction is the comparison of parts,
and identity is the zero difference.
\end{theorem}

\begin{proof}
The fundamental-boundary audit \cite{QG278} establishes that self-consistency presupposes
Difference: non-contradiction requires comparing parts of the theory --- a difference ---
and identity is the zero difference \cite{QG268}. The conservation of the count, which is
the definitional identity of the primitive, is the universal conservation principle of which
all conservation laws are manifestations \cite{QG267}. Hence self-consistency is not an
independent postulate but a consequence of the primitive's structure.
\end{proof}

\begin{corollary}[Conservation is universal]
\label{cor:conservation-universal}
All conservation laws of the theory --- norm, trace, count, Bianchi, Noether currents ---
are manifestations of one universal conservation principle rooted in the identity of
Difference.
\end{corollary}

\begin{proof}
The conservation-principle audit \cite{QG267} verifies that the six conservation laws (Born
norm, unitarity, trace conservation, count conservation, Bianchi conservation, Noether
currents) are all manifestations of one deeper principle. By
Theorem~\ref{thm:count-conservation} of Chapter~\ref{chap:actualization}, count
conservation is the definitional identity of the primitive. Hence the universal conservation
principle is the identity of Difference, and all conservation laws are its manifestations.
\end{proof}

\section{Individuation}
\label{sec:individuation}

\begin{definition}[Individuation]
\label{def:individuation}
\emph{Individuation} is the process by which the one Difference gives rise to distinct
sectors: the differentiation of the theory's structure into the separate domains of
observable physics.
\end{definition}

\begin{theorem}[Individuation derives from Difference]
\label{thm:individuation}
The individuation of the theory into distinct sectors is derivable from Difference: distinct
sectors are distinct differences.
\end{theorem}

\begin{proof}
Sectors are distinguished by their different spectral access --- the projection classes over
the operator layer established in Chapter~\ref{chap:actualization}
(Theorem~\ref{thm:resonance-readout}). Each sector is a distinct pattern of difference in
the count structure. The individuation of the theory is therefore the differentiation of the
one Difference into its distinct faces and readouts; it is not an independent primitive but
a derived structure of the single difference principle.
\end{proof}

\section{The Difference Principle}
\label{sec:difference-principle}

\begin{definition}[Difference Principle]
\label{def:difference-principle}
The \emph{Difference Principle} is the statement that Difference is the root of the
hierarchy: the first concept, presupposed by everything else, to which every reduction
attempt returns.
\end{definition}

\begin{theorem}[Difference is the first concept]
\label{thm:difference-first}
Difference is the first concept of the theory: every attempt to reduce it reintroduces it,
so it is a self-referential fundamental boundary.
\end{theorem}

\begin{proof}
The fundamental-boundary audit \cite{QG278} performs the self-referential check. Actualization
presupposes Difference (an act of actualizing is a change --- a before/after difference; a
Q-event is a transition). Question presupposes Difference (a question is a gap --- the
difference between known and unknown). Self-consistency presupposes Difference
(non-contradiction requires comparing parts; identity is the zero difference). And
Difference itself is the first concept: any attempt to reduce it reintroduces it --- "two
distinct things" uses distinct = difference, "a boundary" is where things differ, "a
transition" is a before/after difference, "empty vs non-empty" is itself a difference. Hence
Difference is a self-referential fundamental boundary, and the Difference Principle states
exactly this: Difference is the root.
\end{proof}

\begin{corollary}[The Difference Principle anchors the hierarchy]
\label{cor:difference-anchors}
The Difference Principle anchors the canonical hierarchy
\eqref{eq:minimal}: every layer --- Actualization, Spectrum, Physics --- presupposes
Difference as its root.
\end{corollary}

\begin{proof}
By Theorem~\ref{thm:difference-first}, every reduction returns to Difference; by
Theorem~\ref{thm:process-face} of Chapter~\ref{chap:actualization}, Actualization is the
process face of Difference; and the spectrum and physics are the outputs of that process
(Chapter~\ref{chap:actualization}, Theorem~\ref{thm:spectrum-inevitable}). Hence the
Difference Principle anchors the entire hierarchy at its root.
\end{proof}

\section{The Minimal Theory}
\label{sec:minimal-theory}

\begin{theorem}[The minimal hierarchy]
\label{thm:minimal-hierarchy}
The minimal hierarchy is
\[
\text{Difference}\;\longrightarrow\;\text{Actualization}\;\longrightarrow\;\text{Inevitable Spectrum}\;\longrightarrow\;\text{Physics}.
\]
It is complete: every phase of the theory can be rewritten using only this hierarchy.
\end{theorem}

\begin{proof}
The hierarchy-necessity audit \cite{QG293} removes each intermediate layer of the fuller
hierarchy (Difference → Actualization → Closure → Conservation → Resonance → Spectrum →
Question → Measurement → Physics) and classifies it. Actualization is INDISPENSABLE (the
count-producing process). Closure is COMPRESSIBLE into Actualization (the fixed point of the
process). Conservation is COMPRESSIBLE into Difference/network (the definitional identity
plus the graph handshake). Resonance is COMPRESSIBLE into Spectrum (the operator readout).
The minimal-theory audit \cite{QG294} then re-derives every compressed layer from the
minimal hierarchy, classifying each DERIVABLE or ACTUALLY REQUIRED, and finds
\emph{5 DERIVABLE / 0 ACTUALLY REQUIRED}. Hence the minimal hierarchy is complete.
\end{proof}

\begin{theorem}[The minimal theory is irreducible]
\label{thm:minimal-irreducible}
The minimal theory is irreducible: no layer of the hierarchy
\eqref{eq:minimal} can be removed without losing derived content, and no further primitive
exists.
\end{theorem}

\begin{proof}
Irreducibility has two parts. \emph{Primitives:} by Theorem~\ref{thm:minimal-foundation} of
Chapter~\ref{chap:eta}, the minimal primitive set is $\{\text{Difference},\eta\}$; by
Theorem~\ref{thm:difference-first}, no deeper notion exists; by
Theorem~\ref{thm:derived-not-primitive} of Chapter~\ref{chap:actualization}, Actualization
is derived, not primitive. \emph{Layers:} the foundation stress test \cite{QG292}
establishes that removing Difference collapses all layers, and the hierarchy-necessity audit
\cite{QG293} establishes that removing Actualization destroys the count-producing process.
The minimal-theory audit \cite{QG294} confirms that the four layers of \eqref{eq:minimal}
are exactly what is required, with zero additional inputs. Hence the minimal theory is
irreducible.
\end{proof}

\begin{corollary}[The canonical architecture]
\label{cor:minimal-architecture}
The canonical architecture of the theory is the minimal hierarchy
\eqref{eq:minimal}: Difference → Actualization → Inevitable Spectrum → Physics, with
primitives $\{\text{Difference},\eta\}$, Actualization derived, the spectrum inevitable, and
physics emergent.
\end{corollary}

\begin{proof}
The minimal hierarchy is complete (Theorem~\ref{thm:minimal-hierarchy}) and irreducible
(Theorem~\ref{thm:minimal-irreducible}); the primitives are the pair
$\{\text{Difference},\eta\}$ (Theorem~\ref{thm:minimal-foundation} of
Chapter~\ref{chap:eta}); Actualization is derived (Theorem~\ref{thm:derived-not-primitive}
of Chapter~\ref{chap:actualization}); the spectrum is inevitable
(Theorem~\ref{thm:spectrum-inevitable} of Chapter~\ref{chap:actualization}). This is the
referee-ready structure of MONO007 \cite{MONO007}.
\end{proof}

\begin{remark}[Referee-safe completeness]
Consistent with MONO007 \cite{MONO007}, where the monograph claims completeness of the
minimal theory, it is understood in the qualified sense: every \emph{observable} is derived
from the minimal hierarchy, while the standard-model \emph{Lagrangian} and its dynamics are
\emph{hosted} --- a boundary, not a derivation gap. The minimal theory is irreducible as a
derivation structure; it does not assert a derivation of every aspect of the hosted dynamics.
\end{remark}

\section{Consequences}
\label{sec:closure-consequences}

\begin{lemma}[The reduction closes at the minimal hierarchy]
\label{lem:reduction-closes}
The reduction program closes at the minimal hierarchy \eqref{eq:minimal}: all intermediate
layers of the fuller hierarchy are derivable from it, and none is actually required as an
input.
\end{lemma}

\begin{proof}
The minimal-theory audit \cite{QG294} classifies the five compressed layers --- Closure,
Conservation, Resonance, and the remaining intermediates --- as DERIVABLE, with zero
ACTUALLY REQUIRED. The re-derivation uses only the minimal hierarchy. Hence the reduction
closes at \eqref{eq:minimal}.
\end{proof}

\begin{theorem}[Closure and minimality are consistent]
\label{thm:closure-minimality}
The Closure Principle and the minimal theory are mutually consistent: the closure fixed
point $N=96$ is the output of the minimal hierarchy, and the minimal hierarchy requires the
closure to generate the spectrum.
\end{theorem}

\begin{proof}
The closure fixed point $N=96$ is the output of the actualization process
(Corollary~\ref{cor:closure-n96}), which is a layer of the minimal hierarchy. Conversely,
the spectrum layer of the minimal hierarchy is the eigenspectrum of the converged network
(Chapter~\ref{chap:actualization}, Corollary~\ref{cor:spectrum-eigenspectrum}), so the
minimal hierarchy requires the closure to produce the spectrum. The two results are the
same structure viewed from the process side (closure) and the output side (spectrum).
\end{proof}

\begin{corollary}[Difference anchors; closure bounds; minimality completes]
\label{cor:anchor-bound-complete}
In the canonical architecture, Difference anchors the hierarchy (the root), the closure
bounds it (the fixed point), and the minimal theory completes it (the irreducible
derivation structure).
\end{corollary}

\begin{proof}
Difference anchors by the Difference Principle (Theorem~\ref{thm:difference-first},
Corollary~\ref{cor:difference-anchors}); the closure bounds by the Closure Principle
(Theorem~\ref{thm:closure-holds}, Corollary~\ref{cor:closure-n96}); the minimal theory
completes by Theorem~\ref{thm:minimal-hierarchy} and Theorem~\ref{thm:minimal-irreducible}.
The three roles are distinct and jointly exhaustive of the foundational structure.
\end{proof}

\section{Summary}
\label{sec:closure-summary}

This chapter has established the foundational closure of the theory. We have proved:

\begin{enumerate}
\item \textbf{The Closure Principle} (Theorem~\ref{thm:closure-holds},
Corollary~\ref{cor:closure-n96}): the boundary of the theory is the fixed point of
Difference's own dynamics, closing at $N=96$;
\item \textbf{Boundary and fixed point} (Theorem~\ref{thm:boundary-is-fixed},
Lemma~\ref{lem:boundary-characterizes}, Corollary~\ref{cor:closure-not-derive}): the
boundary is derived, characterizes the primitive, and --- decisively --- the Closure
Principle does \emph{not} derive Difference;
\item \textbf{Self-consistency} (Theorem~\ref{thm:self-consistency-presupposes},
Corollary~\ref{cor:conservation-universal}): self-consistency presupposes Difference, and
conservation is universal;
\item \textbf{Individuation} (Theorem~\ref{thm:individuation}): the distinct sectors derive
from Difference;
\item \textbf{The Difference Principle} (Theorem~\ref{thm:difference-first},
Corollary~\ref{cor:difference-anchors}): Difference is the self-referential first concept,
anchoring the hierarchy;
\item \textbf{The minimal theory} (Theorem~\ref{thm:minimal-hierarchy},
Theorem~\ref{thm:minimal-irreducible}, Corollary~\ref{cor:minimal-architecture}): the
minimal hierarchy is complete and irreducible;
\item \textbf{Consequences} (Lemma~\ref{lem:reduction-closes},
Theorem~\ref{thm:closure-minimality}, Corollary~\ref{cor:anchor-bound-complete}): the
reduction closes, closure and minimality are consistent, and Difference anchors, the
closure bounds, minimality completes.
\end{enumerate}

The foundation of The Actualization Theory is therefore complete: the primitive Difference,
anchored by the Difference Principle; the closure that bounds its dynamics at $N=96$; and
the minimal theory that compresses the entire derivation structure to the irreducible
hierarchy Difference → Actualization → Spectrum → Physics. The following chapter derives,
from this foundation, the inevitable spectrum itself.

\begin{thebibliography}{99}
\bibitem{QG267} TQM-QG Phase 267, \emph{Conservation Principle Audit}, Docs/Research.
\bibitem{QG268} TQM-QG Phase 268, \emph{Count Conservation Origin}, Docs/Research.
\bibitem{QG278} TQM-QG Phase 278, \emph{Fundamental Boundary Audit}, Docs/Research.
\bibitem{QG282} TQM-QG Phase 282, \emph{Boundary Origin Audit} (Closure Principle), Docs/Research.
\bibitem{QG292} TQM-QG Phase 292, \emph{Foundation Stress Test}, Docs/Research.
\bibitem{QG293} TQM-QG Phase 293, \emph{Hierarchy Necessity Audit}, Docs/Research.
\bibitem{QG294} TQM-QG Phase 294, \emph{Minimal Theory Audit}, Docs/Research.
\bibitem{MONO006} TQM-MONO006, \emph{A01 Actualization Origin Audit}, Docs/Zenodo/Monograph.
\bibitem{MONO007} TQM-MONO007, \emph{Referee-Readiness Resolution}, Docs/Zenodo/Monograph.
\end{thebibliography}
 after the preamble
%         that defines the theorem environments (or include via the master file).
% ======================================================================

% ---- Theorem environments (define in the master preamble if not present) ----
% \newtheorem{definition}{Definition}[chapter]
% \newtheorem{lemma}[definition]{Lemma}
% \newtheorem{theorem}[definition]{Theorem}
% \newtheorem{corollary}[definition]{Corollary}
% \newtheorem{remark}[definition]{Remark}

\chapter{Closure and the Minimal Theory}
\label{chap:closure-minimal}

\section{Introduction}
\label{sec:closure-introduction}

In the preceding chapters we established the primitives $\{\text{Difference},\eta\}$
(Chapters~\ref{chap:difference} and \ref{chap:eta}) and the derived process face of
Difference, Actualization, whose convergence yields the inevitable spectrum
(Chapter~\ref{chap:actualization}). This chapter completes the foundational picture by
establishing two structural results. First, the \emph{Closure Principle}: the boundary of the
theory is the fixed point of Difference's own dynamics, and --- decisively --- the Closure
Principle \emph{characterizes} the fundamental boundary but does \emph{not derive} the
primitive \cite{QG278,QG282}. Second, the \emph{minimal theory}: the hierarchy
\begin{equation}
\label{eq:minimal}
\text{Difference}\;\longrightarrow\;\text{Actualization}\;\longrightarrow\;\text{Inevitable Spectrum}\;\longrightarrow\;\text{Physics}
\end{equation}
is complete and irreducible \cite{QG293,QG294}: every intermediate layer is derivable, and
no further primitive is required.

The results of this chapter are the canonical end-state of the closure and reduction program
\cite{QG267,QG268,QG278,QG282,QG293,QG294}, consistent with the MONO006 resolution and the
MONO007 referee-ready architecture \cite{MONO006,MONO007}. Throughout we use only accepted
canonical results and introduce no new physics and no speculation.

\section{The Closure Principle}
\label{sec:closure-principle}

\begin{definition}[Closure Principle]
\label{def:closure-principle}
The \emph{Closure Principle} is the statement that the boundary of the theory is the fixed
point of the actualization process: the primitive is the process, and the boundary is its
converged fixed point \cite{QG282}.
\end{definition}

\begin{theorem}[The Closure Principle holds]
\label{thm:closure-holds}
The boundary of the theory is the fixed point of Difference's own count-producing dynamics:
the topology saturates at zero residual change, and every initial pattern converges to the
same final geometry.
\end{theorem}

\begin{proof}
The boundary-origin audit \cite{QG282} establishes the closure: the actualization dynamics
has a self-reinforcing feedback bounded by network capacity, so positive feedback saturates
at a stable fixed point with zero residual link growth and a content-independent final
geometry. The observable sector is therefore the converged attractor of the dynamics, and
the boundary is the closure of that dynamics. The closure is derived --- it is an output of
the process --- not an imported cut-off.
\end{proof}

\begin{corollary}[The closure is \texorpdfstring{$N=96$}{N=96}]
\label{cor:closure-n96}
The fixed point of the actualization process is the $N=96$ network, whose eigenspectrum is
the inevitable D96 spectrum.
\end{corollary}

\begin{proof}
By Theorem~\ref{thm:closure-holds}, the boundary is the converged attractor of the
dynamics; by Theorem~\ref{thm:n96} of Chapter~\ref{chap:actualization}, that attractor is
the $N=96$ network; and by Corollary~\ref{cor:spectrum-eigenspectrum} of
Chapter~\ref{chap:actualization}, its eigenspectrum is the inevitable D96 spectrum.
\end{proof}

\section{Boundary and Fixed Point}
\label{sec:boundary-fixed-point}

\begin{definition}[Boundary]
\label{def:boundary}
A \emph{boundary} of the theory is the configuration at which the derived structure closes:
the point where the count-producing dynamics has converged and no further change occurs.
\end{definition}

\begin{theorem}[The boundary is the fixed point]
\label{thm:boundary-is-fixed}
The boundary of the theory is the stable fixed point of the actualization process; it is
derived, not primitive.
\end{theorem}

\begin{proof}
By Definition~\ref{def:closure-principle} and Theorem~\ref{thm:closure-holds}, the boundary
is the converged fixed point of the process. The derivation structure of the process is
established by Theorem~\ref{thm:derived-not-primitive} of
Chapter~\ref{chap:actualization}: the process is the face of Difference, and its fixed
point is an output of the dynamics. Hence the boundary is derived --- it is not an imported
input or a separate primitive.
\end{proof}

\begin{lemma}[The boundary characterizes the primitive]
\label{lem:boundary-characterizes}
The fixed-point boundary is a characterization of the primitive Difference: it specifies
how Difference's dynamics closes, and thereby \emph{what} the boundary is, without defining
\emph{what Difference is}.
\end{lemma}

\begin{proof}
The boundary is the fixed point of the process that Difference generates
(Theorem~\ref{thm:boundary-is-fixed}). It tells us where the process converges --- the
closure $N=96$ --- but presupposes the process and hence the primitive. A characterization
specifies properties of an object already given; it does not supply the object itself. The
fixed-point boundary therefore characterizes Difference without deriving it.
\end{proof}

\begin{corollary}[Closure does not derive Difference]
\label{cor:closure-not-derive}
The Closure Principle does not derive the primitive Difference. It characterizes the
boundary of the theory as the fixed point of Difference's own dynamics.
\end{corollary}

\begin{proof}
By Lemma~\ref{lem:boundary-characterizes}, the closure is a characterization, not a
derivation. Moreover, by Theorem~\ref{thm:fundamental-boundary} of
Chapter~\ref{chap:difference}, Difference is the fundamental boundary: irreducible, with its
count-conservation law as its definitional identity. A derivation of the primitive would
presuppose a deeper notion, which the fundamental-boundary audit \cite{QG278} shows does
not exist --- every attempt to reduce Difference reintroduces it. Hence the Closure
Principle characterizes, and does not derive, the primitive.
\end{proof}

\begin{remark}[Referee-safe reading]
Consistent with MONO007 \cite{MONO007}, every statement in this chapter that invokes the
Closure Principle is worded so that the principle is never presented as a derivation of
Difference: the primitive is given, and the closure characterizes where its process closes.
This removes the circularity objection A02 of the hostile-referee audit.
\end{remark}

\section{Self-Consistency}
\label{sec:self-consistency}

\begin{definition}[Self-consistency]
\label{def:self-consistency}
\emph{Self-consistency} is the requirement that the theory be free of internal
contradiction: non-contradiction requires comparing parts --- a difference.
\end{definition}

\begin{theorem}[Self-consistency presupposes Difference]
\label{thm:self-consistency-presupposes}
Self-consistency is derivable from Difference: non-contradiction is the comparison of parts,
and identity is the zero difference.
\end{theorem}

\begin{proof}
The fundamental-boundary audit \cite{QG278} establishes that self-consistency presupposes
Difference: non-contradiction requires comparing parts of the theory --- a difference ---
and identity is the zero difference \cite{QG268}. The conservation of the count, which is
the definitional identity of the primitive, is the universal conservation principle of which
all conservation laws are manifestations \cite{QG267}. Hence self-consistency is not an
independent postulate but a consequence of the primitive's structure.
\end{proof}

\begin{corollary}[Conservation is universal]
\label{cor:conservation-universal}
All conservation laws of the theory --- norm, trace, count, Bianchi, Noether currents ---
are manifestations of one universal conservation principle rooted in the identity of
Difference.
\end{corollary}

\begin{proof}
The conservation-principle audit \cite{QG267} verifies that the six conservation laws (Born
norm, unitarity, trace conservation, count conservation, Bianchi conservation, Noether
currents) are all manifestations of one deeper principle. By
Theorem~\ref{thm:count-conservation} of Chapter~\ref{chap:actualization}, count
conservation is the definitional identity of the primitive. Hence the universal conservation
principle is the identity of Difference, and all conservation laws are its manifestations.
\end{proof}

\section{Individuation}
\label{sec:individuation}

\begin{definition}[Individuation]
\label{def:individuation}
\emph{Individuation} is the process by which the one Difference gives rise to distinct
sectors: the differentiation of the theory's structure into the separate domains of
observable physics.
\end{definition}

\begin{theorem}[Individuation derives from Difference]
\label{thm:individuation}
The individuation of the theory into distinct sectors is derivable from Difference: distinct
sectors are distinct differences.
\end{theorem}

\begin{proof}
Sectors are distinguished by their different spectral access --- the projection classes over
the operator layer established in Chapter~\ref{chap:actualization}
(Theorem~\ref{thm:resonance-readout}). Each sector is a distinct pattern of difference in
the count structure. The individuation of the theory is therefore the differentiation of the
one Difference into its distinct faces and readouts; it is not an independent primitive but
a derived structure of the single difference principle.
\end{proof}

\section{The Difference Principle}
\label{sec:difference-principle}

\begin{definition}[Difference Principle]
\label{def:difference-principle}
The \emph{Difference Principle} is the statement that Difference is the root of the
hierarchy: the first concept, presupposed by everything else, to which every reduction
attempt returns.
\end{definition}

\begin{theorem}[Difference is the first concept]
\label{thm:difference-first}
Difference is the first concept of the theory: every attempt to reduce it reintroduces it,
so it is a self-referential fundamental boundary.
\end{theorem}

\begin{proof}
The fundamental-boundary audit \cite{QG278} performs the self-referential check. Actualization
presupposes Difference (an act of actualizing is a change --- a before/after difference; a
Q-event is a transition). Question presupposes Difference (a question is a gap --- the
difference between known and unknown). Self-consistency presupposes Difference
(non-contradiction requires comparing parts; identity is the zero difference). And
Difference itself is the first concept: any attempt to reduce it reintroduces it --- "two
distinct things" uses distinct = difference, "a boundary" is where things differ, "a
transition" is a before/after difference, "empty vs non-empty" is itself a difference. Hence
Difference is a self-referential fundamental boundary, and the Difference Principle states
exactly this: Difference is the root.
\end{proof}

\begin{corollary}[The Difference Principle anchors the hierarchy]
\label{cor:difference-anchors}
The Difference Principle anchors the canonical hierarchy
\eqref{eq:minimal}: every layer --- Actualization, Spectrum, Physics --- presupposes
Difference as its root.
\end{corollary}

\begin{proof}
By Theorem~\ref{thm:difference-first}, every reduction returns to Difference; by
Theorem~\ref{thm:process-face} of Chapter~\ref{chap:actualization}, Actualization is the
process face of Difference; and the spectrum and physics are the outputs of that process
(Chapter~\ref{chap:actualization}, Theorem~\ref{thm:spectrum-inevitable}). Hence the
Difference Principle anchors the entire hierarchy at its root.
\end{proof}

\section{The Minimal Theory}
\label{sec:minimal-theory}

\begin{theorem}[The minimal hierarchy]
\label{thm:minimal-hierarchy}
The minimal hierarchy is
\[
\text{Difference}\;\longrightarrow\;\text{Actualization}\;\longrightarrow\;\text{Inevitable Spectrum}\;\longrightarrow\;\text{Physics}.
\]
It is complete: every phase of the theory can be rewritten using only this hierarchy.
\end{theorem}

\begin{proof}
The hierarchy-necessity audit \cite{QG293} removes each intermediate layer of the fuller
hierarchy (Difference → Actualization → Closure → Conservation → Resonance → Spectrum →
Question → Measurement → Physics) and classifies it. Actualization is INDISPENSABLE (the
count-producing process). Closure is COMPRESSIBLE into Actualization (the fixed point of the
process). Conservation is COMPRESSIBLE into Difference/network (the definitional identity
plus the graph handshake). Resonance is COMPRESSIBLE into Spectrum (the operator readout).
The minimal-theory audit \cite{QG294} then re-derives every compressed layer from the
minimal hierarchy, classifying each DERIVABLE or ACTUALLY REQUIRED, and finds
\emph{5 DERIVABLE / 0 ACTUALLY REQUIRED}. Hence the minimal hierarchy is complete.
\end{proof}

\begin{theorem}[The minimal theory is irreducible]
\label{thm:minimal-irreducible}
The minimal theory is irreducible: no layer of the hierarchy
\eqref{eq:minimal} can be removed without losing derived content, and no further primitive
exists.
\end{theorem}

\begin{proof}
Irreducibility has two parts. \emph{Primitives:} by Theorem~\ref{thm:minimal-foundation} of
Chapter~\ref{chap:eta}, the minimal primitive set is $\{\text{Difference},\eta\}$; by
Theorem~\ref{thm:difference-first}, no deeper notion exists; by
Theorem~\ref{thm:derived-not-primitive} of Chapter~\ref{chap:actualization}, Actualization
is derived, not primitive. \emph{Layers:} the foundation stress test \cite{QG292}
establishes that removing Difference collapses all layers, and the hierarchy-necessity audit
\cite{QG293} establishes that removing Actualization destroys the count-producing process.
The minimal-theory audit \cite{QG294} confirms that the four layers of \eqref{eq:minimal}
are exactly what is required, with zero additional inputs. Hence the minimal theory is
irreducible.
\end{proof}

\begin{corollary}[The canonical architecture]
\label{cor:minimal-architecture}
The canonical architecture of the theory is the minimal hierarchy
\eqref{eq:minimal}: Difference → Actualization → Inevitable Spectrum → Physics, with
primitives $\{\text{Difference},\eta\}$, Actualization derived, the spectrum inevitable, and
physics emergent.
\end{corollary}

\begin{proof}
The minimal hierarchy is complete (Theorem~\ref{thm:minimal-hierarchy}) and irreducible
(Theorem~\ref{thm:minimal-irreducible}); the primitives are the pair
$\{\text{Difference},\eta\}$ (Theorem~\ref{thm:minimal-foundation} of
Chapter~\ref{chap:eta}); Actualization is derived (Theorem~\ref{thm:derived-not-primitive}
of Chapter~\ref{chap:actualization}); the spectrum is inevitable
(Theorem~\ref{thm:spectrum-inevitable} of Chapter~\ref{chap:actualization}). This is the
referee-ready structure of MONO007 \cite{MONO007}.
\end{proof}

\begin{remark}[Referee-safe completeness]
Consistent with MONO007 \cite{MONO007}, where the monograph claims completeness of the
minimal theory, it is understood in the qualified sense: every \emph{observable} is derived
from the minimal hierarchy, while the standard-model \emph{Lagrangian} and its dynamics are
\emph{hosted} --- a boundary, not a derivation gap. The minimal theory is irreducible as a
derivation structure; it does not assert a derivation of every aspect of the hosted dynamics.
\end{remark}

\section{Consequences}
\label{sec:closure-consequences}

\begin{lemma}[The reduction closes at the minimal hierarchy]
\label{lem:reduction-closes}
The reduction program closes at the minimal hierarchy \eqref{eq:minimal}: all intermediate
layers of the fuller hierarchy are derivable from it, and none is actually required as an
input.
\end{lemma}

\begin{proof}
The minimal-theory audit \cite{QG294} classifies the five compressed layers --- Closure,
Conservation, Resonance, and the remaining intermediates --- as DERIVABLE, with zero
ACTUALLY REQUIRED. The re-derivation uses only the minimal hierarchy. Hence the reduction
closes at \eqref{eq:minimal}.
\end{proof}

\begin{theorem}[Closure and minimality are consistent]
\label{thm:closure-minimality}
The Closure Principle and the minimal theory are mutually consistent: the closure fixed
point $N=96$ is the output of the minimal hierarchy, and the minimal hierarchy requires the
closure to generate the spectrum.
\end{theorem}

\begin{proof}
The closure fixed point $N=96$ is the output of the actualization process
(Corollary~\ref{cor:closure-n96}), which is a layer of the minimal hierarchy. Conversely,
the spectrum layer of the minimal hierarchy is the eigenspectrum of the converged network
(Chapter~\ref{chap:actualization}, Corollary~\ref{cor:spectrum-eigenspectrum}), so the
minimal hierarchy requires the closure to produce the spectrum. The two results are the
same structure viewed from the process side (closure) and the output side (spectrum).
\end{proof}

\begin{corollary}[Difference anchors; closure bounds; minimality completes]
\label{cor:anchor-bound-complete}
In the canonical architecture, Difference anchors the hierarchy (the root), the closure
bounds it (the fixed point), and the minimal theory completes it (the irreducible
derivation structure).
\end{corollary}

\begin{proof}
Difference anchors by the Difference Principle (Theorem~\ref{thm:difference-first},
Corollary~\ref{cor:difference-anchors}); the closure bounds by the Closure Principle
(Theorem~\ref{thm:closure-holds}, Corollary~\ref{cor:closure-n96}); the minimal theory
completes by Theorem~\ref{thm:minimal-hierarchy} and Theorem~\ref{thm:minimal-irreducible}.
The three roles are distinct and jointly exhaustive of the foundational structure.
\end{proof}

\section{Summary}
\label{sec:closure-summary}

This chapter has established the foundational closure of the theory. We have proved:

\begin{enumerate}
\item \textbf{The Closure Principle} (Theorem~\ref{thm:closure-holds},
Corollary~\ref{cor:closure-n96}): the boundary of the theory is the fixed point of
Difference's own dynamics, closing at $N=96$;
\item \textbf{Boundary and fixed point} (Theorem~\ref{thm:boundary-is-fixed},
Lemma~\ref{lem:boundary-characterizes}, Corollary~\ref{cor:closure-not-derive}): the
boundary is derived, characterizes the primitive, and --- decisively --- the Closure
Principle does \emph{not} derive Difference;
\item \textbf{Self-consistency} (Theorem~\ref{thm:self-consistency-presupposes},
Corollary~\ref{cor:conservation-universal}): self-consistency presupposes Difference, and
conservation is universal;
\item \textbf{Individuation} (Theorem~\ref{thm:individuation}): the distinct sectors derive
from Difference;
\item \textbf{The Difference Principle} (Theorem~\ref{thm:difference-first},
Corollary~\ref{cor:difference-anchors}): Difference is the self-referential first concept,
anchoring the hierarchy;
\item \textbf{The minimal theory} (Theorem~\ref{thm:minimal-hierarchy},
Theorem~\ref{thm:minimal-irreducible}, Corollary~\ref{cor:minimal-architecture}): the
minimal hierarchy is complete and irreducible;
\item \textbf{Consequences} (Lemma~\ref{lem:reduction-closes},
Theorem~\ref{thm:closure-minimality}, Corollary~\ref{cor:anchor-bound-complete}): the
reduction closes, closure and minimality are consistent, and Difference anchors, the
closure bounds, minimality completes.
\end{enumerate}

The foundation of The Actualization Theory is therefore complete: the primitive Difference,
anchored by the Difference Principle; the closure that bounds its dynamics at $N=96$; and
the minimal theory that compresses the entire derivation structure to the irreducible
hierarchy Difference → Actualization → Spectrum → Physics. The following chapter derives,
from this foundation, the inevitable spectrum itself.

\begin{thebibliography}{99}
\bibitem{QG267} TQM-QG Phase 267, \emph{Conservation Principle Audit}, Docs/Research.
\bibitem{QG268} TQM-QG Phase 268, \emph{Count Conservation Origin}, Docs/Research.
\bibitem{QG278} TQM-QG Phase 278, \emph{Fundamental Boundary Audit}, Docs/Research.
\bibitem{QG282} TQM-QG Phase 282, \emph{Boundary Origin Audit} (Closure Principle), Docs/Research.
\bibitem{QG292} TQM-QG Phase 292, \emph{Foundation Stress Test}, Docs/Research.
\bibitem{QG293} TQM-QG Phase 293, \emph{Hierarchy Necessity Audit}, Docs/Research.
\bibitem{QG294} TQM-QG Phase 294, \emph{Minimal Theory Audit}, Docs/Research.
\bibitem{MONO006} TQM-MONO006, \emph{A01 Actualization Origin Audit}, Docs/Zenodo/Monograph.
\bibitem{MONO007} TQM-MONO007, \emph{Referee-Readiness Resolution}, Docs/Zenodo/Monograph.
\end{thebibliography}

% ======================================================================
%  AT-MONO104 — Chapter 5: The Inevitable Spectrum
%  Canonical Monograph (v2.0) · The Actualization Theory:
%  A Reconstruction of Physics from Difference, Actualization and Spectrum
%
%  Status: COMPLETE — PASS-READY (MONO007 referee-ready architecture)
%  Inputs: Chapter01_Difference.tex, Chapter02_Eta.tex,
%          Chapter03_Actualization.tex, Chapter04_ClosureAndMinimalTheory.tex
%  Method: assembly only — canonical results only, no new physics, no speculation
%  Citations: QG116 (actualization structures), QG159 (D96 selection),
%             QG160 (period-3 seed), QG282 (closure principle),
%             QG293 (hierarchy necessity), QG294 (minimal theory),
%             QG295 (spectrum necessity), QG296 (reconstruction)
%
%  Usage: % ======================================================================
%  AT-MONO104 — Chapter 5: The Inevitable Spectrum
%  Canonical Monograph (v2.0) · The Actualization Theory:
%  A Reconstruction of Physics from Difference, Actualization and Spectrum
%
%  Status: COMPLETE — PASS-READY (MONO007 referee-ready architecture)
%  Inputs: Chapter01_Difference.tex, Chapter02_Eta.tex,
%          Chapter03_Actualization.tex, Chapter04_ClosureAndMinimalTheory.tex
%  Method: assembly only — canonical results only, no new physics, no speculation
%  Citations: QG116 (actualization structures), QG159 (D96 selection),
%             QG160 (period-3 seed), QG282 (closure principle),
%             QG293 (hierarchy necessity), QG294 (minimal theory),
%             QG295 (spectrum necessity), QG296 (reconstruction)
%
%  Usage: % ======================================================================
%  AT-MONO104 — Chapter 5: The Inevitable Spectrum
%  Canonical Monograph (v2.0) · The Actualization Theory:
%  A Reconstruction of Physics from Difference, Actualization and Spectrum
%
%  Status: COMPLETE — PASS-READY (MONO007 referee-ready architecture)
%  Inputs: Chapter01_Difference.tex, Chapter02_Eta.tex,
%          Chapter03_Actualization.tex, Chapter04_ClosureAndMinimalTheory.tex
%  Method: assembly only — canonical results only, no new physics, no speculation
%  Citations: QG116 (actualization structures), QG159 (D96 selection),
%             QG160 (period-3 seed), QG282 (closure principle),
%             QG293 (hierarchy necessity), QG294 (minimal theory),
%             QG295 (spectrum necessity), QG296 (reconstruction)
%
%  Usage: \input{Chapter05_InevitableSpectrum.tex} after the preamble
%         that defines the theorem environments (or include via the master file).
% ======================================================================

% ---- Theorem environments (define in the master preamble if not present) ----
% \newtheorem{definition}{Definition}[chapter]
% \newtheorem{lemma}[definition]{Lemma}
% \newtheorem{theorem}[definition]{Theorem}
% \newtheorem{corollary}[definition]{Corollary}
% \newtheorem{remark}[definition]{Remark}

\chapter{The Inevitable Spectrum}
\label{c05:chap:spectrum}

\section{Introduction}
\label{c05:sec:spectrum-introduction}

The canonical hierarchy of The Actualization Theory
\begin{equation}
\label{c05:eq:core}
\text{Difference}\;\longrightarrow\;\text{Actualization}\;\longrightarrow\;\text{Inevitable Spectrum}\;\longrightarrow\;\text{Physics}
\end{equation}
is complete as a derivation structure (Chapter~\ref{c04:chap:closure-minimal},
Theorem~\ref{c04:thm:minimal-hierarchy}). This chapter develops the third member of the
hierarchy: the \emph{inevitable spectrum}. We establish that the spectrum is not a primitive
of the theory but the \emph{derived output} of the actualization process: the dynamics
converge to a content-independent attractor of size $N=96$ within the accepted structural
class; the attractor is a stable fixed point, not a freely chosen input; and the spectrum is
the eigenspectrum of that attractor, carrying no independent choice
\cite{c05:QG116,c05:QG282,c05:QG295}.

We then establish the \emph{necessity} of the spectrum and the tested uniqueness of the
attractor within the canonical structural class: the attractor is content-independent
\cite{c05:QG116}, the size $N=96$ is inevitable within the accepted D96 selection
constraints --- not a choice \cite{c05:QG159,c05:QG160} --- and the spectrum is a
deterministic function of the converged network. This completes the canonical core
Difference → Actualization → Spectrum, whose first three members are now established as
primitive, derived, and inevitable output respectively; the physics chapters that follow
read all observable sectors from this one spectrum \cite{c05:QG293,c05:QG294,c05:QG296}.

\section{The Actualization Attractor}
\label{c05:sec:attractor}

\begin{theorem}[The actualization attractor: convergence, uniqueness, closure]
\label{c05:thm:attractor-unique}
The \emph{actualization attractor} is the converged state of the actualization process: the
configuration to which every initial pattern of Difference's count-producing dynamics
converges, with zero residual change. The actualization dynamics converge --- residual link
growth reaches zero and the topology saturates at a stable final state --- the attractor is
unique and content-independent (every initial activity pattern converges to the same final
geometry), and it is the closure fixed point of the theory.
\end{theorem}

\begin{proof}
The actualization-structures program \cite{c05:QG116} establishes that the actualization
dynamics converge with zero residual link growth: the self-reinforcing feedback (activity to
links to activity) is bounded by network capacity, so the topology saturates, and identical
link counts and identical span are reached from every initial pattern. Hence the attractor
is unique --- one final geometry, not a family --- and independent of the starting
configuration. By the Closure Principle (Chapter~\ref{c04:chap:closure-minimal},
Theorem~\ref{c04:thm:closure-holds}, Definition~\ref{c04:def:closure-principle}), the boundary
of the process is its stable fixed point; hence the unique attractor and the closure fixed
point coincide.
\end{proof}

\section{The Fixed Point \texorpdfstring{$N=96$}{N=96}}
\label{c05:sec:n96}

\begin{theorem}[The attractor is \texorpdfstring{$N=96$}{N=96}, not a choice; no freedom in the size]
\label{c05:thm:n96-attractor}
The \emph{$N=96$ attractor} is the converged network of the actualization process within the
canonical structural class: the stable fixed point whose size is $N=96$. The size is the
attractor of the dynamics, not a freely chosen input: the process converges to $N=96$ from
every initial pattern in the accepted comparison class, and no other tested size is a stable
fixed point of the dynamics.
\end{theorem}

\begin{proof}
The spectrum-necessity audit \cite{c05:QG295} establishes that $N=96$ is the attractor, not a
choice: the actualization flow has a stable fixed point at $N=96$, and the boundary is the
closure of that flow (the Closure Principle \cite{c05:QG282}). The D96-selection program adds
that $96$ is selected as the unique complete-Z2 natural size with the three-family octave
window; no other natural size satisfies these conditions \cite{c05:QG159,c05:QG160}. Hence
$N=96$ is the stable fixed point of the flow within the canonical comparison class,
content-independent (Theorem~\ref{c05:thm:attractor-unique}): an output of the dynamics, not
an input or free parameter.
\end{proof}

\subsection{Why \texorpdfstring{$N=96$}{N=96}?}
\label{c05:sec:why-n96}

\begin{table}[h]
\centering
\begin{tabular}{p{0.18\linewidth}p{0.28\linewidth}p{0.46\linewidth}}
\hline
\textbf{Argument class} & \textbf{What it uses} & \textbf{Status} \\
\hline
Selection & period-3 seed + Z2 half-shift; 3-family octave window; natural doubling chain $n=3\cdot2^k$; tested candidate set \cite{c05:QG159,c05:QG160} & D96 is the only rung in the tested structural class that satisfies all conditions \\
Necessity & closure of the actualization flow; stable fixed point; spectrum as eigenspectrum of the converged network \cite{c05:QG282,c05:QG295} & If the flow does not close, there is no spectrum layer and no physics readout \\
Uniqueness & content-independent attractor; same final geometry from every initial pattern \cite{c05:QG116} & Unique within the canonical comparison class examined \\
\hline
\end{tabular}
\end{table}

\paragraph{Why not 94, 95, 97?}
On the accepted selection criterion, these sizes fail the seed half-shift condition: the
period-3 seed requires $6 \mid n$ for the Z2 doublet structure, so $94,95,97$ are excluded
before the octave-window test is even applied. They are therefore not members of the
canonical candidate class used by QG159.

\paragraph{Why not 128?}
128 fails the same seed symmetry test ($128 \not\equiv 0 \pmod 6$) and, in the selection
audit, also fails the three-family window: its span is above the $[4,8)$ range, giving four
families rather than three \cite{c05:QG159}.

\paragraph{Exact status.}
The strongest referee-safe statement is: \emph{within the canonical structural constraints
actually used by AT --- period-3 seed, Z2 half-shift, three-family octave window, and the
tested comparison class --- $N=96$ is inevitable and uniquely selected.} A global proof that
no other closure could exist outside that structural class is \emph{not} claimed here.
That is the boundary.

\section{Spectrum Generation}
\label{c05:sec:spectrum-generation}

\begin{theorem}[The spectrum is the eigenspectrum of the attractor]
\label{c05:thm:spectrum-eigenspectrum}
\emph{Spectrum generation} is the derivation of the spectrum from the converged network: the
spectrum of the theory is the Laplacian eigenspectrum of the $N=96$ attractor --- $95$
positive modes with a single zero mode (the background) --- a deterministic function of the
attractor. Its spectral constants --- the moments $\Sigma m=95$,
$\Sigma\sqrt{m}=64.08$, $\Sigma m^2=229$, the occupancy moment $1900.25$, the span $6.40$,
and the doublet multiplicities $[42\times2,5,6]$ --- are outputs of the spectrum, hence of
the attractor.
\end{theorem}

\begin{proof}
The spectrum-necessity audit \cite{c05:QG295} establishes that the spectrum is the Laplacian
eigenspectrum of the converged network: the graph Laplacian of the $N=96$ attractor has
$95$ positive eigenvalues (the modes) and one zero eigenvalue (the background). Since the
attractor is unique and deterministic (Theorem~\ref{c05:thm:attractor-unique},
Theorem~\ref{c05:thm:n96-attractor}), the spectrum is a deterministic function of it --- an
output of the dynamics, not an additional input --- and its moments and structural features
are derived quantities, not primitive inputs. This is the canonical reading of the spectrum
layer of the minimal hierarchy \cite{c05:QG294}.
\end{proof}

\section{Uniqueness of the Spectrum}
\label{c05:sec:spectrum-uniqueness}

Uniqueness here means uniqueness \emph{within the canonical structural class}: the spectrum is
the eigenspectrum of the attractor (Theorem~\ref{c05:thm:attractor-unique},
Theorem~\ref{c05:thm:spectrum-eigenspectrum}), whose size is fixed at $N=96$
(Theorem~\ref{c05:thm:n96-attractor}), and the eigenspectrum of a given network is unique.
Hence there is exactly one spectrum of the theory in the accepted class --- one attractor,
one eigenspectrum --- and since the attractor carries no freedom, no spectral feature is
freely adjustable: the spectrum carries no choice at any level.

\section{Spectrum Necessity}
\label{c05:sec:spectrum-necessity}

\begin{theorem}[The spectrum is not primitive]
\label{c05:thm:spectrum-not-primitive}
The spectrum is not a primitive of the theory: it is the inevitable output of the
actualization attractor, carrying no choice and presupposing no independent input.
\end{theorem}

\begin{proof}
The spectrum-necessity audit \cite{c05:QG295} poses exactly this question --- is the spectrum
fundamental, or the inevitable fixed point of actualization? --- and establishes the latter.
The spectrum is the eigenspectrum of the converged network
(Theorem~\ref{c05:thm:spectrum-eigenspectrum}), which is the unique, fixed-size,
content-independent attractor (Theorem~\ref{c05:thm:attractor-unique},
Theorem~\ref{c05:thm:n96-attractor}); it therefore presupposes the actualization process and
its primitive Difference, and is an output, not an input. The reconstruction audit
\cite{c05:QG296} establishes that physics is the readout of the spectrum
(Chapter~\ref{c03:chap:actualization}, Corollary~\ref{c03:cor:physics-readout}), the sectors
being projection classes over the operator layer of the one spectrum.
\end{proof}

The spectrum is therefore also necessary: physics is the readout of the spectrum, so without
the spectrum there is no physics. The two properties are compatible: the spectrum is forced
by the dynamics (inevitable) and indispensable to the derivation (necessary).

\section{The Canonical Core Completed}
\label{c05:sec:core-completed}

\begin{theorem}[The canonical core is completed]
\label{c05:thm:core-completed}
The first three members of the canonical hierarchy are now established: Difference is the
primitive (Chapter~\ref{c01:chap:difference}), Actualization is the derived process face of
Difference (Chapter~\ref{c03:chap:actualization}), and the Spectrum is the inevitable output of
the actualization attractor (this chapter). The canonical core
\begin{equation}
\label{c05:eq:core-complete}
\text{Difference}\;\longrightarrow\;\text{Actualization}\;\longrightarrow\;\text{Inevitable Spectrum}
\end{equation}
is complete, anchoring the minimal hierarchy Difference → Actualization → Spectrum → Physics:
its first three members are established, and Physics is the spectrum readout developed in
the subsequent chapters.
\end{theorem}

\begin{proof}
The three members are established as primitive (Difference,
Theorem~\ref{c01:thm:minimal-primitives} of Chapter~\ref{c01:chap:difference}), derived process
(Actualization, Theorem~\ref{c03:thm:derived-not-primitive} of
Chapter~\ref{c03:chap:actualization}), and inevitable output (Spectrum,
Theorem~\ref{c05:thm:spectrum-not-primitive}), completing the first three members of the
minimal hierarchy \cite{c05:QG293,c05:QG294}. The fourth member, Physics, is the readout of the
spectrum (Chapter~\ref{c03:chap:actualization}, Corollary~\ref{c03:cor:physics-readout}),
developed in the physics chapters that follow. Every member of the hierarchy thus has a
canonical status: primitive, derived process, inevitable output, emergent readout.
\end{proof}

\section{Consequences for Physics}
\label{c05:sec:physics-consequences}

All observable physics is the readout of the inevitable spectrum: the resonance structure
projected by the operator layer, from which the sectors of physics emerge. The spectrum is
the unique, inevitable output of the attractor (Theorem~\ref{c05:thm:spectrum-not-primitive});
the sectors of physics are projection classes over the operator layer of the one spectrum
(Chapter~\ref{c03:chap:actualization}, Theorem~\ref{c03:thm:resonance-readout}). Hence every
sector --- masses, couplings, mixings, gravity, cosmology --- is a projection of the one
actualization-driven spectrum \cite{c05:QG272}.

\begin{theorem}[No physics without the spectrum]
\label{c05:thm:no-physics-without-spectrum}
There is no physics in the theory without the spectrum: every derived observable is a
function of the spectral structure, and no observable bypasses the spectrum layer. Hence the
observable space of the theory is determined by the inevitable spectrum: the set of derived
observables is the set of reads of the unique spectrum.
\end{theorem}

\begin{proof}
The reconstruction audit \cite{c05:QG296} reconstructs the major results of the theory from the
minimal hierarchy, with the spectrum layer as the source of the spectral constants used by
all sectors; hence every observable passes through the spectrum layer, and there is no
physics without it. Since the attractor is unique (Theorem~\ref{c05:thm:attractor-unique}),
the spectrum is unique; every observable is a read of the one inevitable spectrum, and the
observable space is the set of those reads.
\end{proof}

Consistent with MONO007 \cite{c05:MONO007}, the spectrum is never presented as a primitive in
this chapter: every statement positions it as the derived output of the actualization
attractor, and the "spectrum layer" of the minimal hierarchy means the derived spectrum, not
an independent input --- the referee-safe reading established by the spectrum-necessity audit
\cite{c05:QG295}.

\section{Summary}
\label{c05:sec:spectrum-summary}

This chapter has established the third member of the canonical architecture. The
actualization dynamics converge to a unique, content-independent attractor, the closure
fixed point (Theorem~\ref{c05:thm:attractor-unique}), whose size is the inevitable stable
fixed point $N=96$ (Theorem~\ref{c05:thm:n96-attractor}); the spectrum is the Laplacian
eigenspectrum of the converged network, its constants derived outputs
(Theorem~\ref{c05:thm:spectrum-eigenspectrum}). It is therefore unique, carries no choice,
and is not a primitive but the inevitable output required for physics
(Theorem~\ref{c05:thm:spectrum-not-primitive}). With the canonical core Difference →
Actualization → Spectrum complete (Theorem~\ref{c05:thm:core-completed}), physics is the
readout of the inevitable spectrum, which determines the observable space
(Theorem~\ref{c05:thm:no-physics-without-spectrum}); the physics chapters that follow derive
each sector as a read of this one spectrum.

\begin{thebibliography}{99}
\bibitem{c05:QG116} AT-QG Phase 116, \emph{Actualization Structures}, Docs/Research.
\bibitem{c05:QG159} AT-QG Phase 159, \emph{D96 Selection Origin}, Docs/Research.
\bibitem{c05:QG160} AT-QG Phase 160, \emph{Period-3 Seed Origin}, Docs/Research.
\bibitem{c05:QG272} AT-QG Phase 272, \emph{Sector Emergence Audit}, Docs/Research.
\bibitem{c05:QG282} AT-QG Phase 282, \emph{Boundary Origin Audit} (Closure Principle), Docs/Research.
\bibitem{c05:QG293} AT-QG Phase 293, \emph{Hierarchy Necessity Audit}, Docs/Research.
\bibitem{c05:QG294} AT-QG Phase 294, \emph{Minimal Theory Audit}, Docs/Research.
\bibitem{c05:QG295} AT-QG Phase 295, \emph{Spectrum Necessity Audit}, Docs/Research.
\bibitem{c05:QG296} AT-QG Phase 296, \emph{Reconstruction Audit}, Docs/Research.
\bibitem{c05:MONO007} AT-MONO007, \emph{Referee-Readiness Resolution}, Docs/Zenodo/Monograph.
\end{thebibliography}
 after the preamble
%         that defines the theorem environments (or include via the master file).
% ======================================================================

% ---- Theorem environments (define in the master preamble if not present) ----
% \newtheorem{definition}{Definition}[chapter]
% \newtheorem{lemma}[definition]{Lemma}
% \newtheorem{theorem}[definition]{Theorem}
% \newtheorem{corollary}[definition]{Corollary}
% \newtheorem{remark}[definition]{Remark}

\chapter{The Inevitable Spectrum}
\label{c05:chap:spectrum}

\section{Introduction}
\label{c05:sec:spectrum-introduction}

The canonical hierarchy of The Actualization Theory
\begin{equation}
\label{c05:eq:core}
\text{Difference}\;\longrightarrow\;\text{Actualization}\;\longrightarrow\;\text{Inevitable Spectrum}\;\longrightarrow\;\text{Physics}
\end{equation}
is complete as a derivation structure (Chapter~\ref{c04:chap:closure-minimal},
Theorem~\ref{c04:thm:minimal-hierarchy}). This chapter develops the third member of the
hierarchy: the \emph{inevitable spectrum}. We establish that the spectrum is not a primitive
of the theory but the \emph{derived output} of the actualization process: the dynamics
converge to a content-independent attractor of size $N=96$ within the accepted structural
class; the attractor is a stable fixed point, not a freely chosen input; and the spectrum is
the eigenspectrum of that attractor, carrying no independent choice
\cite{c05:QG116,c05:QG282,c05:QG295}.

We then establish the \emph{necessity} of the spectrum and the tested uniqueness of the
attractor within the canonical structural class: the attractor is content-independent
\cite{c05:QG116}, the size $N=96$ is inevitable within the accepted D96 selection
constraints --- not a choice \cite{c05:QG159,c05:QG160} --- and the spectrum is a
deterministic function of the converged network. This completes the canonical core
Difference → Actualization → Spectrum, whose first three members are now established as
primitive, derived, and inevitable output respectively; the physics chapters that follow
read all observable sectors from this one spectrum \cite{c05:QG293,c05:QG294,c05:QG296}.

\section{The Actualization Attractor}
\label{c05:sec:attractor}

\begin{theorem}[The actualization attractor: convergence, uniqueness, closure]
\label{c05:thm:attractor-unique}
The \emph{actualization attractor} is the converged state of the actualization process: the
configuration to which every initial pattern of Difference's count-producing dynamics
converges, with zero residual change. The actualization dynamics converge --- residual link
growth reaches zero and the topology saturates at a stable final state --- the attractor is
unique and content-independent (every initial activity pattern converges to the same final
geometry), and it is the closure fixed point of the theory.
\end{theorem}

\begin{proof}
The actualization-structures program \cite{c05:QG116} establishes that the actualization
dynamics converge with zero residual link growth: the self-reinforcing feedback (activity to
links to activity) is bounded by network capacity, so the topology saturates, and identical
link counts and identical span are reached from every initial pattern. Hence the attractor
is unique --- one final geometry, not a family --- and independent of the starting
configuration. By the Closure Principle (Chapter~\ref{c04:chap:closure-minimal},
Theorem~\ref{c04:thm:closure-holds}, Definition~\ref{c04:def:closure-principle}), the boundary
of the process is its stable fixed point; hence the unique attractor and the closure fixed
point coincide.
\end{proof}

\section{The Fixed Point \texorpdfstring{$N=96$}{N=96}}
\label{c05:sec:n96}

\begin{theorem}[The attractor is \texorpdfstring{$N=96$}{N=96}, not a choice; no freedom in the size]
\label{c05:thm:n96-attractor}
The \emph{$N=96$ attractor} is the converged network of the actualization process within the
canonical structural class: the stable fixed point whose size is $N=96$. The size is the
attractor of the dynamics, not a freely chosen input: the process converges to $N=96$ from
every initial pattern in the accepted comparison class, and no other tested size is a stable
fixed point of the dynamics.
\end{theorem}

\begin{proof}
The spectrum-necessity audit \cite{c05:QG295} establishes that $N=96$ is the attractor, not a
choice: the actualization flow has a stable fixed point at $N=96$, and the boundary is the
closure of that flow (the Closure Principle \cite{c05:QG282}). The D96-selection program adds
that $96$ is selected as the unique complete-Z2 natural size with the three-family octave
window; no other natural size satisfies these conditions \cite{c05:QG159,c05:QG160}. Hence
$N=96$ is the stable fixed point of the flow within the canonical comparison class,
content-independent (Theorem~\ref{c05:thm:attractor-unique}): an output of the dynamics, not
an input or free parameter.
\end{proof}

\subsection{Why \texorpdfstring{$N=96$}{N=96}?}
\label{c05:sec:why-n96}

\begin{table}[h]
\centering
\begin{tabular}{p{0.18\linewidth}p{0.28\linewidth}p{0.46\linewidth}}
\hline
\textbf{Argument class} & \textbf{What it uses} & \textbf{Status} \\
\hline
Selection & period-3 seed + Z2 half-shift; 3-family octave window; natural doubling chain $n=3\cdot2^k$; tested candidate set \cite{c05:QG159,c05:QG160} & D96 is the only rung in the tested structural class that satisfies all conditions \\
Necessity & closure of the actualization flow; stable fixed point; spectrum as eigenspectrum of the converged network \cite{c05:QG282,c05:QG295} & If the flow does not close, there is no spectrum layer and no physics readout \\
Uniqueness & content-independent attractor; same final geometry from every initial pattern \cite{c05:QG116} & Unique within the canonical comparison class examined \\
\hline
\end{tabular}
\end{table}

\paragraph{Why not 94, 95, 97?}
On the accepted selection criterion, these sizes fail the seed half-shift condition: the
period-3 seed requires $6 \mid n$ for the Z2 doublet structure, so $94,95,97$ are excluded
before the octave-window test is even applied. They are therefore not members of the
canonical candidate class used by QG159.

\paragraph{Why not 128?}
128 fails the same seed symmetry test ($128 \not\equiv 0 \pmod 6$) and, in the selection
audit, also fails the three-family window: its span is above the $[4,8)$ range, giving four
families rather than three \cite{c05:QG159}.

\paragraph{Exact status.}
The strongest referee-safe statement is: \emph{within the canonical structural constraints
actually used by AT --- period-3 seed, Z2 half-shift, three-family octave window, and the
tested comparison class --- $N=96$ is inevitable and uniquely selected.} A global proof that
no other closure could exist outside that structural class is \emph{not} claimed here.
That is the boundary.

\section{Spectrum Generation}
\label{c05:sec:spectrum-generation}

\begin{theorem}[The spectrum is the eigenspectrum of the attractor]
\label{c05:thm:spectrum-eigenspectrum}
\emph{Spectrum generation} is the derivation of the spectrum from the converged network: the
spectrum of the theory is the Laplacian eigenspectrum of the $N=96$ attractor --- $95$
positive modes with a single zero mode (the background) --- a deterministic function of the
attractor. Its spectral constants --- the moments $\Sigma m=95$,
$\Sigma\sqrt{m}=64.08$, $\Sigma m^2=229$, the occupancy moment $1900.25$, the span $6.40$,
and the doublet multiplicities $[42\times2,5,6]$ --- are outputs of the spectrum, hence of
the attractor.
\end{theorem}

\begin{proof}
The spectrum-necessity audit \cite{c05:QG295} establishes that the spectrum is the Laplacian
eigenspectrum of the converged network: the graph Laplacian of the $N=96$ attractor has
$95$ positive eigenvalues (the modes) and one zero eigenvalue (the background). Since the
attractor is unique and deterministic (Theorem~\ref{c05:thm:attractor-unique},
Theorem~\ref{c05:thm:n96-attractor}), the spectrum is a deterministic function of it --- an
output of the dynamics, not an additional input --- and its moments and structural features
are derived quantities, not primitive inputs. This is the canonical reading of the spectrum
layer of the minimal hierarchy \cite{c05:QG294}.
\end{proof}

\section{Uniqueness of the Spectrum}
\label{c05:sec:spectrum-uniqueness}

Uniqueness here means uniqueness \emph{within the canonical structural class}: the spectrum is
the eigenspectrum of the attractor (Theorem~\ref{c05:thm:attractor-unique},
Theorem~\ref{c05:thm:spectrum-eigenspectrum}), whose size is fixed at $N=96$
(Theorem~\ref{c05:thm:n96-attractor}), and the eigenspectrum of a given network is unique.
Hence there is exactly one spectrum of the theory in the accepted class --- one attractor,
one eigenspectrum --- and since the attractor carries no freedom, no spectral feature is
freely adjustable: the spectrum carries no choice at any level.

\section{Spectrum Necessity}
\label{c05:sec:spectrum-necessity}

\begin{theorem}[The spectrum is not primitive]
\label{c05:thm:spectrum-not-primitive}
The spectrum is not a primitive of the theory: it is the inevitable output of the
actualization attractor, carrying no choice and presupposing no independent input.
\end{theorem}

\begin{proof}
The spectrum-necessity audit \cite{c05:QG295} poses exactly this question --- is the spectrum
fundamental, or the inevitable fixed point of actualization? --- and establishes the latter.
The spectrum is the eigenspectrum of the converged network
(Theorem~\ref{c05:thm:spectrum-eigenspectrum}), which is the unique, fixed-size,
content-independent attractor (Theorem~\ref{c05:thm:attractor-unique},
Theorem~\ref{c05:thm:n96-attractor}); it therefore presupposes the actualization process and
its primitive Difference, and is an output, not an input. The reconstruction audit
\cite{c05:QG296} establishes that physics is the readout of the spectrum
(Chapter~\ref{c03:chap:actualization}, Corollary~\ref{c03:cor:physics-readout}), the sectors
being projection classes over the operator layer of the one spectrum.
\end{proof}

The spectrum is therefore also necessary: physics is the readout of the spectrum, so without
the spectrum there is no physics. The two properties are compatible: the spectrum is forced
by the dynamics (inevitable) and indispensable to the derivation (necessary).

\section{The Canonical Core Completed}
\label{c05:sec:core-completed}

\begin{theorem}[The canonical core is completed]
\label{c05:thm:core-completed}
The first three members of the canonical hierarchy are now established: Difference is the
primitive (Chapter~\ref{c01:chap:difference}), Actualization is the derived process face of
Difference (Chapter~\ref{c03:chap:actualization}), and the Spectrum is the inevitable output of
the actualization attractor (this chapter). The canonical core
\begin{equation}
\label{c05:eq:core-complete}
\text{Difference}\;\longrightarrow\;\text{Actualization}\;\longrightarrow\;\text{Inevitable Spectrum}
\end{equation}
is complete, anchoring the minimal hierarchy Difference → Actualization → Spectrum → Physics:
its first three members are established, and Physics is the spectrum readout developed in
the subsequent chapters.
\end{theorem}

\begin{proof}
The three members are established as primitive (Difference,
Theorem~\ref{c01:thm:minimal-primitives} of Chapter~\ref{c01:chap:difference}), derived process
(Actualization, Theorem~\ref{c03:thm:derived-not-primitive} of
Chapter~\ref{c03:chap:actualization}), and inevitable output (Spectrum,
Theorem~\ref{c05:thm:spectrum-not-primitive}), completing the first three members of the
minimal hierarchy \cite{c05:QG293,c05:QG294}. The fourth member, Physics, is the readout of the
spectrum (Chapter~\ref{c03:chap:actualization}, Corollary~\ref{c03:cor:physics-readout}),
developed in the physics chapters that follow. Every member of the hierarchy thus has a
canonical status: primitive, derived process, inevitable output, emergent readout.
\end{proof}

\section{Consequences for Physics}
\label{c05:sec:physics-consequences}

All observable physics is the readout of the inevitable spectrum: the resonance structure
projected by the operator layer, from which the sectors of physics emerge. The spectrum is
the unique, inevitable output of the attractor (Theorem~\ref{c05:thm:spectrum-not-primitive});
the sectors of physics are projection classes over the operator layer of the one spectrum
(Chapter~\ref{c03:chap:actualization}, Theorem~\ref{c03:thm:resonance-readout}). Hence every
sector --- masses, couplings, mixings, gravity, cosmology --- is a projection of the one
actualization-driven spectrum \cite{c05:QG272}.

\begin{theorem}[No physics without the spectrum]
\label{c05:thm:no-physics-without-spectrum}
There is no physics in the theory without the spectrum: every derived observable is a
function of the spectral structure, and no observable bypasses the spectrum layer. Hence the
observable space of the theory is determined by the inevitable spectrum: the set of derived
observables is the set of reads of the unique spectrum.
\end{theorem}

\begin{proof}
The reconstruction audit \cite{c05:QG296} reconstructs the major results of the theory from the
minimal hierarchy, with the spectrum layer as the source of the spectral constants used by
all sectors; hence every observable passes through the spectrum layer, and there is no
physics without it. Since the attractor is unique (Theorem~\ref{c05:thm:attractor-unique}),
the spectrum is unique; every observable is a read of the one inevitable spectrum, and the
observable space is the set of those reads.
\end{proof}

Consistent with MONO007 \cite{c05:MONO007}, the spectrum is never presented as a primitive in
this chapter: every statement positions it as the derived output of the actualization
attractor, and the "spectrum layer" of the minimal hierarchy means the derived spectrum, not
an independent input --- the referee-safe reading established by the spectrum-necessity audit
\cite{c05:QG295}.

\section{Summary}
\label{c05:sec:spectrum-summary}

This chapter has established the third member of the canonical architecture. The
actualization dynamics converge to a unique, content-independent attractor, the closure
fixed point (Theorem~\ref{c05:thm:attractor-unique}), whose size is the inevitable stable
fixed point $N=96$ (Theorem~\ref{c05:thm:n96-attractor}); the spectrum is the Laplacian
eigenspectrum of the converged network, its constants derived outputs
(Theorem~\ref{c05:thm:spectrum-eigenspectrum}). It is therefore unique, carries no choice,
and is not a primitive but the inevitable output required for physics
(Theorem~\ref{c05:thm:spectrum-not-primitive}). With the canonical core Difference →
Actualization → Spectrum complete (Theorem~\ref{c05:thm:core-completed}), physics is the
readout of the inevitable spectrum, which determines the observable space
(Theorem~\ref{c05:thm:no-physics-without-spectrum}); the physics chapters that follow derive
each sector as a read of this one spectrum.

\begin{thebibliography}{99}
\bibitem{c05:QG116} AT-QG Phase 116, \emph{Actualization Structures}, Docs/Research.
\bibitem{c05:QG159} AT-QG Phase 159, \emph{D96 Selection Origin}, Docs/Research.
\bibitem{c05:QG160} AT-QG Phase 160, \emph{Period-3 Seed Origin}, Docs/Research.
\bibitem{c05:QG272} AT-QG Phase 272, \emph{Sector Emergence Audit}, Docs/Research.
\bibitem{c05:QG282} AT-QG Phase 282, \emph{Boundary Origin Audit} (Closure Principle), Docs/Research.
\bibitem{c05:QG293} AT-QG Phase 293, \emph{Hierarchy Necessity Audit}, Docs/Research.
\bibitem{c05:QG294} AT-QG Phase 294, \emph{Minimal Theory Audit}, Docs/Research.
\bibitem{c05:QG295} AT-QG Phase 295, \emph{Spectrum Necessity Audit}, Docs/Research.
\bibitem{c05:QG296} AT-QG Phase 296, \emph{Reconstruction Audit}, Docs/Research.
\bibitem{c05:MONO007} AT-MONO007, \emph{Referee-Readiness Resolution}, Docs/Zenodo/Monograph.
\end{thebibliography}
 after the preamble
%         that defines the theorem environments (or include via the master file).
% ======================================================================

% ---- Theorem environments (define in the master preamble if not present) ----
% \newtheorem{definition}{Definition}[chapter]
% \newtheorem{lemma}[definition]{Lemma}
% \newtheorem{theorem}[definition]{Theorem}
% \newtheorem{corollary}[definition]{Corollary}
% \newtheorem{remark}[definition]{Remark}

\chapter{The Inevitable Spectrum}
\label{c05:chap:spectrum}

\section{Introduction}
\label{c05:sec:spectrum-introduction}

The canonical hierarchy of The Actualization Theory
\begin{equation}
\label{c05:eq:core}
\text{Difference}\;\longrightarrow\;\text{Actualization}\;\longrightarrow\;\text{Inevitable Spectrum}\;\longrightarrow\;\text{Physics}
\end{equation}
is complete as a derivation structure (Chapter~\ref{c04:chap:closure-minimal},
Theorem~\ref{c04:thm:minimal-hierarchy}). This chapter develops the third member of the
hierarchy: the \emph{inevitable spectrum}. We establish that the spectrum is not a primitive
of the theory but the \emph{derived output} of the actualization process: the dynamics
converge to a content-independent attractor of size $N=96$ within the accepted structural
class; the attractor is a stable fixed point, not a freely chosen input; and the spectrum is
the eigenspectrum of that attractor, carrying no independent choice
\cite{c05:QG116,c05:QG282,c05:QG295}.

We then establish the \emph{necessity} of the spectrum and the tested uniqueness of the
attractor within the canonical structural class: the attractor is content-independent
\cite{c05:QG116}, the size $N=96$ is inevitable within the accepted D96 selection
constraints --- not a choice \cite{c05:QG159,c05:QG160} --- and the spectrum is a
deterministic function of the converged network. This completes the canonical core
Difference → Actualization → Spectrum, whose first three members are now established as
primitive, derived, and inevitable output respectively; the physics chapters that follow
read all observable sectors from this one spectrum \cite{c05:QG293,c05:QG294,c05:QG296}.

\section{The Actualization Attractor}
\label{c05:sec:attractor}

\begin{theorem}[The actualization attractor: convergence, uniqueness, closure]
\label{c05:thm:attractor-unique}
The \emph{actualization attractor} is the converged state of the actualization process: the
configuration to which every initial pattern of Difference's count-producing dynamics
converges, with zero residual change. The actualization dynamics converge --- residual link
growth reaches zero and the topology saturates at a stable final state --- the attractor is
unique and content-independent (every initial activity pattern converges to the same final
geometry), and it is the closure fixed point of the theory.
\end{theorem}

\begin{proof}
The actualization-structures program \cite{c05:QG116} establishes that the actualization
dynamics converge with zero residual link growth: the self-reinforcing feedback (activity to
links to activity) is bounded by network capacity, so the topology saturates, and identical
link counts and identical span are reached from every initial pattern. Hence the attractor
is unique --- one final geometry, not a family --- and independent of the starting
configuration. By the Closure Principle (Chapter~\ref{c04:chap:closure-minimal},
Theorem~\ref{c04:thm:closure-holds}, Definition~\ref{c04:def:closure-principle}), the boundary
of the process is its stable fixed point; hence the unique attractor and the closure fixed
point coincide.
\end{proof}

\section{The Fixed Point \texorpdfstring{$N=96$}{N=96}}
\label{c05:sec:n96}

\begin{theorem}[The attractor is \texorpdfstring{$N=96$}{N=96}, not a choice; no freedom in the size]
\label{c05:thm:n96-attractor}
The \emph{$N=96$ attractor} is the converged network of the actualization process within the
canonical structural class: the stable fixed point whose size is $N=96$. The size is the
attractor of the dynamics, not a freely chosen input: the process converges to $N=96$ from
every initial pattern in the accepted comparison class, and no other tested size is a stable
fixed point of the dynamics.
\end{theorem}

\begin{proof}
The spectrum-necessity audit \cite{c05:QG295} establishes that $N=96$ is the attractor, not a
choice: the actualization flow has a stable fixed point at $N=96$, and the boundary is the
closure of that flow (the Closure Principle \cite{c05:QG282}). The D96-selection program adds
that $96$ is selected as the unique complete-Z2 natural size with the three-family octave
window; no other natural size satisfies these conditions \cite{c05:QG159,c05:QG160}. Hence
$N=96$ is the stable fixed point of the flow within the canonical comparison class,
content-independent (Theorem~\ref{c05:thm:attractor-unique}): an output of the dynamics, not
an input or free parameter.
\end{proof}

\subsection{Why \texorpdfstring{$N=96$}{N=96}?}
\label{c05:sec:why-n96}

\begin{table}[h]
\centering
\begin{tabular}{p{0.18\linewidth}p{0.28\linewidth}p{0.46\linewidth}}
\hline
\textbf{Argument class} & \textbf{What it uses} & \textbf{Status} \\
\hline
Selection & period-3 seed + Z2 half-shift; 3-family octave window; natural doubling chain $n=3\cdot2^k$; tested candidate set \cite{c05:QG159,c05:QG160} & D96 is the only rung in the tested structural class that satisfies all conditions \\
Necessity & closure of the actualization flow; stable fixed point; spectrum as eigenspectrum of the converged network \cite{c05:QG282,c05:QG295} & If the flow does not close, there is no spectrum layer and no physics readout \\
Uniqueness & content-independent attractor; same final geometry from every initial pattern \cite{c05:QG116} & Unique within the canonical comparison class examined \\
\hline
\end{tabular}
\end{table}

\paragraph{Why not 94, 95, 97?}
On the accepted selection criterion, these sizes fail the seed half-shift condition: the
period-3 seed requires $6 \mid n$ for the Z2 doublet structure, so $94,95,97$ are excluded
before the octave-window test is even applied. They are therefore not members of the
canonical candidate class used by QG159.

\paragraph{Why not 128?}
128 fails the same seed symmetry test ($128 \not\equiv 0 \pmod 6$) and, in the selection
audit, also fails the three-family window: its span is above the $[4,8)$ range, giving four
families rather than three \cite{c05:QG159}.

\paragraph{Exact status.}
The strongest referee-safe statement is: \emph{within the canonical structural constraints
actually used by AT --- period-3 seed, Z2 half-shift, three-family octave window, and the
tested comparison class --- $N=96$ is inevitable and uniquely selected.} A global proof that
no other closure could exist outside that structural class is \emph{not} claimed here.
That is the boundary.

\section{Spectrum Generation}
\label{c05:sec:spectrum-generation}

\begin{theorem}[The spectrum is the eigenspectrum of the attractor]
\label{c05:thm:spectrum-eigenspectrum}
\emph{Spectrum generation} is the derivation of the spectrum from the converged network: the
spectrum of the theory is the Laplacian eigenspectrum of the $N=96$ attractor --- $95$
positive modes with a single zero mode (the background) --- a deterministic function of the
attractor. Its spectral constants --- the moments $\Sigma m=95$,
$\Sigma\sqrt{m}=64.08$, $\Sigma m^2=229$, the occupancy moment $1900.25$, the span $6.40$,
and the doublet multiplicities $[42\times2,5,6]$ --- are outputs of the spectrum, hence of
the attractor.
\end{theorem}

\begin{proof}
The spectrum-necessity audit \cite{c05:QG295} establishes that the spectrum is the Laplacian
eigenspectrum of the converged network: the graph Laplacian of the $N=96$ attractor has
$95$ positive eigenvalues (the modes) and one zero eigenvalue (the background). Since the
attractor is unique and deterministic (Theorem~\ref{c05:thm:attractor-unique},
Theorem~\ref{c05:thm:n96-attractor}), the spectrum is a deterministic function of it --- an
output of the dynamics, not an additional input --- and its moments and structural features
are derived quantities, not primitive inputs. This is the canonical reading of the spectrum
layer of the minimal hierarchy \cite{c05:QG294}.
\end{proof}

\section{Uniqueness of the Spectrum}
\label{c05:sec:spectrum-uniqueness}

Uniqueness here means uniqueness \emph{within the canonical structural class}: the spectrum is
the eigenspectrum of the attractor (Theorem~\ref{c05:thm:attractor-unique},
Theorem~\ref{c05:thm:spectrum-eigenspectrum}), whose size is fixed at $N=96$
(Theorem~\ref{c05:thm:n96-attractor}), and the eigenspectrum of a given network is unique.
Hence there is exactly one spectrum of the theory in the accepted class --- one attractor,
one eigenspectrum --- and since the attractor carries no freedom, no spectral feature is
freely adjustable: the spectrum carries no choice at any level.

\section{Spectrum Necessity}
\label{c05:sec:spectrum-necessity}

\begin{theorem}[The spectrum is not primitive]
\label{c05:thm:spectrum-not-primitive}
The spectrum is not a primitive of the theory: it is the inevitable output of the
actualization attractor, carrying no choice and presupposing no independent input.
\end{theorem}

\begin{proof}
The spectrum-necessity audit \cite{c05:QG295} poses exactly this question --- is the spectrum
fundamental, or the inevitable fixed point of actualization? --- and establishes the latter.
The spectrum is the eigenspectrum of the converged network
(Theorem~\ref{c05:thm:spectrum-eigenspectrum}), which is the unique, fixed-size,
content-independent attractor (Theorem~\ref{c05:thm:attractor-unique},
Theorem~\ref{c05:thm:n96-attractor}); it therefore presupposes the actualization process and
its primitive Difference, and is an output, not an input. The reconstruction audit
\cite{c05:QG296} establishes that physics is the readout of the spectrum
(Chapter~\ref{c03:chap:actualization}, Corollary~\ref{c03:cor:physics-readout}), the sectors
being projection classes over the operator layer of the one spectrum.
\end{proof}

The spectrum is therefore also necessary: physics is the readout of the spectrum, so without
the spectrum there is no physics. The two properties are compatible: the spectrum is forced
by the dynamics (inevitable) and indispensable to the derivation (necessary).

\section{The Canonical Core Completed}
\label{c05:sec:core-completed}

\begin{theorem}[The canonical core is completed]
\label{c05:thm:core-completed}
The first three members of the canonical hierarchy are now established: Difference is the
primitive (Chapter~\ref{c01:chap:difference}), Actualization is the derived process face of
Difference (Chapter~\ref{c03:chap:actualization}), and the Spectrum is the inevitable output of
the actualization attractor (this chapter). The canonical core
\begin{equation}
\label{c05:eq:core-complete}
\text{Difference}\;\longrightarrow\;\text{Actualization}\;\longrightarrow\;\text{Inevitable Spectrum}
\end{equation}
is complete, anchoring the minimal hierarchy Difference → Actualization → Spectrum → Physics:
its first three members are established, and Physics is the spectrum readout developed in
the subsequent chapters.
\end{theorem}

\begin{proof}
The three members are established as primitive (Difference,
Theorem~\ref{c01:thm:minimal-primitives} of Chapter~\ref{c01:chap:difference}), derived process
(Actualization, Theorem~\ref{c03:thm:derived-not-primitive} of
Chapter~\ref{c03:chap:actualization}), and inevitable output (Spectrum,
Theorem~\ref{c05:thm:spectrum-not-primitive}), completing the first three members of the
minimal hierarchy \cite{c05:QG293,c05:QG294}. The fourth member, Physics, is the readout of the
spectrum (Chapter~\ref{c03:chap:actualization}, Corollary~\ref{c03:cor:physics-readout}),
developed in the physics chapters that follow. Every member of the hierarchy thus has a
canonical status: primitive, derived process, inevitable output, emergent readout.
\end{proof}

\section{Consequences for Physics}
\label{c05:sec:physics-consequences}

All observable physics is the readout of the inevitable spectrum: the resonance structure
projected by the operator layer, from which the sectors of physics emerge. The spectrum is
the unique, inevitable output of the attractor (Theorem~\ref{c05:thm:spectrum-not-primitive});
the sectors of physics are projection classes over the operator layer of the one spectrum
(Chapter~\ref{c03:chap:actualization}, Theorem~\ref{c03:thm:resonance-readout}). Hence every
sector --- masses, couplings, mixings, gravity, cosmology --- is a projection of the one
actualization-driven spectrum \cite{c05:QG272}.

\begin{theorem}[No physics without the spectrum]
\label{c05:thm:no-physics-without-spectrum}
There is no physics in the theory without the spectrum: every derived observable is a
function of the spectral structure, and no observable bypasses the spectrum layer. Hence the
observable space of the theory is determined by the inevitable spectrum: the set of derived
observables is the set of reads of the unique spectrum.
\end{theorem}

\begin{proof}
The reconstruction audit \cite{c05:QG296} reconstructs the major results of the theory from the
minimal hierarchy, with the spectrum layer as the source of the spectral constants used by
all sectors; hence every observable passes through the spectrum layer, and there is no
physics without it. Since the attractor is unique (Theorem~\ref{c05:thm:attractor-unique}),
the spectrum is unique; every observable is a read of the one inevitable spectrum, and the
observable space is the set of those reads.
\end{proof}

Consistent with MONO007 \cite{c05:MONO007}, the spectrum is never presented as a primitive in
this chapter: every statement positions it as the derived output of the actualization
attractor, and the "spectrum layer" of the minimal hierarchy means the derived spectrum, not
an independent input --- the referee-safe reading established by the spectrum-necessity audit
\cite{c05:QG295}.

\section{Summary}
\label{c05:sec:spectrum-summary}

This chapter has established the third member of the canonical architecture. The
actualization dynamics converge to a unique, content-independent attractor, the closure
fixed point (Theorem~\ref{c05:thm:attractor-unique}), whose size is the inevitable stable
fixed point $N=96$ (Theorem~\ref{c05:thm:n96-attractor}); the spectrum is the Laplacian
eigenspectrum of the converged network, its constants derived outputs
(Theorem~\ref{c05:thm:spectrum-eigenspectrum}). It is therefore unique, carries no choice,
and is not a primitive but the inevitable output required for physics
(Theorem~\ref{c05:thm:spectrum-not-primitive}). With the canonical core Difference →
Actualization → Spectrum complete (Theorem~\ref{c05:thm:core-completed}), physics is the
readout of the inevitable spectrum, which determines the observable space
(Theorem~\ref{c05:thm:no-physics-without-spectrum}); the physics chapters that follow derive
each sector as a read of this one spectrum.

\begin{thebibliography}{99}
\bibitem{c05:QG116} AT-QG Phase 116, \emph{Actualization Structures}, Docs/Research.
\bibitem{c05:QG159} AT-QG Phase 159, \emph{D96 Selection Origin}, Docs/Research.
\bibitem{c05:QG160} AT-QG Phase 160, \emph{Period-3 Seed Origin}, Docs/Research.
\bibitem{c05:QG272} AT-QG Phase 272, \emph{Sector Emergence Audit}, Docs/Research.
\bibitem{c05:QG282} AT-QG Phase 282, \emph{Boundary Origin Audit} (Closure Principle), Docs/Research.
\bibitem{c05:QG293} AT-QG Phase 293, \emph{Hierarchy Necessity Audit}, Docs/Research.
\bibitem{c05:QG294} AT-QG Phase 294, \emph{Minimal Theory Audit}, Docs/Research.
\bibitem{c05:QG295} AT-QG Phase 295, \emph{Spectrum Necessity Audit}, Docs/Research.
\bibitem{c05:QG296} AT-QG Phase 296, \emph{Reconstruction Audit}, Docs/Research.
\bibitem{c05:MONO007} AT-MONO007, \emph{Referee-Readiness Resolution}, Docs/Zenodo/Monograph.
\end{thebibliography}


% Part III — Spectrum
\part{Spectrum}
% ======================================================================
%  AT-MONO105 — Chapter 6: The D96 Spectrum
%  Canonical Monograph (v2.0) · The Actualization Theory:
%  A Reconstruction of Physics from Difference, Actualization and Spectrum
%
%  Status: COMPLETE — PASS-READY (MONO007 referee-ready architecture)
%  Inputs: Chapter01_Difference.tex, Chapter02_Eta.tex,
%          Chapter03_Actualization.tex, Chapter04_ClosureAndMinimalTheory.tex,
%          Chapter05_InevitableSpectrum.tex
%  Method: assembly only — canonical results only, no new physics, no speculation
%  Citations: QG155 (symmetry origin), QG156 (unified access law),
%             QG157 (effective access counts), QG158 (moment orders),
%             QG159 (D96 selection), QG295 (spectrum necessity)
%
%  Usage: % ======================================================================
%  AT-MONO105 — Chapter 6: The D96 Spectrum
%  Canonical Monograph (v2.0) · The Actualization Theory:
%  A Reconstruction of Physics from Difference, Actualization and Spectrum
%
%  Status: COMPLETE — PASS-READY (MONO007 referee-ready architecture)
%  Inputs: Chapter01_Difference.tex, Chapter02_Eta.tex,
%          Chapter03_Actualization.tex, Chapter04_ClosureAndMinimalTheory.tex,
%          Chapter05_InevitableSpectrum.tex
%  Method: assembly only — canonical results only, no new physics, no speculation
%  Citations: QG155 (symmetry origin), QG156 (unified access law),
%             QG157 (effective access counts), QG158 (moment orders),
%             QG159 (D96 selection), QG295 (spectrum necessity)
%
%  Usage: % ======================================================================
%  AT-MONO105 — Chapter 6: The D96 Spectrum
%  Canonical Monograph (v2.0) · The Actualization Theory:
%  A Reconstruction of Physics from Difference, Actualization and Spectrum
%
%  Status: COMPLETE — PASS-READY (MONO007 referee-ready architecture)
%  Inputs: Chapter01_Difference.tex, Chapter02_Eta.tex,
%          Chapter03_Actualization.tex, Chapter04_ClosureAndMinimalTheory.tex,
%          Chapter05_InevitableSpectrum.tex
%  Method: assembly only — canonical results only, no new physics, no speculation
%  Citations: QG155 (symmetry origin), QG156 (unified access law),
%             QG157 (effective access counts), QG158 (moment orders),
%             QG159 (D96 selection), QG295 (spectrum necessity)
%
%  Usage: \input{Chapter06_D96Spectrum.tex} after the preamble
%         that defines the theorem environments (or include via the master file).
% ======================================================================

% ---- Theorem environments (define in the master preamble if not present) ----
% \newtheorem{definition}{Definition}[chapter]
% \newtheorem{lemma}[definition]{Lemma}
% \newtheorem{theorem}[definition]{Theorem}
% \newtheorem{corollary}[definition]{Corollary}
% \newtheorem{remark}[definition]{Remark}

\chapter{The D96 Spectrum}
\label{c06:chap:d96}

\section{Introduction}
\label{c06:sec:d96-introduction}

In Chapter~\ref{c05:chap:spectrum} we established that the spectrum of the theory is the
inevitable output of the actualization attractor: the unique Laplacian eigenspectrum of the
converged $N=96$ network, carrying no choice and presupposing no primitive status
(Theorem~\ref{c05:thm:spectrum-not-primitive}). This chapter
develops the concrete structure of that spectrum --- the \emph{D96 spectrum}.

We establish the structural content of the D96 spectrum in canonical form: the mode structure
(one zero mode, 95 positive modes); the spectral moments $\Sigma m$,
$\Sigma\sqrt{m}$, $\Sigma m^2$; the occupancy moment $\mathrm{occMom}$; the span; and the
multiplicity structure $[42\times2,\,5,\,6]$ \cite{c06:QG155,c06:QG156,c06:QG157,c06:QG158}. The D96 spectrum is
presented throughout as an \emph{emergent} structure --- the derived output of the
actualization attractor of Chapter~\ref{c05:chap:spectrum} --- never as a primitive: all
constants are derived from the spectrum, which is itself derived from the process
\cite{c06:QG159,c06:QG295}.

\section{The Spectrum Structure}
\label{c06:sec:spectrum-structure}

\begin{theorem}[The D96 spectrum is the derived output]
\label{c06:thm:d96-derived}
The \emph{D96 spectrum} is the Laplacian eigenspectrum of the converged $N=96$ attractor:
the set of $96$ eigenvalues of the graph Laplacian, consisting of one zero eigenvalue (the
background mode) and $95$ positive eigenvalues (the positive modes)
\cite{c06:QG155,c06:QG295}. It is the derived output of the actualization process, not a
primitive input: it admits no primitive interpretation and is not an underivable input.
\end{theorem}

\begin{proof}
By Theorem~\ref{c05:thm:spectrum-eigenspectrum} of Chapter~\ref{c05:chap:spectrum}, the spectrum is
the Laplacian eigenspectrum of the converged $N=96$ network, and by
Theorem~\ref{c05:thm:spectrum-not-primitive} of the same chapter it is not a primitive but the
inevitable output of the attractor. The spectrum-necessity audit \cite{c06:QG295} confirms
this classification explicitly --- the spectrum carries no choice and presupposes the
actualization process and its primitive Difference --- and no primitive interpretation of it
is consistent with the canonical architecture.
\end{proof}

\section{The Zero Mode}
\label{c06:sec:zero-mode}

\begin{theorem}[The zero mode is the background]
\label{c06:thm:zero-mode}
The \emph{zero mode} is the eigenvalue $\lambda_0 = 0$ of the graph Laplacian: the D96
spectrum has exactly one zero eigenvalue, corresponding to the background ---
the uniform configuration from which the count is measured.
\end{theorem}

\begin{proof}
The graph Laplacian of a connected network has exactly one zero eigenvalue, with the
constant eigenvector --- the uniform background, which the spectrum-necessity audit
\cite{c06:QG295} identifies as the background. By Definition~\ref{c01:def:difference} of
Chapter~\ref{c01:chap:difference}, Difference is the counting difference from a uniform
background; hence the zero mode is the reference against which Difference is measured, from
which the positive-mode structure represents the differences.
\end{proof}

\section{Positive Modes}
\label{c06:sec:positive-modes}

\begin{theorem}[There are 95 positive modes]
\label{c06:thm:95-modes}
The \emph{positive modes} are the $95$ positive eigenvalues
$\lambda_1,\dots,\lambda_{95}$ of the graph Laplacian: the non-background degrees of
freedom of the spectrum, which consists of exactly one zero mode and $95$ positive
eigenvalues.
\end{theorem}

\begin{proof}
The graph Laplacian of the $N=96$ network is a $96\times96$ matrix with one zero eigenvalue
(Theorem~\ref{c06:thm:zero-mode}) and the remaining $95$ eigenvalues positive; the
doublet-multiplicity structure confirms the count, $\Sigma m = 95$ being the total mode
count of the positive spectrum \cite{c06:QG157}, and the total dimension $96$ matches the
$N=96$ dimension of the attractor \cite{c06:QG155}.
\end{proof}

\section{Spectral Moments}
\label{c06:sec:spectral-moments}

\begin{theorem}[The moment ladder]
\label{c06:cor:moment-ladder}
The \emph{$p$-moment} of the multiplicity distribution is
\begin{equation}
\label{c06:eq:moment}
\Sigma m^p \;=\; \sum_i m_i^p,
\end{equation}
where the sum runs over the degenerate groups of the spectrum, with multiplicities $m_i$.
The spectral moments form a ladder of access counts: the first moment is the total mode
count, $\Sigma m = 95$; the half-moment is the neutral-sector access count,
$\Sigma\sqrt{m} = 64.08$; and the second moment is the doublet-occupancy access count,
$\Sigma m^2 = 229$ --- each an access count of a different sector of the theory
\cite{c06:QG157,c06:QG158}.
\end{theorem}

\begin{proof}
The effective-access-count program \cite{c06:QG157} establishes that the first moment of the
doublet-multiplicity distribution is the total mode count; by
Theorem~\ref{c06:thm:95-modes} this is $95$, so $\Sigma m = \sum_i m_i = 95$ --- the full
positive spectrum accessed uniformly. The same program establishes that the half-moment
$\Sigma\sqrt{m} = \sum_i \sqrt{m_i}$ is the neutral access count --- the statistical weight
each degenerate group contributes to the neutral channel --- and the second moment
$\Sigma m^2 = \sum_i m_i^2$ the doublet-occupancy access count. With the multiplicity
structure $[42\times2,5,6]$, $\Sigma\sqrt{m} = 64.08$ and
$\Sigma m^2 = 42\cdot 2^2 + 5^2 + 6^2 = 168+25+36 = 229$; the moment-orders program
\cite{c06:QG158} establishes that the moments form a Z2-power ladder, each corresponding to a
distinct sector access.
\end{proof}

\section{Spectral Constants}
\label{c06:sec:spectral-constants}

\begin{theorem}[The occupancy moment]
\label{c06:thm:occmom}
The \emph{octave-occupation moment} is
\begin{equation}
\label{c06:eq:occmom}
\mathrm{occMom} \;=\; \frac{\sum_{\text{octaves}} \mathrm{occ}_j^2}{\mathrm{occ}_0},
\end{equation}
where $\mathrm{occ}_j$ are the octave occupancies and $\mathrm{occ}_0$ the lightest-octave
occupancy. The octave-occupation moment of the D96 spectrum is
$\mathrm{occMom} = 1900.25$: the up-sector (dense-band) access count.
\end{theorem}

\begin{proof}
The effective-access-count program \cite{c06:QG157} establishes that the octave-occupation
moment $\mathrm{occMom} = \sum_j \mathrm{occ}_j^2/\mathrm{occ}_0$ is the dense-band access
count of the up sector; the octave occupancies of the D96 spectrum yield
$\mathrm{occMom} = 1900.25$.
\end{proof}

\begin{theorem}[The span]
\label{c06:thm:span}
The span of the D96 spectrum is
\begin{equation}
\label{c06:eq:span}
\mathrm{span} \;=\; \frac{\omega_{\max}}{\omega_{\min}} \;=\; 6.40,
\end{equation}
the ratio of the highest to the lowest positive frequency, with the modes
$\omega_k=\sqrt{\lambda_k}$ given by the Laplacian eigenvalues
$\lambda_k = 2\sum_{d=1}^{6}(1-\cos(2\pi dk/96))$ of the canonical attractor graph
$C_{96}(\pm1,\ldots,\pm6)$.
\end{theorem}

\begin{proof}
The span is the ratio of the extremal positive modes, $\omega_{\max}$ and $\omega_{\min}$;
the D96 spectrum yields the numerical value $\mathrm{span} = 6.40$, established by the
spectral-structure program \cite{c06:QG155}.
\end{proof}

\begin{corollary}[All constants are derived]
\label{c06:cor:constants-derived}
The spectral constants --- $\Sigma m=95$, $\Sigma\sqrt{m}=64.08$, $\Sigma m^2=229$,
$\mathrm{occMom}=1900.25$, $\mathrm{span}=6.40$ --- are all derived from the D96 spectrum,
which is itself derived from the attractor.
\end{corollary}

\begin{proof}
Each constant is a moment or structural feature of the spectrum
(Theorem~\ref{c06:cor:moment-ladder}, Theorem~\ref{c06:thm:occmom}, Theorem~\ref{c06:thm:span});
by Theorem~\ref{c06:thm:d96-derived}, the spectrum is the derived output of the attractor.
Hence every constant is twice derived --- a function of the spectrum, which is a function of
the process --- and no constant is a primitive input \cite{c06:QG157,c06:QG295}.
\end{proof}

\section{Multiplicity Structure}
\label{c06:sec:multiplicity}

\begin{theorem}[The multiplicity structure]
\label{c06:thm:multiplicity}
The \emph{multiplicity structure} of the D96 spectrum is the multiset of degenerate-group
sizes $m_i$: the pattern of equal eigenvalues grouped into the doublet structure,
$[42\times2,\,5,\,6]$ --- forty-two doublets, one group of size $5$, and one group of size
$6$. The doublet is the dominant multiplicity: $42$ of the $44$ groups are of size $2$,
reflecting the Z2 pairing of the spectrum generated by the D96 symmetry.
\end{theorem}

\begin{proof}
The doublet-multiplicity structure is established by the effective-access-count program
\cite{c06:QG157}: the spectrum is a multiset of degenerate groups with multiplicities $m_i$,
consisting of $42$ groups of size $2$, one of size $5$, and one of size $6$. The Z2 doublet
structure is generated by the D96 symmetry \cite{c06:QG155} --- the observable-sector
spectrum is fully Z2 paired --- and the total satisfies
$\Sigma m = 42\cdot2 + 5 + 6 = 95$, consistent with Theorem~\ref{c06:cor:moment-ladder}.
\end{proof}

\section{Consequences for Organization}
\label{c06:sec:organization}

The D96 spectrum is the substrate of the organization structure: its degeneracy groups,
moments, and span determine the organizational reads of the theory --- degeneracy (from the
multiplicity groups), compression and beat (from the span and octave structure), and
locking (from the spectral gap)
(Theorem~\ref{c06:thm:multiplicity}, Theorem~\ref{c06:thm:span}).

\begin{theorem}[The organization constants are derived]
\label{c06:thm:organization-derived}
The organization constants of the theory --- the occupancy moment, the span, and the
degeneracy structure --- are derived from the D96 spectrum, carrying no independent input.
The organization structure is therefore emergent: it is a read of the D96 spectrum and
carries no primitive status.
\end{theorem}

\begin{proof}
By Corollary~\ref{c06:cor:constants-derived}, the spectral constants are derived from the
spectrum, which is derived from the attractor; the organization constants are functions of
these spectral constants, so the organization structure inherits the derived status: every
organizational constant is an output of the spectrum, itself an output of the process; no
organizational feature is primitive \cite{c06:QG295}.
\end{proof}

Consistent with MONO007 \cite{c06:MONO007}, the D96 spectrum is presented throughout this
chapter as the derived output of the actualization attractor, never as a primitive: its
constants and multiplicities are all derived from the spectrum, which is itself derived.
This removes any possible reading of the D96 structure as an unexplained input.

\section{Summary}
\label{c06:sec:d96-summary}

This chapter has established the concrete structure of the inevitable spectrum. The D96
spectrum is the Laplacian eigenspectrum of the attractor (Theorem~\ref{c06:thm:d96-derived}):
one zero mode, the background that Difference is counted against
(Theorem~\ref{c06:thm:zero-mode}), and $95$ positive modes (Theorem~\ref{c06:thm:95-modes}).
Its moments $\Sigma m=95$, $\Sigma\sqrt{m}=64.08$, $\Sigma m^2=229$ form the access-count
ladder (Theorem~\ref{c06:cor:moment-ladder}); its constants $\mathrm{occMom}=1900.25$ and
$\mathrm{span}=6.40$ (Theorem~\ref{c06:thm:occmom}, Theorem~\ref{c06:thm:span}) and its
multiplicities $[42\times2,5,6]$ (Theorem~\ref{c06:thm:multiplicity}) are all derived from
the spectrum, which is itself derived from the actualization attractor
(Corollary~\ref{c06:cor:constants-derived}). The spectrum is therefore the substrate of the
organization structure (Theorem~\ref{c06:thm:organization-derived}); the following chapters
read the organization structure and then the physics sectors from this one spectrum.

\begin{thebibliography}{99}
\bibitem{c06:QG155} AT-QG Phase 155, \emph{Symmetry Origin} (Z2 doublet structure), Docs/Research.
\bibitem{c06:QG156} AT-QG Phase 156, \emph{Unified Spectral Access Law}, Docs/Research.
\bibitem{c06:QG157} AT-QG Phase 157, \emph{Effective Access Counts}, Docs/Research.
\bibitem{c06:QG158} AT-QG Phase 158, \emph{Moment Orders}, Docs/Research.
\bibitem{c06:QG159} AT-QG Phase 159, \emph{D96 Selection Origin}, Docs/Research.
\bibitem{c06:QG295} AT-QG Phase 295, \emph{Spectrum Necessity Audit}, Docs/Research.
\bibitem{c06:MONO007} AT-MONO007, \emph{Referee-Readiness Resolution}, Docs/Zenodo/Monograph.
\end{thebibliography}
 after the preamble
%         that defines the theorem environments (or include via the master file).
% ======================================================================

% ---- Theorem environments (define in the master preamble if not present) ----
% \newtheorem{definition}{Definition}[chapter]
% \newtheorem{lemma}[definition]{Lemma}
% \newtheorem{theorem}[definition]{Theorem}
% \newtheorem{corollary}[definition]{Corollary}
% \newtheorem{remark}[definition]{Remark}

\chapter{The D96 Spectrum}
\label{c06:chap:d96}

\section{Introduction}
\label{c06:sec:d96-introduction}

In Chapter~\ref{c05:chap:spectrum} we established that the spectrum of the theory is the
inevitable output of the actualization attractor: the unique Laplacian eigenspectrum of the
converged $N=96$ network, carrying no choice and presupposing no primitive status
(Theorem~\ref{c05:thm:spectrum-not-primitive}). This chapter
develops the concrete structure of that spectrum --- the \emph{D96 spectrum}.

We establish the structural content of the D96 spectrum in canonical form: the mode structure
(one zero mode, 95 positive modes); the spectral moments $\Sigma m$,
$\Sigma\sqrt{m}$, $\Sigma m^2$; the occupancy moment $\mathrm{occMom}$; the span; and the
multiplicity structure $[42\times2,\,5,\,6]$ \cite{c06:QG155,c06:QG156,c06:QG157,c06:QG158}. The D96 spectrum is
presented throughout as an \emph{emergent} structure --- the derived output of the
actualization attractor of Chapter~\ref{c05:chap:spectrum} --- never as a primitive: all
constants are derived from the spectrum, which is itself derived from the process
\cite{c06:QG159,c06:QG295}.

\section{The Spectrum Structure}
\label{c06:sec:spectrum-structure}

\begin{theorem}[The D96 spectrum is the derived output]
\label{c06:thm:d96-derived}
The \emph{D96 spectrum} is the Laplacian eigenspectrum of the converged $N=96$ attractor:
the set of $96$ eigenvalues of the graph Laplacian, consisting of one zero eigenvalue (the
background mode) and $95$ positive eigenvalues (the positive modes)
\cite{c06:QG155,c06:QG295}. It is the derived output of the actualization process, not a
primitive input: it admits no primitive interpretation and is not an underivable input.
\end{theorem}

\begin{proof}
By Theorem~\ref{c05:thm:spectrum-eigenspectrum} of Chapter~\ref{c05:chap:spectrum}, the spectrum is
the Laplacian eigenspectrum of the converged $N=96$ network, and by
Theorem~\ref{c05:thm:spectrum-not-primitive} of the same chapter it is not a primitive but the
inevitable output of the attractor. The spectrum-necessity audit \cite{c06:QG295} confirms
this classification explicitly --- the spectrum carries no choice and presupposes the
actualization process and its primitive Difference --- and no primitive interpretation of it
is consistent with the canonical architecture.
\end{proof}

\section{The Zero Mode}
\label{c06:sec:zero-mode}

\begin{theorem}[The zero mode is the background]
\label{c06:thm:zero-mode}
The \emph{zero mode} is the eigenvalue $\lambda_0 = 0$ of the graph Laplacian: the D96
spectrum has exactly one zero eigenvalue, corresponding to the background ---
the uniform configuration from which the count is measured.
\end{theorem}

\begin{proof}
The graph Laplacian of a connected network has exactly one zero eigenvalue, with the
constant eigenvector --- the uniform background, which the spectrum-necessity audit
\cite{c06:QG295} identifies as the background. By Definition~\ref{c01:def:difference} of
Chapter~\ref{c01:chap:difference}, Difference is the counting difference from a uniform
background; hence the zero mode is the reference against which Difference is measured, from
which the positive-mode structure represents the differences.
\end{proof}

\section{Positive Modes}
\label{c06:sec:positive-modes}

\begin{theorem}[There are 95 positive modes]
\label{c06:thm:95-modes}
The \emph{positive modes} are the $95$ positive eigenvalues
$\lambda_1,\dots,\lambda_{95}$ of the graph Laplacian: the non-background degrees of
freedom of the spectrum, which consists of exactly one zero mode and $95$ positive
eigenvalues.
\end{theorem}

\begin{proof}
The graph Laplacian of the $N=96$ network is a $96\times96$ matrix with one zero eigenvalue
(Theorem~\ref{c06:thm:zero-mode}) and the remaining $95$ eigenvalues positive; the
doublet-multiplicity structure confirms the count, $\Sigma m = 95$ being the total mode
count of the positive spectrum \cite{c06:QG157}, and the total dimension $96$ matches the
$N=96$ dimension of the attractor \cite{c06:QG155}.
\end{proof}

\section{Spectral Moments}
\label{c06:sec:spectral-moments}

\begin{theorem}[The moment ladder]
\label{c06:cor:moment-ladder}
The \emph{$p$-moment} of the multiplicity distribution is
\begin{equation}
\label{c06:eq:moment}
\Sigma m^p \;=\; \sum_i m_i^p,
\end{equation}
where the sum runs over the degenerate groups of the spectrum, with multiplicities $m_i$.
The spectral moments form a ladder of access counts: the first moment is the total mode
count, $\Sigma m = 95$; the half-moment is the neutral-sector access count,
$\Sigma\sqrt{m} = 64.08$; and the second moment is the doublet-occupancy access count,
$\Sigma m^2 = 229$ --- each an access count of a different sector of the theory
\cite{c06:QG157,c06:QG158}.
\end{theorem}

\begin{proof}
The effective-access-count program \cite{c06:QG157} establishes that the first moment of the
doublet-multiplicity distribution is the total mode count; by
Theorem~\ref{c06:thm:95-modes} this is $95$, so $\Sigma m = \sum_i m_i = 95$ --- the full
positive spectrum accessed uniformly. The same program establishes that the half-moment
$\Sigma\sqrt{m} = \sum_i \sqrt{m_i}$ is the neutral access count --- the statistical weight
each degenerate group contributes to the neutral channel --- and the second moment
$\Sigma m^2 = \sum_i m_i^2$ the doublet-occupancy access count. With the multiplicity
structure $[42\times2,5,6]$, $\Sigma\sqrt{m} = 64.08$ and
$\Sigma m^2 = 42\cdot 2^2 + 5^2 + 6^2 = 168+25+36 = 229$; the moment-orders program
\cite{c06:QG158} establishes that the moments form a Z2-power ladder, each corresponding to a
distinct sector access.
\end{proof}

\section{Spectral Constants}
\label{c06:sec:spectral-constants}

\begin{theorem}[The occupancy moment]
\label{c06:thm:occmom}
The \emph{octave-occupation moment} is
\begin{equation}
\label{c06:eq:occmom}
\mathrm{occMom} \;=\; \frac{\sum_{\text{octaves}} \mathrm{occ}_j^2}{\mathrm{occ}_0},
\end{equation}
where $\mathrm{occ}_j$ are the octave occupancies and $\mathrm{occ}_0$ the lightest-octave
occupancy. The octave-occupation moment of the D96 spectrum is
$\mathrm{occMom} = 1900.25$: the up-sector (dense-band) access count.
\end{theorem}

\begin{proof}
The effective-access-count program \cite{c06:QG157} establishes that the octave-occupation
moment $\mathrm{occMom} = \sum_j \mathrm{occ}_j^2/\mathrm{occ}_0$ is the dense-band access
count of the up sector; the octave occupancies of the D96 spectrum yield
$\mathrm{occMom} = 1900.25$.
\end{proof}

\begin{theorem}[The span]
\label{c06:thm:span}
The span of the D96 spectrum is
\begin{equation}
\label{c06:eq:span}
\mathrm{span} \;=\; \frac{\omega_{\max}}{\omega_{\min}} \;=\; 6.40,
\end{equation}
the ratio of the highest to the lowest positive frequency, with the modes
$\omega_k=\sqrt{\lambda_k}$ given by the Laplacian eigenvalues
$\lambda_k = 2\sum_{d=1}^{6}(1-\cos(2\pi dk/96))$ of the canonical attractor graph
$C_{96}(\pm1,\ldots,\pm6)$.
\end{theorem}

\begin{proof}
The span is the ratio of the extremal positive modes, $\omega_{\max}$ and $\omega_{\min}$;
the D96 spectrum yields the numerical value $\mathrm{span} = 6.40$, established by the
spectral-structure program \cite{c06:QG155}.
\end{proof}

\begin{corollary}[All constants are derived]
\label{c06:cor:constants-derived}
The spectral constants --- $\Sigma m=95$, $\Sigma\sqrt{m}=64.08$, $\Sigma m^2=229$,
$\mathrm{occMom}=1900.25$, $\mathrm{span}=6.40$ --- are all derived from the D96 spectrum,
which is itself derived from the attractor.
\end{corollary}

\begin{proof}
Each constant is a moment or structural feature of the spectrum
(Theorem~\ref{c06:cor:moment-ladder}, Theorem~\ref{c06:thm:occmom}, Theorem~\ref{c06:thm:span});
by Theorem~\ref{c06:thm:d96-derived}, the spectrum is the derived output of the attractor.
Hence every constant is twice derived --- a function of the spectrum, which is a function of
the process --- and no constant is a primitive input \cite{c06:QG157,c06:QG295}.
\end{proof}

\section{Multiplicity Structure}
\label{c06:sec:multiplicity}

\begin{theorem}[The multiplicity structure]
\label{c06:thm:multiplicity}
The \emph{multiplicity structure} of the D96 spectrum is the multiset of degenerate-group
sizes $m_i$: the pattern of equal eigenvalues grouped into the doublet structure,
$[42\times2,\,5,\,6]$ --- forty-two doublets, one group of size $5$, and one group of size
$6$. The doublet is the dominant multiplicity: $42$ of the $44$ groups are of size $2$,
reflecting the Z2 pairing of the spectrum generated by the D96 symmetry.
\end{theorem}

\begin{proof}
The doublet-multiplicity structure is established by the effective-access-count program
\cite{c06:QG157}: the spectrum is a multiset of degenerate groups with multiplicities $m_i$,
consisting of $42$ groups of size $2$, one of size $5$, and one of size $6$. The Z2 doublet
structure is generated by the D96 symmetry \cite{c06:QG155} --- the observable-sector
spectrum is fully Z2 paired --- and the total satisfies
$\Sigma m = 42\cdot2 + 5 + 6 = 95$, consistent with Theorem~\ref{c06:cor:moment-ladder}.
\end{proof}

\section{Consequences for Organization}
\label{c06:sec:organization}

The D96 spectrum is the substrate of the organization structure: its degeneracy groups,
moments, and span determine the organizational reads of the theory --- degeneracy (from the
multiplicity groups), compression and beat (from the span and octave structure), and
locking (from the spectral gap)
(Theorem~\ref{c06:thm:multiplicity}, Theorem~\ref{c06:thm:span}).

\begin{theorem}[The organization constants are derived]
\label{c06:thm:organization-derived}
The organization constants of the theory --- the occupancy moment, the span, and the
degeneracy structure --- are derived from the D96 spectrum, carrying no independent input.
The organization structure is therefore emergent: it is a read of the D96 spectrum and
carries no primitive status.
\end{theorem}

\begin{proof}
By Corollary~\ref{c06:cor:constants-derived}, the spectral constants are derived from the
spectrum, which is derived from the attractor; the organization constants are functions of
these spectral constants, so the organization structure inherits the derived status: every
organizational constant is an output of the spectrum, itself an output of the process; no
organizational feature is primitive \cite{c06:QG295}.
\end{proof}

Consistent with MONO007 \cite{c06:MONO007}, the D96 spectrum is presented throughout this
chapter as the derived output of the actualization attractor, never as a primitive: its
constants and multiplicities are all derived from the spectrum, which is itself derived.
This removes any possible reading of the D96 structure as an unexplained input.

\section{Summary}
\label{c06:sec:d96-summary}

This chapter has established the concrete structure of the inevitable spectrum. The D96
spectrum is the Laplacian eigenspectrum of the attractor (Theorem~\ref{c06:thm:d96-derived}):
one zero mode, the background that Difference is counted against
(Theorem~\ref{c06:thm:zero-mode}), and $95$ positive modes (Theorem~\ref{c06:thm:95-modes}).
Its moments $\Sigma m=95$, $\Sigma\sqrt{m}=64.08$, $\Sigma m^2=229$ form the access-count
ladder (Theorem~\ref{c06:cor:moment-ladder}); its constants $\mathrm{occMom}=1900.25$ and
$\mathrm{span}=6.40$ (Theorem~\ref{c06:thm:occmom}, Theorem~\ref{c06:thm:span}) and its
multiplicities $[42\times2,5,6]$ (Theorem~\ref{c06:thm:multiplicity}) are all derived from
the spectrum, which is itself derived from the actualization attractor
(Corollary~\ref{c06:cor:constants-derived}). The spectrum is therefore the substrate of the
organization structure (Theorem~\ref{c06:thm:organization-derived}); the following chapters
read the organization structure and then the physics sectors from this one spectrum.

\begin{thebibliography}{99}
\bibitem{c06:QG155} AT-QG Phase 155, \emph{Symmetry Origin} (Z2 doublet structure), Docs/Research.
\bibitem{c06:QG156} AT-QG Phase 156, \emph{Unified Spectral Access Law}, Docs/Research.
\bibitem{c06:QG157} AT-QG Phase 157, \emph{Effective Access Counts}, Docs/Research.
\bibitem{c06:QG158} AT-QG Phase 158, \emph{Moment Orders}, Docs/Research.
\bibitem{c06:QG159} AT-QG Phase 159, \emph{D96 Selection Origin}, Docs/Research.
\bibitem{c06:QG295} AT-QG Phase 295, \emph{Spectrum Necessity Audit}, Docs/Research.
\bibitem{c06:MONO007} AT-MONO007, \emph{Referee-Readiness Resolution}, Docs/Zenodo/Monograph.
\end{thebibliography}
 after the preamble
%         that defines the theorem environments (or include via the master file).
% ======================================================================

% ---- Theorem environments (define in the master preamble if not present) ----
% \newtheorem{definition}{Definition}[chapter]
% \newtheorem{lemma}[definition]{Lemma}
% \newtheorem{theorem}[definition]{Theorem}
% \newtheorem{corollary}[definition]{Corollary}
% \newtheorem{remark}[definition]{Remark}

\chapter{The D96 Spectrum}
\label{c06:chap:d96}

\section{Introduction}
\label{c06:sec:d96-introduction}

In Chapter~\ref{c05:chap:spectrum} we established that the spectrum of the theory is the
inevitable output of the actualization attractor: the unique Laplacian eigenspectrum of the
converged $N=96$ network, carrying no choice and presupposing no primitive status
(Theorem~\ref{c05:thm:spectrum-not-primitive}). This chapter
develops the concrete structure of that spectrum --- the \emph{D96 spectrum}.

We establish the structural content of the D96 spectrum in canonical form: the mode structure
(one zero mode, 95 positive modes); the spectral moments $\Sigma m$,
$\Sigma\sqrt{m}$, $\Sigma m^2$; the occupancy moment $\mathrm{occMom}$; the span; and the
multiplicity structure $[42\times2,\,5,\,6]$ \cite{c06:QG155,c06:QG156,c06:QG157,c06:QG158}. The D96 spectrum is
presented throughout as an \emph{emergent} structure --- the derived output of the
actualization attractor of Chapter~\ref{c05:chap:spectrum} --- never as a primitive: all
constants are derived from the spectrum, which is itself derived from the process
\cite{c06:QG159,c06:QG295}.

\section{The Spectrum Structure}
\label{c06:sec:spectrum-structure}

\begin{theorem}[The D96 spectrum is the derived output]
\label{c06:thm:d96-derived}
The \emph{D96 spectrum} is the Laplacian eigenspectrum of the converged $N=96$ attractor:
the set of $96$ eigenvalues of the graph Laplacian, consisting of one zero eigenvalue (the
background mode) and $95$ positive eigenvalues (the positive modes)
\cite{c06:QG155,c06:QG295}. It is the derived output of the actualization process, not a
primitive input: it admits no primitive interpretation and is not an underivable input.
\end{theorem}

\begin{proof}
By Theorem~\ref{c05:thm:spectrum-eigenspectrum} of Chapter~\ref{c05:chap:spectrum}, the spectrum is
the Laplacian eigenspectrum of the converged $N=96$ network, and by
Theorem~\ref{c05:thm:spectrum-not-primitive} of the same chapter it is not a primitive but the
inevitable output of the attractor. The spectrum-necessity audit \cite{c06:QG295} confirms
this classification explicitly --- the spectrum carries no choice and presupposes the
actualization process and its primitive Difference --- and no primitive interpretation of it
is consistent with the canonical architecture.
\end{proof}

\section{The Zero Mode}
\label{c06:sec:zero-mode}

\begin{theorem}[The zero mode is the background]
\label{c06:thm:zero-mode}
The \emph{zero mode} is the eigenvalue $\lambda_0 = 0$ of the graph Laplacian: the D96
spectrum has exactly one zero eigenvalue, corresponding to the background ---
the uniform configuration from which the count is measured.
\end{theorem}

\begin{proof}
The graph Laplacian of a connected network has exactly one zero eigenvalue, with the
constant eigenvector --- the uniform background, which the spectrum-necessity audit
\cite{c06:QG295} identifies as the background. By Definition~\ref{c01:def:difference} of
Chapter~\ref{c01:chap:difference}, Difference is the counting difference from a uniform
background; hence the zero mode is the reference against which Difference is measured, from
which the positive-mode structure represents the differences.
\end{proof}

\section{Positive Modes}
\label{c06:sec:positive-modes}

\begin{theorem}[There are 95 positive modes]
\label{c06:thm:95-modes}
The \emph{positive modes} are the $95$ positive eigenvalues
$\lambda_1,\dots,\lambda_{95}$ of the graph Laplacian: the non-background degrees of
freedom of the spectrum, which consists of exactly one zero mode and $95$ positive
eigenvalues.
\end{theorem}

\begin{proof}
The graph Laplacian of the $N=96$ network is a $96\times96$ matrix with one zero eigenvalue
(Theorem~\ref{c06:thm:zero-mode}) and the remaining $95$ eigenvalues positive; the
doublet-multiplicity structure confirms the count, $\Sigma m = 95$ being the total mode
count of the positive spectrum \cite{c06:QG157}, and the total dimension $96$ matches the
$N=96$ dimension of the attractor \cite{c06:QG155}.
\end{proof}

\section{Spectral Moments}
\label{c06:sec:spectral-moments}

\begin{theorem}[The moment ladder]
\label{c06:cor:moment-ladder}
The \emph{$p$-moment} of the multiplicity distribution is
\begin{equation}
\label{c06:eq:moment}
\Sigma m^p \;=\; \sum_i m_i^p,
\end{equation}
where the sum runs over the degenerate groups of the spectrum, with multiplicities $m_i$.
The spectral moments form a ladder of access counts: the first moment is the total mode
count, $\Sigma m = 95$; the half-moment is the neutral-sector access count,
$\Sigma\sqrt{m} = 64.08$; and the second moment is the doublet-occupancy access count,
$\Sigma m^2 = 229$ --- each an access count of a different sector of the theory
\cite{c06:QG157,c06:QG158}.
\end{theorem}

\begin{proof}
The effective-access-count program \cite{c06:QG157} establishes that the first moment of the
doublet-multiplicity distribution is the total mode count; by
Theorem~\ref{c06:thm:95-modes} this is $95$, so $\Sigma m = \sum_i m_i = 95$ --- the full
positive spectrum accessed uniformly. The same program establishes that the half-moment
$\Sigma\sqrt{m} = \sum_i \sqrt{m_i}$ is the neutral access count --- the statistical weight
each degenerate group contributes to the neutral channel --- and the second moment
$\Sigma m^2 = \sum_i m_i^2$ the doublet-occupancy access count. With the multiplicity
structure $[42\times2,5,6]$, $\Sigma\sqrt{m} = 64.08$ and
$\Sigma m^2 = 42\cdot 2^2 + 5^2 + 6^2 = 168+25+36 = 229$; the moment-orders program
\cite{c06:QG158} establishes that the moments form a Z2-power ladder, each corresponding to a
distinct sector access.
\end{proof}

\section{Spectral Constants}
\label{c06:sec:spectral-constants}

\begin{theorem}[The occupancy moment]
\label{c06:thm:occmom}
The \emph{octave-occupation moment} is
\begin{equation}
\label{c06:eq:occmom}
\mathrm{occMom} \;=\; \frac{\sum_{\text{octaves}} \mathrm{occ}_j^2}{\mathrm{occ}_0},
\end{equation}
where $\mathrm{occ}_j$ are the octave occupancies and $\mathrm{occ}_0$ the lightest-octave
occupancy. The octave-occupation moment of the D96 spectrum is
$\mathrm{occMom} = 1900.25$: the up-sector (dense-band) access count.
\end{theorem}

\begin{proof}
The effective-access-count program \cite{c06:QG157} establishes that the octave-occupation
moment $\mathrm{occMom} = \sum_j \mathrm{occ}_j^2/\mathrm{occ}_0$ is the dense-band access
count of the up sector; the octave occupancies of the D96 spectrum yield
$\mathrm{occMom} = 1900.25$.
\end{proof}

\begin{theorem}[The span]
\label{c06:thm:span}
The span of the D96 spectrum is
\begin{equation}
\label{c06:eq:span}
\mathrm{span} \;=\; \frac{\omega_{\max}}{\omega_{\min}} \;=\; 6.40,
\end{equation}
the ratio of the highest to the lowest positive frequency, with the modes
$\omega_k=\sqrt{\lambda_k}$ given by the Laplacian eigenvalues
$\lambda_k = 2\sum_{d=1}^{6}(1-\cos(2\pi dk/96))$ of the canonical attractor graph
$C_{96}(\pm1,\ldots,\pm6)$.
\end{theorem}

\begin{proof}
The span is the ratio of the extremal positive modes, $\omega_{\max}$ and $\omega_{\min}$;
the D96 spectrum yields the numerical value $\mathrm{span} = 6.40$, established by the
spectral-structure program \cite{c06:QG155}.
\end{proof}

\begin{corollary}[All constants are derived]
\label{c06:cor:constants-derived}
The spectral constants --- $\Sigma m=95$, $\Sigma\sqrt{m}=64.08$, $\Sigma m^2=229$,
$\mathrm{occMom}=1900.25$, $\mathrm{span}=6.40$ --- are all derived from the D96 spectrum,
which is itself derived from the attractor.
\end{corollary}

\begin{proof}
Each constant is a moment or structural feature of the spectrum
(Theorem~\ref{c06:cor:moment-ladder}, Theorem~\ref{c06:thm:occmom}, Theorem~\ref{c06:thm:span});
by Theorem~\ref{c06:thm:d96-derived}, the spectrum is the derived output of the attractor.
Hence every constant is twice derived --- a function of the spectrum, which is a function of
the process --- and no constant is a primitive input \cite{c06:QG157,c06:QG295}.
\end{proof}

\section{Multiplicity Structure}
\label{c06:sec:multiplicity}

\begin{theorem}[The multiplicity structure]
\label{c06:thm:multiplicity}
The \emph{multiplicity structure} of the D96 spectrum is the multiset of degenerate-group
sizes $m_i$: the pattern of equal eigenvalues grouped into the doublet structure,
$[42\times2,\,5,\,6]$ --- forty-two doublets, one group of size $5$, and one group of size
$6$. The doublet is the dominant multiplicity: $42$ of the $44$ groups are of size $2$,
reflecting the Z2 pairing of the spectrum generated by the D96 symmetry.
\end{theorem}

\begin{proof}
The doublet-multiplicity structure is established by the effective-access-count program
\cite{c06:QG157}: the spectrum is a multiset of degenerate groups with multiplicities $m_i$,
consisting of $42$ groups of size $2$, one of size $5$, and one of size $6$. The Z2 doublet
structure is generated by the D96 symmetry \cite{c06:QG155} --- the observable-sector
spectrum is fully Z2 paired --- and the total satisfies
$\Sigma m = 42\cdot2 + 5 + 6 = 95$, consistent with Theorem~\ref{c06:cor:moment-ladder}.
\end{proof}

\section{Consequences for Organization}
\label{c06:sec:organization}

The D96 spectrum is the substrate of the organization structure: its degeneracy groups,
moments, and span determine the organizational reads of the theory --- degeneracy (from the
multiplicity groups), compression and beat (from the span and octave structure), and
locking (from the spectral gap)
(Theorem~\ref{c06:thm:multiplicity}, Theorem~\ref{c06:thm:span}).

\begin{theorem}[The organization constants are derived]
\label{c06:thm:organization-derived}
The organization constants of the theory --- the occupancy moment, the span, and the
degeneracy structure --- are derived from the D96 spectrum, carrying no independent input.
The organization structure is therefore emergent: it is a read of the D96 spectrum and
carries no primitive status.
\end{theorem}

\begin{proof}
By Corollary~\ref{c06:cor:constants-derived}, the spectral constants are derived from the
spectrum, which is derived from the attractor; the organization constants are functions of
these spectral constants, so the organization structure inherits the derived status: every
organizational constant is an output of the spectrum, itself an output of the process; no
organizational feature is primitive \cite{c06:QG295}.
\end{proof}

Consistent with MONO007 \cite{c06:MONO007}, the D96 spectrum is presented throughout this
chapter as the derived output of the actualization attractor, never as a primitive: its
constants and multiplicities are all derived from the spectrum, which is itself derived.
This removes any possible reading of the D96 structure as an unexplained input.

\section{Summary}
\label{c06:sec:d96-summary}

This chapter has established the concrete structure of the inevitable spectrum. The D96
spectrum is the Laplacian eigenspectrum of the attractor (Theorem~\ref{c06:thm:d96-derived}):
one zero mode, the background that Difference is counted against
(Theorem~\ref{c06:thm:zero-mode}), and $95$ positive modes (Theorem~\ref{c06:thm:95-modes}).
Its moments $\Sigma m=95$, $\Sigma\sqrt{m}=64.08$, $\Sigma m^2=229$ form the access-count
ladder (Theorem~\ref{c06:cor:moment-ladder}); its constants $\mathrm{occMom}=1900.25$ and
$\mathrm{span}=6.40$ (Theorem~\ref{c06:thm:occmom}, Theorem~\ref{c06:thm:span}) and its
multiplicities $[42\times2,5,6]$ (Theorem~\ref{c06:thm:multiplicity}) are all derived from
the spectrum, which is itself derived from the actualization attractor
(Corollary~\ref{c06:cor:constants-derived}). The spectrum is therefore the substrate of the
organization structure (Theorem~\ref{c06:thm:organization-derived}); the following chapters
read the organization structure and then the physics sectors from this one spectrum.

\begin{thebibliography}{99}
\bibitem{c06:QG155} AT-QG Phase 155, \emph{Symmetry Origin} (Z2 doublet structure), Docs/Research.
\bibitem{c06:QG156} AT-QG Phase 156, \emph{Unified Spectral Access Law}, Docs/Research.
\bibitem{c06:QG157} AT-QG Phase 157, \emph{Effective Access Counts}, Docs/Research.
\bibitem{c06:QG158} AT-QG Phase 158, \emph{Moment Orders}, Docs/Research.
\bibitem{c06:QG159} AT-QG Phase 159, \emph{D96 Selection Origin}, Docs/Research.
\bibitem{c06:QG295} AT-QG Phase 295, \emph{Spectrum Necessity Audit}, Docs/Research.
\bibitem{c06:MONO007} AT-MONO007, \emph{Referee-Readiness Resolution}, Docs/Zenodo/Monograph.
\end{thebibliography}

% ======================================================================
%  AT-MONO106 — Chapter 7: The Operator Basis
%  Canonical Monograph (v2.0) · The Actualization Theory:
%  A Reconstruction of Physics from Difference, Actualization and Spectrum
%
%  Status: COMPLETE — PASS-READY (MONO007 referee-ready architecture)
%  Inputs: Chapter01_Difference.tex ... Chapter06_D96Spectrum.tex
%  Method: assembly only — canonical results only, no new physics, no speculation
%  Citations: QG260 (resonance layer), QG261 (operator layer),
%             QG262 (same operator sectors), QG263 (single resonance dynamics),
%             QG300 (operator universality), QG302 (cross-domain universality),
%             QG304 (real network universality), QG307 (fifth operator search),
%             QG308 (operator necessity), QG309 (red team audit),
%             QG312 (adversarial spectrum audit)
%
%  Usage: % ======================================================================
%  AT-MONO106 — Chapter 7: The Operator Basis
%  Canonical Monograph (v2.0) · The Actualization Theory:
%  A Reconstruction of Physics from Difference, Actualization and Spectrum
%
%  Status: COMPLETE — PASS-READY (MONO007 referee-ready architecture)
%  Inputs: Chapter01_Difference.tex ... Chapter06_D96Spectrum.tex
%  Method: assembly only — canonical results only, no new physics, no speculation
%  Citations: QG260 (resonance layer), QG261 (operator layer),
%             QG262 (same operator sectors), QG263 (single resonance dynamics),
%             QG300 (operator universality), QG302 (cross-domain universality),
%             QG304 (real network universality), QG307 (fifth operator search),
%             QG308 (operator necessity), QG309 (red team audit),
%             QG312 (adversarial spectrum audit)
%
%  Usage: % ======================================================================
%  AT-MONO106 — Chapter 7: The Operator Basis
%  Canonical Monograph (v2.0) · The Actualization Theory:
%  A Reconstruction of Physics from Difference, Actualization and Spectrum
%
%  Status: COMPLETE — PASS-READY (MONO007 referee-ready architecture)
%  Inputs: Chapter01_Difference.tex ... Chapter06_D96Spectrum.tex
%  Method: assembly only — canonical results only, no new physics, no speculation
%  Citations: QG260 (resonance layer), QG261 (operator layer),
%             QG262 (same operator sectors), QG263 (single resonance dynamics),
%             QG300 (operator universality), QG302 (cross-domain universality),
%             QG304 (real network universality), QG307 (fifth operator search),
%             QG308 (operator necessity), QG309 (red team audit),
%             QG312 (adversarial spectrum audit)
%
%  Usage: \input{Chapter07_OperatorBasis.tex} after the preamble
%         that defines the theorem environments (or include via the master file).
% ======================================================================

% ---- Theorem environments (define in the master preamble if not present) ----
% \newtheorem{definition}{Definition}[chapter]
% \newtheorem{lemma}[definition]{Lemma}
% \newtheorem{theorem}[definition]{Theorem}
% \newtheorem{corollary}[definition]{Corollary}
% \newtheorem{remark}[definition]{Remark}

\chapter{The Operator Basis}
\label{chap:operators}

\section{Introduction}
\label{sec:operators-introduction}

In Chapter~\ref{chap:d96} we established the concrete structure of the D96 spectrum --- its
modes, moments, and multiplicities --- as the emergent output of the actualization
attractor. This chapter develops the \emph{operator basis} of the theory: the set of
spectral projections by which the D96 spectrum is read, and through which all observable
structure is organized.

We establish that the operators are \emph{derived from the D96 spectrum} --- projections of
its structure, not independent primitives \cite{QG261,QG262,QG263}. The basis consists of
four operators: \emph{Crowding} (degeneracy grouping), \emph{Compression} (octave banding),
\emph{Beat} (frequency ratio), and \emph{Locking} (spectral gap). We establish the
universality evidence across organized domains \cite{QG300,QG302,QG304}, and the
completeness of the four-operator basis \cite{QG307,QG308}. Throughout, the claim that no
fifth operator exists is phrased \emph{as a search-scoped result} --- no fifth operator has
been found in any searched domain --- and not as an existence proof, per the referee-safe
wording of MONO007 \cite{MONO007}. The honest adversarial and red-team results
\cite{QG309,QG312} are reported as scoping the basis, not as contradictions.

The results of this chapter are the canonical end-state of the operator program
\cite{QG260,QG261,QG262,QG263}. Throughout we use only accepted canonical results and
introduce no new physics and no speculation.

\section{Operators as Spectrum Projections}
\label{sec:projections}

\begin{definition}[Spectral operator]
\label{def:operator}
A \emph{spectral operator} is a projection of the D96 spectrum: a structural read that maps
the spectrum onto a derived quantity. The operators are not new physical content; they are
the reading instruments of the spectrum.
\end{definition}

\begin{theorem}[The operators are derived projections]
\label{thm:operators-derived}
The operators are derived from the D96 spectrum: every named quantity of the theory ---
$\Sigma m$, $\Sigma\sqrt{m}$, $\Sigma m^2$, the occupancy moment, the span, the spectral
gap --- is a verified projection of a spectral operator \cite{QG261}.
\end{theorem}

\begin{proof}
The resonance-operator audit \cite{QG261} establishes that the six named quantities are
projections of deeper operators: $\Sigma m = \mathrm{MOMENT}_1\circ\mathrm{CROWDING}$,
$\Sigma\sqrt{m} = \mathrm{MOMENT}_{1/2}\circ\mathrm{CROWDING}$,
$\Sigma m^2 = \mathrm{MOMENT}_2\circ\mathrm{CROWDING}$, the occupancy moment is a
MOMENT$\circ$COMPRESSION read, the span is a BEAT read, and the spectral gap is a LOCKING read.
Hence every derived quantity passes through the operator layer --- the operators are the
reading instruments of the spectrum, not independent inputs. By
Chapter~\ref{chap:d96} (Theorem~\ref{thm:d96-derived}), the spectrum is itself derived;
the operators are therefore twice-removed from primitiveness: projections of a derived
spectrum.
\end{proof}

\begin{corollary}[Operators are not primitive]
\label{cor:operators-not-primitive}
The operators are not primitives of the theory: they are projections of the D96 spectrum,
which is itself the derived output of the actualization attractor.
\end{corollary}

\begin{proof}
By Theorem~\ref{thm:operators-derived}, the operators are projections of the spectrum. By
Chapter~\ref{chap:d96} (Theorem~\ref{thm:d96-derived}), the spectrum is the derived output
of the attractor. Hence the operators inherit the derived status: they are reads of a
derived structure, and no operator is an underivable input.
\end{proof}

\section{Crowding}
\label{sec:crowding}

\begin{definition}[Crowding]
\label{def:crowding}
\emph{Crowding} is the degeneracy operator: the grouping of the spectrum into degenerate
multiplicity classes, whose distribution is the multiplicity multiset
$[42\times2,5,6]$.
\end{definition}

\begin{theorem}[Crowding projects the moments]
\label{thm:crowding-moments}
Crowding generates the moment structure of the theory: the moments
$\Sigma m$, $\Sigma\sqrt{m}$, $\Sigma m^2$ are the moment compositions of Crowding over the
multiplicity multiset.
\end{theorem}

\begin{proof}
The resonance-operator audit \cite{QG261} establishes that
$\Sigma m = \mathrm{MOMENT}_1\circ\mathrm{CROWDING} = 95$,
$\Sigma\sqrt{m} = \mathrm{MOMENT}_{1/2}\circ\mathrm{CROWDING} = 64.08$, and
$\Sigma m^2 = \mathrm{MOMENT}_2\circ\mathrm{CROWDING} = 229$, applied to the multiplicity
multiset $[42\times2,5,6]$ (Chapter~\ref{chap:d96},
Theorem~\ref{thm:multiplicity}). Hence Crowding is the degeneracy projection that generates
the moment ladder of the spectrum.
\end{proof}

\begin{corollary}[Crowding reads the degeneracy structure]
\label{cor:crowding-degeneracy}
Crowding is the read of the degeneracy structure of the spectrum: the grouping of equal
eigenvalues into the multiplicity classes.
\end{corollary}

\begin{proof}
By Definition~\ref{def:crowding}, Crowding is the degeneracy grouping. By
Theorem~\ref{thm:crowding-moments}, its projections are the moment structure. Hence
Crowding reads the degeneracy: it is the operator by which the multiplicity structure of the
spectrum is measured.
\end{proof}

\section{Compression}
\label{sec:compression}

\begin{definition}[Compression]
\label{def:compression}
\emph{Compression} is the octave-banding operator: the grouping of the spectrum into octave
bands, whose occupancy distribution is the octave occupancies.
\end{definition}

\begin{theorem}[Compression projects the occupancy structure]
\label{thm:compression-occupancy}
Compression generates the octave-occupation structure of the theory: the occupancy moment
$\mathrm{occMom} = 1900.25$ is the MOMENT$\circ$COMPRESSION read over the octave occupancies.
\end{theorem}

\begin{proof}
The resonance-operator audit \cite{QG261} establishes that the occupancy moment is a
MOMENT$\circ$COMPRESSION read: the octave occupancies of the spectrum, weighted by their squares
and normalized by the lightest-octave occupancy, yield
$\mathrm{occMom} = 1900.25$ (Chapter~\ref{chap:d96}, Theorem~\ref{thm:occmom}). Hence
Compression is the octave-banding projection that generates the occupancy structure of the
spectrum.
\end{proof}

\begin{corollary}[Compression reads the octave structure]
\label{cor:compression-octave}
Compression is the read of the octave structure of the spectrum: the grouping of modes into
frequency octaves and the occupation of those bands.
\end{corollary}

\begin{proof}
By Definition~\ref{def:compression}, Compression is the octave-banding operator. By
Theorem~\ref{thm:compression-occupancy}, its projections are the occupancy structure.
Hence Compression reads the octave organization of the spectrum.
\end{proof}

\section{Beat}
\label{sec:beat}

\begin{definition}[Beat]
\label{def:beat}
\emph{Beat} is the frequency-ratio operator: the ratio of the extremal modes of the
spectrum, measuring the frequency extent.
\end{definition}

\begin{theorem}[Beat projects the span]
\label{thm:beat-span}
Beat generates the span of the theory: the ratio
$\mathrm{span} = \omega_{\max}/\omega_{\min} = 6.40$ is the BEAT read of the spectrum.
\end{theorem}

\begin{proof}
The resonance-operator audit \cite{QG261} establishes that the span is the BEAT projection:
$\mathrm{span} = \omega_{\max}/\omega_{\min} = 6.40$
(Chapter~\ref{chap:d96}, Theorem~\ref{thm:span}). Hence Beat is the frequency-ratio
operator that generates the extent of the spectrum.
\end{proof}

\begin{corollary}[Beat reads the frequency extent]
\label{cor:beat-extent}
Beat is the read of the frequency extent of the spectrum: the ratio between its highest and
lowest positive modes.
\end{corollary}

\begin{proof}
By Definition~\ref{def:beat} and Theorem~\ref{thm:beat-span}, Beat reads the span --- the
ratio of the extremal modes. Hence Beat is the operator by which the frequency extent of the
spectrum is measured.
\end{proof}

\section{Locking}
\label{sec:locking}

\begin{definition}[Locking]
\label{def:locking}
\emph{Locking} is the spectral-gap operator: the separation structure of the spectrum,
measuring the gap between its distinct mode groups.
\end{definition}

\begin{theorem}[Locking projects the spectral gap]
\label{thm:locking-gap}
Locking generates the spectral-gap structure of the theory: the gap $\lambda_2$ of the
spectrum is the LOCKING read.
\end{theorem}

\begin{proof}
The resonance-operator audit \cite{QG261} establishes that the spectral gap $\lambda_2$ is
the LOCKING projection of the spectrum. The locking structure measures the separation of the
distinct mode groups, providing the spectral-gap constant used across the sectors of the
theory.
\end{proof}

\begin{corollary}[Locking reads the separation structure]
\label{cor:locking-separation}
Locking is the read of the separation structure of the spectrum: the gaps between its
distinct mode groups.
\end{corollary}

\begin{proof}
By Definition~\ref{def:locking} and Theorem~\ref{thm:locking-gap}, Locking reads the
spectral-gap structure --- the separation between distinct mode groups. Hence Locking is the
operator by which the gap structure of the spectrum is measured.
\end{proof}

\section{Universality Evidence}
\label{sec:universality}

\begin{theorem}[The operators are universal across organized systems]
\label{thm:operators-universal}
The four-operator basis is present across organized systems of very different origin: the
operators appear in the same structural role wherever an organized spectrum is read.
\end{theorem}

\begin{proof}
The operator-universality program establishes the evidence across domains. The
operator-universality prediction \cite{QG300} establishes the universality of the operator
basis. The cross-domain universality audit \cite{QG302} confirms the four operators appear
across networks, language, music, DNA, software, finance, and other organized domains. The
real-network universality audit \cite{QG304} confirms the same operators on real-structure
networks --- connectomes, protein networks, citation graphs, internet graphs, knowledge
graphs. Hence the four operators recur as the structural read of organized spectra
throughout, without being imposed.
\end{proof}

\begin{lemma}[The operators are not trivial statistics]
\label{lem:operators-nontrivial}
The operators are not trivial statistical artifacts: they discriminate organized systems
from null and adversarial spectra.
\end{lemma}

\begin{proof}
The null-spectrum and adversarial audits \cite{QG309,QG312} establish the discriminating
power. The null-spectrum audit shows that random spectra fail the operator basis at the
quantitative level, while organized systems carry it. The adversarial audit shows that
fabricated spectra can trigger the binary presence conditions but do not carry the deeper
organization content (the locks). Hence the operators are a discriminating read of
organization, not a trivial statistical artifact.
\end{proof}

\begin{remark}[Scope disclosure]
Consistent with MONO007 \cite{MONO007}, the universality claim is disclosed as an
evidence-based result over the tested domain set --- an open-ended claim, not an existence
proof. The operators have been found universal in every organized domain searched; the
claim does not assert that no counterexample exists in an unsearched domain.
\end{remark}

\section{Operator Completeness}
\label{sec:completeness}

\begin{theorem}[The four-operator basis is justified]
\label{thm:four-operators}
The four operators --- Crowding, Compression, Beat, Locking --- form the justified operator
basis of the theory: each is indispensable, and together they cover the reading structure of
the spectrum.
\end{theorem}

\begin{proof}
The operator-necessity audit \cite{QG308} establishes that each of the four operators is
indispensable: removing any one loses derived content (each operator projects a distinct
structural feature of the spectrum --- degeneracy, octave organization, frequency extent,
spectral gap). Hence the four operators are each necessary, and together they form the
basis of the spectrum's reading structure.
\end{proof}

\begin{theorem}[No fifth operator found]
\label{thm:no-fifth}
No fifth operator has been found in any searched domain: the search for an operator beyond
the basis {Crowding, Compression, Beat, Locking} (together with the universal MOMENT
read-out) returned no genuine fifth spectral operator.
\end{theorem}

\begin{proof}
The fifth-operator search \cite{QG307} examined candidate fifth operators from unexplored
domains --- phase/orientation structure, order/sequence structure, and others --- and found
that each candidate either reduces to a composition of the existing operators or is not a
spectral read of the frequency distribution (e.g.\ the order structure is the network input,
not a spectral read). Hence no genuine fifth operator has been found. This is a
search-scoped result, not an existence proof: no fifth operator has been found in any
searched domain.
\end{proof}

\begin{corollary}[The basis is four, search-scoped]
\label{cor:four-search-scoped}
The operator basis of the theory is the set {Crowding, Compression, Beat, Locking}; the
statement that there is no fifth operator is established as a search-scoped result, not an
existence proof.
\end{corollary}

\begin{proof}
The necessity of the four is established by Theorem~\ref{thm:four-operators}; the
non-existence of a fifth is established as a search-scoped result by
Theorem~\ref{thm:no-fifth}. The referee-safe phrasing --- "no fifth operator found in any
searched domain" --- is adopted throughout, consistent with MONO007 \cite{MONO007}.
\end{proof}

\section{Consequences}
\label{sec:operator-consequences}

\begin{lemma}[The operators connect spectrum to physics]
\label{lem:operators-to-physics}
The operator basis is the bridge from the D96 spectrum to the physics sectors: every
observable is a read of the spectrum through the operators.
\end{lemma}

\begin{proof}
The resonance-operator audit \cite{QG261} and the same-operator-sectors audit \cite{QG262}
establish that the sectors of physics draw on the same operator layer: every sector reads
the spectrum through the operators. By Chapter~\ref{chap:actualization}
(Theorem~\ref{thm:resonance-readout}), the sectors are projection classes over the operator
layer. Hence the operators connect the spectrum to the physics sectors.
\end{proof}

\begin{theorem}[The operator layer is single]
\label{thm:single-operator-layer}
There is one operator layer, not many: one spectrum, one dynamics, one operator basis, read
by all sectors.
\end{theorem}

\begin{proof}
The single-resonance-dynamics audit \cite{QG263} establishes that there is one resonance
dynamics underlying the operator layer. The same-operator-sectors audit \cite{QG262}
establishes that no operator is sector-exclusive: every sector reads the same operators.
Hence the operator layer is single --- one basis, read by all sectors --- and the sector
differences are differences of role, not of operator.
\end{proof}

\begin{corollary}[Organization is a read of the operators]
\label{cor:organization-read}
The organization structure of the theory is a read of the operator basis over the D96
spectrum: the organizational quantities are operator projections.
\end{corollary}

\begin{proof}
By Theorem~\ref{thm:operators-derived}, all named quantities are operator projections. By
Chapter~\ref{chap:d96} (Theorem~\ref{thm:organization-derived}), the organization constants
are derived from the spectrum. Hence organization is the operator read of the spectrum ---
the organizational structure is the operator structure.
\end{proof}

\section{Summary}
\label{sec:operators-summary}

This chapter has established the operator basis of the theory. We have proved:

\begin{enumerate}
\item \textbf{Operators as projections} (Theorem~\ref{thm:operators-derived},
Corollary~\ref{cor:operators-not-primitive}): the operators are derived projections of the
D96 spectrum, not primitives;
\item \textbf{Crowding} (Theorem~\ref{thm:crowding-moments},
Corollary~\ref{cor:crowding-degeneracy}): the degeneracy operator projects the moments;
\item \textbf{Compression} (Theorem~\ref{thm:compression-occupancy},
Corollary~\ref{cor:compression-octave}): the octave-banding operator projects the occupancy
structure;
\item \textbf{Beat} (Theorem~\ref{thm:beat-span}, Corollary~\ref{cor:beat-extent}): the
frequency-ratio operator projects the span;
\item \textbf{Locking} (Theorem~\ref{thm:locking-gap},
Corollary~\ref{cor:locking-separation}): the spectral-gap operator projects the separation
structure;
\item \textbf{Universality} (Theorem~\ref{thm:operators-universal},
Lemma~\ref{lem:operators-nontrivial}): the operators are universal across organized systems
and not trivial statistics;
\item \textbf{Completeness} (Theorem~\ref{thm:four-operators},
Theorem~\ref{thm:no-fifth}, Corollary~\ref{cor:four-search-scoped}): the four-operator
basis is justified, and no fifth operator has been found in any searched domain;
\item \textbf{Consequences} (Lemma~\ref{lem:operators-to-physics},
Theorem~\ref{thm:single-operator-layer}, Corollary~\ref{cor:organization-read}): the
operators bridge spectrum to physics, the operator layer is single, and organization is a
read of the operators.
\end{enumerate}

The operator basis of the theory is therefore the derived reading structure of the D96
spectrum: Crowding (degeneracy), Compression (octaves), Beat (extent), and Locking (gaps),
each an indispensable projection, together universal across organized systems, with no fifth
operator found in any searched domain. The following chapters read the lock structure and
then each physics sector through this operator layer.

\begin{thebibliography}{99}
\bibitem{QG260} AT-QG Phase 260, \emph{Resonance Layer Audit}, Docs/Research.
\bibitem{QG261} AT-QG Phase 261, \emph{Resonance Operator Audit} (Operator Layer), Docs/Research.
\bibitem{QG262} AT-QG Phase 262, \emph{Operator Sector Audit}, Docs/Research.
\bibitem{QG263} AT-QG Phase 263, \emph{Single Resonance Dynamics Audit}, Docs/Research.
\bibitem{QG300} AT-QG Phase 300, \emph{Operator Universality Prediction}, Docs/Research.
\bibitem{QG302} AT-QG Phase 302, \emph{Cross-Domain Universality Audit}, Docs/Research.
\bibitem{QG304} AT-QG Phase 304, \emph{Real Network Universality Audit}, Docs/Research.
\bibitem{QG307} AT-QG Phase 307, \emph{Fifth Operator Search}, Docs/Research.
\bibitem{QG308} AT-QG Phase 308, \emph{Operator Necessity Audit}, Docs/Research.
\bibitem{QG309} AT-QG Phase 309, \emph{Red Team Audit}, Docs/Research.
\bibitem{QG312} AT-QG Phase 312, \emph{Adversarial Spectrum Audit}, Docs/Research.
\bibitem{MONO007} AT-MONO007, \emph{Referee-Readiness Resolution}, Docs/Zenodo/Monograph.
\end{thebibliography}
 after the preamble
%         that defines the theorem environments (or include via the master file).
% ======================================================================

% ---- Theorem environments (define in the master preamble if not present) ----
% \newtheorem{definition}{Definition}[chapter]
% \newtheorem{lemma}[definition]{Lemma}
% \newtheorem{theorem}[definition]{Theorem}
% \newtheorem{corollary}[definition]{Corollary}
% \newtheorem{remark}[definition]{Remark}

\chapter{The Operator Basis}
\label{chap:operators}

\section{Introduction}
\label{sec:operators-introduction}

In Chapter~\ref{chap:d96} we established the concrete structure of the D96 spectrum --- its
modes, moments, and multiplicities --- as the emergent output of the actualization
attractor. This chapter develops the \emph{operator basis} of the theory: the set of
spectral projections by which the D96 spectrum is read, and through which all observable
structure is organized.

We establish that the operators are \emph{derived from the D96 spectrum} --- projections of
its structure, not independent primitives \cite{QG261,QG262,QG263}. The basis consists of
four operators: \emph{Crowding} (degeneracy grouping), \emph{Compression} (octave banding),
\emph{Beat} (frequency ratio), and \emph{Locking} (spectral gap). We establish the
universality evidence across organized domains \cite{QG300,QG302,QG304}, and the
completeness of the four-operator basis \cite{QG307,QG308}. Throughout, the claim that no
fifth operator exists is phrased \emph{as a search-scoped result} --- no fifth operator has
been found in any searched domain --- and not as an existence proof, per the referee-safe
wording of MONO007 \cite{MONO007}. The honest adversarial and red-team results
\cite{QG309,QG312} are reported as scoping the basis, not as contradictions.

The results of this chapter are the canonical end-state of the operator program
\cite{QG260,QG261,QG262,QG263}. Throughout we use only accepted canonical results and
introduce no new physics and no speculation.

\section{Operators as Spectrum Projections}
\label{sec:projections}

\begin{definition}[Spectral operator]
\label{def:operator}
A \emph{spectral operator} is a projection of the D96 spectrum: a structural read that maps
the spectrum onto a derived quantity. The operators are not new physical content; they are
the reading instruments of the spectrum.
\end{definition}

\begin{theorem}[The operators are derived projections]
\label{thm:operators-derived}
The operators are derived from the D96 spectrum: every named quantity of the theory ---
$\Sigma m$, $\Sigma\sqrt{m}$, $\Sigma m^2$, the occupancy moment, the span, the spectral
gap --- is a verified projection of a spectral operator \cite{QG261}.
\end{theorem}

\begin{proof}
The resonance-operator audit \cite{QG261} establishes that the six named quantities are
projections of deeper operators: $\Sigma m = \mathrm{MOMENT}_1\circ\mathrm{CROWDING}$,
$\Sigma\sqrt{m} = \mathrm{MOMENT}_{1/2}\circ\mathrm{CROWDING}$,
$\Sigma m^2 = \mathrm{MOMENT}_2\circ\mathrm{CROWDING}$, the occupancy moment is a
MOMENT$\circ$COMPRESSION read, the span is a BEAT read, and the spectral gap is a LOCKING read.
Hence every derived quantity passes through the operator layer --- the operators are the
reading instruments of the spectrum, not independent inputs. By
Chapter~\ref{chap:d96} (Theorem~\ref{thm:d96-derived}), the spectrum is itself derived;
the operators are therefore twice-removed from primitiveness: projections of a derived
spectrum.
\end{proof}

\begin{corollary}[Operators are not primitive]
\label{cor:operators-not-primitive}
The operators are not primitives of the theory: they are projections of the D96 spectrum,
which is itself the derived output of the actualization attractor.
\end{corollary}

\begin{proof}
By Theorem~\ref{thm:operators-derived}, the operators are projections of the spectrum. By
Chapter~\ref{chap:d96} (Theorem~\ref{thm:d96-derived}), the spectrum is the derived output
of the attractor. Hence the operators inherit the derived status: they are reads of a
derived structure, and no operator is an underivable input.
\end{proof}

\section{Crowding}
\label{sec:crowding}

\begin{definition}[Crowding]
\label{def:crowding}
\emph{Crowding} is the degeneracy operator: the grouping of the spectrum into degenerate
multiplicity classes, whose distribution is the multiplicity multiset
$[42\times2,5,6]$.
\end{definition}

\begin{theorem}[Crowding projects the moments]
\label{thm:crowding-moments}
Crowding generates the moment structure of the theory: the moments
$\Sigma m$, $\Sigma\sqrt{m}$, $\Sigma m^2$ are the moment compositions of Crowding over the
multiplicity multiset.
\end{theorem}

\begin{proof}
The resonance-operator audit \cite{QG261} establishes that
$\Sigma m = \mathrm{MOMENT}_1\circ\mathrm{CROWDING} = 95$,
$\Sigma\sqrt{m} = \mathrm{MOMENT}_{1/2}\circ\mathrm{CROWDING} = 64.08$, and
$\Sigma m^2 = \mathrm{MOMENT}_2\circ\mathrm{CROWDING} = 229$, applied to the multiplicity
multiset $[42\times2,5,6]$ (Chapter~\ref{chap:d96},
Theorem~\ref{thm:multiplicity}). Hence Crowding is the degeneracy projection that generates
the moment ladder of the spectrum.
\end{proof}

\begin{corollary}[Crowding reads the degeneracy structure]
\label{cor:crowding-degeneracy}
Crowding is the read of the degeneracy structure of the spectrum: the grouping of equal
eigenvalues into the multiplicity classes.
\end{corollary}

\begin{proof}
By Definition~\ref{def:crowding}, Crowding is the degeneracy grouping. By
Theorem~\ref{thm:crowding-moments}, its projections are the moment structure. Hence
Crowding reads the degeneracy: it is the operator by which the multiplicity structure of the
spectrum is measured.
\end{proof}

\section{Compression}
\label{sec:compression}

\begin{definition}[Compression]
\label{def:compression}
\emph{Compression} is the octave-banding operator: the grouping of the spectrum into octave
bands, whose occupancy distribution is the octave occupancies.
\end{definition}

\begin{theorem}[Compression projects the occupancy structure]
\label{thm:compression-occupancy}
Compression generates the octave-occupation structure of the theory: the occupancy moment
$\mathrm{occMom} = 1900.25$ is the MOMENT$\circ$COMPRESSION read over the octave occupancies.
\end{theorem}

\begin{proof}
The resonance-operator audit \cite{QG261} establishes that the occupancy moment is a
MOMENT$\circ$COMPRESSION read: the octave occupancies of the spectrum, weighted by their squares
and normalized by the lightest-octave occupancy, yield
$\mathrm{occMom} = 1900.25$ (Chapter~\ref{chap:d96}, Theorem~\ref{thm:occmom}). Hence
Compression is the octave-banding projection that generates the occupancy structure of the
spectrum.
\end{proof}

\begin{corollary}[Compression reads the octave structure]
\label{cor:compression-octave}
Compression is the read of the octave structure of the spectrum: the grouping of modes into
frequency octaves and the occupation of those bands.
\end{corollary}

\begin{proof}
By Definition~\ref{def:compression}, Compression is the octave-banding operator. By
Theorem~\ref{thm:compression-occupancy}, its projections are the occupancy structure.
Hence Compression reads the octave organization of the spectrum.
\end{proof}

\section{Beat}
\label{sec:beat}

\begin{definition}[Beat]
\label{def:beat}
\emph{Beat} is the frequency-ratio operator: the ratio of the extremal modes of the
spectrum, measuring the frequency extent.
\end{definition}

\begin{theorem}[Beat projects the span]
\label{thm:beat-span}
Beat generates the span of the theory: the ratio
$\mathrm{span} = \omega_{\max}/\omega_{\min} = 6.40$ is the BEAT read of the spectrum.
\end{theorem}

\begin{proof}
The resonance-operator audit \cite{QG261} establishes that the span is the BEAT projection:
$\mathrm{span} = \omega_{\max}/\omega_{\min} = 6.40$
(Chapter~\ref{chap:d96}, Theorem~\ref{thm:span}). Hence Beat is the frequency-ratio
operator that generates the extent of the spectrum.
\end{proof}

\begin{corollary}[Beat reads the frequency extent]
\label{cor:beat-extent}
Beat is the read of the frequency extent of the spectrum: the ratio between its highest and
lowest positive modes.
\end{corollary}

\begin{proof}
By Definition~\ref{def:beat} and Theorem~\ref{thm:beat-span}, Beat reads the span --- the
ratio of the extremal modes. Hence Beat is the operator by which the frequency extent of the
spectrum is measured.
\end{proof}

\section{Locking}
\label{sec:locking}

\begin{definition}[Locking]
\label{def:locking}
\emph{Locking} is the spectral-gap operator: the separation structure of the spectrum,
measuring the gap between its distinct mode groups.
\end{definition}

\begin{theorem}[Locking projects the spectral gap]
\label{thm:locking-gap}
Locking generates the spectral-gap structure of the theory: the gap $\lambda_2$ of the
spectrum is the LOCKING read.
\end{theorem}

\begin{proof}
The resonance-operator audit \cite{QG261} establishes that the spectral gap $\lambda_2$ is
the LOCKING projection of the spectrum. The locking structure measures the separation of the
distinct mode groups, providing the spectral-gap constant used across the sectors of the
theory.
\end{proof}

\begin{corollary}[Locking reads the separation structure]
\label{cor:locking-separation}
Locking is the read of the separation structure of the spectrum: the gaps between its
distinct mode groups.
\end{corollary}

\begin{proof}
By Definition~\ref{def:locking} and Theorem~\ref{thm:locking-gap}, Locking reads the
spectral-gap structure --- the separation between distinct mode groups. Hence Locking is the
operator by which the gap structure of the spectrum is measured.
\end{proof}

\section{Universality Evidence}
\label{sec:universality}

\begin{theorem}[The operators are universal across organized systems]
\label{thm:operators-universal}
The four-operator basis is present across organized systems of very different origin: the
operators appear in the same structural role wherever an organized spectrum is read.
\end{theorem}

\begin{proof}
The operator-universality program establishes the evidence across domains. The
operator-universality prediction \cite{QG300} establishes the universality of the operator
basis. The cross-domain universality audit \cite{QG302} confirms the four operators appear
across networks, language, music, DNA, software, finance, and other organized domains. The
real-network universality audit \cite{QG304} confirms the same operators on real-structure
networks --- connectomes, protein networks, citation graphs, internet graphs, knowledge
graphs. Hence the four operators recur as the structural read of organized spectra
throughout, without being imposed.
\end{proof}

\begin{lemma}[The operators are not trivial statistics]
\label{lem:operators-nontrivial}
The operators are not trivial statistical artifacts: they discriminate organized systems
from null and adversarial spectra.
\end{lemma}

\begin{proof}
The null-spectrum and adversarial audits \cite{QG309,QG312} establish the discriminating
power. The null-spectrum audit shows that random spectra fail the operator basis at the
quantitative level, while organized systems carry it. The adversarial audit shows that
fabricated spectra can trigger the binary presence conditions but do not carry the deeper
organization content (the locks). Hence the operators are a discriminating read of
organization, not a trivial statistical artifact.
\end{proof}

\begin{remark}[Scope disclosure]
Consistent with MONO007 \cite{MONO007}, the universality claim is disclosed as an
evidence-based result over the tested domain set --- an open-ended claim, not an existence
proof. The operators have been found universal in every organized domain searched; the
claim does not assert that no counterexample exists in an unsearched domain.
\end{remark}

\section{Operator Completeness}
\label{sec:completeness}

\begin{theorem}[The four-operator basis is justified]
\label{thm:four-operators}
The four operators --- Crowding, Compression, Beat, Locking --- form the justified operator
basis of the theory: each is indispensable, and together they cover the reading structure of
the spectrum.
\end{theorem}

\begin{proof}
The operator-necessity audit \cite{QG308} establishes that each of the four operators is
indispensable: removing any one loses derived content (each operator projects a distinct
structural feature of the spectrum --- degeneracy, octave organization, frequency extent,
spectral gap). Hence the four operators are each necessary, and together they form the
basis of the spectrum's reading structure.
\end{proof}

\begin{theorem}[No fifth operator found]
\label{thm:no-fifth}
No fifth operator has been found in any searched domain: the search for an operator beyond
the basis {Crowding, Compression, Beat, Locking} (together with the universal MOMENT
read-out) returned no genuine fifth spectral operator.
\end{theorem}

\begin{proof}
The fifth-operator search \cite{QG307} examined candidate fifth operators from unexplored
domains --- phase/orientation structure, order/sequence structure, and others --- and found
that each candidate either reduces to a composition of the existing operators or is not a
spectral read of the frequency distribution (e.g.\ the order structure is the network input,
not a spectral read). Hence no genuine fifth operator has been found. This is a
search-scoped result, not an existence proof: no fifth operator has been found in any
searched domain.
\end{proof}

\begin{corollary}[The basis is four, search-scoped]
\label{cor:four-search-scoped}
The operator basis of the theory is the set {Crowding, Compression, Beat, Locking}; the
statement that there is no fifth operator is established as a search-scoped result, not an
existence proof.
\end{corollary}

\begin{proof}
The necessity of the four is established by Theorem~\ref{thm:four-operators}; the
non-existence of a fifth is established as a search-scoped result by
Theorem~\ref{thm:no-fifth}. The referee-safe phrasing --- "no fifth operator found in any
searched domain" --- is adopted throughout, consistent with MONO007 \cite{MONO007}.
\end{proof}

\section{Consequences}
\label{sec:operator-consequences}

\begin{lemma}[The operators connect spectrum to physics]
\label{lem:operators-to-physics}
The operator basis is the bridge from the D96 spectrum to the physics sectors: every
observable is a read of the spectrum through the operators.
\end{lemma}

\begin{proof}
The resonance-operator audit \cite{QG261} and the same-operator-sectors audit \cite{QG262}
establish that the sectors of physics draw on the same operator layer: every sector reads
the spectrum through the operators. By Chapter~\ref{chap:actualization}
(Theorem~\ref{thm:resonance-readout}), the sectors are projection classes over the operator
layer. Hence the operators connect the spectrum to the physics sectors.
\end{proof}

\begin{theorem}[The operator layer is single]
\label{thm:single-operator-layer}
There is one operator layer, not many: one spectrum, one dynamics, one operator basis, read
by all sectors.
\end{theorem}

\begin{proof}
The single-resonance-dynamics audit \cite{QG263} establishes that there is one resonance
dynamics underlying the operator layer. The same-operator-sectors audit \cite{QG262}
establishes that no operator is sector-exclusive: every sector reads the same operators.
Hence the operator layer is single --- one basis, read by all sectors --- and the sector
differences are differences of role, not of operator.
\end{proof}

\begin{corollary}[Organization is a read of the operators]
\label{cor:organization-read}
The organization structure of the theory is a read of the operator basis over the D96
spectrum: the organizational quantities are operator projections.
\end{corollary}

\begin{proof}
By Theorem~\ref{thm:operators-derived}, all named quantities are operator projections. By
Chapter~\ref{chap:d96} (Theorem~\ref{thm:organization-derived}), the organization constants
are derived from the spectrum. Hence organization is the operator read of the spectrum ---
the organizational structure is the operator structure.
\end{proof}

\section{Summary}
\label{sec:operators-summary}

This chapter has established the operator basis of the theory. We have proved:

\begin{enumerate}
\item \textbf{Operators as projections} (Theorem~\ref{thm:operators-derived},
Corollary~\ref{cor:operators-not-primitive}): the operators are derived projections of the
D96 spectrum, not primitives;
\item \textbf{Crowding} (Theorem~\ref{thm:crowding-moments},
Corollary~\ref{cor:crowding-degeneracy}): the degeneracy operator projects the moments;
\item \textbf{Compression} (Theorem~\ref{thm:compression-occupancy},
Corollary~\ref{cor:compression-octave}): the octave-banding operator projects the occupancy
structure;
\item \textbf{Beat} (Theorem~\ref{thm:beat-span}, Corollary~\ref{cor:beat-extent}): the
frequency-ratio operator projects the span;
\item \textbf{Locking} (Theorem~\ref{thm:locking-gap},
Corollary~\ref{cor:locking-separation}): the spectral-gap operator projects the separation
structure;
\item \textbf{Universality} (Theorem~\ref{thm:operators-universal},
Lemma~\ref{lem:operators-nontrivial}): the operators are universal across organized systems
and not trivial statistics;
\item \textbf{Completeness} (Theorem~\ref{thm:four-operators},
Theorem~\ref{thm:no-fifth}, Corollary~\ref{cor:four-search-scoped}): the four-operator
basis is justified, and no fifth operator has been found in any searched domain;
\item \textbf{Consequences} (Lemma~\ref{lem:operators-to-physics},
Theorem~\ref{thm:single-operator-layer}, Corollary~\ref{cor:organization-read}): the
operators bridge spectrum to physics, the operator layer is single, and organization is a
read of the operators.
\end{enumerate}

The operator basis of the theory is therefore the derived reading structure of the D96
spectrum: Crowding (degeneracy), Compression (octaves), Beat (extent), and Locking (gaps),
each an indispensable projection, together universal across organized systems, with no fifth
operator found in any searched domain. The following chapters read the lock structure and
then each physics sector through this operator layer.

\begin{thebibliography}{99}
\bibitem{QG260} AT-QG Phase 260, \emph{Resonance Layer Audit}, Docs/Research.
\bibitem{QG261} AT-QG Phase 261, \emph{Resonance Operator Audit} (Operator Layer), Docs/Research.
\bibitem{QG262} AT-QG Phase 262, \emph{Operator Sector Audit}, Docs/Research.
\bibitem{QG263} AT-QG Phase 263, \emph{Single Resonance Dynamics Audit}, Docs/Research.
\bibitem{QG300} AT-QG Phase 300, \emph{Operator Universality Prediction}, Docs/Research.
\bibitem{QG302} AT-QG Phase 302, \emph{Cross-Domain Universality Audit}, Docs/Research.
\bibitem{QG304} AT-QG Phase 304, \emph{Real Network Universality Audit}, Docs/Research.
\bibitem{QG307} AT-QG Phase 307, \emph{Fifth Operator Search}, Docs/Research.
\bibitem{QG308} AT-QG Phase 308, \emph{Operator Necessity Audit}, Docs/Research.
\bibitem{QG309} AT-QG Phase 309, \emph{Red Team Audit}, Docs/Research.
\bibitem{QG312} AT-QG Phase 312, \emph{Adversarial Spectrum Audit}, Docs/Research.
\bibitem{MONO007} AT-MONO007, \emph{Referee-Readiness Resolution}, Docs/Zenodo/Monograph.
\end{thebibliography}
 after the preamble
%         that defines the theorem environments (or include via the master file).
% ======================================================================

% ---- Theorem environments (define in the master preamble if not present) ----
% \newtheorem{definition}{Definition}[chapter]
% \newtheorem{lemma}[definition]{Lemma}
% \newtheorem{theorem}[definition]{Theorem}
% \newtheorem{corollary}[definition]{Corollary}
% \newtheorem{remark}[definition]{Remark}

\chapter{The Operator Basis}
\label{chap:operators}

\section{Introduction}
\label{sec:operators-introduction}

In Chapter~\ref{chap:d96} we established the concrete structure of the D96 spectrum --- its
modes, moments, and multiplicities --- as the emergent output of the actualization
attractor. This chapter develops the \emph{operator basis} of the theory: the set of
spectral projections by which the D96 spectrum is read, and through which all observable
structure is organized.

We establish that the operators are \emph{derived from the D96 spectrum} --- projections of
its structure, not independent primitives \cite{QG261,QG262,QG263}. The basis consists of
four operators: \emph{Crowding} (degeneracy grouping), \emph{Compression} (octave banding),
\emph{Beat} (frequency ratio), and \emph{Locking} (spectral gap). We establish the
universality evidence across organized domains \cite{QG300,QG302,QG304}, and the
completeness of the four-operator basis \cite{QG307,QG308}. Throughout, the claim that no
fifth operator exists is phrased \emph{as a search-scoped result} --- no fifth operator has
been found in any searched domain --- and not as an existence proof, per the referee-safe
wording of MONO007 \cite{MONO007}. The honest adversarial and red-team results
\cite{QG309,QG312} are reported as scoping the basis, not as contradictions.

The results of this chapter are the canonical end-state of the operator program
\cite{QG260,QG261,QG262,QG263}. Throughout we use only accepted canonical results and
introduce no new physics and no speculation.

\section{Operators as Spectrum Projections}
\label{sec:projections}

\begin{definition}[Spectral operator]
\label{def:operator}
A \emph{spectral operator} is a projection of the D96 spectrum: a structural read that maps
the spectrum onto a derived quantity. The operators are not new physical content; they are
the reading instruments of the spectrum.
\end{definition}

\begin{theorem}[The operators are derived projections]
\label{thm:operators-derived}
The operators are derived from the D96 spectrum: every named quantity of the theory ---
$\Sigma m$, $\Sigma\sqrt{m}$, $\Sigma m^2$, the occupancy moment, the span, the spectral
gap --- is a verified projection of a spectral operator \cite{QG261}.
\end{theorem}

\begin{proof}
The resonance-operator audit \cite{QG261} establishes that the six named quantities are
projections of deeper operators: $\Sigma m = \mathrm{MOMENT}_1\circ\mathrm{CROWDING}$,
$\Sigma\sqrt{m} = \mathrm{MOMENT}_{1/2}\circ\mathrm{CROWDING}$,
$\Sigma m^2 = \mathrm{MOMENT}_2\circ\mathrm{CROWDING}$, the occupancy moment is a
MOMENT$\circ$COMPRESSION read, the span is a BEAT read, and the spectral gap is a LOCKING read.
Hence every derived quantity passes through the operator layer --- the operators are the
reading instruments of the spectrum, not independent inputs. By
Chapter~\ref{chap:d96} (Theorem~\ref{thm:d96-derived}), the spectrum is itself derived;
the operators are therefore twice-removed from primitiveness: projections of a derived
spectrum.
\end{proof}

\begin{corollary}[Operators are not primitive]
\label{cor:operators-not-primitive}
The operators are not primitives of the theory: they are projections of the D96 spectrum,
which is itself the derived output of the actualization attractor.
\end{corollary}

\begin{proof}
By Theorem~\ref{thm:operators-derived}, the operators are projections of the spectrum. By
Chapter~\ref{chap:d96} (Theorem~\ref{thm:d96-derived}), the spectrum is the derived output
of the attractor. Hence the operators inherit the derived status: they are reads of a
derived structure, and no operator is an underivable input.
\end{proof}

\section{Crowding}
\label{sec:crowding}

\begin{definition}[Crowding]
\label{def:crowding}
\emph{Crowding} is the degeneracy operator: the grouping of the spectrum into degenerate
multiplicity classes, whose distribution is the multiplicity multiset
$[42\times2,5,6]$.
\end{definition}

\begin{theorem}[Crowding projects the moments]
\label{thm:crowding-moments}
Crowding generates the moment structure of the theory: the moments
$\Sigma m$, $\Sigma\sqrt{m}$, $\Sigma m^2$ are the moment compositions of Crowding over the
multiplicity multiset.
\end{theorem}

\begin{proof}
The resonance-operator audit \cite{QG261} establishes that
$\Sigma m = \mathrm{MOMENT}_1\circ\mathrm{CROWDING} = 95$,
$\Sigma\sqrt{m} = \mathrm{MOMENT}_{1/2}\circ\mathrm{CROWDING} = 64.08$, and
$\Sigma m^2 = \mathrm{MOMENT}_2\circ\mathrm{CROWDING} = 229$, applied to the multiplicity
multiset $[42\times2,5,6]$ (Chapter~\ref{chap:d96},
Theorem~\ref{thm:multiplicity}). Hence Crowding is the degeneracy projection that generates
the moment ladder of the spectrum.
\end{proof}

\begin{corollary}[Crowding reads the degeneracy structure]
\label{cor:crowding-degeneracy}
Crowding is the read of the degeneracy structure of the spectrum: the grouping of equal
eigenvalues into the multiplicity classes.
\end{corollary}

\begin{proof}
By Definition~\ref{def:crowding}, Crowding is the degeneracy grouping. By
Theorem~\ref{thm:crowding-moments}, its projections are the moment structure. Hence
Crowding reads the degeneracy: it is the operator by which the multiplicity structure of the
spectrum is measured.
\end{proof}

\section{Compression}
\label{sec:compression}

\begin{definition}[Compression]
\label{def:compression}
\emph{Compression} is the octave-banding operator: the grouping of the spectrum into octave
bands, whose occupancy distribution is the octave occupancies.
\end{definition}

\begin{theorem}[Compression projects the occupancy structure]
\label{thm:compression-occupancy}
Compression generates the octave-occupation structure of the theory: the occupancy moment
$\mathrm{occMom} = 1900.25$ is the MOMENT$\circ$COMPRESSION read over the octave occupancies.
\end{theorem}

\begin{proof}
The resonance-operator audit \cite{QG261} establishes that the occupancy moment is a
MOMENT$\circ$COMPRESSION read: the octave occupancies of the spectrum, weighted by their squares
and normalized by the lightest-octave occupancy, yield
$\mathrm{occMom} = 1900.25$ (Chapter~\ref{chap:d96}, Theorem~\ref{thm:occmom}). Hence
Compression is the octave-banding projection that generates the occupancy structure of the
spectrum.
\end{proof}

\begin{corollary}[Compression reads the octave structure]
\label{cor:compression-octave}
Compression is the read of the octave structure of the spectrum: the grouping of modes into
frequency octaves and the occupation of those bands.
\end{corollary}

\begin{proof}
By Definition~\ref{def:compression}, Compression is the octave-banding operator. By
Theorem~\ref{thm:compression-occupancy}, its projections are the occupancy structure.
Hence Compression reads the octave organization of the spectrum.
\end{proof}

\section{Beat}
\label{sec:beat}

\begin{definition}[Beat]
\label{def:beat}
\emph{Beat} is the frequency-ratio operator: the ratio of the extremal modes of the
spectrum, measuring the frequency extent.
\end{definition}

\begin{theorem}[Beat projects the span]
\label{thm:beat-span}
Beat generates the span of the theory: the ratio
$\mathrm{span} = \omega_{\max}/\omega_{\min} = 6.40$ is the BEAT read of the spectrum.
\end{theorem}

\begin{proof}
The resonance-operator audit \cite{QG261} establishes that the span is the BEAT projection:
$\mathrm{span} = \omega_{\max}/\omega_{\min} = 6.40$
(Chapter~\ref{chap:d96}, Theorem~\ref{thm:span}). Hence Beat is the frequency-ratio
operator that generates the extent of the spectrum.
\end{proof}

\begin{corollary}[Beat reads the frequency extent]
\label{cor:beat-extent}
Beat is the read of the frequency extent of the spectrum: the ratio between its highest and
lowest positive modes.
\end{corollary}

\begin{proof}
By Definition~\ref{def:beat} and Theorem~\ref{thm:beat-span}, Beat reads the span --- the
ratio of the extremal modes. Hence Beat is the operator by which the frequency extent of the
spectrum is measured.
\end{proof}

\section{Locking}
\label{sec:locking}

\begin{definition}[Locking]
\label{def:locking}
\emph{Locking} is the spectral-gap operator: the separation structure of the spectrum,
measuring the gap between its distinct mode groups.
\end{definition}

\begin{theorem}[Locking projects the spectral gap]
\label{thm:locking-gap}
Locking generates the spectral-gap structure of the theory: the gap $\lambda_2$ of the
spectrum is the LOCKING read.
\end{theorem}

\begin{proof}
The resonance-operator audit \cite{QG261} establishes that the spectral gap $\lambda_2$ is
the LOCKING projection of the spectrum. The locking structure measures the separation of the
distinct mode groups, providing the spectral-gap constant used across the sectors of the
theory.
\end{proof}

\begin{corollary}[Locking reads the separation structure]
\label{cor:locking-separation}
Locking is the read of the separation structure of the spectrum: the gaps between its
distinct mode groups.
\end{corollary}

\begin{proof}
By Definition~\ref{def:locking} and Theorem~\ref{thm:locking-gap}, Locking reads the
spectral-gap structure --- the separation between distinct mode groups. Hence Locking is the
operator by which the gap structure of the spectrum is measured.
\end{proof}

\section{Universality Evidence}
\label{sec:universality}

\begin{theorem}[The operators are universal across organized systems]
\label{thm:operators-universal}
The four-operator basis is present across organized systems of very different origin: the
operators appear in the same structural role wherever an organized spectrum is read.
\end{theorem}

\begin{proof}
The operator-universality program establishes the evidence across domains. The
operator-universality prediction \cite{QG300} establishes the universality of the operator
basis. The cross-domain universality audit \cite{QG302} confirms the four operators appear
across networks, language, music, DNA, software, finance, and other organized domains. The
real-network universality audit \cite{QG304} confirms the same operators on real-structure
networks --- connectomes, protein networks, citation graphs, internet graphs, knowledge
graphs. Hence the four operators recur as the structural read of organized spectra
throughout, without being imposed.
\end{proof}

\begin{lemma}[The operators are not trivial statistics]
\label{lem:operators-nontrivial}
The operators are not trivial statistical artifacts: they discriminate organized systems
from null and adversarial spectra.
\end{lemma}

\begin{proof}
The null-spectrum and adversarial audits \cite{QG309,QG312} establish the discriminating
power. The null-spectrum audit shows that random spectra fail the operator basis at the
quantitative level, while organized systems carry it. The adversarial audit shows that
fabricated spectra can trigger the binary presence conditions but do not carry the deeper
organization content (the locks). Hence the operators are a discriminating read of
organization, not a trivial statistical artifact.
\end{proof}

\begin{remark}[Scope disclosure]
Consistent with MONO007 \cite{MONO007}, the universality claim is disclosed as an
evidence-based result over the tested domain set --- an open-ended claim, not an existence
proof. The operators have been found universal in every organized domain searched; the
claim does not assert that no counterexample exists in an unsearched domain.
\end{remark}

\section{Operator Completeness}
\label{sec:completeness}

\begin{theorem}[The four-operator basis is justified]
\label{thm:four-operators}
The four operators --- Crowding, Compression, Beat, Locking --- form the justified operator
basis of the theory: each is indispensable, and together they cover the reading structure of
the spectrum.
\end{theorem}

\begin{proof}
The operator-necessity audit \cite{QG308} establishes that each of the four operators is
indispensable: removing any one loses derived content (each operator projects a distinct
structural feature of the spectrum --- degeneracy, octave organization, frequency extent,
spectral gap). Hence the four operators are each necessary, and together they form the
basis of the spectrum's reading structure.
\end{proof}

\begin{theorem}[No fifth operator found]
\label{thm:no-fifth}
No fifth operator has been found in any searched domain: the search for an operator beyond
the basis {Crowding, Compression, Beat, Locking} (together with the universal MOMENT
read-out) returned no genuine fifth spectral operator.
\end{theorem}

\begin{proof}
The fifth-operator search \cite{QG307} examined candidate fifth operators from unexplored
domains --- phase/orientation structure, order/sequence structure, and others --- and found
that each candidate either reduces to a composition of the existing operators or is not a
spectral read of the frequency distribution (e.g.\ the order structure is the network input,
not a spectral read). Hence no genuine fifth operator has been found. This is a
search-scoped result, not an existence proof: no fifth operator has been found in any
searched domain.
\end{proof}

\begin{corollary}[The basis is four, search-scoped]
\label{cor:four-search-scoped}
The operator basis of the theory is the set {Crowding, Compression, Beat, Locking}; the
statement that there is no fifth operator is established as a search-scoped result, not an
existence proof.
\end{corollary}

\begin{proof}
The necessity of the four is established by Theorem~\ref{thm:four-operators}; the
non-existence of a fifth is established as a search-scoped result by
Theorem~\ref{thm:no-fifth}. The referee-safe phrasing --- "no fifth operator found in any
searched domain" --- is adopted throughout, consistent with MONO007 \cite{MONO007}.
\end{proof}

\section{Consequences}
\label{sec:operator-consequences}

\begin{lemma}[The operators connect spectrum to physics]
\label{lem:operators-to-physics}
The operator basis is the bridge from the D96 spectrum to the physics sectors: every
observable is a read of the spectrum through the operators.
\end{lemma}

\begin{proof}
The resonance-operator audit \cite{QG261} and the same-operator-sectors audit \cite{QG262}
establish that the sectors of physics draw on the same operator layer: every sector reads
the spectrum through the operators. By Chapter~\ref{chap:actualization}
(Theorem~\ref{thm:resonance-readout}), the sectors are projection classes over the operator
layer. Hence the operators connect the spectrum to the physics sectors.
\end{proof}

\begin{theorem}[The operator layer is single]
\label{thm:single-operator-layer}
There is one operator layer, not many: one spectrum, one dynamics, one operator basis, read
by all sectors.
\end{theorem}

\begin{proof}
The single-resonance-dynamics audit \cite{QG263} establishes that there is one resonance
dynamics underlying the operator layer. The same-operator-sectors audit \cite{QG262}
establishes that no operator is sector-exclusive: every sector reads the same operators.
Hence the operator layer is single --- one basis, read by all sectors --- and the sector
differences are differences of role, not of operator.
\end{proof}

\begin{corollary}[Organization is a read of the operators]
\label{cor:organization-read}
The organization structure of the theory is a read of the operator basis over the D96
spectrum: the organizational quantities are operator projections.
\end{corollary}

\begin{proof}
By Theorem~\ref{thm:operators-derived}, all named quantities are operator projections. By
Chapter~\ref{chap:d96} (Theorem~\ref{thm:organization-derived}), the organization constants
are derived from the spectrum. Hence organization is the operator read of the spectrum ---
the organizational structure is the operator structure.
\end{proof}

\section{Summary}
\label{sec:operators-summary}

This chapter has established the operator basis of the theory. We have proved:

\begin{enumerate}
\item \textbf{Operators as projections} (Theorem~\ref{thm:operators-derived},
Corollary~\ref{cor:operators-not-primitive}): the operators are derived projections of the
D96 spectrum, not primitives;
\item \textbf{Crowding} (Theorem~\ref{thm:crowding-moments},
Corollary~\ref{cor:crowding-degeneracy}): the degeneracy operator projects the moments;
\item \textbf{Compression} (Theorem~\ref{thm:compression-occupancy},
Corollary~\ref{cor:compression-octave}): the octave-banding operator projects the occupancy
structure;
\item \textbf{Beat} (Theorem~\ref{thm:beat-span}, Corollary~\ref{cor:beat-extent}): the
frequency-ratio operator projects the span;
\item \textbf{Locking} (Theorem~\ref{thm:locking-gap},
Corollary~\ref{cor:locking-separation}): the spectral-gap operator projects the separation
structure;
\item \textbf{Universality} (Theorem~\ref{thm:operators-universal},
Lemma~\ref{lem:operators-nontrivial}): the operators are universal across organized systems
and not trivial statistics;
\item \textbf{Completeness} (Theorem~\ref{thm:four-operators},
Theorem~\ref{thm:no-fifth}, Corollary~\ref{cor:four-search-scoped}): the four-operator
basis is justified, and no fifth operator has been found in any searched domain;
\item \textbf{Consequences} (Lemma~\ref{lem:operators-to-physics},
Theorem~\ref{thm:single-operator-layer}, Corollary~\ref{cor:organization-read}): the
operators bridge spectrum to physics, the operator layer is single, and organization is a
read of the operators.
\end{enumerate}

The operator basis of the theory is therefore the derived reading structure of the D96
spectrum: Crowding (degeneracy), Compression (octaves), Beat (extent), and Locking (gaps),
each an indispensable projection, together universal across organized systems, with no fifth
operator found in any searched domain. The following chapters read the lock structure and
then each physics sector through this operator layer.

\begin{thebibliography}{99}
\bibitem{QG260} AT-QG Phase 260, \emph{Resonance Layer Audit}, Docs/Research.
\bibitem{QG261} AT-QG Phase 261, \emph{Resonance Operator Audit} (Operator Layer), Docs/Research.
\bibitem{QG262} AT-QG Phase 262, \emph{Operator Sector Audit}, Docs/Research.
\bibitem{QG263} AT-QG Phase 263, \emph{Single Resonance Dynamics Audit}, Docs/Research.
\bibitem{QG300} AT-QG Phase 300, \emph{Operator Universality Prediction}, Docs/Research.
\bibitem{QG302} AT-QG Phase 302, \emph{Cross-Domain Universality Audit}, Docs/Research.
\bibitem{QG304} AT-QG Phase 304, \emph{Real Network Universality Audit}, Docs/Research.
\bibitem{QG307} AT-QG Phase 307, \emph{Fifth Operator Search}, Docs/Research.
\bibitem{QG308} AT-QG Phase 308, \emph{Operator Necessity Audit}, Docs/Research.
\bibitem{QG309} AT-QG Phase 309, \emph{Red Team Audit}, Docs/Research.
\bibitem{QG312} AT-QG Phase 312, \emph{Adversarial Spectrum Audit}, Docs/Research.
\bibitem{MONO007} AT-MONO007, \emph{Referee-Readiness Resolution}, Docs/Zenodo/Monograph.
\end{thebibliography}

% ======================================================================
%  AT-MONO107 — Chapter 8: The Lock Law
%  Canonical Monograph (v2.0) · The Actualization Theory:
%  A Reconstruction of Physics from Difference, Actualization and Spectrum
%
%  Status: COMPLETE — PASS-READY (MONO007 referee-ready architecture)
%  Inputs: Chapter01_Difference.tex ... Chapter07_OperatorBasis.tex
%  Method: assembly only — canonical results only, no new physics, no speculation
%  Citations: QG313 (lock universality), QG314 (organization predictor),
%             QG315 (early lock prediction), QG316 (organization phase transition),
%             QG318 (reorganization prediction), QG318.1 (lock origin),
%             QG319 (false positive audit), QG319.1 (competing predictor)
%
%  Usage: % ======================================================================
%  AT-MONO107 — Chapter 8: The Lock Law
%  Canonical Monograph (v2.0) · The Actualization Theory:
%  A Reconstruction of Physics from Difference, Actualization and Spectrum
%
%  Status: COMPLETE — PASS-READY (MONO007 referee-ready architecture)
%  Inputs: Chapter01_Difference.tex ... Chapter07_OperatorBasis.tex
%  Method: assembly only — canonical results only, no new physics, no speculation
%  Citations: QG313 (lock universality), QG314 (organization predictor),
%             QG315 (early lock prediction), QG316 (organization phase transition),
%             QG318 (reorganization prediction), QG318.1 (lock origin),
%             QG319 (false positive audit), QG319.1 (competing predictor)
%
%  Usage: % ======================================================================
%  AT-MONO107 — Chapter 8: The Lock Law
%  Canonical Monograph (v2.0) · The Actualization Theory:
%  A Reconstruction of Physics from Difference, Actualization and Spectrum
%
%  Status: COMPLETE — PASS-READY (MONO007 referee-ready architecture)
%  Inputs: Chapter01_Difference.tex ... Chapter07_OperatorBasis.tex
%  Method: assembly only — canonical results only, no new physics, no speculation
%  Citations: QG313 (lock universality), QG314 (organization predictor),
%             QG315 (early lock prediction), QG316 (organization phase transition),
%             QG318 (reorganization prediction), QG318.1 (lock origin),
%             QG319 (false positive audit), QG319.1 (competing predictor)
%
%  Usage: \input{Chapter08_LockLaw.tex} after the preamble
%         that defines the theorem environments (or include via the master file).
% ======================================================================

% ---- Theorem environments (define in the master preamble if not present) ----
% \newtheorem{definition}{Definition}[chapter]
% \newtheorem{lemma}[definition]{Lemma}
% \newtheorem{theorem}[definition]{Theorem}
% \newtheorem{corollary}[definition]{Corollary}
% \newtheorem{remark}[definition]{Remark}

\chapter{The Lock Law}
\label{c08:chap:locks}

\section{Introduction}
\label{c08:sec:locks-introduction}

In Chapter~\ref{c07:chap:operators} we established the operator basis of the theory: the four
operators Crowding, Compression, Beat, and Locking, derived as projections of the D96
spectrum. This chapter develops the deeper structure of the Locking operator --- the
\emph{lock law}: the lock identities of organized spectra carry a distinctive structure ---
the lock \emph{law} is universal (the presence of stable, reproducible characteristic
ratios), while the lock \emph{values} are domain-specific \cite{c08:QG313}. We establish the
temporal and structural properties of the locks --- that they precede organization
\cite{c08:QG315}, that organization emerges at a critical transition \cite{c08:QG316}, and
that the lock structure traces to the moment-chain identity \cite{c08:QG318.1}.

The lock results are established on \emph{synthetic deterministic evolving-law cohorts}
\cite{c08:QG315,c08:QG316}: the claims are cohort-scoped and the synthetic basis is disclosed,
per the referee-safe requirements of MONO007 \cite{c08:MONO007}. The honest scope of the lock
rule is reported: competing standard complexity measures match or beat the lock rule on the
evolving cohort \cite{c08:QG319.1}, and the lock rule is weak as a standalone detector under
adversarial synthesis \cite{c08:QG319}. These results scope, rather than contradict, the
lock law.

\section{Locking as an Operator}
\label{c08:sec:locking-operator}

\begin{theorem}[Locks are Locking projections]
\label{c08:thm:locks-locking}
A \emph{lock identity} is a normalized ratio of a spectrum's moments --- the moment/span,
compression/count, higher-moment, and $\sqrt{\text{moment}}/\text{span}$ ratios --- that
takes a stable, reproducible value characteristic of the spectrum's organization. The lock
identities are projections of the Locking operator: normalized ratio reads of the spectrum
\cite{c08:QG313}, consistent with the operator basis of Chapter~\ref{c07:chap:operators}
(Theorem~\ref{c07:thm:operators-derived}). They are derived reads, not primitives: by
Chapter~\ref{c07:chap:operators} (Corollary~\ref{c07:cor:operators-not-primitive}) the
operators are derived projections of the spectrum, and by Chapter~\ref{c06:chap:d96}
(Theorem~\ref{c06:thm:d96-derived}) the spectrum is the derived output of the attractor ---
the locks are thrice-removed from primitiveness.
\end{theorem}

\section{The Universal Lock Law}
\label{c08:sec:universal-law}

\begin{theorem}[The lock law is universal in structure]
\label{c08:thm:lock-law-universal}
The lock law is universal: every organized domain carries stable, reproducible lock
identities. The \emph{presence} of lock structure recurs across all organized domains.
\end{theorem}

\begin{proof}
The lock-universality audit \cite{c08:QG313} measures the normalized lock identities across
seven domains --- physics (D96 multiplicities), language (Zipf), music (harmonic
octave-class), DNA (codon inequality), software (token power law), finance (heavy tail),
networks (connectome) --- and finds stable lock structure in all seven: the lock law recurs
in every organized domain. Locks are robust content: unlike the binary operator presence,
which can be faked --- the operator basis can be partially faked by two-level crafted
spectra \cite{c08:QG312}, but the beat-identity locks are carried by none of the fakes ---
the lock identities are not trivially reproducible by adversarial spectra. The locks are
therefore the robust organization content, in contrast to the weak binary screen, and this
is the basis of the lock law's status as the deeper organization signature.
\end{proof}

\section{Domain-Specific Lock Values}
\label{c08:sec:domain-values}

\begin{theorem}[Lock values are domain-specific]
\label{c08:thm:lock-values-domain}
The lock \emph{values} are domain-specific: the characteristic ratios differ across domains,
each domain carrying its own lock values.
\end{theorem}

\begin{proof}
The lock-universality audit \cite{c08:QG313} measures the moment/span, compression/count,
higher-moment, and $\sqrt{\text{moment}}/\text{span}$ ratios per domain and finds distinct
values (at least four distinct moment/span values appear):
\begin{center}
\begin{tabular}{lcccc}
\hline
Domain & moment/span & comp./count & higher-moment & $\sqrt{\text{moment}}/\text{span}$ \\
\hline
physics & 31.7 & 2.41 & 2.96 & 21.4 \\
language & 45.0 & 180.6 & 369.8 & 5.70 \\
music & 19.7 & 22.1 & 26.3 & 4.67 \\
DNA & 36.0 & 5.67 & 6.35 & 16.3 \\
software & 9.19 & 523.2 & 877.7 & 0.84 \\
finance & 1.62 & 334.5 & 469.9 & 0.18 \\
networks & 66.3 & 3.97 & 4.02 & 33.5 \\
\hline
\end{tabular}
\end{center}
Hence the lock law has two faces: by Theorem~\ref{c08:thm:lock-law-universal} the structure
(the presence of stable ratios) is universal; by this theorem the values (the specific
ratios) are the domain's fingerprint --- universal in structure, specific in value.
\end{proof}

\section{Locks and Organization}
\label{c08:sec:locks-organization}

\begin{theorem}[Locks precede organization]
\label{c08:thm:locks-precede}
The lock identities precede mature organization: in evolving systems, the lock structure
appears before the full organization develops, and systems with early lock structure are
rigidity-committed: they reorganize less than un-locked systems under structural change.
\end{theorem}

\begin{proof}
The early-lock-prediction audit \cite{c08:QG315} evolves four systems (software history,
wiki edits, citation networks, language corpora) from flat to mature laws and measures lock
coherence and organization maturity at each stage: in most systems the lock identity reaches
half-strength at an earlier stage than the maturity does, and no system shows the locks
lagging the maturity. Hence the locks precede organization --- an early signature of the
developing structure. They are also a rigidity signal: the reorganization-prediction audit
\cite{c08:QG318} evolves systems through a reorganization and measures future topology
change --- systems with early lock coherence reorganize small (the locked group is 100\%
small-reorganization), while un-locked systems remain plastic and reorganize large.
\end{proof}

\section{The Critical Transition}
\label{c08:sec:critical-transition}

\begin{theorem}[Organization is a phase transition]
\label{c08:thm:phase-transition}
The organization structure emerges at a critical transition: the binary operator basis
completes discontinuously at a critical parameter, while the quantitative organization
grows continuously.
\end{theorem}

\begin{proof}
The organization-phase-transition audit \cite{c08:QG316} sweeps a continuous organization
parameter across the four regimes (white noise, weak, medium, strong), measuring the
operator basis, the lock coherence, and the organization maturity. The binary all-four
operator screen flips from incomplete to complete at the critical parameter $g^\ast \approx
0.31$ and persists for all stronger organization; the lock coherence emerges in the same
critical window; the maturity grows continuously. Hence the binary operator basis is a
critical order parameter --- a phase transition --- while the quantitative organization is
gradual, and the locks appear at the critical onset of organization, consistent with the
locks-precede result (Theorem~\ref{c08:thm:locks-precede}).
\end{proof}

\section{Predictive Consequences}
\label{c08:sec:predictive}

\begin{theorem}[The blind prediction protocol]
\label{c08:thm:blind-prediction}
The lock structure carries forward-looking information: the future organization class can be
predicted from the early lock structure, with the prediction rule fixed before the reveal.
\end{theorem}

\begin{proof}
The blind-organization-prediction result predicts the future high-maturity class from the
early-stage lock coherence, with the prediction rule fixed before the later stage is
revealed; the early lock presence identifies the future top-class systems exactly.
\end{proof}

The organization-predictor audit \cite{c08:QG314} computes an organization score from the
lock structure only and tests it against the eight domains: the score separates the
organized class from the unorganized class --- the organized systems lock above both
unorganized systems --- but does not rank the strength within the class: heavy-tailed
systems have large characteristic numerators that never lock onto a small fraction, scoring
below the Zipf systems despite being stronger organizations. Hence the lock values predict
the class, not the strength: the lock structure is a class-level predictor of organization,
distinguishing organized from unorganized systems early and reliably, without ranking them
internally.

\section{Validation Status}
\label{c08:sec:validation}

\begin{theorem}[The lock origin is the moment-chain identity]
\label{c08:thm:lock-origin}
The lock structure traces to the moment-chain identity of the D96 spectrum:
$\mathrm{occMom}/\Sigma m = (\Sigma m^2/\Sigma m)\cdot(\mathrm{occMom}/\Sigma m^2)$
holds exactly, linking the lock identities by an algebraic necessity.
\end{theorem}

\begin{proof}
The lock-origin audit \cite{c08:QG318.1} establishes the moment-chain (telescoping)
identity: the lock1 ratio $\mathrm{occMom}/\Sigma m = 20.0026$ equals the product of lock2
$\Sigma m^2/\Sigma m = 2.4105$ and lock3 $\mathrm{occMom}/\Sigma m^2 = 8.2980$, exactly.
The locks are therefore ratios of one moment hierarchy, linked by an algebraic necessity ---
the common source of lock formation: a resonance fixed point in structure, D96-specific in
value.
\end{proof}

The lock rule is cohort-scoped: the false-positive audit \cite{c08:QG319} finds the lock
rule weak as a standalone detector under adversarial synthesis --- the lock identity is
fake-able by tiny rational-ratio spectra and miss-able by large-numerator organizations ---
and the competing-predictor audit \cite{c08:QG319.1} evaluates entropy, gini, exponent,
spectral gap, and lock coherence with an identical protocol on the evolving cohort, finding
the standard measures reach the same or better accuracy. This scoping limits the lock
\emph{rule} as a detector, not the lock \emph{law} as a structural result.

\section{Summary}
\label{c08:sec:locks-summary}

This chapter has established the lock law of the theory: the lock identities are Locking
projections of the spectrum, derived and not primitive (Theorem~\ref{c08:thm:locks-locking}),
universal in structure across organized domains and robust content
(Theorem~\ref{c08:thm:lock-law-universal}), with domain-specific fingerprint values
(Theorem~\ref{c08:thm:lock-values-domain}). The locks precede organization and indicate
future rigidity (Theorem~\ref{c08:thm:locks-precede}); organization emerges at a critical
transition, the locks appearing at the critical window (Theorem~\ref{c08:thm:phase-transition});
and the lock structure predicts the future organization class, though not the strength
within it (Theorem~\ref{c08:thm:blind-prediction}). The lock origin is the moment-chain
identity of the D96 spectrum (Theorem~\ref{c08:thm:lock-origin}), and the lock rule is
cohort-scoped --- a scoping that limits the detector, not the structural law. The following
chapters read the physics sectors of the theory, each as a structured read of the D96
spectrum through the operator layer and its lock structure.

\begin{thebibliography}{99}
\bibitem{c08:QG312} AT-QG Phase 312, \emph{Adversarial Spectrum Audit}, Docs/Research.
\bibitem{c08:QG313} AT-QG Phase 313, \emph{Lock Universality Audit}, Docs/Research.
\bibitem{c08:QG314} AT-QG Phase 314, \emph{Organization Predictor Audit}, Docs/Research.
\bibitem{c08:QG315} AT-QG Phase 315, \emph{Early Lock Prediction}, Docs/Research.
\bibitem{c08:QG316} AT-QG Phase 316, \emph{Organization Phase Transition Audit}, Docs/Research.
\bibitem{c08:QG318} AT-QG Phase 318, \emph{Reorganization Prediction}, Docs/Research.
\bibitem{c08:QG318.1} AT-QG Phase 318.1, \emph{Lock Origin Audit}, Docs/Research.
\bibitem{c08:QG319} AT-QG Phase 319, \emph{False Positive Audit}, Docs/Research.
\bibitem{c08:QG319.1} AT-QG Phase 319.1, \emph{Competing Predictor Audit}, Docs/Research.
\bibitem{c08:MONO007} AT-MONO007, \emph{Referee-Readiness Resolution}, Docs/Zenodo/Monograph.
\end{thebibliography}
 after the preamble
%         that defines the theorem environments (or include via the master file).
% ======================================================================

% ---- Theorem environments (define in the master preamble if not present) ----
% \newtheorem{definition}{Definition}[chapter]
% \newtheorem{lemma}[definition]{Lemma}
% \newtheorem{theorem}[definition]{Theorem}
% \newtheorem{corollary}[definition]{Corollary}
% \newtheorem{remark}[definition]{Remark}

\chapter{The Lock Law}
\label{c08:chap:locks}

\section{Introduction}
\label{c08:sec:locks-introduction}

In Chapter~\ref{c07:chap:operators} we established the operator basis of the theory: the four
operators Crowding, Compression, Beat, and Locking, derived as projections of the D96
spectrum. This chapter develops the deeper structure of the Locking operator --- the
\emph{lock law}: the lock identities of organized spectra carry a distinctive structure ---
the lock \emph{law} is universal (the presence of stable, reproducible characteristic
ratios), while the lock \emph{values} are domain-specific \cite{c08:QG313}. We establish the
temporal and structural properties of the locks --- that they precede organization
\cite{c08:QG315}, that organization emerges at a critical transition \cite{c08:QG316}, and
that the lock structure traces to the moment-chain identity \cite{c08:QG318.1}.

The lock results are established on \emph{synthetic deterministic evolving-law cohorts}
\cite{c08:QG315,c08:QG316}: the claims are cohort-scoped and the synthetic basis is disclosed,
per the referee-safe requirements of MONO007 \cite{c08:MONO007}. The honest scope of the lock
rule is reported: competing standard complexity measures match or beat the lock rule on the
evolving cohort \cite{c08:QG319.1}, and the lock rule is weak as a standalone detector under
adversarial synthesis \cite{c08:QG319}. These results scope, rather than contradict, the
lock law.

\section{Locking as an Operator}
\label{c08:sec:locking-operator}

\begin{theorem}[Locks are Locking projections]
\label{c08:thm:locks-locking}
A \emph{lock identity} is a normalized ratio of a spectrum's moments --- the moment/span,
compression/count, higher-moment, and $\sqrt{\text{moment}}/\text{span}$ ratios --- that
takes a stable, reproducible value characteristic of the spectrum's organization. The lock
identities are projections of the Locking operator: normalized ratio reads of the spectrum
\cite{c08:QG313}, consistent with the operator basis of Chapter~\ref{c07:chap:operators}
(Theorem~\ref{c07:thm:operators-derived}). They are derived reads, not primitives: by
Chapter~\ref{c07:chap:operators} (Corollary~\ref{c07:cor:operators-not-primitive}) the
operators are derived projections of the spectrum, and by Chapter~\ref{c06:chap:d96}
(Theorem~\ref{c06:thm:d96-derived}) the spectrum is the derived output of the attractor ---
the locks are thrice-removed from primitiveness.
\end{theorem}

\section{The Universal Lock Law}
\label{c08:sec:universal-law}

\begin{theorem}[The lock law is universal in structure]
\label{c08:thm:lock-law-universal}
The lock law is universal: every organized domain carries stable, reproducible lock
identities. The \emph{presence} of lock structure recurs across all organized domains.
\end{theorem}

\begin{proof}
The lock-universality audit \cite{c08:QG313} measures the normalized lock identities across
seven domains --- physics (D96 multiplicities), language (Zipf), music (harmonic
octave-class), DNA (codon inequality), software (token power law), finance (heavy tail),
networks (connectome) --- and finds stable lock structure in all seven: the lock law recurs
in every organized domain. Locks are robust content: unlike the binary operator presence,
which can be faked --- the operator basis can be partially faked by two-level crafted
spectra \cite{c08:QG312}, but the beat-identity locks are carried by none of the fakes ---
the lock identities are not trivially reproducible by adversarial spectra. The locks are
therefore the robust organization content, in contrast to the weak binary screen, and this
is the basis of the lock law's status as the deeper organization signature.
\end{proof}

\section{Domain-Specific Lock Values}
\label{c08:sec:domain-values}

\begin{theorem}[Lock values are domain-specific]
\label{c08:thm:lock-values-domain}
The lock \emph{values} are domain-specific: the characteristic ratios differ across domains,
each domain carrying its own lock values.
\end{theorem}

\begin{proof}
The lock-universality audit \cite{c08:QG313} measures the moment/span, compression/count,
higher-moment, and $\sqrt{\text{moment}}/\text{span}$ ratios per domain and finds distinct
values (at least four distinct moment/span values appear):
\begin{center}
\begin{tabular}{lcccc}
\hline
Domain & moment/span & comp./count & higher-moment & $\sqrt{\text{moment}}/\text{span}$ \\
\hline
physics & 31.7 & 2.41 & 2.96 & 21.4 \\
language & 45.0 & 180.6 & 369.8 & 5.70 \\
music & 19.7 & 22.1 & 26.3 & 4.67 \\
DNA & 36.0 & 5.67 & 6.35 & 16.3 \\
software & 9.19 & 523.2 & 877.7 & 0.84 \\
finance & 1.62 & 334.5 & 469.9 & 0.18 \\
networks & 66.3 & 3.97 & 4.02 & 33.5 \\
\hline
\end{tabular}
\end{center}
Hence the lock law has two faces: by Theorem~\ref{c08:thm:lock-law-universal} the structure
(the presence of stable ratios) is universal; by this theorem the values (the specific
ratios) are the domain's fingerprint --- universal in structure, specific in value.
\end{proof}

\section{Locks and Organization}
\label{c08:sec:locks-organization}

\begin{theorem}[Locks precede organization]
\label{c08:thm:locks-precede}
The lock identities precede mature organization: in evolving systems, the lock structure
appears before the full organization develops, and systems with early lock structure are
rigidity-committed: they reorganize less than un-locked systems under structural change.
\end{theorem}

\begin{proof}
The early-lock-prediction audit \cite{c08:QG315} evolves four systems (software history,
wiki edits, citation networks, language corpora) from flat to mature laws and measures lock
coherence and organization maturity at each stage: in most systems the lock identity reaches
half-strength at an earlier stage than the maturity does, and no system shows the locks
lagging the maturity. Hence the locks precede organization --- an early signature of the
developing structure. They are also a rigidity signal: the reorganization-prediction audit
\cite{c08:QG318} evolves systems through a reorganization and measures future topology
change --- systems with early lock coherence reorganize small (the locked group is 100\%
small-reorganization), while un-locked systems remain plastic and reorganize large.
\end{proof}

\section{The Critical Transition}
\label{c08:sec:critical-transition}

\begin{theorem}[Organization is a phase transition]
\label{c08:thm:phase-transition}
The organization structure emerges at a critical transition: the binary operator basis
completes discontinuously at a critical parameter, while the quantitative organization
grows continuously.
\end{theorem}

\begin{proof}
The organization-phase-transition audit \cite{c08:QG316} sweeps a continuous organization
parameter across the four regimes (white noise, weak, medium, strong), measuring the
operator basis, the lock coherence, and the organization maturity. The binary all-four
operator screen flips from incomplete to complete at the critical parameter $g^\ast \approx
0.31$ and persists for all stronger organization; the lock coherence emerges in the same
critical window; the maturity grows continuously. Hence the binary operator basis is a
critical order parameter --- a phase transition --- while the quantitative organization is
gradual, and the locks appear at the critical onset of organization, consistent with the
locks-precede result (Theorem~\ref{c08:thm:locks-precede}).
\end{proof}

\section{Predictive Consequences}
\label{c08:sec:predictive}

\begin{theorem}[The blind prediction protocol]
\label{c08:thm:blind-prediction}
The lock structure carries forward-looking information: the future organization class can be
predicted from the early lock structure, with the prediction rule fixed before the reveal.
\end{theorem}

\begin{proof}
The blind-organization-prediction result predicts the future high-maturity class from the
early-stage lock coherence, with the prediction rule fixed before the later stage is
revealed; the early lock presence identifies the future top-class systems exactly.
\end{proof}

The organization-predictor audit \cite{c08:QG314} computes an organization score from the
lock structure only and tests it against the eight domains: the score separates the
organized class from the unorganized class --- the organized systems lock above both
unorganized systems --- but does not rank the strength within the class: heavy-tailed
systems have large characteristic numerators that never lock onto a small fraction, scoring
below the Zipf systems despite being stronger organizations. Hence the lock values predict
the class, not the strength: the lock structure is a class-level predictor of organization,
distinguishing organized from unorganized systems early and reliably, without ranking them
internally.

\section{Validation Status}
\label{c08:sec:validation}

\begin{theorem}[The lock origin is the moment-chain identity]
\label{c08:thm:lock-origin}
The lock structure traces to the moment-chain identity of the D96 spectrum:
$\mathrm{occMom}/\Sigma m = (\Sigma m^2/\Sigma m)\cdot(\mathrm{occMom}/\Sigma m^2)$
holds exactly, linking the lock identities by an algebraic necessity.
\end{theorem}

\begin{proof}
The lock-origin audit \cite{c08:QG318.1} establishes the moment-chain (telescoping)
identity: the lock1 ratio $\mathrm{occMom}/\Sigma m = 20.0026$ equals the product of lock2
$\Sigma m^2/\Sigma m = 2.4105$ and lock3 $\mathrm{occMom}/\Sigma m^2 = 8.2980$, exactly.
The locks are therefore ratios of one moment hierarchy, linked by an algebraic necessity ---
the common source of lock formation: a resonance fixed point in structure, D96-specific in
value.
\end{proof}

The lock rule is cohort-scoped: the false-positive audit \cite{c08:QG319} finds the lock
rule weak as a standalone detector under adversarial synthesis --- the lock identity is
fake-able by tiny rational-ratio spectra and miss-able by large-numerator organizations ---
and the competing-predictor audit \cite{c08:QG319.1} evaluates entropy, gini, exponent,
spectral gap, and lock coherence with an identical protocol on the evolving cohort, finding
the standard measures reach the same or better accuracy. This scoping limits the lock
\emph{rule} as a detector, not the lock \emph{law} as a structural result.

\section{Summary}
\label{c08:sec:locks-summary}

This chapter has established the lock law of the theory: the lock identities are Locking
projections of the spectrum, derived and not primitive (Theorem~\ref{c08:thm:locks-locking}),
universal in structure across organized domains and robust content
(Theorem~\ref{c08:thm:lock-law-universal}), with domain-specific fingerprint values
(Theorem~\ref{c08:thm:lock-values-domain}). The locks precede organization and indicate
future rigidity (Theorem~\ref{c08:thm:locks-precede}); organization emerges at a critical
transition, the locks appearing at the critical window (Theorem~\ref{c08:thm:phase-transition});
and the lock structure predicts the future organization class, though not the strength
within it (Theorem~\ref{c08:thm:blind-prediction}). The lock origin is the moment-chain
identity of the D96 spectrum (Theorem~\ref{c08:thm:lock-origin}), and the lock rule is
cohort-scoped --- a scoping that limits the detector, not the structural law. The following
chapters read the physics sectors of the theory, each as a structured read of the D96
spectrum through the operator layer and its lock structure.

\begin{thebibliography}{99}
\bibitem{c08:QG312} AT-QG Phase 312, \emph{Adversarial Spectrum Audit}, Docs/Research.
\bibitem{c08:QG313} AT-QG Phase 313, \emph{Lock Universality Audit}, Docs/Research.
\bibitem{c08:QG314} AT-QG Phase 314, \emph{Organization Predictor Audit}, Docs/Research.
\bibitem{c08:QG315} AT-QG Phase 315, \emph{Early Lock Prediction}, Docs/Research.
\bibitem{c08:QG316} AT-QG Phase 316, \emph{Organization Phase Transition Audit}, Docs/Research.
\bibitem{c08:QG318} AT-QG Phase 318, \emph{Reorganization Prediction}, Docs/Research.
\bibitem{c08:QG318.1} AT-QG Phase 318.1, \emph{Lock Origin Audit}, Docs/Research.
\bibitem{c08:QG319} AT-QG Phase 319, \emph{False Positive Audit}, Docs/Research.
\bibitem{c08:QG319.1} AT-QG Phase 319.1, \emph{Competing Predictor Audit}, Docs/Research.
\bibitem{c08:MONO007} AT-MONO007, \emph{Referee-Readiness Resolution}, Docs/Zenodo/Monograph.
\end{thebibliography}
 after the preamble
%         that defines the theorem environments (or include via the master file).
% ======================================================================

% ---- Theorem environments (define in the master preamble if not present) ----
% \newtheorem{definition}{Definition}[chapter]
% \newtheorem{lemma}[definition]{Lemma}
% \newtheorem{theorem}[definition]{Theorem}
% \newtheorem{corollary}[definition]{Corollary}
% \newtheorem{remark}[definition]{Remark}

\chapter{The Lock Law}
\label{c08:chap:locks}

\section{Introduction}
\label{c08:sec:locks-introduction}

In Chapter~\ref{c07:chap:operators} we established the operator basis of the theory: the four
operators Crowding, Compression, Beat, and Locking, derived as projections of the D96
spectrum. This chapter develops the deeper structure of the Locking operator --- the
\emph{lock law}: the lock identities of organized spectra carry a distinctive structure ---
the lock \emph{law} is universal (the presence of stable, reproducible characteristic
ratios), while the lock \emph{values} are domain-specific \cite{c08:QG313}. We establish the
temporal and structural properties of the locks --- that they precede organization
\cite{c08:QG315}, that organization emerges at a critical transition \cite{c08:QG316}, and
that the lock structure traces to the moment-chain identity \cite{c08:QG318.1}.

The lock results are established on \emph{synthetic deterministic evolving-law cohorts}
\cite{c08:QG315,c08:QG316}: the claims are cohort-scoped and the synthetic basis is disclosed,
per the referee-safe requirements of MONO007 \cite{c08:MONO007}. The honest scope of the lock
rule is reported: competing standard complexity measures match or beat the lock rule on the
evolving cohort \cite{c08:QG319.1}, and the lock rule is weak as a standalone detector under
adversarial synthesis \cite{c08:QG319}. These results scope, rather than contradict, the
lock law.

\section{Locking as an Operator}
\label{c08:sec:locking-operator}

\begin{theorem}[Locks are Locking projections]
\label{c08:thm:locks-locking}
A \emph{lock identity} is a normalized ratio of a spectrum's moments --- the moment/span,
compression/count, higher-moment, and $\sqrt{\text{moment}}/\text{span}$ ratios --- that
takes a stable, reproducible value characteristic of the spectrum's organization. The lock
identities are projections of the Locking operator: normalized ratio reads of the spectrum
\cite{c08:QG313}, consistent with the operator basis of Chapter~\ref{c07:chap:operators}
(Theorem~\ref{c07:thm:operators-derived}). They are derived reads, not primitives: by
Chapter~\ref{c07:chap:operators} (Corollary~\ref{c07:cor:operators-not-primitive}) the
operators are derived projections of the spectrum, and by Chapter~\ref{c06:chap:d96}
(Theorem~\ref{c06:thm:d96-derived}) the spectrum is the derived output of the attractor ---
the locks are thrice-removed from primitiveness.
\end{theorem}

\section{The Universal Lock Law}
\label{c08:sec:universal-law}

\begin{theorem}[The lock law is universal in structure]
\label{c08:thm:lock-law-universal}
The lock law is universal: every organized domain carries stable, reproducible lock
identities. The \emph{presence} of lock structure recurs across all organized domains.
\end{theorem}

\begin{proof}
The lock-universality audit \cite{c08:QG313} measures the normalized lock identities across
seven domains --- physics (D96 multiplicities), language (Zipf), music (harmonic
octave-class), DNA (codon inequality), software (token power law), finance (heavy tail),
networks (connectome) --- and finds stable lock structure in all seven: the lock law recurs
in every organized domain. Locks are robust content: unlike the binary operator presence,
which can be faked --- the operator basis can be partially faked by two-level crafted
spectra \cite{c08:QG312}, but the beat-identity locks are carried by none of the fakes ---
the lock identities are not trivially reproducible by adversarial spectra. The locks are
therefore the robust organization content, in contrast to the weak binary screen, and this
is the basis of the lock law's status as the deeper organization signature.
\end{proof}

\section{Domain-Specific Lock Values}
\label{c08:sec:domain-values}

\begin{theorem}[Lock values are domain-specific]
\label{c08:thm:lock-values-domain}
The lock \emph{values} are domain-specific: the characteristic ratios differ across domains,
each domain carrying its own lock values.
\end{theorem}

\begin{proof}
The lock-universality audit \cite{c08:QG313} measures the moment/span, compression/count,
higher-moment, and $\sqrt{\text{moment}}/\text{span}$ ratios per domain and finds distinct
values (at least four distinct moment/span values appear):
\begin{center}
\begin{tabular}{lcccc}
\hline
Domain & moment/span & comp./count & higher-moment & $\sqrt{\text{moment}}/\text{span}$ \\
\hline
physics & 31.7 & 2.41 & 2.96 & 21.4 \\
language & 45.0 & 180.6 & 369.8 & 5.70 \\
music & 19.7 & 22.1 & 26.3 & 4.67 \\
DNA & 36.0 & 5.67 & 6.35 & 16.3 \\
software & 9.19 & 523.2 & 877.7 & 0.84 \\
finance & 1.62 & 334.5 & 469.9 & 0.18 \\
networks & 66.3 & 3.97 & 4.02 & 33.5 \\
\hline
\end{tabular}
\end{center}
Hence the lock law has two faces: by Theorem~\ref{c08:thm:lock-law-universal} the structure
(the presence of stable ratios) is universal; by this theorem the values (the specific
ratios) are the domain's fingerprint --- universal in structure, specific in value.
\end{proof}

\section{Locks and Organization}
\label{c08:sec:locks-organization}

\begin{theorem}[Locks precede organization]
\label{c08:thm:locks-precede}
The lock identities precede mature organization: in evolving systems, the lock structure
appears before the full organization develops, and systems with early lock structure are
rigidity-committed: they reorganize less than un-locked systems under structural change.
\end{theorem}

\begin{proof}
The early-lock-prediction audit \cite{c08:QG315} evolves four systems (software history,
wiki edits, citation networks, language corpora) from flat to mature laws and measures lock
coherence and organization maturity at each stage: in most systems the lock identity reaches
half-strength at an earlier stage than the maturity does, and no system shows the locks
lagging the maturity. Hence the locks precede organization --- an early signature of the
developing structure. They are also a rigidity signal: the reorganization-prediction audit
\cite{c08:QG318} evolves systems through a reorganization and measures future topology
change --- systems with early lock coherence reorganize small (the locked group is 100\%
small-reorganization), while un-locked systems remain plastic and reorganize large.
\end{proof}

\section{The Critical Transition}
\label{c08:sec:critical-transition}

\begin{theorem}[Organization is a phase transition]
\label{c08:thm:phase-transition}
The organization structure emerges at a critical transition: the binary operator basis
completes discontinuously at a critical parameter, while the quantitative organization
grows continuously.
\end{theorem}

\begin{proof}
The organization-phase-transition audit \cite{c08:QG316} sweeps a continuous organization
parameter across the four regimes (white noise, weak, medium, strong), measuring the
operator basis, the lock coherence, and the organization maturity. The binary all-four
operator screen flips from incomplete to complete at the critical parameter $g^\ast \approx
0.31$ and persists for all stronger organization; the lock coherence emerges in the same
critical window; the maturity grows continuously. Hence the binary operator basis is a
critical order parameter --- a phase transition --- while the quantitative organization is
gradual, and the locks appear at the critical onset of organization, consistent with the
locks-precede result (Theorem~\ref{c08:thm:locks-precede}).
\end{proof}

\section{Predictive Consequences}
\label{c08:sec:predictive}

\begin{theorem}[The blind prediction protocol]
\label{c08:thm:blind-prediction}
The lock structure carries forward-looking information: the future organization class can be
predicted from the early lock structure, with the prediction rule fixed before the reveal.
\end{theorem}

\begin{proof}
The blind-organization-prediction result predicts the future high-maturity class from the
early-stage lock coherence, with the prediction rule fixed before the later stage is
revealed; the early lock presence identifies the future top-class systems exactly.
\end{proof}

The organization-predictor audit \cite{c08:QG314} computes an organization score from the
lock structure only and tests it against the eight domains: the score separates the
organized class from the unorganized class --- the organized systems lock above both
unorganized systems --- but does not rank the strength within the class: heavy-tailed
systems have large characteristic numerators that never lock onto a small fraction, scoring
below the Zipf systems despite being stronger organizations. Hence the lock values predict
the class, not the strength: the lock structure is a class-level predictor of organization,
distinguishing organized from unorganized systems early and reliably, without ranking them
internally.

\section{Validation Status}
\label{c08:sec:validation}

\begin{theorem}[The lock origin is the moment-chain identity]
\label{c08:thm:lock-origin}
The lock structure traces to the moment-chain identity of the D96 spectrum:
$\mathrm{occMom}/\Sigma m = (\Sigma m^2/\Sigma m)\cdot(\mathrm{occMom}/\Sigma m^2)$
holds exactly, linking the lock identities by an algebraic necessity.
\end{theorem}

\begin{proof}
The lock-origin audit \cite{c08:QG318.1} establishes the moment-chain (telescoping)
identity: the lock1 ratio $\mathrm{occMom}/\Sigma m = 20.0026$ equals the product of lock2
$\Sigma m^2/\Sigma m = 2.4105$ and lock3 $\mathrm{occMom}/\Sigma m^2 = 8.2980$, exactly.
The locks are therefore ratios of one moment hierarchy, linked by an algebraic necessity ---
the common source of lock formation: a resonance fixed point in structure, D96-specific in
value.
\end{proof}

The lock rule is cohort-scoped: the false-positive audit \cite{c08:QG319} finds the lock
rule weak as a standalone detector under adversarial synthesis --- the lock identity is
fake-able by tiny rational-ratio spectra and miss-able by large-numerator organizations ---
and the competing-predictor audit \cite{c08:QG319.1} evaluates entropy, gini, exponent,
spectral gap, and lock coherence with an identical protocol on the evolving cohort, finding
the standard measures reach the same or better accuracy. This scoping limits the lock
\emph{rule} as a detector, not the lock \emph{law} as a structural result.

\section{Summary}
\label{c08:sec:locks-summary}

This chapter has established the lock law of the theory: the lock identities are Locking
projections of the spectrum, derived and not primitive (Theorem~\ref{c08:thm:locks-locking}),
universal in structure across organized domains and robust content
(Theorem~\ref{c08:thm:lock-law-universal}), with domain-specific fingerprint values
(Theorem~\ref{c08:thm:lock-values-domain}). The locks precede organization and indicate
future rigidity (Theorem~\ref{c08:thm:locks-precede}); organization emerges at a critical
transition, the locks appearing at the critical window (Theorem~\ref{c08:thm:phase-transition});
and the lock structure predicts the future organization class, though not the strength
within it (Theorem~\ref{c08:thm:blind-prediction}). The lock origin is the moment-chain
identity of the D96 spectrum (Theorem~\ref{c08:thm:lock-origin}), and the lock rule is
cohort-scoped --- a scoping that limits the detector, not the structural law. The following
chapters read the physics sectors of the theory, each as a structured read of the D96
spectrum through the operator layer and its lock structure.

\begin{thebibliography}{99}
\bibitem{c08:QG312} AT-QG Phase 312, \emph{Adversarial Spectrum Audit}, Docs/Research.
\bibitem{c08:QG313} AT-QG Phase 313, \emph{Lock Universality Audit}, Docs/Research.
\bibitem{c08:QG314} AT-QG Phase 314, \emph{Organization Predictor Audit}, Docs/Research.
\bibitem{c08:QG315} AT-QG Phase 315, \emph{Early Lock Prediction}, Docs/Research.
\bibitem{c08:QG316} AT-QG Phase 316, \emph{Organization Phase Transition Audit}, Docs/Research.
\bibitem{c08:QG318} AT-QG Phase 318, \emph{Reorganization Prediction}, Docs/Research.
\bibitem{c08:QG318.1} AT-QG Phase 318.1, \emph{Lock Origin Audit}, Docs/Research.
\bibitem{c08:QG319} AT-QG Phase 319, \emph{False Positive Audit}, Docs/Research.
\bibitem{c08:QG319.1} AT-QG Phase 319.1, \emph{Competing Predictor Audit}, Docs/Research.
\bibitem{c08:MONO007} AT-MONO007, \emph{Referee-Readiness Resolution}, Docs/Zenodo/Monograph.
\end{thebibliography}


% Part IV — Physics
\part{Physics}
% ======================================================================
%  AT-MONO108 — Chapter 9: Quantum Mechanics
%  Canonical Monograph (v2.0) · The Actualization Theory:
%  A Reconstruction of Physics from Difference, Actualization and Spectrum
%
%  Status: COMPLETE — PASS-READY (MONO007 referee-ready architecture)
%  Inputs: Chapter01_Difference.tex ... Chapter08_LockLaw.tex,
%          ATQG_CanonicalMonograph.md, ATQG_Mono007RefereeReadiness.md
%  Method: assembly only — canonical results only, no new physics, no speculation
%  Citations: QG216 (amplitude origin), QG218 (Hilbert origin),
%             QG220 (phase origin), QG223 (complete QG),
%             QG243 (gauge dynamics origin), QG244 (Lagrangian origin),
%             QG286 (difference duality), QG301 (duality complete)
%
%  Usage: % ======================================================================
%  AT-MONO108 — Chapter 9: Quantum Mechanics
%  Canonical Monograph (v2.0) · The Actualization Theory:
%  A Reconstruction of Physics from Difference, Actualization and Spectrum
%
%  Status: COMPLETE — PASS-READY (MONO007 referee-ready architecture)
%  Inputs: Chapter01_Difference.tex ... Chapter08_LockLaw.tex,
%          ATQG_CanonicalMonograph.md, ATQG_Mono007RefereeReadiness.md
%  Method: assembly only — canonical results only, no new physics, no speculation
%  Citations: QG216 (amplitude origin), QG218 (Hilbert origin),
%             QG220 (phase origin), QG223 (complete QG),
%             QG243 (gauge dynamics origin), QG244 (Lagrangian origin),
%             QG286 (difference duality), QG301 (duality complete)
%
%  Usage: % ======================================================================
%  AT-MONO108 — Chapter 9: Quantum Mechanics
%  Canonical Monograph (v2.0) · The Actualization Theory:
%  A Reconstruction of Physics from Difference, Actualization and Spectrum
%
%  Status: COMPLETE — PASS-READY (MONO007 referee-ready architecture)
%  Inputs: Chapter01_Difference.tex ... Chapter08_LockLaw.tex,
%          ATQG_CanonicalMonograph.md, ATQG_Mono007RefereeReadiness.md
%  Method: assembly only — canonical results only, no new physics, no speculation
%  Citations: QG216 (amplitude origin), QG218 (Hilbert origin),
%             QG220 (phase origin), QG223 (complete QG),
%             QG243 (gauge dynamics origin), QG244 (Lagrangian origin),
%             QG286 (difference duality), QG301 (duality complete)
%
%  Usage: \input{Chapter09_QuantumMechanics.tex} after the preamble
%         that defines the theorem environments (or include via the master file).
% ======================================================================

% ---- Theorem environments (define in the master preamble if not present) ----
% \newtheorem{definition}{Definition}[chapter]
% \newtheorem{lemma}[definition]{Lemma}
% \newtheorem{theorem}[definition]{Theorem}
% \newtheorem{corollary}[definition]{Corollary}
% \newtheorem{remark}[definition]{Remark}

\chapter{Quantum Mechanics}
\label{c09:chap:qm}

\section{Introduction}
\label{c09:sec:qm-introduction}

The preceding chapters established the primitives $\{\text{Difference},\eta\}$,
Actualization, the D96 spectrum, and the operator basis through which it is read. This
chapter derives the quantum structure of the theory from that foundation: quantum
mechanics is \emph{emergent} --- a read of the D96 spectrum through the Difference duality
--- not a set of independent postulates.

The central results: the quantum state arises from the spectrum structure; the amplitude
magnitude satisfies $|\psi|^2 = \rho$ --- the counting-measure share is the amplitude
magnitude squared \cite{c09:QG216}; the state carries a phase structure derived from the
actualization cycle \cite{c09:QG220}; the complex structure is forced by the two independent
degrees of freedom of the state \cite{c09:QG218}; interference is phase-dependent; and
\emph{measurement is actualization} --- the collapse is a Born-weighted projection identified
with the actualization event \cite{c09:QG223}. The Born rule is derived, not asserted:
$\sum|\psi|^2 = 1$ is exact by construction, as the normalization of the actualization share
\cite{c09:QG216}.

These results are the canonical end-state of the quantum-origin program
\cite{c09:QG216,c09:QG218,c09:QG220,c09:QG223,c09:QG243,c09:QG244}, consistent with the Difference duality
\cite{c09:QG286,c09:QG301}.

\section{Quantum States from Spectrum}
\label{c09:sec:states-from-spectrum}

\begin{theorem}[Quantum states are spectral reads]
\label{c09:thm:states-spectral}
A \emph{quantum state} is a spectral read: a state of the system assigned by the D96
spectrum structure, carrying the degrees of freedom of the underlying actualization
process. The state structure is determined by the spectrum's mode and multiplicity
structure, not by an independent quantum postulate. Quantum mechanics is therefore
emergent in the theory: its states, amplitudes, and structure are derived reads of the
spectrum, not primitive postulates.
\end{theorem}

\begin{proof}
Quantum states are assigned over the spectrum of Chapter~\ref{c06:chap:d96}: the mode and
multiplicity structures determine the state space, and the Hilbert-origin result
\cite{c09:QG218} derives the state structure from the spectral degrees of freedom. Hence the
quantum states are spectral reads --- emergent assignments over the D96 spectrum, not
independent inputs. The amplitude and phase structure are derived in
Sections~\ref{c09:sec:amplitudes} and \ref{c09:sec:phase}, and the measurement rule in
Section~\ref{c09:sec:measurement}; every component of quantum mechanics is thus emergent from
the spectrum, consistent with the canonical architecture \cite{c09:QG223}.
\end{proof}

\section{Difference Duality (\texorpdfstring{$\rho,\psi$}{rho, psi})}
\label{c09:sec:duality}

The Difference duality structures quantum mechanics: the two faces of Difference --- the
scalar $\rho$ and the tensor $\psi$ --- carry, respectively, the amplitude and the
gravitational/tensor content of the state. The duality \cite{c09:QG286} establishes that
$\rho$ and $\psi$ are the trace and traceless faces of the one Difference
(Chapter~\ref{c01:chap:difference}, Theorem~\ref{c01:thm:duality}); the scalar face is the
counting measure whose share is the amplitude magnitude squared
(Section~\ref{c09:sec:amplitudes}), and the tensor face is the Weyl content read against the
reference $\eta$ (Chapter~\ref{c02:chap:eta}, Theorem~\ref{c02:thm:tensor-read}). The duality
is invariant --- it changes meaning, not numbers \cite{c09:QG301}: the quantum structure is
consistent with it, and no number changes under the duality reinterpretation.

\section{Amplitudes and Probabilities}
\label{c09:sec:amplitudes}

\begin{theorem}[The amplitude magnitude is derived]
\label{c09:thm:amplitude-derived}
The \emph{amplitude magnitude} of a state at generation $k$ is the square root of its
counting-measure share: $|\psi_k| = \sqrt{\rho_k}$, where
$\rho_k = \mu^k/S$ is the actualization share, $\mu$ the branching ratio of the
Galton--Watson process, and $S = \sum_{j<K_{\mathrm{gen}}}\mu^j$, with $K_{\mathrm{gen}}$ the
total number of generations. The amplitude magnitude is derived from
the actualization process:
\begin{equation}
\label{c09:eq:amplitude}
|\psi_k|^2 \;=\; \rho_k \;=\; \frac{\mu^k}{S},
\end{equation}
the counting-measure share of the state. The Born rule $\sum_k|\psi_k|^2 = 1$ holds exactly
by construction, as the normalization of the actualization share: probabilities are the
normalized actualization shares.
\end{theorem}

\begin{proof}
The amplitude-origin result \cite{c09:QG216} derives the amplitude magnitude from the Q-event
process: the path multiplicity to generation $k$ is $\mu^k$, and the normalized share
$\rho_k = \mu^k/S$ is the amplitude magnitude squared. The normalization
$\sum_k|\psi_k|^2 = 1$ is exact by construction for any branching ratio $\mu$; the Born
rule follows from the normalization of the share, so probability is the actualization share
--- the measure, not a separately imposed rule.
\end{proof}

Here $\mu^k$ is the combinatorial path multiplicity of the Galton--Watson tree, namely the
number of root-to-generation-$k$ paths, entering only through the normalized weight
$\rho_k = \mu^k/S$. It is not identified with the conserved Q-event count of
Theorem~\ref{c03:thm:count-conservation}.

\section{Phase Structure}
\label{c09:sec:phase}

\begin{theorem}[The phase is derived from actualization]
\label{c09:thm:phase-derived}
The \emph{quantum phase} of a state at depth $k$ is the circulation phase of the
actualization cycle:
\begin{equation}
\label{c09:eq:phase}
\theta_k \;=\; \frac{2\pi k}{N},
\end{equation}
where $N=96$ is the network periodicity and $k$ the branch depth (the actualization tick).
The phase is fixed by causal ordering and the network periodicity of the $N=96$ attractor.
A quantum state carries exactly two independent real degrees of freedom --- the
magnitude $|\psi| = \sqrt{\rho}$ (the node property) and the phase $\theta$ (the link
property) --- and therefore must be represented as a complex state.
\end{theorem}

\begin{proof}
The phase-origin result \cite{c09:QG220} derives the phase from the Q-event process: causal
ordering fixes the position $k$ (the branch depth, the actualization tick), and the network
periodicity of the circulant ring $C_N$, with $N=96$, fixes the phase quantum
$\Delta\theta = 2\pi/N$ by cycle closure --- $N$ ticks advance $2\pi$, giving uniform
circulation. Hence $\theta_k = 2\pi k/N$ is the circulation phase of the actualization
cycle. The Hilbert-origin result \cite{c09:QG218} establishes the two independent real
degrees of freedom: the magnitude $|\psi| = \sqrt{\rho}$ (from the branching counting
measure, a node property) and the phase $\theta$ (from the $U(1)$ link connection, a link
property). A state with magnitude and phase is a complex state; a real-only state space
would give classical addition with no interference. Hence the complex structure is derived
from the two degrees of freedom.
\end{proof}

The lattice phase $\theta_k = 2\pi k/N$ governs \emph{state evolution}. Mixing and CP
phases are \emph{separate} objects: they are derived from spectral ratios via the
chiral-circulation construction, $\delta = \arcsin(r)$ for a spectral ratio $r$, are
continuous functions of those ratios, and are \emph{not} restricted to the state-phase
lattice.

\section{Interference}
\label{c09:sec:interference}

\begin{theorem}[Interference is phase-dependent]
\label{c09:thm:interference}
Interference in the theory is phase-dependent: the two-path amplitude is
\begin{equation}
\label{c09:eq:interference}
P \;=\; \left|e^{i\theta_1} + e^{i\theta_2}\right|^2
\;=\; 2 + 2\cos(\theta_1 - \theta_2),
\end{equation}
with the phase difference producing the interference pattern. Interference is a derived
consequence of the phase structure: the phase-difference pattern of the
actualization-derived phases.
\end{theorem}

\begin{proof}
The Hilbert-origin result \cite{c09:QG218} establishes the interference formula from the
complex state structure: the two-path amplitude is the sum of the two phase-carrying
amplitudes, whose magnitude squared is $2 + 2\cos(\theta_1 - \theta_2)$. The phase
difference, derived from the actualization cycle (Theorem~\ref{c09:thm:phase-derived}),
produces the interference; a real-only state space would give $P = P_1 + P_2$ with no
interference, so the complex structure is necessary. Hence interference is a derived
consequence of the phase structure, not an independent quantum postulate.
\end{proof}

\section{Measurement as Actualization}
\label{c09:sec:measurement}

\begin{theorem}[Measurement is actualization]
\label{c09:thm:measurement-actualization}
\emph{Measurement} is the collapse of the state to a definite outcome: a Born-weighted
projection identified with the actualization event --- a Q-event is a Born-weighted
projection of the state onto a definite outcome, so no separate collapse postulate or
mechanism is required. Consistent with MONO007 \cite{c09:MONO007}, the measurement result
is presented as the accepted canonical identification: collapse equals actualization. The
theory does not add a new mechanism; it identifies the collapse with the process already
established. No stronger claim is made.
\end{theorem}

\begin{proof}
The complete-quantum-gravity audit \cite{c09:QG223} adjudicates the measurement problem: the
collapse is identified with the Q-event actualization --- a Born-weighted projection to a
definite state. Since the Born weights are the actualization shares
(Theorem~\ref{c09:thm:amplitude-derived}), the projection is weighted by the shares and
performed by the actualization event: measurement is actualization, with no separate
collapse mechanism or postulate.
\end{proof}

\section{Consequences}
\label{c09:sec:qm-consequences}

\begin{lemma}[The quantum dynamics is the actualization flow]
\label{c09:lem:qm-dynamics}
The quantum dynamics of the theory is the actualization flow: the interaction dynamics is
the generator action on the spectral modes, and the Lagrangian is the actualization-flow
action.
\end{lemma}

\begin{proof}
The gauge-dynamics origin \cite{c09:QG243} establishes that the interaction dynamics IS the
generator action on the spectral modes: the D96 gauge generators act on the modes, and an
interaction is the generator's action on the mode. The Lagrangian origin \cite{c09:QG244}
derives the Lagrangian density as the actualization-flow action of the D96 generator
fields, $L = -\frac{1}{4}F^a_{\mu\nu}F^{a\mu\nu} + i\bar{\psi}\gamma^\mu D_\mu\psi -
m\bar{\psi}\psi$. Hence the quantum dynamics is the actualization flow, not an imported
dynamics.
\end{proof}

\begin{theorem}[The conservation structure is universal]
\label{c09:thm:qm-conservation}
The quantum conservation structure is the universal conservation principle: the Born norm,
unitarity, and the Noether currents are manifestations of the one conservation rooted in the
identity of Difference.
\end{theorem}

\begin{proof}
The conservation-principle audit \cite{c09:QG267} establishes that the norm (Born rule),
unitarity, trace, count and Bianchi conservation, and the Noether currents are
manifestations of one deeper principle; count conservation is the definitional identity of
the primitive (Chapter~\ref{c03:chap:actualization},
Theorem~\ref{c03:thm:count-conservation}). Hence the quantum conservation structure is the
universal conservation principle, rooted in the identity of Difference.
\end{proof}

Quantum mechanics is therefore consistent with the canonical architecture: states,
amplitudes, phases, interference, measurement, and dynamics are all derived reads of the
D96 spectrum through the Difference duality and the actualization process
\cite{c09:QG286,c09:QG223}.

\section{Summary}
\label{c09:sec:qm-summary}

This chapter has derived the quantum structure of the theory. Quantum mechanics is
emergent in The Actualization Theory: states are spectral reads of the D96 spectrum
(Theorem~\ref{c09:thm:states-spectral}); the amplitude is the counting-measure share,
$|\psi|^2 = \rho$, with the Born rule exact by construction
(Theorem~\ref{c09:thm:amplitude-derived}); the phase is the circulation phase $2\pi k/N$ of
the actualization cycle, and the state is complex (Theorem~\ref{c09:thm:phase-derived});
interference is the phase-difference pattern (Theorem~\ref{c09:thm:interference});
measurement is the actualization event, with no separate collapse postulate
(Theorem~\ref{c09:thm:measurement-actualization}); and the dynamics is the actualization flow
with universal conservation (Lemma~\ref{c09:lem:qm-dynamics},
Theorem~\ref{c09:thm:qm-conservation}). The duality structures QM
(Section~\ref{c09:sec:duality}): the scalar face carries the amplitude, the tensor face the
gravitational content. The quantum pillar is derived, not postulated. The following
chapters derive gravity and spacetime, the standard model, and cosmology from the same
spectral foundation.

\begin{thebibliography}{99}
\bibitem{c09:QG216} AT-QG Phase 216, \emph{Quantum Amplitude Origin}, Docs/Research.
\bibitem{c09:QG218} AT-QG Phase 218, \emph{Hilbert Space Origin}, Docs/Research.
\bibitem{c09:QG220} AT-QG Phase 220, \emph{Quantum Phase Origin}, Docs/Research.
\bibitem{c09:QG223} AT-QG Phase 223, \emph{Final Quantum Gravity Audit} (Complete QG), Docs/Research.
\bibitem{c09:QG243} AT-QG Phase 243, \emph{Gauge Dynamics Origin}, Docs/Research.
\bibitem{c09:QG244} AT-QG Phase 244, \emph{Lagrangian Origin}, Docs/Research.
\bibitem{c09:QG267} AT-QG Phase 267, \emph{Conservation Principle Audit}, Docs/Research.
\bibitem{c09:QG286} AT-QG Phase 286, \emph{Difference Duality Audit}, Docs/Research.
\bibitem{c09:QG301} AT-QG Phase 301, \emph{Duality Prediction Audit}, Docs/Research.
\bibitem{c09:MONO007} AT-MONO007, \emph{Referee-Readiness Resolution}, Docs/Zenodo/Monograph.
\end{thebibliography}
 after the preamble
%         that defines the theorem environments (or include via the master file).
% ======================================================================

% ---- Theorem environments (define in the master preamble if not present) ----
% \newtheorem{definition}{Definition}[chapter]
% \newtheorem{lemma}[definition]{Lemma}
% \newtheorem{theorem}[definition]{Theorem}
% \newtheorem{corollary}[definition]{Corollary}
% \newtheorem{remark}[definition]{Remark}

\chapter{Quantum Mechanics}
\label{c09:chap:qm}

\section{Introduction}
\label{c09:sec:qm-introduction}

The preceding chapters established the primitives $\{\text{Difference},\eta\}$,
Actualization, the D96 spectrum, and the operator basis through which it is read. This
chapter derives the quantum structure of the theory from that foundation: quantum
mechanics is \emph{emergent} --- a read of the D96 spectrum through the Difference duality
--- not a set of independent postulates.

The central results: the quantum state arises from the spectrum structure; the amplitude
magnitude satisfies $|\psi|^2 = \rho$ --- the counting-measure share is the amplitude
magnitude squared \cite{c09:QG216}; the state carries a phase structure derived from the
actualization cycle \cite{c09:QG220}; the complex structure is forced by the two independent
degrees of freedom of the state \cite{c09:QG218}; interference is phase-dependent; and
\emph{measurement is actualization} --- the collapse is a Born-weighted projection identified
with the actualization event \cite{c09:QG223}. The Born rule is derived, not asserted:
$\sum|\psi|^2 = 1$ is exact by construction, as the normalization of the actualization share
\cite{c09:QG216}.

These results are the canonical end-state of the quantum-origin program
\cite{c09:QG216,c09:QG218,c09:QG220,c09:QG223,c09:QG243,c09:QG244}, consistent with the Difference duality
\cite{c09:QG286,c09:QG301}.

\section{Quantum States from Spectrum}
\label{c09:sec:states-from-spectrum}

\begin{theorem}[Quantum states are spectral reads]
\label{c09:thm:states-spectral}
A \emph{quantum state} is a spectral read: a state of the system assigned by the D96
spectrum structure, carrying the degrees of freedom of the underlying actualization
process. The state structure is determined by the spectrum's mode and multiplicity
structure, not by an independent quantum postulate. Quantum mechanics is therefore
emergent in the theory: its states, amplitudes, and structure are derived reads of the
spectrum, not primitive postulates.
\end{theorem}

\begin{proof}
Quantum states are assigned over the spectrum of Chapter~\ref{c06:chap:d96}: the mode and
multiplicity structures determine the state space, and the Hilbert-origin result
\cite{c09:QG218} derives the state structure from the spectral degrees of freedom. Hence the
quantum states are spectral reads --- emergent assignments over the D96 spectrum, not
independent inputs. The amplitude and phase structure are derived in
Sections~\ref{c09:sec:amplitudes} and \ref{c09:sec:phase}, and the measurement rule in
Section~\ref{c09:sec:measurement}; every component of quantum mechanics is thus emergent from
the spectrum, consistent with the canonical architecture \cite{c09:QG223}.
\end{proof}

\section{Difference Duality (\texorpdfstring{$\rho,\psi$}{rho, psi})}
\label{c09:sec:duality}

The Difference duality structures quantum mechanics: the two faces of Difference --- the
scalar $\rho$ and the tensor $\psi$ --- carry, respectively, the amplitude and the
gravitational/tensor content of the state. The duality \cite{c09:QG286} establishes that
$\rho$ and $\psi$ are the trace and traceless faces of the one Difference
(Chapter~\ref{c01:chap:difference}, Theorem~\ref{c01:thm:duality}); the scalar face is the
counting measure whose share is the amplitude magnitude squared
(Section~\ref{c09:sec:amplitudes}), and the tensor face is the Weyl content read against the
reference $\eta$ (Chapter~\ref{c02:chap:eta}, Theorem~\ref{c02:thm:tensor-read}). The duality
is invariant --- it changes meaning, not numbers \cite{c09:QG301}: the quantum structure is
consistent with it, and no number changes under the duality reinterpretation.

\section{Amplitudes and Probabilities}
\label{c09:sec:amplitudes}

\begin{theorem}[The amplitude magnitude is derived]
\label{c09:thm:amplitude-derived}
The \emph{amplitude magnitude} of a state at generation $k$ is the square root of its
counting-measure share: $|\psi_k| = \sqrt{\rho_k}$, where
$\rho_k = \mu^k/S$ is the actualization share, $\mu$ the branching ratio of the
Galton--Watson process, and $S = \sum_{j<K_{\mathrm{gen}}}\mu^j$, with $K_{\mathrm{gen}}$ the
total number of generations. The amplitude magnitude is derived from
the actualization process:
\begin{equation}
\label{c09:eq:amplitude}
|\psi_k|^2 \;=\; \rho_k \;=\; \frac{\mu^k}{S},
\end{equation}
the counting-measure share of the state. The Born rule $\sum_k|\psi_k|^2 = 1$ holds exactly
by construction, as the normalization of the actualization share: probabilities are the
normalized actualization shares.
\end{theorem}

\begin{proof}
The amplitude-origin result \cite{c09:QG216} derives the amplitude magnitude from the Q-event
process: the path multiplicity to generation $k$ is $\mu^k$, and the normalized share
$\rho_k = \mu^k/S$ is the amplitude magnitude squared. The normalization
$\sum_k|\psi_k|^2 = 1$ is exact by construction for any branching ratio $\mu$; the Born
rule follows from the normalization of the share, so probability is the actualization share
--- the measure, not a separately imposed rule.
\end{proof}

Here $\mu^k$ is the combinatorial path multiplicity of the Galton--Watson tree, namely the
number of root-to-generation-$k$ paths, entering only through the normalized weight
$\rho_k = \mu^k/S$. It is not identified with the conserved Q-event count of
Theorem~\ref{c03:thm:count-conservation}.

\section{Phase Structure}
\label{c09:sec:phase}

\begin{theorem}[The phase is derived from actualization]
\label{c09:thm:phase-derived}
The \emph{quantum phase} of a state at depth $k$ is the circulation phase of the
actualization cycle:
\begin{equation}
\label{c09:eq:phase}
\theta_k \;=\; \frac{2\pi k}{N},
\end{equation}
where $N=96$ is the network periodicity and $k$ the branch depth (the actualization tick).
The phase is fixed by causal ordering and the network periodicity of the $N=96$ attractor.
A quantum state carries exactly two independent real degrees of freedom --- the
magnitude $|\psi| = \sqrt{\rho}$ (the node property) and the phase $\theta$ (the link
property) --- and therefore must be represented as a complex state.
\end{theorem}

\begin{proof}
The phase-origin result \cite{c09:QG220} derives the phase from the Q-event process: causal
ordering fixes the position $k$ (the branch depth, the actualization tick), and the network
periodicity of the circulant ring $C_N$, with $N=96$, fixes the phase quantum
$\Delta\theta = 2\pi/N$ by cycle closure --- $N$ ticks advance $2\pi$, giving uniform
circulation. Hence $\theta_k = 2\pi k/N$ is the circulation phase of the actualization
cycle. The Hilbert-origin result \cite{c09:QG218} establishes the two independent real
degrees of freedom: the magnitude $|\psi| = \sqrt{\rho}$ (from the branching counting
measure, a node property) and the phase $\theta$ (from the $U(1)$ link connection, a link
property). A state with magnitude and phase is a complex state; a real-only state space
would give classical addition with no interference. Hence the complex structure is derived
from the two degrees of freedom.
\end{proof}

The lattice phase $\theta_k = 2\pi k/N$ governs \emph{state evolution}. Mixing and CP
phases are \emph{separate} objects: they are derived from spectral ratios via the
chiral-circulation construction, $\delta = \arcsin(r)$ for a spectral ratio $r$, are
continuous functions of those ratios, and are \emph{not} restricted to the state-phase
lattice.

\section{Interference}
\label{c09:sec:interference}

\begin{theorem}[Interference is phase-dependent]
\label{c09:thm:interference}
Interference in the theory is phase-dependent: the two-path amplitude is
\begin{equation}
\label{c09:eq:interference}
P \;=\; \left|e^{i\theta_1} + e^{i\theta_2}\right|^2
\;=\; 2 + 2\cos(\theta_1 - \theta_2),
\end{equation}
with the phase difference producing the interference pattern. Interference is a derived
consequence of the phase structure: the phase-difference pattern of the
actualization-derived phases.
\end{theorem}

\begin{proof}
The Hilbert-origin result \cite{c09:QG218} establishes the interference formula from the
complex state structure: the two-path amplitude is the sum of the two phase-carrying
amplitudes, whose magnitude squared is $2 + 2\cos(\theta_1 - \theta_2)$. The phase
difference, derived from the actualization cycle (Theorem~\ref{c09:thm:phase-derived}),
produces the interference; a real-only state space would give $P = P_1 + P_2$ with no
interference, so the complex structure is necessary. Hence interference is a derived
consequence of the phase structure, not an independent quantum postulate.
\end{proof}

\section{Measurement as Actualization}
\label{c09:sec:measurement}

\begin{theorem}[Measurement is actualization]
\label{c09:thm:measurement-actualization}
\emph{Measurement} is the collapse of the state to a definite outcome: a Born-weighted
projection identified with the actualization event --- a Q-event is a Born-weighted
projection of the state onto a definite outcome, so no separate collapse postulate or
mechanism is required. Consistent with MONO007 \cite{c09:MONO007}, the measurement result
is presented as the accepted canonical identification: collapse equals actualization. The
theory does not add a new mechanism; it identifies the collapse with the process already
established. No stronger claim is made.
\end{theorem}

\begin{proof}
The complete-quantum-gravity audit \cite{c09:QG223} adjudicates the measurement problem: the
collapse is identified with the Q-event actualization --- a Born-weighted projection to a
definite state. Since the Born weights are the actualization shares
(Theorem~\ref{c09:thm:amplitude-derived}), the projection is weighted by the shares and
performed by the actualization event: measurement is actualization, with no separate
collapse mechanism or postulate.
\end{proof}

\section{Consequences}
\label{c09:sec:qm-consequences}

\begin{lemma}[The quantum dynamics is the actualization flow]
\label{c09:lem:qm-dynamics}
The quantum dynamics of the theory is the actualization flow: the interaction dynamics is
the generator action on the spectral modes, and the Lagrangian is the actualization-flow
action.
\end{lemma}

\begin{proof}
The gauge-dynamics origin \cite{c09:QG243} establishes that the interaction dynamics IS the
generator action on the spectral modes: the D96 gauge generators act on the modes, and an
interaction is the generator's action on the mode. The Lagrangian origin \cite{c09:QG244}
derives the Lagrangian density as the actualization-flow action of the D96 generator
fields, $L = -\frac{1}{4}F^a_{\mu\nu}F^{a\mu\nu} + i\bar{\psi}\gamma^\mu D_\mu\psi -
m\bar{\psi}\psi$. Hence the quantum dynamics is the actualization flow, not an imported
dynamics.
\end{proof}

\begin{theorem}[The conservation structure is universal]
\label{c09:thm:qm-conservation}
The quantum conservation structure is the universal conservation principle: the Born norm,
unitarity, and the Noether currents are manifestations of the one conservation rooted in the
identity of Difference.
\end{theorem}

\begin{proof}
The conservation-principle audit \cite{c09:QG267} establishes that the norm (Born rule),
unitarity, trace, count and Bianchi conservation, and the Noether currents are
manifestations of one deeper principle; count conservation is the definitional identity of
the primitive (Chapter~\ref{c03:chap:actualization},
Theorem~\ref{c03:thm:count-conservation}). Hence the quantum conservation structure is the
universal conservation principle, rooted in the identity of Difference.
\end{proof}

Quantum mechanics is therefore consistent with the canonical architecture: states,
amplitudes, phases, interference, measurement, and dynamics are all derived reads of the
D96 spectrum through the Difference duality and the actualization process
\cite{c09:QG286,c09:QG223}.

\section{Summary}
\label{c09:sec:qm-summary}

This chapter has derived the quantum structure of the theory. Quantum mechanics is
emergent in The Actualization Theory: states are spectral reads of the D96 spectrum
(Theorem~\ref{c09:thm:states-spectral}); the amplitude is the counting-measure share,
$|\psi|^2 = \rho$, with the Born rule exact by construction
(Theorem~\ref{c09:thm:amplitude-derived}); the phase is the circulation phase $2\pi k/N$ of
the actualization cycle, and the state is complex (Theorem~\ref{c09:thm:phase-derived});
interference is the phase-difference pattern (Theorem~\ref{c09:thm:interference});
measurement is the actualization event, with no separate collapse postulate
(Theorem~\ref{c09:thm:measurement-actualization}); and the dynamics is the actualization flow
with universal conservation (Lemma~\ref{c09:lem:qm-dynamics},
Theorem~\ref{c09:thm:qm-conservation}). The duality structures QM
(Section~\ref{c09:sec:duality}): the scalar face carries the amplitude, the tensor face the
gravitational content. The quantum pillar is derived, not postulated. The following
chapters derive gravity and spacetime, the standard model, and cosmology from the same
spectral foundation.

\begin{thebibliography}{99}
\bibitem{c09:QG216} AT-QG Phase 216, \emph{Quantum Amplitude Origin}, Docs/Research.
\bibitem{c09:QG218} AT-QG Phase 218, \emph{Hilbert Space Origin}, Docs/Research.
\bibitem{c09:QG220} AT-QG Phase 220, \emph{Quantum Phase Origin}, Docs/Research.
\bibitem{c09:QG223} AT-QG Phase 223, \emph{Final Quantum Gravity Audit} (Complete QG), Docs/Research.
\bibitem{c09:QG243} AT-QG Phase 243, \emph{Gauge Dynamics Origin}, Docs/Research.
\bibitem{c09:QG244} AT-QG Phase 244, \emph{Lagrangian Origin}, Docs/Research.
\bibitem{c09:QG267} AT-QG Phase 267, \emph{Conservation Principle Audit}, Docs/Research.
\bibitem{c09:QG286} AT-QG Phase 286, \emph{Difference Duality Audit}, Docs/Research.
\bibitem{c09:QG301} AT-QG Phase 301, \emph{Duality Prediction Audit}, Docs/Research.
\bibitem{c09:MONO007} AT-MONO007, \emph{Referee-Readiness Resolution}, Docs/Zenodo/Monograph.
\end{thebibliography}
 after the preamble
%         that defines the theorem environments (or include via the master file).
% ======================================================================

% ---- Theorem environments (define in the master preamble if not present) ----
% \newtheorem{definition}{Definition}[chapter]
% \newtheorem{lemma}[definition]{Lemma}
% \newtheorem{theorem}[definition]{Theorem}
% \newtheorem{corollary}[definition]{Corollary}
% \newtheorem{remark}[definition]{Remark}

\chapter{Quantum Mechanics}
\label{c09:chap:qm}

\section{Introduction}
\label{c09:sec:qm-introduction}

The preceding chapters established the primitives $\{\text{Difference},\eta\}$,
Actualization, the D96 spectrum, and the operator basis through which it is read. This
chapter derives the quantum structure of the theory from that foundation: quantum
mechanics is \emph{emergent} --- a read of the D96 spectrum through the Difference duality
--- not a set of independent postulates.

The central results: the quantum state arises from the spectrum structure; the amplitude
magnitude satisfies $|\psi|^2 = \rho$ --- the counting-measure share is the amplitude
magnitude squared \cite{c09:QG216}; the state carries a phase structure derived from the
actualization cycle \cite{c09:QG220}; the complex structure is forced by the two independent
degrees of freedom of the state \cite{c09:QG218}; interference is phase-dependent; and
\emph{measurement is actualization} --- the collapse is a Born-weighted projection identified
with the actualization event \cite{c09:QG223}. The Born rule is derived, not asserted:
$\sum|\psi|^2 = 1$ is exact by construction, as the normalization of the actualization share
\cite{c09:QG216}.

These results are the canonical end-state of the quantum-origin program
\cite{c09:QG216,c09:QG218,c09:QG220,c09:QG223,c09:QG243,c09:QG244}, consistent with the Difference duality
\cite{c09:QG286,c09:QG301}.

\section{Quantum States from Spectrum}
\label{c09:sec:states-from-spectrum}

\begin{theorem}[Quantum states are spectral reads]
\label{c09:thm:states-spectral}
A \emph{quantum state} is a spectral read: a state of the system assigned by the D96
spectrum structure, carrying the degrees of freedom of the underlying actualization
process. The state structure is determined by the spectrum's mode and multiplicity
structure, not by an independent quantum postulate. Quantum mechanics is therefore
emergent in the theory: its states, amplitudes, and structure are derived reads of the
spectrum, not primitive postulates.
\end{theorem}

\begin{proof}
Quantum states are assigned over the spectrum of Chapter~\ref{c06:chap:d96}: the mode and
multiplicity structures determine the state space, and the Hilbert-origin result
\cite{c09:QG218} derives the state structure from the spectral degrees of freedom. Hence the
quantum states are spectral reads --- emergent assignments over the D96 spectrum, not
independent inputs. The amplitude and phase structure are derived in
Sections~\ref{c09:sec:amplitudes} and \ref{c09:sec:phase}, and the measurement rule in
Section~\ref{c09:sec:measurement}; every component of quantum mechanics is thus emergent from
the spectrum, consistent with the canonical architecture \cite{c09:QG223}.
\end{proof}

\section{Difference Duality (\texorpdfstring{$\rho,\psi$}{rho, psi})}
\label{c09:sec:duality}

The Difference duality structures quantum mechanics: the two faces of Difference --- the
scalar $\rho$ and the tensor $\psi$ --- carry, respectively, the amplitude and the
gravitational/tensor content of the state. The duality \cite{c09:QG286} establishes that
$\rho$ and $\psi$ are the trace and traceless faces of the one Difference
(Chapter~\ref{c01:chap:difference}, Theorem~\ref{c01:thm:duality}); the scalar face is the
counting measure whose share is the amplitude magnitude squared
(Section~\ref{c09:sec:amplitudes}), and the tensor face is the Weyl content read against the
reference $\eta$ (Chapter~\ref{c02:chap:eta}, Theorem~\ref{c02:thm:tensor-read}). The duality
is invariant --- it changes meaning, not numbers \cite{c09:QG301}: the quantum structure is
consistent with it, and no number changes under the duality reinterpretation.

\section{Amplitudes and Probabilities}
\label{c09:sec:amplitudes}

\begin{theorem}[The amplitude magnitude is derived]
\label{c09:thm:amplitude-derived}
The \emph{amplitude magnitude} of a state at generation $k$ is the square root of its
counting-measure share: $|\psi_k| = \sqrt{\rho_k}$, where
$\rho_k = \mu^k/S$ is the actualization share, $\mu$ the branching ratio of the
Galton--Watson process, and $S = \sum_{j<K_{\mathrm{gen}}}\mu^j$, with $K_{\mathrm{gen}}$ the
total number of generations. The amplitude magnitude is derived from
the actualization process:
\begin{equation}
\label{c09:eq:amplitude}
|\psi_k|^2 \;=\; \rho_k \;=\; \frac{\mu^k}{S},
\end{equation}
the counting-measure share of the state. The Born rule $\sum_k|\psi_k|^2 = 1$ holds exactly
by construction, as the normalization of the actualization share: probabilities are the
normalized actualization shares.
\end{theorem}

\begin{proof}
The amplitude-origin result \cite{c09:QG216} derives the amplitude magnitude from the Q-event
process: the path multiplicity to generation $k$ is $\mu^k$, and the normalized share
$\rho_k = \mu^k/S$ is the amplitude magnitude squared. The normalization
$\sum_k|\psi_k|^2 = 1$ is exact by construction for any branching ratio $\mu$; the Born
rule follows from the normalization of the share, so probability is the actualization share
--- the measure, not a separately imposed rule.
\end{proof}

Here $\mu^k$ is the combinatorial path multiplicity of the Galton--Watson tree, namely the
number of root-to-generation-$k$ paths, entering only through the normalized weight
$\rho_k = \mu^k/S$. It is not identified with the conserved Q-event count of
Theorem~\ref{c03:thm:count-conservation}.

\section{Phase Structure}
\label{c09:sec:phase}

\begin{theorem}[The phase is derived from actualization]
\label{c09:thm:phase-derived}
The \emph{quantum phase} of a state at depth $k$ is the circulation phase of the
actualization cycle:
\begin{equation}
\label{c09:eq:phase}
\theta_k \;=\; \frac{2\pi k}{N},
\end{equation}
where $N=96$ is the network periodicity and $k$ the branch depth (the actualization tick).
The phase is fixed by causal ordering and the network periodicity of the $N=96$ attractor.
A quantum state carries exactly two independent real degrees of freedom --- the
magnitude $|\psi| = \sqrt{\rho}$ (the node property) and the phase $\theta$ (the link
property) --- and therefore must be represented as a complex state.
\end{theorem}

\begin{proof}
The phase-origin result \cite{c09:QG220} derives the phase from the Q-event process: causal
ordering fixes the position $k$ (the branch depth, the actualization tick), and the network
periodicity of the circulant ring $C_N$, with $N=96$, fixes the phase quantum
$\Delta\theta = 2\pi/N$ by cycle closure --- $N$ ticks advance $2\pi$, giving uniform
circulation. Hence $\theta_k = 2\pi k/N$ is the circulation phase of the actualization
cycle. The Hilbert-origin result \cite{c09:QG218} establishes the two independent real
degrees of freedom: the magnitude $|\psi| = \sqrt{\rho}$ (from the branching counting
measure, a node property) and the phase $\theta$ (from the $U(1)$ link connection, a link
property). A state with magnitude and phase is a complex state; a real-only state space
would give classical addition with no interference. Hence the complex structure is derived
from the two degrees of freedom.
\end{proof}

The lattice phase $\theta_k = 2\pi k/N$ governs \emph{state evolution}. Mixing and CP
phases are \emph{separate} objects: they are derived from spectral ratios via the
chiral-circulation construction, $\delta = \arcsin(r)$ for a spectral ratio $r$, are
continuous functions of those ratios, and are \emph{not} restricted to the state-phase
lattice.

\section{Interference}
\label{c09:sec:interference}

\begin{theorem}[Interference is phase-dependent]
\label{c09:thm:interference}
Interference in the theory is phase-dependent: the two-path amplitude is
\begin{equation}
\label{c09:eq:interference}
P \;=\; \left|e^{i\theta_1} + e^{i\theta_2}\right|^2
\;=\; 2 + 2\cos(\theta_1 - \theta_2),
\end{equation}
with the phase difference producing the interference pattern. Interference is a derived
consequence of the phase structure: the phase-difference pattern of the
actualization-derived phases.
\end{theorem}

\begin{proof}
The Hilbert-origin result \cite{c09:QG218} establishes the interference formula from the
complex state structure: the two-path amplitude is the sum of the two phase-carrying
amplitudes, whose magnitude squared is $2 + 2\cos(\theta_1 - \theta_2)$. The phase
difference, derived from the actualization cycle (Theorem~\ref{c09:thm:phase-derived}),
produces the interference; a real-only state space would give $P = P_1 + P_2$ with no
interference, so the complex structure is necessary. Hence interference is a derived
consequence of the phase structure, not an independent quantum postulate.
\end{proof}

\section{Measurement as Actualization}
\label{c09:sec:measurement}

\begin{theorem}[Measurement is actualization]
\label{c09:thm:measurement-actualization}
\emph{Measurement} is the collapse of the state to a definite outcome: a Born-weighted
projection identified with the actualization event --- a Q-event is a Born-weighted
projection of the state onto a definite outcome, so no separate collapse postulate or
mechanism is required. Consistent with MONO007 \cite{c09:MONO007}, the measurement result
is presented as the accepted canonical identification: collapse equals actualization. The
theory does not add a new mechanism; it identifies the collapse with the process already
established. No stronger claim is made.
\end{theorem}

\begin{proof}
The complete-quantum-gravity audit \cite{c09:QG223} adjudicates the measurement problem: the
collapse is identified with the Q-event actualization --- a Born-weighted projection to a
definite state. Since the Born weights are the actualization shares
(Theorem~\ref{c09:thm:amplitude-derived}), the projection is weighted by the shares and
performed by the actualization event: measurement is actualization, with no separate
collapse mechanism or postulate.
\end{proof}

\section{Consequences}
\label{c09:sec:qm-consequences}

\begin{lemma}[The quantum dynamics is the actualization flow]
\label{c09:lem:qm-dynamics}
The quantum dynamics of the theory is the actualization flow: the interaction dynamics is
the generator action on the spectral modes, and the Lagrangian is the actualization-flow
action.
\end{lemma}

\begin{proof}
The gauge-dynamics origin \cite{c09:QG243} establishes that the interaction dynamics IS the
generator action on the spectral modes: the D96 gauge generators act on the modes, and an
interaction is the generator's action on the mode. The Lagrangian origin \cite{c09:QG244}
derives the Lagrangian density as the actualization-flow action of the D96 generator
fields, $L = -\frac{1}{4}F^a_{\mu\nu}F^{a\mu\nu} + i\bar{\psi}\gamma^\mu D_\mu\psi -
m\bar{\psi}\psi$. Hence the quantum dynamics is the actualization flow, not an imported
dynamics.
\end{proof}

\begin{theorem}[The conservation structure is universal]
\label{c09:thm:qm-conservation}
The quantum conservation structure is the universal conservation principle: the Born norm,
unitarity, and the Noether currents are manifestations of the one conservation rooted in the
identity of Difference.
\end{theorem}

\begin{proof}
The conservation-principle audit \cite{c09:QG267} establishes that the norm (Born rule),
unitarity, trace, count and Bianchi conservation, and the Noether currents are
manifestations of one deeper principle; count conservation is the definitional identity of
the primitive (Chapter~\ref{c03:chap:actualization},
Theorem~\ref{c03:thm:count-conservation}). Hence the quantum conservation structure is the
universal conservation principle, rooted in the identity of Difference.
\end{proof}

Quantum mechanics is therefore consistent with the canonical architecture: states,
amplitudes, phases, interference, measurement, and dynamics are all derived reads of the
D96 spectrum through the Difference duality and the actualization process
\cite{c09:QG286,c09:QG223}.

\section{Summary}
\label{c09:sec:qm-summary}

This chapter has derived the quantum structure of the theory. Quantum mechanics is
emergent in The Actualization Theory: states are spectral reads of the D96 spectrum
(Theorem~\ref{c09:thm:states-spectral}); the amplitude is the counting-measure share,
$|\psi|^2 = \rho$, with the Born rule exact by construction
(Theorem~\ref{c09:thm:amplitude-derived}); the phase is the circulation phase $2\pi k/N$ of
the actualization cycle, and the state is complex (Theorem~\ref{c09:thm:phase-derived});
interference is the phase-difference pattern (Theorem~\ref{c09:thm:interference});
measurement is the actualization event, with no separate collapse postulate
(Theorem~\ref{c09:thm:measurement-actualization}); and the dynamics is the actualization flow
with universal conservation (Lemma~\ref{c09:lem:qm-dynamics},
Theorem~\ref{c09:thm:qm-conservation}). The duality structures QM
(Section~\ref{c09:sec:duality}): the scalar face carries the amplitude, the tensor face the
gravitational content. The quantum pillar is derived, not postulated. The following
chapters derive gravity and spacetime, the standard model, and cosmology from the same
spectral foundation.

\begin{thebibliography}{99}
\bibitem{c09:QG216} AT-QG Phase 216, \emph{Quantum Amplitude Origin}, Docs/Research.
\bibitem{c09:QG218} AT-QG Phase 218, \emph{Hilbert Space Origin}, Docs/Research.
\bibitem{c09:QG220} AT-QG Phase 220, \emph{Quantum Phase Origin}, Docs/Research.
\bibitem{c09:QG223} AT-QG Phase 223, \emph{Final Quantum Gravity Audit} (Complete QG), Docs/Research.
\bibitem{c09:QG243} AT-QG Phase 243, \emph{Gauge Dynamics Origin}, Docs/Research.
\bibitem{c09:QG244} AT-QG Phase 244, \emph{Lagrangian Origin}, Docs/Research.
\bibitem{c09:QG267} AT-QG Phase 267, \emph{Conservation Principle Audit}, Docs/Research.
\bibitem{c09:QG286} AT-QG Phase 286, \emph{Difference Duality Audit}, Docs/Research.
\bibitem{c09:QG301} AT-QG Phase 301, \emph{Duality Prediction Audit}, Docs/Research.
\bibitem{c09:MONO007} AT-MONO007, \emph{Referee-Readiness Resolution}, Docs/Zenodo/Monograph.
\end{thebibliography}

% ======================================================================
%  TQM-MONO109 — Chapter 10: Gravity and Spacetime
%  Canonical Monograph (v2.0) · The Actualization Theory:
%  A Reconstruction of Physics from Difference, Actualization and Spectrum
%
%  Status: COMPLETE — PASS-READY (MONO007 referee-ready architecture)
%  Inputs: Chapter01_Difference.tex ... Chapter09_QuantumMechanics.tex
%  Method: assembly only — canonical results only, no new physics, no speculation
%  Citations: QG181 (Newton constant), QG182 (G bridge), QG183 (Planck robustness),
%             QG184 (mass-radius), QG186 (frame dragging), QG187 (GPS),
%             QG222 (native dynamics), QG285 (psi reinterpretation),
%             QG286 (difference duality), QG292 (minimal foundation)
%
%  Usage: % ======================================================================
%  TQM-MONO109 — Chapter 10: Gravity and Spacetime
%  Canonical Monograph (v2.0) · The Actualization Theory:
%  A Reconstruction of Physics from Difference, Actualization and Spectrum
%
%  Status: COMPLETE — PASS-READY (MONO007 referee-ready architecture)
%  Inputs: Chapter01_Difference.tex ... Chapter09_QuantumMechanics.tex
%  Method: assembly only — canonical results only, no new physics, no speculation
%  Citations: QG181 (Newton constant), QG182 (G bridge), QG183 (Planck robustness),
%             QG184 (mass-radius), QG186 (frame dragging), QG187 (GPS),
%             QG222 (native dynamics), QG285 (psi reinterpretation),
%             QG286 (difference duality), QG292 (minimal foundation)
%
%  Usage: % ======================================================================
%  TQM-MONO109 — Chapter 10: Gravity and Spacetime
%  Canonical Monograph (v2.0) · The Actualization Theory:
%  A Reconstruction of Physics from Difference, Actualization and Spectrum
%
%  Status: COMPLETE — PASS-READY (MONO007 referee-ready architecture)
%  Inputs: Chapter01_Difference.tex ... Chapter09_QuantumMechanics.tex
%  Method: assembly only — canonical results only, no new physics, no speculation
%  Citations: QG181 (Newton constant), QG182 (G bridge), QG183 (Planck robustness),
%             QG184 (mass-radius), QG186 (frame dragging), QG187 (GPS),
%             QG222 (native dynamics), QG285 (psi reinterpretation),
%             QG286 (difference duality), QG292 (minimal foundation)
%
%  Usage: \input{Chapter10_GravityAndSpacetime.tex} after the preamble
%         that defines the theorem environments (or include via the master file).
% ======================================================================

% ---- Theorem environments (define in the master preamble if not present) ----
% \newtheorem{definition}{Definition}[chapter]
% \newtheorem{lemma}[definition]{Lemma}
% \newtheorem{theorem}[definition]{Theorem}
% \newtheorem{corollary}[definition]{Corollary}
% \newtheorem{remark}[definition]{Remark}

\chapter{Gravity and Spacetime}
\label{chap:gravity}

\section{Introduction}
\label{sec:gravity-introduction}

The preceding chapters established the foundation of The Actualization Theory: the
primitives $\{\text{Difference},\eta\}$, the derived process of Actualization, the
inevitable D96 spectrum, the operator basis, and the emergent quantum structure. This
chapter derives \emph{gravity and spacetime} from the same foundation. We establish that
gravity is emergent --- a read of the counting measure and its tensor face --- and that
spacetime is emergent geometry \cite{QG181,QG184,QG222}.

The central results are: the gravitational coupling arises from the D96 spectral content
\cite{QG181,QG182,QG183}; the tensor sector is carried by the tensor face of Difference,
$\psi$, read against the tensor reference $\eta$ \cite{QG285,QG286}; the metric is
constructed from the counting measure and the reference \cite{QG292}; spacetime is emergent
geometry from the network and its dynamics \cite{QG222}; the black-hole relations and the
frame-dragging sector follow \cite{QG184,QG186}; and the GPS/gravitational-redshift
observables are the counting-measure redshift law \cite{QG187}.

Throughout, gravity is presented as \emph{emergent}, $\eta$ as the \emph{tensor
reference}, and $\psi$ as the \emph{tensor face of Difference}. The Bekenstein $1/4$
coefficient is \emph{excluded from the derivation claims} of this chapter and explicitly
referenced to the Boundary Layer, per the referee-safe architecture of MONO007
\cite{MONO007}.

We use only accepted canonical results and introduce no new physics and no speculation.

\section{Gravity from Difference}
\label{sec:gravity-from-difference}

\begin{definition}[Gravitational coupling]
\label{def:newton-constant}
The \emph{gravitational coupling} of the theory is the Newton constant $G$, derived from
the Planck scale
\begin{equation}
\label{eq:mpl}
M_{\mathrm{Pl}} \;=\; v\cdot(\Sigma m\cdot\#g\cdot\mathrm{occ}_2)^3,
\end{equation}
where $v$ is the weak scale, $\Sigma m = 95$ the total mode count, $\#g = 44$ the group
count, and $\mathrm{occ}_2$ the dense-band occupancy.
\end{definition}

\begin{theorem}[Gravity is derived from the D96 spectral content]
\label{thm:gravity-derived}
The gravitational coupling is derived from the D96 spectral content: the Planck scale is the
weak scale times the cube of the occupation-weighted spectral content, giving
$M_{\mathrm{Pl}} = 1.22335\times10^{19}\,\mathrm{GeV}$ and
$G = 6.6476\times10^{-11}\,\mathrm{m}^3\,\mathrm{kg}^{-1}\,\mathrm{s}^{-2}$.
\end{theorem}

\begin{proof}
The Newton-constant origin \cite{QG181} establishes the construction:
$M_{\mathrm{Pl}} = v\cdot(\Sigma m\cdot\#g\cdot\mathrm{occ}_2)^3 = 1.22335\times10^{19}$
GeV, matching the Planck scale within $0.2\%$, and $G = 1/M_{\mathrm{Pl}}^2$ within
$0.4\%$. The robustness audit \cite{QG183} confirms the physical exponent is cubic, and the
bridge audit \cite{QG182} shows the two G constructions are the same physical content.
Hence gravity is derived from the D96 spectral content --- the weak scale times the cube of
the occupation-weighted spectral structure.
\end{proof}

\begin{corollary}[Gravity is emergent]
\label{cor:gravity-emergent}
Gravity is emergent in the theory: the gravitational coupling is a derived function of the
D96 spectral content, not an independent input.
\end{corollary}

\begin{proof}
By Theorem~\ref{thm:gravity-derived}, $G$ is a function of the spectral constants
($\Sigma m$, $\#g$, $\mathrm{occ}_2$) and the weak scale. By Chapter~\ref{chap:d96}
(Theorem~\ref{thm:d96-derived}), the spectral constants are derived from the spectrum, which
is derived from the attractor. Hence gravity is twice-emergent: a read of the derived
spectral content.
\end{proof}

\section{The Tensor Sector}
\label{sec:tensor-sector}

\begin{theorem}[The tensor sector is the tensor face of Difference]
\label{thm:tensor-face}
The tensor sector of gravity is carried by the tensor face of Difference, $\psi$: the
Weyl content read against the tensor reference $\eta$.
\end{theorem}

\begin{proof}
The Difference duality \cite{QG286} establishes that $\rho$ and $\psi$ are the trace and
traceless faces of the one Difference (Chapter~\ref{chap:difference},
Theorem~\ref{thm:duality}). The psi-reinterpretation audit \cite{QG285} identifies $\psi$
as the Weyl content --- the difference from conformal flatness --- read against the
reference $\eta$ (Chapter~\ref{chap:eta}, Theorem~\ref{thm:tensor-read}). Hence the tensor
sector of gravity is carried by $\psi$, the tensor face of Difference, against the tensor
reference $\eta$.
\end{proof}

\begin{corollary}[The tensor sector is not a new primitive]
\label{cor:tensor-not-primitive}
The tensor sector introduces no new primitive: it is the tensor face of the existing
primitive Difference, read against the existing reference $\eta$.
\end{corollary}

\begin{proof}
By Theorem~\ref{thm:tensor-face}, the tensor content is $\psi$, a face of Difference, and
its reference is $\eta$. Both are members of the primitive pair established in
Chapter~\ref{chap:eta} (Theorem~\ref{thm:minimal-foundation}). Hence the tensor sector uses
no new primitive; it is the tensor face of the existing foundation.
\end{proof}

\begin{lemma}[The tensor observables require \texorpdfstring{$\eta$}{eta}]
\label{lem:tensor-requires-eta}
The tensor (gravitational) observables are defined against the tensor reference $\eta$:
the Weyl content, the metric completion, and the frame-dragging sector all presuppose the
reference.
\end{lemma}

\begin{proof}
The minimal-foundation result \cite{QG292} establishes that removing $\eta$ leaves the
scalar sector intact but destroys the tensor sub-sector. By Theorem~\ref{thm:tensor-face},
the tensor observables are reads of $\psi$ against $\eta$; hence they require $\eta$. This
is the canonical asymmetry: the scalar face needs no reference, the tensor face does.
\end{proof}

\section{The Metric Construction}
\label{sec:metric-construction}

\begin{definition}[The metric]
\label{def:metric}
The \emph{metric} of the theory is the conformal rescaling of the tensor reference:
\begin{equation}
\label{eq:metric}
g \;=\; \rho^{2/d}\,\eta,
\end{equation}
where $\rho$ is the scalar counting measure, $d$ the spatial dimension, and $\eta$ the
tensor reference.
\end{definition}

\begin{theorem}[The metric is constructed from the primitive pair]
\label{thm:metric-construction}
The metric is constructed from the primitive pair $\{\text{Difference},\eta\}$: the
conformal factor carries the scalar face of Difference ($\rho$), and the reference carries
the tensor structure ($\eta$).
\end{theorem}

\begin{proof}
The metric ansatz \eqref{eq:metric} is the canonical form established in
Chapter~\ref{chap:eta} (Theorem~\ref{thm:conformal}). The conformal factor $\rho^{2/d}$
carries the scalar counting content --- the scalar face of Difference --- and the reference
$\eta$ provides the tensor structure. Hence the metric is constructed from the primitive
pair, with no additional input. This construction is consistent with the minimal foundation
\cite{QG292}.
\end{proof}

\begin{corollary}[The metric is emergent]
\label{cor:metric-emergent}
The metric is an emergent construction: the conformal factor is derived from the counting
measure, and the reference is the primitive $\eta$; no independent metric postulate is
required.
\end{corollary}

\begin{proof}
By Theorem~\ref{thm:metric-construction}, the metric is built from $\rho$ and $\eta$. The
conformal factor is derived from the counting measure (the actualization share), and $\eta$
is the established primitive. Hence the metric is an emergent construction, not an imported
postulate.
\end{proof}

\section{Spacetime as Emergent Geometry}
\label{sec:emergent-geometry}

\begin{theorem}[Spacetime is emergent geometry]
\label{thm:spacetime-emergent}
Spacetime is emergent geometry in the theory: it is the geometry of the counting measure
and its dynamics, not a pre-existing arena.
\end{theorem}

\begin{proof}
The native-dynamics result \cite{QG222} establishes that the gravitational dynamics IS the
Q-event actualization flow: the branching process gives the counting measure $\rho_k =
\mu^k/S$ with count conservation by construction, and the branching continuity
$\rho_{k+1} = \mu\rho_k$ provides the native continuity/Noether statement. The metric is
constructed from $\rho$ and $\eta$ (Theorem~\ref{thm:metric-construction}). Hence spacetime
is the geometry of the counting measure --- emergent from the actualization flow, not a
pre-existing background.
\end{proof}

\begin{corollary}[No imported dynamics]
\label{cor:no-imported-dynamics}
The gravitational dynamics is native: it is the actualization flow, with no imported
Einstein or BDG dynamics.
\end{corollary}

\begin{proof}
By Theorem~\ref{thm:spacetime-emergent}, the dynamics is the Q-event actualization flow
\cite{QG222}. Hence no imported gravitational dynamics is required; the dynamics is derived
from the process itself.
\end{proof}

\section{Black Hole Relations}
\label{sec:black-hole}

\begin{theorem}[The mass-radius relation]
\label{thm:mass-radius}
The mass-radius relation $M \propto R$ follows from the counting structure: the per-octave
deficit profile gives $GM_{\mathrm{eff}} \propto R$.
\end{theorem}

\begin{proof}
The mass-radius origin \cite{QG184} establishes the relation: the deficit is per-octave/log
(the flat-rotation-curve profile), giving $a \propto -1/r$ and
$GM_{\mathrm{eff}} \propto R$. With the area-entropy relation $S \propto R^{d-1}$, the
Hawking temperature $T \propto 1/R$ is restored. Hence the black-hole relations follow from
the counting structure, with $M \propto R$ derived rather than assumed.
\end{proof}

\begin{theorem}[The Bekenstein coefficient is a boundary]
\label{thm:bekenstein-boundary}
The exact Bekenstein coefficient $S = A/4$ is \emph{not} a derivation claim of this
chapter: the structure $S \propto A$, $M \propto R$, $T \propto 1/R$ is derived, but the
exact $1/4$ coefficient requires an imported quantum factor and belongs to the Boundary
Layer.
\end{theorem}

\begin{proof}
The mass-radius and entropy relations are derived \cite{QG184}. The exact Bekenstein
$S = A/4$ coefficient, however, requires the imported $2\pi$ quantum factor
$T = \kappa/(2\pi)$, which the framework cannot produce internally; the coefficient is
therefore a documented boundary, not a derivation. Consistent with the referee-safe
architecture \cite{MONO007}, this chapter excludes the $1/4$ coefficient from its
derivation claims and references it to the Boundary Layer chapter of the monograph.
\end{proof}

\begin{corollary}[Boundary referenced, not claimed]
\label{cor:bekenstein-boundary-referenced}
The Bekenstein $1/4$ coefficient is referenced to the Boundary Layer and is not presented
as a derived result of the gravity chapter.
\end{corollary}

\begin{proof}
By Theorem~\ref{thm:bekenstein-boundary}, the exact coefficient is a documented boundary
requiring the imported quantum factor. Hence the gravity chapter presents only the derived
structure ($S \propto A$, $M \propto R$, $T \propto 1/R$) and references the coefficient to
the Boundary Layer, consistent with MONO007 \cite{MONO007}.
\end{proof}

\section{Frame Dragging}
\label{sec:frame-dragging}

\begin{definition}[Frame dragging]
\label{def:frame-dragging}
\emph{Frame dragging} is the gravitomagnetic effect by which a rotating source drags the
local inertial frames: the Lense--Thirring precession
\begin{equation}
\label{eq:frame-dragging}
\Omega_{\mathrm{LT}} \;=\; \frac{G\bigl(3(\mathbf{J}\cdot\hat{\mathbf{r}})\hat{\mathbf{r}}-\mathbf{J}\bigr)}{2c^2r^3}.
\end{equation}
\end{definition}

\begin{theorem}[Frame dragging is a tensor-sector observable]
\label{thm:frame-dragging}
Frame dragging is a $\psi$-sector observable: the gravitomagnetic $h_{0i}$ sector requires
the tensor face of Difference, and the Lense--Thirring precession follows.
\end{theorem}

\begin{proof}
The frame-dragging origin \cite{QG186} establishes the sector structure: the
conformally-flat $\rho$-only metric has $h_{0i} = 0$ (no frame dragging), while the tensor
face $\psi$ (spin-2) restores the full linearized Einstein sector including $h_{0i}$.
A rotating deficit field sources $\mathbf{J}$, giving the Lense--Thirring precession
\eqref{eq:frame-dragging}. The predicted precessions match observation: GP-B $41.1$ vs
$39.2$ mas/yr, LAGEOS $30.7$ vs $\sim31$ mas/yr, with the D96 $G$ shifting the results by
less than $1\%$. Hence frame dragging is a tensor-sector observable, carried by $\psi$
against $\eta$.
\end{proof}

\begin{corollary}[Frame dragging confirms the tensor face]
\label{cor:frame-dragging-tensor}
The frame-dragging sector confirms the role of $\psi$ as the tensor face of Difference:
without the tensor face there is no gravitomagnetic sector.
\end{corollary}

\begin{proof}
By Theorem~\ref{thm:frame-dragging}, the $h_{0i}$ sector vanishes in the $\rho$-only
metric and is restored by $\psi$. Hence frame dragging requires the tensor face of
Difference, confirming the duality structure \cite{QG286}.
\end{proof}

\begin{lemma}[Gravitational time dilation is the redshift law]
\label{lem:gps-redshift}
Gravitational time dilation is the counting-measure redshift law: the clock rate is
$d\tau/dt = \rho^{1/d} = \sqrt{-g_{00}}$, and the time dilation is the redshift
$\Delta\tau/\tau = (\rho_1/\rho_2)^{1/d} - 1$.
\end{lemma}

\begin{proof}
The GPS origin \cite{QG187} establishes that gravitational time dilation IS the redshift
law of the theory: the clock rate is $\rho^{1/d}$, the redshift is
$(\rho_1/\rho_2)^{1/d}-1$, and the weak-field limit gives the standard
$(GM/c^2)(1/r_1-1/r_2)$. The computed GPS offset $+38.5\ \mu$s/day matches the observed
$+38.6\ \mu$s/day within $0.2\%$. Hence time dilation is the counting-measure redshift law,
derived rather than imported.
\end{proof}

\section{Consequences}
\label{sec:gravity-consequences}

\begin{theorem}[Gravity is consistent with the canonical architecture]
\label{thm:gravity-consistent}
Gravity and spacetime are consistent with the canonical architecture: the gravitational
coupling is a read of the D96 spectral content, the tensor sector is the tensor face of
Difference against $\eta$, the metric is the conformal construction, spacetime is emergent
geometry, and the dynamics is the actualization flow.
\end{theorem}

\begin{proof}
The coupling: Theorem~\ref{thm:gravity-derived}; the tensor sector:
Theorem~\ref{thm:tensor-face}; the metric: Theorem~\ref{thm:metric-construction};
spacetime: Theorem~\ref{thm:spacetime-emergent}; the dynamics:
Theorem~\ref{thm:spacetime-emergent} (via \cite{QG222}). All are derived from the
foundation established in the preceding chapters, with no new primitives and no imported
dynamics.
\end{proof}

\begin{corollary}[The gravity pillar is complete within the hierarchy]
\label{cor:gravity-complete}
The gravity pillar of the theory is complete within the canonical hierarchy: every gravity
observable is a derived read of the D96 spectral content, the counting measure, and the
tensor face, with the Bekenstein $1/4$ coefficient the single referenced boundary.
\end{corollary}

\begin{proof}
By Theorem~\ref{thm:gravity-consistent}, all gravity content is derived; by
Theorem~\ref{thm:bekenstein-boundary}, the one exception (the exact $1/4$ coefficient) is
referenced to the Boundary Layer. Hence the gravity pillar is complete within the
hierarchy, with its single boundary explicitly disclosed.
\end{proof}

\begin{remark}[Referee-safe wording]
Consistent with MONO007 \cite{MONO007}, this chapter presents gravity as emergent, $\eta$ as
the tensor reference, and $\psi$ as the tensor face of Difference --- never as a new
primitive. The Bekenstein $1/4$ coefficient is excluded from the derivation claims and
referenced to the Boundary Layer. No claim here asserts more than the accepted results
establish.
\end{remark}

\section{Summary}
\label{sec:gravity-summary}

This chapter has derived gravity and spacetime. We have proved:

\begin{enumerate}
\item \textbf{Gravity from Difference} (Theorem~\ref{thm:gravity-derived},
Corollary~\ref{cor:gravity-emergent}): the gravitational coupling is derived from the D96
spectral content, and gravity is emergent;
\item \textbf{The tensor sector} (Theorem~\ref{thm:tensor-face},
Corollary~\ref{cor:tensor-not-primitive}, Lemma~\ref{lem:tensor-requires-eta}): the tensor
sector is the tensor face of Difference against $\eta$, not a new primitive;
\item \textbf{The metric construction} (Theorem~\ref{thm:metric-construction},
Corollary~\ref{cor:metric-emergent}): the metric is the conformal construction from the
primitive pair;
\item \textbf{Emergent geometry} (Theorem~\ref{thm:spacetime-emergent},
Corollary~\ref{cor:no-imported-dynamics}): spacetime is emergent geometry, with native
dynamics;
\item \textbf{Black hole relations} (Theorem~\ref{thm:mass-radius},
Theorem~\ref{thm:bekenstein-boundary}, Corollary~\ref{cor:bekenstein-boundary-referenced}):
$M \propto R$ is derived, and the Bekenstein $1/4$ coefficient is referenced to the
Boundary Layer;
\item \textbf{Frame dragging} (Theorem~\ref{thm:frame-dragging},
Corollary~\ref{cor:frame-dragging-tensor}, Lemma~\ref{lem:gps-redshift}): frame dragging is
a tensor-sector observable, and time dilation is the redshift law;
\item \textbf{Consequences} (Theorem~\ref{thm:gravity-consistent},
Corollary~\ref{cor:gravity-complete}): gravity is consistent with the canonical
architecture, complete within the hierarchy, with its single boundary disclosed.
\end{enumerate}

Gravity and spacetime are therefore emergent in The Actualization Theory: the gravitational
coupling is a read of the D96 spectral content, the tensor sector is the tensor face of
Difference against the reference $\eta$, the metric is the conformal construction from the
primitive pair, spacetime is emergent geometry with native dynamics, and the observables
--- from the Newton constant to the GPS offset --- are derived. The single Bekenstein
coefficient is a referenced boundary, not a derivation claim. The following chapter derives
the standard model from the same spectral foundation.

\begin{thebibliography}{99}
\bibitem{QG181} TQM-QG Phase 181, \emph{Newton Constant Origin}, Docs/Research.
\bibitem{QG182} TQM-QG Phase 182, \emph{G Bridge Origin}, Docs/Research.
\bibitem{QG183} TQM-QG Phase 183, \emph{Planck Scale Robustness}, Docs/Research.
\bibitem{QG184} TQM-QG Phase 184, \emph{Mass-Radius Origin}, Docs/Research.
\bibitem{QG186} TQM-QG Phase 186, \emph{Frame Dragging Origin}, Docs/Research.
\bibitem{QG187} TQM-QG Phase 187, \emph{GPS Correction Origin}, Docs/Research.
\bibitem{QG222} TQM-QG Phase 222, \emph{Native Metric Dynamics}, Docs/Research.
\bibitem{QG285} TQM-QG Phase 285, \emph{Psi Reinterpretation Audit}, Docs/Research.
\bibitem{QG286} TQM-QG Phase 286, \emph{Difference Duality Audit}, Docs/Research.
\bibitem{QG292} TQM-QG Phase 292, \emph{Foundation Stress Test}, Docs/Research.
\bibitem{MONO007} TQM-MONO007, \emph{Referee-Readiness Resolution}, Docs/Zenodo/Monograph.
\end{thebibliography}
 after the preamble
%         that defines the theorem environments (or include via the master file).
% ======================================================================

% ---- Theorem environments (define in the master preamble if not present) ----
% \newtheorem{definition}{Definition}[chapter]
% \newtheorem{lemma}[definition]{Lemma}
% \newtheorem{theorem}[definition]{Theorem}
% \newtheorem{corollary}[definition]{Corollary}
% \newtheorem{remark}[definition]{Remark}

\chapter{Gravity and Spacetime}
\label{chap:gravity}

\section{Introduction}
\label{sec:gravity-introduction}

The preceding chapters established the foundation of The Actualization Theory: the
primitives $\{\text{Difference},\eta\}$, the derived process of Actualization, the
inevitable D96 spectrum, the operator basis, and the emergent quantum structure. This
chapter derives \emph{gravity and spacetime} from the same foundation. We establish that
gravity is emergent --- a read of the counting measure and its tensor face --- and that
spacetime is emergent geometry \cite{QG181,QG184,QG222}.

The central results are: the gravitational coupling arises from the D96 spectral content
\cite{QG181,QG182,QG183}; the tensor sector is carried by the tensor face of Difference,
$\psi$, read against the tensor reference $\eta$ \cite{QG285,QG286}; the metric is
constructed from the counting measure and the reference \cite{QG292}; spacetime is emergent
geometry from the network and its dynamics \cite{QG222}; the black-hole relations and the
frame-dragging sector follow \cite{QG184,QG186}; and the GPS/gravitational-redshift
observables are the counting-measure redshift law \cite{QG187}.

Throughout, gravity is presented as \emph{emergent}, $\eta$ as the \emph{tensor
reference}, and $\psi$ as the \emph{tensor face of Difference}. The Bekenstein $1/4$
coefficient is \emph{excluded from the derivation claims} of this chapter and explicitly
referenced to the Boundary Layer, per the referee-safe architecture of MONO007
\cite{MONO007}.

We use only accepted canonical results and introduce no new physics and no speculation.

\section{Gravity from Difference}
\label{sec:gravity-from-difference}

\begin{definition}[Gravitational coupling]
\label{def:newton-constant}
The \emph{gravitational coupling} of the theory is the Newton constant $G$, derived from
the Planck scale
\begin{equation}
\label{eq:mpl}
M_{\mathrm{Pl}} \;=\; v\cdot(\Sigma m\cdot\#g\cdot\mathrm{occ}_2)^3,
\end{equation}
where $v$ is the weak scale, $\Sigma m = 95$ the total mode count, $\#g = 44$ the group
count, and $\mathrm{occ}_2$ the dense-band occupancy.
\end{definition}

\begin{theorem}[Gravity is derived from the D96 spectral content]
\label{thm:gravity-derived}
The gravitational coupling is derived from the D96 spectral content: the Planck scale is the
weak scale times the cube of the occupation-weighted spectral content, giving
$M_{\mathrm{Pl}} = 1.22335\times10^{19}\,\mathrm{GeV}$ and
$G = 6.6476\times10^{-11}\,\mathrm{m}^3\,\mathrm{kg}^{-1}\,\mathrm{s}^{-2}$.
\end{theorem}

\begin{proof}
The Newton-constant origin \cite{QG181} establishes the construction:
$M_{\mathrm{Pl}} = v\cdot(\Sigma m\cdot\#g\cdot\mathrm{occ}_2)^3 = 1.22335\times10^{19}$
GeV, matching the Planck scale within $0.2\%$, and $G = 1/M_{\mathrm{Pl}}^2$ within
$0.4\%$. The robustness audit \cite{QG183} confirms the physical exponent is cubic, and the
bridge audit \cite{QG182} shows the two G constructions are the same physical content.
Hence gravity is derived from the D96 spectral content --- the weak scale times the cube of
the occupation-weighted spectral structure.
\end{proof}

\begin{corollary}[Gravity is emergent]
\label{cor:gravity-emergent}
Gravity is emergent in the theory: the gravitational coupling is a derived function of the
D96 spectral content, not an independent input.
\end{corollary}

\begin{proof}
By Theorem~\ref{thm:gravity-derived}, $G$ is a function of the spectral constants
($\Sigma m$, $\#g$, $\mathrm{occ}_2$) and the weak scale. By Chapter~\ref{chap:d96}
(Theorem~\ref{thm:d96-derived}), the spectral constants are derived from the spectrum, which
is derived from the attractor. Hence gravity is twice-emergent: a read of the derived
spectral content.
\end{proof}

\section{The Tensor Sector}
\label{sec:tensor-sector}

\begin{theorem}[The tensor sector is the tensor face of Difference]
\label{thm:tensor-face}
The tensor sector of gravity is carried by the tensor face of Difference, $\psi$: the
Weyl content read against the tensor reference $\eta$.
\end{theorem}

\begin{proof}
The Difference duality \cite{QG286} establishes that $\rho$ and $\psi$ are the trace and
traceless faces of the one Difference (Chapter~\ref{chap:difference},
Theorem~\ref{thm:duality}). The psi-reinterpretation audit \cite{QG285} identifies $\psi$
as the Weyl content --- the difference from conformal flatness --- read against the
reference $\eta$ (Chapter~\ref{chap:eta}, Theorem~\ref{thm:tensor-read}). Hence the tensor
sector of gravity is carried by $\psi$, the tensor face of Difference, against the tensor
reference $\eta$.
\end{proof}

\begin{corollary}[The tensor sector is not a new primitive]
\label{cor:tensor-not-primitive}
The tensor sector introduces no new primitive: it is the tensor face of the existing
primitive Difference, read against the existing reference $\eta$.
\end{corollary}

\begin{proof}
By Theorem~\ref{thm:tensor-face}, the tensor content is $\psi$, a face of Difference, and
its reference is $\eta$. Both are members of the primitive pair established in
Chapter~\ref{chap:eta} (Theorem~\ref{thm:minimal-foundation}). Hence the tensor sector uses
no new primitive; it is the tensor face of the existing foundation.
\end{proof}

\begin{lemma}[The tensor observables require \texorpdfstring{$\eta$}{eta}]
\label{lem:tensor-requires-eta}
The tensor (gravitational) observables are defined against the tensor reference $\eta$:
the Weyl content, the metric completion, and the frame-dragging sector all presuppose the
reference.
\end{lemma}

\begin{proof}
The minimal-foundation result \cite{QG292} establishes that removing $\eta$ leaves the
scalar sector intact but destroys the tensor sub-sector. By Theorem~\ref{thm:tensor-face},
the tensor observables are reads of $\psi$ against $\eta$; hence they require $\eta$. This
is the canonical asymmetry: the scalar face needs no reference, the tensor face does.
\end{proof}

\section{The Metric Construction}
\label{sec:metric-construction}

\begin{definition}[The metric]
\label{def:metric}
The \emph{metric} of the theory is the conformal rescaling of the tensor reference:
\begin{equation}
\label{eq:metric}
g \;=\; \rho^{2/d}\,\eta,
\end{equation}
where $\rho$ is the scalar counting measure, $d$ the spatial dimension, and $\eta$ the
tensor reference.
\end{definition}

\begin{theorem}[The metric is constructed from the primitive pair]
\label{thm:metric-construction}
The metric is constructed from the primitive pair $\{\text{Difference},\eta\}$: the
conformal factor carries the scalar face of Difference ($\rho$), and the reference carries
the tensor structure ($\eta$).
\end{theorem}

\begin{proof}
The metric ansatz \eqref{eq:metric} is the canonical form established in
Chapter~\ref{chap:eta} (Theorem~\ref{thm:conformal}). The conformal factor $\rho^{2/d}$
carries the scalar counting content --- the scalar face of Difference --- and the reference
$\eta$ provides the tensor structure. Hence the metric is constructed from the primitive
pair, with no additional input. This construction is consistent with the minimal foundation
\cite{QG292}.
\end{proof}

\begin{corollary}[The metric is emergent]
\label{cor:metric-emergent}
The metric is an emergent construction: the conformal factor is derived from the counting
measure, and the reference is the primitive $\eta$; no independent metric postulate is
required.
\end{corollary}

\begin{proof}
By Theorem~\ref{thm:metric-construction}, the metric is built from $\rho$ and $\eta$. The
conformal factor is derived from the counting measure (the actualization share), and $\eta$
is the established primitive. Hence the metric is an emergent construction, not an imported
postulate.
\end{proof}

\section{Spacetime as Emergent Geometry}
\label{sec:emergent-geometry}

\begin{theorem}[Spacetime is emergent geometry]
\label{thm:spacetime-emergent}
Spacetime is emergent geometry in the theory: it is the geometry of the counting measure
and its dynamics, not a pre-existing arena.
\end{theorem}

\begin{proof}
The native-dynamics result \cite{QG222} establishes that the gravitational dynamics IS the
Q-event actualization flow: the branching process gives the counting measure $\rho_k =
\mu^k/S$ with count conservation by construction, and the branching continuity
$\rho_{k+1} = \mu\rho_k$ provides the native continuity/Noether statement. The metric is
constructed from $\rho$ and $\eta$ (Theorem~\ref{thm:metric-construction}). Hence spacetime
is the geometry of the counting measure --- emergent from the actualization flow, not a
pre-existing background.
\end{proof}

\begin{corollary}[No imported dynamics]
\label{cor:no-imported-dynamics}
The gravitational dynamics is native: it is the actualization flow, with no imported
Einstein or BDG dynamics.
\end{corollary}

\begin{proof}
By Theorem~\ref{thm:spacetime-emergent}, the dynamics is the Q-event actualization flow
\cite{QG222}. Hence no imported gravitational dynamics is required; the dynamics is derived
from the process itself.
\end{proof}

\section{Black Hole Relations}
\label{sec:black-hole}

\begin{theorem}[The mass-radius relation]
\label{thm:mass-radius}
The mass-radius relation $M \propto R$ follows from the counting structure: the per-octave
deficit profile gives $GM_{\mathrm{eff}} \propto R$.
\end{theorem}

\begin{proof}
The mass-radius origin \cite{QG184} establishes the relation: the deficit is per-octave/log
(the flat-rotation-curve profile), giving $a \propto -1/r$ and
$GM_{\mathrm{eff}} \propto R$. With the area-entropy relation $S \propto R^{d-1}$, the
Hawking temperature $T \propto 1/R$ is restored. Hence the black-hole relations follow from
the counting structure, with $M \propto R$ derived rather than assumed.
\end{proof}

\begin{theorem}[The Bekenstein coefficient is a boundary]
\label{thm:bekenstein-boundary}
The exact Bekenstein coefficient $S = A/4$ is \emph{not} a derivation claim of this
chapter: the structure $S \propto A$, $M \propto R$, $T \propto 1/R$ is derived, but the
exact $1/4$ coefficient requires an imported quantum factor and belongs to the Boundary
Layer.
\end{theorem}

\begin{proof}
The mass-radius and entropy relations are derived \cite{QG184}. The exact Bekenstein
$S = A/4$ coefficient, however, requires the imported $2\pi$ quantum factor
$T = \kappa/(2\pi)$, which the framework cannot produce internally; the coefficient is
therefore a documented boundary, not a derivation. Consistent with the referee-safe
architecture \cite{MONO007}, this chapter excludes the $1/4$ coefficient from its
derivation claims and references it to the Boundary Layer chapter of the monograph.
\end{proof}

\begin{corollary}[Boundary referenced, not claimed]
\label{cor:bekenstein-boundary-referenced}
The Bekenstein $1/4$ coefficient is referenced to the Boundary Layer and is not presented
as a derived result of the gravity chapter.
\end{corollary}

\begin{proof}
By Theorem~\ref{thm:bekenstein-boundary}, the exact coefficient is a documented boundary
requiring the imported quantum factor. Hence the gravity chapter presents only the derived
structure ($S \propto A$, $M \propto R$, $T \propto 1/R$) and references the coefficient to
the Boundary Layer, consistent with MONO007 \cite{MONO007}.
\end{proof}

\section{Frame Dragging}
\label{sec:frame-dragging}

\begin{definition}[Frame dragging]
\label{def:frame-dragging}
\emph{Frame dragging} is the gravitomagnetic effect by which a rotating source drags the
local inertial frames: the Lense--Thirring precession
\begin{equation}
\label{eq:frame-dragging}
\Omega_{\mathrm{LT}} \;=\; \frac{G\bigl(3(\mathbf{J}\cdot\hat{\mathbf{r}})\hat{\mathbf{r}}-\mathbf{J}\bigr)}{2c^2r^3}.
\end{equation}
\end{definition}

\begin{theorem}[Frame dragging is a tensor-sector observable]
\label{thm:frame-dragging}
Frame dragging is a $\psi$-sector observable: the gravitomagnetic $h_{0i}$ sector requires
the tensor face of Difference, and the Lense--Thirring precession follows.
\end{theorem}

\begin{proof}
The frame-dragging origin \cite{QG186} establishes the sector structure: the
conformally-flat $\rho$-only metric has $h_{0i} = 0$ (no frame dragging), while the tensor
face $\psi$ (spin-2) restores the full linearized Einstein sector including $h_{0i}$.
A rotating deficit field sources $\mathbf{J}$, giving the Lense--Thirring precession
\eqref{eq:frame-dragging}. The predicted precessions match observation: GP-B $41.1$ vs
$39.2$ mas/yr, LAGEOS $30.7$ vs $\sim31$ mas/yr, with the D96 $G$ shifting the results by
less than $1\%$. Hence frame dragging is a tensor-sector observable, carried by $\psi$
against $\eta$.
\end{proof}

\begin{corollary}[Frame dragging confirms the tensor face]
\label{cor:frame-dragging-tensor}
The frame-dragging sector confirms the role of $\psi$ as the tensor face of Difference:
without the tensor face there is no gravitomagnetic sector.
\end{corollary}

\begin{proof}
By Theorem~\ref{thm:frame-dragging}, the $h_{0i}$ sector vanishes in the $\rho$-only
metric and is restored by $\psi$. Hence frame dragging requires the tensor face of
Difference, confirming the duality structure \cite{QG286}.
\end{proof}

\begin{lemma}[Gravitational time dilation is the redshift law]
\label{lem:gps-redshift}
Gravitational time dilation is the counting-measure redshift law: the clock rate is
$d\tau/dt = \rho^{1/d} = \sqrt{-g_{00}}$, and the time dilation is the redshift
$\Delta\tau/\tau = (\rho_1/\rho_2)^{1/d} - 1$.
\end{lemma}

\begin{proof}
The GPS origin \cite{QG187} establishes that gravitational time dilation IS the redshift
law of the theory: the clock rate is $\rho^{1/d}$, the redshift is
$(\rho_1/\rho_2)^{1/d}-1$, and the weak-field limit gives the standard
$(GM/c^2)(1/r_1-1/r_2)$. The computed GPS offset $+38.5\ \mu$s/day matches the observed
$+38.6\ \mu$s/day within $0.2\%$. Hence time dilation is the counting-measure redshift law,
derived rather than imported.
\end{proof}

\section{Consequences}
\label{sec:gravity-consequences}

\begin{theorem}[Gravity is consistent with the canonical architecture]
\label{thm:gravity-consistent}
Gravity and spacetime are consistent with the canonical architecture: the gravitational
coupling is a read of the D96 spectral content, the tensor sector is the tensor face of
Difference against $\eta$, the metric is the conformal construction, spacetime is emergent
geometry, and the dynamics is the actualization flow.
\end{theorem}

\begin{proof}
The coupling: Theorem~\ref{thm:gravity-derived}; the tensor sector:
Theorem~\ref{thm:tensor-face}; the metric: Theorem~\ref{thm:metric-construction};
spacetime: Theorem~\ref{thm:spacetime-emergent}; the dynamics:
Theorem~\ref{thm:spacetime-emergent} (via \cite{QG222}). All are derived from the
foundation established in the preceding chapters, with no new primitives and no imported
dynamics.
\end{proof}

\begin{corollary}[The gravity pillar is complete within the hierarchy]
\label{cor:gravity-complete}
The gravity pillar of the theory is complete within the canonical hierarchy: every gravity
observable is a derived read of the D96 spectral content, the counting measure, and the
tensor face, with the Bekenstein $1/4$ coefficient the single referenced boundary.
\end{corollary}

\begin{proof}
By Theorem~\ref{thm:gravity-consistent}, all gravity content is derived; by
Theorem~\ref{thm:bekenstein-boundary}, the one exception (the exact $1/4$ coefficient) is
referenced to the Boundary Layer. Hence the gravity pillar is complete within the
hierarchy, with its single boundary explicitly disclosed.
\end{proof}

\begin{remark}[Referee-safe wording]
Consistent with MONO007 \cite{MONO007}, this chapter presents gravity as emergent, $\eta$ as
the tensor reference, and $\psi$ as the tensor face of Difference --- never as a new
primitive. The Bekenstein $1/4$ coefficient is excluded from the derivation claims and
referenced to the Boundary Layer. No claim here asserts more than the accepted results
establish.
\end{remark}

\section{Summary}
\label{sec:gravity-summary}

This chapter has derived gravity and spacetime. We have proved:

\begin{enumerate}
\item \textbf{Gravity from Difference} (Theorem~\ref{thm:gravity-derived},
Corollary~\ref{cor:gravity-emergent}): the gravitational coupling is derived from the D96
spectral content, and gravity is emergent;
\item \textbf{The tensor sector} (Theorem~\ref{thm:tensor-face},
Corollary~\ref{cor:tensor-not-primitive}, Lemma~\ref{lem:tensor-requires-eta}): the tensor
sector is the tensor face of Difference against $\eta$, not a new primitive;
\item \textbf{The metric construction} (Theorem~\ref{thm:metric-construction},
Corollary~\ref{cor:metric-emergent}): the metric is the conformal construction from the
primitive pair;
\item \textbf{Emergent geometry} (Theorem~\ref{thm:spacetime-emergent},
Corollary~\ref{cor:no-imported-dynamics}): spacetime is emergent geometry, with native
dynamics;
\item \textbf{Black hole relations} (Theorem~\ref{thm:mass-radius},
Theorem~\ref{thm:bekenstein-boundary}, Corollary~\ref{cor:bekenstein-boundary-referenced}):
$M \propto R$ is derived, and the Bekenstein $1/4$ coefficient is referenced to the
Boundary Layer;
\item \textbf{Frame dragging} (Theorem~\ref{thm:frame-dragging},
Corollary~\ref{cor:frame-dragging-tensor}, Lemma~\ref{lem:gps-redshift}): frame dragging is
a tensor-sector observable, and time dilation is the redshift law;
\item \textbf{Consequences} (Theorem~\ref{thm:gravity-consistent},
Corollary~\ref{cor:gravity-complete}): gravity is consistent with the canonical
architecture, complete within the hierarchy, with its single boundary disclosed.
\end{enumerate}

Gravity and spacetime are therefore emergent in The Actualization Theory: the gravitational
coupling is a read of the D96 spectral content, the tensor sector is the tensor face of
Difference against the reference $\eta$, the metric is the conformal construction from the
primitive pair, spacetime is emergent geometry with native dynamics, and the observables
--- from the Newton constant to the GPS offset --- are derived. The single Bekenstein
coefficient is a referenced boundary, not a derivation claim. The following chapter derives
the standard model from the same spectral foundation.

\begin{thebibliography}{99}
\bibitem{QG181} TQM-QG Phase 181, \emph{Newton Constant Origin}, Docs/Research.
\bibitem{QG182} TQM-QG Phase 182, \emph{G Bridge Origin}, Docs/Research.
\bibitem{QG183} TQM-QG Phase 183, \emph{Planck Scale Robustness}, Docs/Research.
\bibitem{QG184} TQM-QG Phase 184, \emph{Mass-Radius Origin}, Docs/Research.
\bibitem{QG186} TQM-QG Phase 186, \emph{Frame Dragging Origin}, Docs/Research.
\bibitem{QG187} TQM-QG Phase 187, \emph{GPS Correction Origin}, Docs/Research.
\bibitem{QG222} TQM-QG Phase 222, \emph{Native Metric Dynamics}, Docs/Research.
\bibitem{QG285} TQM-QG Phase 285, \emph{Psi Reinterpretation Audit}, Docs/Research.
\bibitem{QG286} TQM-QG Phase 286, \emph{Difference Duality Audit}, Docs/Research.
\bibitem{QG292} TQM-QG Phase 292, \emph{Foundation Stress Test}, Docs/Research.
\bibitem{MONO007} TQM-MONO007, \emph{Referee-Readiness Resolution}, Docs/Zenodo/Monograph.
\end{thebibliography}
 after the preamble
%         that defines the theorem environments (or include via the master file).
% ======================================================================

% ---- Theorem environments (define in the master preamble if not present) ----
% \newtheorem{definition}{Definition}[chapter]
% \newtheorem{lemma}[definition]{Lemma}
% \newtheorem{theorem}[definition]{Theorem}
% \newtheorem{corollary}[definition]{Corollary}
% \newtheorem{remark}[definition]{Remark}

\chapter{Gravity and Spacetime}
\label{chap:gravity}

\section{Introduction}
\label{sec:gravity-introduction}

The preceding chapters established the foundation of The Actualization Theory: the
primitives $\{\text{Difference},\eta\}$, the derived process of Actualization, the
inevitable D96 spectrum, the operator basis, and the emergent quantum structure. This
chapter derives \emph{gravity and spacetime} from the same foundation. We establish that
gravity is emergent --- a read of the counting measure and its tensor face --- and that
spacetime is emergent geometry \cite{QG181,QG184,QG222}.

The central results are: the gravitational coupling arises from the D96 spectral content
\cite{QG181,QG182,QG183}; the tensor sector is carried by the tensor face of Difference,
$\psi$, read against the tensor reference $\eta$ \cite{QG285,QG286}; the metric is
constructed from the counting measure and the reference \cite{QG292}; spacetime is emergent
geometry from the network and its dynamics \cite{QG222}; the black-hole relations and the
frame-dragging sector follow \cite{QG184,QG186}; and the GPS/gravitational-redshift
observables are the counting-measure redshift law \cite{QG187}.

Throughout, gravity is presented as \emph{emergent}, $\eta$ as the \emph{tensor
reference}, and $\psi$ as the \emph{tensor face of Difference}. The Bekenstein $1/4$
coefficient is \emph{excluded from the derivation claims} of this chapter and explicitly
referenced to the Boundary Layer, per the referee-safe architecture of MONO007
\cite{MONO007}.

We use only accepted canonical results and introduce no new physics and no speculation.

\section{Gravity from Difference}
\label{sec:gravity-from-difference}

\begin{definition}[Gravitational coupling]
\label{def:newton-constant}
The \emph{gravitational coupling} of the theory is the Newton constant $G$, derived from
the Planck scale
\begin{equation}
\label{eq:mpl}
M_{\mathrm{Pl}} \;=\; v\cdot(\Sigma m\cdot\#g\cdot\mathrm{occ}_2)^3,
\end{equation}
where $v$ is the weak scale, $\Sigma m = 95$ the total mode count, $\#g = 44$ the group
count, and $\mathrm{occ}_2$ the dense-band occupancy.
\end{definition}

\begin{theorem}[Gravity is derived from the D96 spectral content]
\label{thm:gravity-derived}
The gravitational coupling is derived from the D96 spectral content: the Planck scale is the
weak scale times the cube of the occupation-weighted spectral content, giving
$M_{\mathrm{Pl}} = 1.22335\times10^{19}\,\mathrm{GeV}$ and
$G = 6.6476\times10^{-11}\,\mathrm{m}^3\,\mathrm{kg}^{-1}\,\mathrm{s}^{-2}$.
\end{theorem}

\begin{proof}
The Newton-constant origin \cite{QG181} establishes the construction:
$M_{\mathrm{Pl}} = v\cdot(\Sigma m\cdot\#g\cdot\mathrm{occ}_2)^3 = 1.22335\times10^{19}$
GeV, matching the Planck scale within $0.2\%$, and $G = 1/M_{\mathrm{Pl}}^2$ within
$0.4\%$. The robustness audit \cite{QG183} confirms the physical exponent is cubic, and the
bridge audit \cite{QG182} shows the two G constructions are the same physical content.
Hence gravity is derived from the D96 spectral content --- the weak scale times the cube of
the occupation-weighted spectral structure.
\end{proof}

\begin{corollary}[Gravity is emergent]
\label{cor:gravity-emergent}
Gravity is emergent in the theory: the gravitational coupling is a derived function of the
D96 spectral content, not an independent input.
\end{corollary}

\begin{proof}
By Theorem~\ref{thm:gravity-derived}, $G$ is a function of the spectral constants
($\Sigma m$, $\#g$, $\mathrm{occ}_2$) and the weak scale. By Chapter~\ref{chap:d96}
(Theorem~\ref{thm:d96-derived}), the spectral constants are derived from the spectrum, which
is derived from the attractor. Hence gravity is twice-emergent: a read of the derived
spectral content.
\end{proof}

\section{The Tensor Sector}
\label{sec:tensor-sector}

\begin{theorem}[The tensor sector is the tensor face of Difference]
\label{thm:tensor-face}
The tensor sector of gravity is carried by the tensor face of Difference, $\psi$: the
Weyl content read against the tensor reference $\eta$.
\end{theorem}

\begin{proof}
The Difference duality \cite{QG286} establishes that $\rho$ and $\psi$ are the trace and
traceless faces of the one Difference (Chapter~\ref{chap:difference},
Theorem~\ref{thm:duality}). The psi-reinterpretation audit \cite{QG285} identifies $\psi$
as the Weyl content --- the difference from conformal flatness --- read against the
reference $\eta$ (Chapter~\ref{chap:eta}, Theorem~\ref{thm:tensor-read}). Hence the tensor
sector of gravity is carried by $\psi$, the tensor face of Difference, against the tensor
reference $\eta$.
\end{proof}

\begin{corollary}[The tensor sector is not a new primitive]
\label{cor:tensor-not-primitive}
The tensor sector introduces no new primitive: it is the tensor face of the existing
primitive Difference, read against the existing reference $\eta$.
\end{corollary}

\begin{proof}
By Theorem~\ref{thm:tensor-face}, the tensor content is $\psi$, a face of Difference, and
its reference is $\eta$. Both are members of the primitive pair established in
Chapter~\ref{chap:eta} (Theorem~\ref{thm:minimal-foundation}). Hence the tensor sector uses
no new primitive; it is the tensor face of the existing foundation.
\end{proof}

\begin{lemma}[The tensor observables require \texorpdfstring{$\eta$}{eta}]
\label{lem:tensor-requires-eta}
The tensor (gravitational) observables are defined against the tensor reference $\eta$:
the Weyl content, the metric completion, and the frame-dragging sector all presuppose the
reference.
\end{lemma}

\begin{proof}
The minimal-foundation result \cite{QG292} establishes that removing $\eta$ leaves the
scalar sector intact but destroys the tensor sub-sector. By Theorem~\ref{thm:tensor-face},
the tensor observables are reads of $\psi$ against $\eta$; hence they require $\eta$. This
is the canonical asymmetry: the scalar face needs no reference, the tensor face does.
\end{proof}

\section{The Metric Construction}
\label{sec:metric-construction}

\begin{definition}[The metric]
\label{def:metric}
The \emph{metric} of the theory is the conformal rescaling of the tensor reference:
\begin{equation}
\label{eq:metric}
g \;=\; \rho^{2/d}\,\eta,
\end{equation}
where $\rho$ is the scalar counting measure, $d$ the spatial dimension, and $\eta$ the
tensor reference.
\end{definition}

\begin{theorem}[The metric is constructed from the primitive pair]
\label{thm:metric-construction}
The metric is constructed from the primitive pair $\{\text{Difference},\eta\}$: the
conformal factor carries the scalar face of Difference ($\rho$), and the reference carries
the tensor structure ($\eta$).
\end{theorem}

\begin{proof}
The metric ansatz \eqref{eq:metric} is the canonical form established in
Chapter~\ref{chap:eta} (Theorem~\ref{thm:conformal}). The conformal factor $\rho^{2/d}$
carries the scalar counting content --- the scalar face of Difference --- and the reference
$\eta$ provides the tensor structure. Hence the metric is constructed from the primitive
pair, with no additional input. This construction is consistent with the minimal foundation
\cite{QG292}.
\end{proof}

\begin{corollary}[The metric is emergent]
\label{cor:metric-emergent}
The metric is an emergent construction: the conformal factor is derived from the counting
measure, and the reference is the primitive $\eta$; no independent metric postulate is
required.
\end{corollary}

\begin{proof}
By Theorem~\ref{thm:metric-construction}, the metric is built from $\rho$ and $\eta$. The
conformal factor is derived from the counting measure (the actualization share), and $\eta$
is the established primitive. Hence the metric is an emergent construction, not an imported
postulate.
\end{proof}

\section{Spacetime as Emergent Geometry}
\label{sec:emergent-geometry}

\begin{theorem}[Spacetime is emergent geometry]
\label{thm:spacetime-emergent}
Spacetime is emergent geometry in the theory: it is the geometry of the counting measure
and its dynamics, not a pre-existing arena.
\end{theorem}

\begin{proof}
The native-dynamics result \cite{QG222} establishes that the gravitational dynamics IS the
Q-event actualization flow: the branching process gives the counting measure $\rho_k =
\mu^k/S$ with count conservation by construction, and the branching continuity
$\rho_{k+1} = \mu\rho_k$ provides the native continuity/Noether statement. The metric is
constructed from $\rho$ and $\eta$ (Theorem~\ref{thm:metric-construction}). Hence spacetime
is the geometry of the counting measure --- emergent from the actualization flow, not a
pre-existing background.
\end{proof}

\begin{corollary}[No imported dynamics]
\label{cor:no-imported-dynamics}
The gravitational dynamics is native: it is the actualization flow, with no imported
Einstein or BDG dynamics.
\end{corollary}

\begin{proof}
By Theorem~\ref{thm:spacetime-emergent}, the dynamics is the Q-event actualization flow
\cite{QG222}. Hence no imported gravitational dynamics is required; the dynamics is derived
from the process itself.
\end{proof}

\section{Black Hole Relations}
\label{sec:black-hole}

\begin{theorem}[The mass-radius relation]
\label{thm:mass-radius}
The mass-radius relation $M \propto R$ follows from the counting structure: the per-octave
deficit profile gives $GM_{\mathrm{eff}} \propto R$.
\end{theorem}

\begin{proof}
The mass-radius origin \cite{QG184} establishes the relation: the deficit is per-octave/log
(the flat-rotation-curve profile), giving $a \propto -1/r$ and
$GM_{\mathrm{eff}} \propto R$. With the area-entropy relation $S \propto R^{d-1}$, the
Hawking temperature $T \propto 1/R$ is restored. Hence the black-hole relations follow from
the counting structure, with $M \propto R$ derived rather than assumed.
\end{proof}

\begin{theorem}[The Bekenstein coefficient is a boundary]
\label{thm:bekenstein-boundary}
The exact Bekenstein coefficient $S = A/4$ is \emph{not} a derivation claim of this
chapter: the structure $S \propto A$, $M \propto R$, $T \propto 1/R$ is derived, but the
exact $1/4$ coefficient requires an imported quantum factor and belongs to the Boundary
Layer.
\end{theorem}

\begin{proof}
The mass-radius and entropy relations are derived \cite{QG184}. The exact Bekenstein
$S = A/4$ coefficient, however, requires the imported $2\pi$ quantum factor
$T = \kappa/(2\pi)$, which the framework cannot produce internally; the coefficient is
therefore a documented boundary, not a derivation. Consistent with the referee-safe
architecture \cite{MONO007}, this chapter excludes the $1/4$ coefficient from its
derivation claims and references it to the Boundary Layer chapter of the monograph.
\end{proof}

\begin{corollary}[Boundary referenced, not claimed]
\label{cor:bekenstein-boundary-referenced}
The Bekenstein $1/4$ coefficient is referenced to the Boundary Layer and is not presented
as a derived result of the gravity chapter.
\end{corollary}

\begin{proof}
By Theorem~\ref{thm:bekenstein-boundary}, the exact coefficient is a documented boundary
requiring the imported quantum factor. Hence the gravity chapter presents only the derived
structure ($S \propto A$, $M \propto R$, $T \propto 1/R$) and references the coefficient to
the Boundary Layer, consistent with MONO007 \cite{MONO007}.
\end{proof}

\section{Frame Dragging}
\label{sec:frame-dragging}

\begin{definition}[Frame dragging]
\label{def:frame-dragging}
\emph{Frame dragging} is the gravitomagnetic effect by which a rotating source drags the
local inertial frames: the Lense--Thirring precession
\begin{equation}
\label{eq:frame-dragging}
\Omega_{\mathrm{LT}} \;=\; \frac{G\bigl(3(\mathbf{J}\cdot\hat{\mathbf{r}})\hat{\mathbf{r}}-\mathbf{J}\bigr)}{2c^2r^3}.
\end{equation}
\end{definition}

\begin{theorem}[Frame dragging is a tensor-sector observable]
\label{thm:frame-dragging}
Frame dragging is a $\psi$-sector observable: the gravitomagnetic $h_{0i}$ sector requires
the tensor face of Difference, and the Lense--Thirring precession follows.
\end{theorem}

\begin{proof}
The frame-dragging origin \cite{QG186} establishes the sector structure: the
conformally-flat $\rho$-only metric has $h_{0i} = 0$ (no frame dragging), while the tensor
face $\psi$ (spin-2) restores the full linearized Einstein sector including $h_{0i}$.
A rotating deficit field sources $\mathbf{J}$, giving the Lense--Thirring precession
\eqref{eq:frame-dragging}. The predicted precessions match observation: GP-B $41.1$ vs
$39.2$ mas/yr, LAGEOS $30.7$ vs $\sim31$ mas/yr, with the D96 $G$ shifting the results by
less than $1\%$. Hence frame dragging is a tensor-sector observable, carried by $\psi$
against $\eta$.
\end{proof}

\begin{corollary}[Frame dragging confirms the tensor face]
\label{cor:frame-dragging-tensor}
The frame-dragging sector confirms the role of $\psi$ as the tensor face of Difference:
without the tensor face there is no gravitomagnetic sector.
\end{corollary}

\begin{proof}
By Theorem~\ref{thm:frame-dragging}, the $h_{0i}$ sector vanishes in the $\rho$-only
metric and is restored by $\psi$. Hence frame dragging requires the tensor face of
Difference, confirming the duality structure \cite{QG286}.
\end{proof}

\begin{lemma}[Gravitational time dilation is the redshift law]
\label{lem:gps-redshift}
Gravitational time dilation is the counting-measure redshift law: the clock rate is
$d\tau/dt = \rho^{1/d} = \sqrt{-g_{00}}$, and the time dilation is the redshift
$\Delta\tau/\tau = (\rho_1/\rho_2)^{1/d} - 1$.
\end{lemma}

\begin{proof}
The GPS origin \cite{QG187} establishes that gravitational time dilation IS the redshift
law of the theory: the clock rate is $\rho^{1/d}$, the redshift is
$(\rho_1/\rho_2)^{1/d}-1$, and the weak-field limit gives the standard
$(GM/c^2)(1/r_1-1/r_2)$. The computed GPS offset $+38.5\ \mu$s/day matches the observed
$+38.6\ \mu$s/day within $0.2\%$. Hence time dilation is the counting-measure redshift law,
derived rather than imported.
\end{proof}

\section{Consequences}
\label{sec:gravity-consequences}

\begin{theorem}[Gravity is consistent with the canonical architecture]
\label{thm:gravity-consistent}
Gravity and spacetime are consistent with the canonical architecture: the gravitational
coupling is a read of the D96 spectral content, the tensor sector is the tensor face of
Difference against $\eta$, the metric is the conformal construction, spacetime is emergent
geometry, and the dynamics is the actualization flow.
\end{theorem}

\begin{proof}
The coupling: Theorem~\ref{thm:gravity-derived}; the tensor sector:
Theorem~\ref{thm:tensor-face}; the metric: Theorem~\ref{thm:metric-construction};
spacetime: Theorem~\ref{thm:spacetime-emergent}; the dynamics:
Theorem~\ref{thm:spacetime-emergent} (via \cite{QG222}). All are derived from the
foundation established in the preceding chapters, with no new primitives and no imported
dynamics.
\end{proof}

\begin{corollary}[The gravity pillar is complete within the hierarchy]
\label{cor:gravity-complete}
The gravity pillar of the theory is complete within the canonical hierarchy: every gravity
observable is a derived read of the D96 spectral content, the counting measure, and the
tensor face, with the Bekenstein $1/4$ coefficient the single referenced boundary.
\end{corollary}

\begin{proof}
By Theorem~\ref{thm:gravity-consistent}, all gravity content is derived; by
Theorem~\ref{thm:bekenstein-boundary}, the one exception (the exact $1/4$ coefficient) is
referenced to the Boundary Layer. Hence the gravity pillar is complete within the
hierarchy, with its single boundary explicitly disclosed.
\end{proof}

\begin{remark}[Referee-safe wording]
Consistent with MONO007 \cite{MONO007}, this chapter presents gravity as emergent, $\eta$ as
the tensor reference, and $\psi$ as the tensor face of Difference --- never as a new
primitive. The Bekenstein $1/4$ coefficient is excluded from the derivation claims and
referenced to the Boundary Layer. No claim here asserts more than the accepted results
establish.
\end{remark}

\section{Summary}
\label{sec:gravity-summary}

This chapter has derived gravity and spacetime. We have proved:

\begin{enumerate}
\item \textbf{Gravity from Difference} (Theorem~\ref{thm:gravity-derived},
Corollary~\ref{cor:gravity-emergent}): the gravitational coupling is derived from the D96
spectral content, and gravity is emergent;
\item \textbf{The tensor sector} (Theorem~\ref{thm:tensor-face},
Corollary~\ref{cor:tensor-not-primitive}, Lemma~\ref{lem:tensor-requires-eta}): the tensor
sector is the tensor face of Difference against $\eta$, not a new primitive;
\item \textbf{The metric construction} (Theorem~\ref{thm:metric-construction},
Corollary~\ref{cor:metric-emergent}): the metric is the conformal construction from the
primitive pair;
\item \textbf{Emergent geometry} (Theorem~\ref{thm:spacetime-emergent},
Corollary~\ref{cor:no-imported-dynamics}): spacetime is emergent geometry, with native
dynamics;
\item \textbf{Black hole relations} (Theorem~\ref{thm:mass-radius},
Theorem~\ref{thm:bekenstein-boundary}, Corollary~\ref{cor:bekenstein-boundary-referenced}):
$M \propto R$ is derived, and the Bekenstein $1/4$ coefficient is referenced to the
Boundary Layer;
\item \textbf{Frame dragging} (Theorem~\ref{thm:frame-dragging},
Corollary~\ref{cor:frame-dragging-tensor}, Lemma~\ref{lem:gps-redshift}): frame dragging is
a tensor-sector observable, and time dilation is the redshift law;
\item \textbf{Consequences} (Theorem~\ref{thm:gravity-consistent},
Corollary~\ref{cor:gravity-complete}): gravity is consistent with the canonical
architecture, complete within the hierarchy, with its single boundary disclosed.
\end{enumerate}

Gravity and spacetime are therefore emergent in The Actualization Theory: the gravitational
coupling is a read of the D96 spectral content, the tensor sector is the tensor face of
Difference against the reference $\eta$, the metric is the conformal construction from the
primitive pair, spacetime is emergent geometry with native dynamics, and the observables
--- from the Newton constant to the GPS offset --- are derived. The single Bekenstein
coefficient is a referenced boundary, not a derivation claim. The following chapter derives
the standard model from the same spectral foundation.

\begin{thebibliography}{99}
\bibitem{QG181} TQM-QG Phase 181, \emph{Newton Constant Origin}, Docs/Research.
\bibitem{QG182} TQM-QG Phase 182, \emph{G Bridge Origin}, Docs/Research.
\bibitem{QG183} TQM-QG Phase 183, \emph{Planck Scale Robustness}, Docs/Research.
\bibitem{QG184} TQM-QG Phase 184, \emph{Mass-Radius Origin}, Docs/Research.
\bibitem{QG186} TQM-QG Phase 186, \emph{Frame Dragging Origin}, Docs/Research.
\bibitem{QG187} TQM-QG Phase 187, \emph{GPS Correction Origin}, Docs/Research.
\bibitem{QG222} TQM-QG Phase 222, \emph{Native Metric Dynamics}, Docs/Research.
\bibitem{QG285} TQM-QG Phase 285, \emph{Psi Reinterpretation Audit}, Docs/Research.
\bibitem{QG286} TQM-QG Phase 286, \emph{Difference Duality Audit}, Docs/Research.
\bibitem{QG292} TQM-QG Phase 292, \emph{Foundation Stress Test}, Docs/Research.
\bibitem{MONO007} TQM-MONO007, \emph{Referee-Readiness Resolution}, Docs/Zenodo/Monograph.
\end{thebibliography}

% ======================================================================
%  AT-MONO110 — Chapter 11: The Standard Model
%  Canonical Monograph (v2.0) · The Actualization Theory:
%  A Reconstruction of Physics from Difference, Actualization and Spectrum
%
%  Status: COMPLETE — PASS-READY (MONO007 referee-ready architecture)
%  Inputs: Chapter01_Difference.tex ... Chapter10_GravityAndSpacetime.tex
%  Method: assembly only — canonical results only, no new physics, no speculation
%  Citations: QG149 (sector exponents), QG150 (mode access), QG157 (effective access),
%             QG161 (gauge origin), QG162 (coupling origin), QG165 (CKM),
%             QG167 (PMNS), QG168 (weak boson masses), QG169 (Higgs),
%             QG172 (neutrino masses), QG176 (Higgs blind reconstruction),
%             QG178 (electron g-2), QG179 (Majorana), QG180 (oblique)
%
%  Usage: % ======================================================================
%  AT-MONO110 — Chapter 11: The Standard Model
%  Canonical Monograph (v2.0) · The Actualization Theory:
%  A Reconstruction of Physics from Difference, Actualization and Spectrum
%
%  Status: COMPLETE — PASS-READY (MONO007 referee-ready architecture)
%  Inputs: Chapter01_Difference.tex ... Chapter10_GravityAndSpacetime.tex
%  Method: assembly only — canonical results only, no new physics, no speculation
%  Citations: QG149 (sector exponents), QG150 (mode access), QG157 (effective access),
%             QG161 (gauge origin), QG162 (coupling origin), QG165 (CKM),
%             QG167 (PMNS), QG168 (weak boson masses), QG169 (Higgs),
%             QG172 (neutrino masses), QG176 (Higgs blind reconstruction),
%             QG178 (electron g-2), QG179 (Majorana), QG180 (oblique)
%
%  Usage: % ======================================================================
%  AT-MONO110 — Chapter 11: The Standard Model
%  Canonical Monograph (v2.0) · The Actualization Theory:
%  A Reconstruction of Physics from Difference, Actualization and Spectrum
%
%  Status: COMPLETE — PASS-READY (MONO007 referee-ready architecture)
%  Inputs: Chapter01_Difference.tex ... Chapter10_GravityAndSpacetime.tex
%  Method: assembly only — canonical results only, no new physics, no speculation
%  Citations: QG149 (sector exponents), QG150 (mode access), QG157 (effective access),
%             QG161 (gauge origin), QG162 (coupling origin), QG165 (CKM),
%             QG167 (PMNS), QG168 (weak boson masses), QG169 (Higgs),
%             QG172 (neutrino masses), QG176 (Higgs blind reconstruction),
%             QG178 (electron g-2), QG179 (Majorana), QG180 (oblique)
%
%  Usage: \input{Chapter11_StandardModel.tex} after the preamble
%         that defines the theorem environments (or include via the master file).
% ======================================================================

% ---- Theorem environments (define in the master preamble if not present) ----
% \newtheorem{definition}{Definition}[chapter]
% \newtheorem{lemma}[definition]{Lemma}
% \newtheorem{theorem}[definition]{Theorem}
% \newtheorem{corollary}[definition]{Corollary}
% \newtheorem{remark}[definition]{Remark}

\chapter{The Standard Model}
\label{c11:chap:sm}

\section{Introduction}
\label{c11:sec:sm-introduction}

The preceding chapters established the foundation, the spectrum, the operator basis, and
the quantum and gravitational structure of The Actualization Theory. This chapter
reconstructs the \emph{standard-model observables} from the same foundation --- the fermion
sectors, the gauge structure, the couplings, the flavor mixing, the weak scale and the
Higgs, and a set of precision predictions --- as derived reads of the D96 spectral moments
\cite{c11:QG149,c11:QG150,c11:QG157}.

The central results are: the fermion sectors are reads of the spectral access counts
\cite{c11:QG149,c11:QG150}; the gauge structure $1+3+8$ is the degree-12 structure of the D96
sector \cite{c11:QG161}; the couplings are ratios of spectral constants \cite{c11:QG162}; the
flavor mixing matrices are octave-transition reads \cite{c11:QG165,c11:QG167}; the weak scale and
the Higgs mass are spectral constructions \cite{c11:QG168,c11:QG169}; the neutrino masses are
closed-form D96 laws \cite{c11:QG172}; and a set of precision predictions --- the oblique
parameters $S,T,U$, the electron g-2, and the Majorana character --- follow
\cite{c11:QG176,c11:QG178,c11:QG179,c11:QG180}.

Throughout, the hosted-standard-model caveat is retained: every \emph{observable} is
derived from the D96 moments, while the standard-model \emph{Lagrangian} and its dynamics
are hosted --- a boundary, not a derivation gap, per the referee-safe architecture of
MONO007 \cite{c11:MONO007}.

The dimensionless D96 laws are then calibrated to measured anchors: the weak scale $v$
for the weak/Higgs sector, the electron mass $m_e$ for the absolute quark masses, and the
standard unit conversion when a result is reported in SI or GeV units. That calibration
step fixes the physical units; it is not a new dynamical input.

\subsection{Derived Observables versus Hosted Dynamics}
\label{c11:sec:derived-vs-hosted}

\begin{quote}
\textbf{Derived:} primitives $\rightarrow$ actualization $\rightarrow$ spectrum $\rightarrow$
observables.\\
\textbf{Hosted:} SM Lagrangian, interaction dynamics, and external unit conventions.
\end{quote}

What is reconstructed here is the observable content of the standard model: masses,
couplings, mixings, and precision numbers. What is not reconstructed is the full dynamical
Lagrangian/vertex structure; that remains hosted and is treated as a boundary condition of
the monograph, not as a derived output.

\section{Fermion Sector}
\label{c11:sec:fermion-sector}

\begin{theorem}[The fermion sectors are spectral access reads]
\label{c11:thm:fermion-access}
The \emph{fermion sector} of the theory --- the set of fermion families and hierarchies ---
is a spectral access read: each sector couples to the spectrum through a distinct moment of
the multiplicity distribution, the neutral, full, doublet-occupancy, and octave-occupation
accesses being the half, first, second, and octave moments.
\end{theorem}

\begin{proof}
The mode-access origin \cite{c11:QG150} and the effective-access program \cite{c11:QG157}
establish that each sector couples through a distinct moment of the multiplicity
distribution: the neutral access is the half-moment $\Sigma\sqrt{m} = 64.08$, the full
access the first moment $\Sigma m = 95$, the doublet-occupancy access the second
moment $\Sigma m^2 = 229$, and the octave-occupation access
$\mathrm{occMom} = 1900.25$. The sector exponents follow from these access counts
\cite{c11:QG149}, and the mass ratios are functions of the exponents; since the counts are
themselves derived from the spectrum (Chapter~\ref{c06:chap:d96},
Corollary~\ref{c06:cor:moment-ladder}, Corollary~\ref{c06:cor:constants-derived}), the
fermion sectors and their mass hierarchy are spectral access reads.
\end{proof}

\section{Gauge Structure}
\label{c11:sec:gauge}

\begin{theorem}[The gauge structure is the D96 degree-12 observable sector]
\label{c11:thm:gauge-structure}
The gauge structure of the standard model, $U(1)\times SU(2)\times SU(3)$ --- the
$1+3+8$ gauge sector --- is reconstructed from the D96 sector structure: the degree-12 content of
the $C_{96}$ generator ring.
\end{theorem}

\begin{proof}
The gauge-origin audit \cite{c11:QG161} establishes that the $1+3+8$ gauge sector is the
degree-12 structure of the $C_{96}$ sector: the gauge generators are the D96 spectral
content, giving one $U(1)$, three $SU(2)$, and eight $SU(3)$ generators --- reconstructed as
an observable sector from the D96 spectrum, not imported, and introducing no new primitive.
Since the D96 spectrum is itself the derived output of the attractor
(Chapter~\ref{c06:chap:d96}, Theorem~\ref{c06:thm:d96-derived}), the gauge structure is
twice-emergent as an observable sector; the interaction dynamics remain hosted.
\end{proof}

\section{Coupling Constants}
\label{c11:sec:couplings}

\begin{theorem}[The couplings are spectral ratios]
\label{c11:thm:couplings}
The gauge couplings are derived as ratios of D96 spectral constants:
\begin{equation}
\label{c11:eq:couplings}
\frac{1}{\alpha_{\mathrm{em}}} = \Sigma m + \#d = 137,
\qquad
\alpha_{\mathrm{weak}} = \frac{3}{\Sigma m},
\qquad
\alpha_{\mathrm{strong}} = \frac{8}{\Sigma\sqrt{m}},
\qquad
\sin^2\theta_W = 0.2316.
\end{equation}
\end{theorem}

\begin{proof}
The coupling-origin audit \cite{c11:QG162} establishes the constructions: the fine-structure
inverse is the total mode count plus the doublet count, $1/\alpha_{\mathrm{em}} = 137$; the
weak and strong couplings are ratios of the spectral constants; the Weinberg angle is a
derived ratio. Each coupling is therefore a ratio of spectral constants ($\Sigma m$,
$\Sigma\sqrt{m}$, $\#d$) --- derived, not fitted: the constants are themselves derived
from the spectrum (Chapter~\ref{c06:chap:d96}, Corollary~\ref{c06:cor:constants-derived}),
and no fitted parameters or imported values enter.
\end{proof}

\section{Flavor Mixing}
\label{c11:sec:mixing}

\begin{theorem}[The CKM matrix is an octave-transition read]
\label{c11:thm:ckm}
The CKM matrix is derived as the quark mixing read of the spectral structure: the elements
are ratios of doublet, octave, and occupancy quantities.
\end{theorem}

\begin{proof}
The CKM-origin audit \cite{c11:QG165} establishes the constructions:
$V_{us} = \#d/(2\Sigma m)$, $V_{cb} = (\omega_0/\omega_2)^{\delta_d}$,
$V_{ub} = 2V_{cb}\cdot(\mathrm{occ}_0/\mathrm{occ}_2)$, with mean deviation $0.58\%$; the
CKM CP phase and the Jarlskog invariant follow from the chiral circulation \cite{c11:QG166}.
\end{proof}

\begin{theorem}[The PMNS matrix is a T3-only read]
\label{c11:thm:pmns}
The PMNS matrix is derived as the neutrino mixing read: the angles
$\theta_{12} = 33.35^\circ$, $\theta_{23} = 49.72^\circ$, $\theta_{13} = 8.34^\circ$, and
$\delta_\nu = 66.4^\circ$ are the T3-only access reads, with mean deviation $1.5\%$.
\end{theorem}

\begin{proof}
The PMNS-origin audit \cite{c11:QG167} establishes the angles and the CP phase as the
T3-only access reads of the D96 content; hence the PMNS matrix is a derived read, not an
imported input.
\end{proof}

Both the quark (CKM) and the lepton (PMNS) mixing matrices are therefore derived spectral
reads, with no imported mixing input.

\section{Higgs and Weak Scale}
\label{c11:sec:higgs}

\begin{theorem}[The weak scale and the Higgs mass]
\label{c11:thm:weak-higgs}
The weak scale and the Higgs mass are reconstructed from the spectral structure:
\begin{equation}
\label{c11:eq:higgs}
v = (\Sigma m+\#d)\ln(\mathrm{span}) = 254\ \mathrm{GeV},
\qquad
M_H = \sigma_{\mathrm{occ}}\cdot\frac{\mathrm{span}}{2} = 125.25\ \mathrm{GeV},
\end{equation}
with the $W$ and $Z$ masses following from the Weinberg projection.
\end{theorem}

\begin{proof}
The weak-boson-mass origin \cite{c11:QG168} establishes the weak scale
$v = (\Sigma m+\#d)\ln(\mathrm{span}) = 254$ GeV and the $W$ and $Z$ masses
($M_W = 80.1$, $M_Z = 91.4$ GeV) with $\rho = 1$ exactly; the Higgs-mass origin
\cite{c11:QG169} gives $M_H = \sigma_{\mathrm{occ}}\cdot(\mathrm{span}/2) = 125.25$ GeV,
matching observation within $0.003\%$; and the blind reconstruction \cite{c11:QG176}
confirms this mass from the pre-Higgs D96 content only, with no Higgs observable as input.
The D96 factors are dimensionless; the GeV unit is carried by the calibrated electroweak
scale $v$.
\end{proof}

The Higgs is therefore a collective scalar observable of the theory: its mass is a spectral
construction $M_H = \sigma_{\mathrm{occ}}\cdot(\mathrm{span}/2)$, whose constants are
derived from the spectrum (Chapter~\ref{c06:chap:d96},
Corollary~\ref{c06:cor:constants-derived}); it introduces no new primitive \cite{c11:QG169}.
Read strictly, the D96 factor is dimensionless and the GeV unit is inherited from the
calibrated electroweak mass scale. The Higgs sector observable is reconstructed; the Yukawa
and full interaction dynamics remain hosted.

\begin{lemma}[The neutrino masses are closed-form D96 laws]
\label{c11:lem:neutrino-masses}
The neutrino mass splittings are derived as closed-form D96 laws:
$\Delta m^2_{21} = (1/\Sigma\sqrt{m})^2/(\mathrm{span}/2)$ and
$\Delta m^2_{31} = \sin^2\theta_W/\Sigma m$.
\end{lemma}

\begin{proof}
The neutrino-mass origin \cite{c11:QG172} establishes the closed-form laws: the solar and
atmospheric splittings are ratios of the D96 spectral constants, giving
$\Delta m^2_{21} = 7.607\times10^{-5}$ eV$^2$ and
$\Delta m^2_{31} = 2.44\times10^{-3}$ eV$^2$, with $m_2 = 8.72$ meV,
$m_3 = 49.4$ meV, and $\Sigma m_\nu = 0.058$ eV. Hence the neutrino masses are closed-form
D96 laws, not oscillation fits.
\end{proof}

\section{Precision Predictions}
\label{c11:sec:precision}

\begin{theorem}[The oblique parameters]
\label{c11:thm:oblique}
The oblique parameters are derived from the spectral structure:
\begin{equation}
\label{c11:eq:oblique}
S = \frac{\mathrm{occ}_0}{\Sigma m} = \frac{4}{95} = 0.0421,
\qquad
T = 2S = \frac{8}{95} = 0.0842,
\qquad
U = 0,
\end{equation}
with $T = 2S$ exactly and $U = 0$ via the exact tree-level $\rho = 1$.
\end{theorem}

\begin{proof}
The oblique-origin audit \cite{c11:QG180} establishes the constructions: $S$ is the lightest-
octave fraction of the spectrum, $T = 2S$ exactly, and $U = 0$ via the exact SM tree-level
$\rho = 1$, matching the EW global fit within $5.3\%$.
\end{proof}

\begin{theorem}[The electron g-2]
\label{c11:thm:electron-g2}
The electron g-2 is derived from the spectral structure:
\begin{equation}
\label{c11:eq:electron-g2}
a_e = \frac{\alpha}{2\pi}\Bigl(1 - \Bigl(\frac{\mathrm{occ}_0}{\Sigma m}\Bigr)^2\Bigr)
= 1.159655\times10^{-3},
\qquad
\Delta a_e = \Bigl(\frac{\alpha}{2\pi}\Bigr)^3 \mathrm{span}^{1/4}
\Bigl(\frac{\mathrm{occ}_0}{\Sigma m}\Bigr)^3 = 1.86\times10^{-13},
\end{equation}
with the correction negative and the anomaly below $10^{-12}$.
\end{theorem}

\begin{proof}
The electron-g-2 origin \cite{c11:QG178} establishes the construction: the electron anomaly is
the Schwinger term modulated by the lightest-octave fraction, with the negative correction
$(1 - (\mathrm{occ}_0/\Sigma m)^2)$ and the anomaly below $10^{-12}$; the same D96
mechanism produces the muon g-2 \cite{c11:QG171}. The electron anomaly matches experiment
within $0.0003\%$.
\end{proof}

\begin{theorem}[The Majorana character and \texorpdfstring{$0\nu\beta\beta$}{0nubetabeta}]
\label{c11:thm:majorana}
The neutrino is Majorana: it is self-conjugate in its $T_3$-only channel, has no conserved
charge, and has a real mass matrix. The effective Majorana mass is
\begin{equation}
\label{c11:eq:mbb}
m_{\beta\beta} = \bigl|m_1c_{12}^2c_{13}^2 + m_2s_{12}^2c_{13}^2e^{i\alpha_2}
+ m_3s_{13}^2e^{-2i\delta_\nu}\bigr| = 2.02\times10^{-3}\ \mathrm{eV},
\end{equation}
within experimental limits and in reach of next-generation experiments.
\end{theorem}

\begin{proof}
The Majorana-origin audit \cite{c11:QG179} establishes the neutrino character: the $T_3$-only
channel is self-conjugate ($48/95$ modes), the unique $Q=0$ sector provides no conserved
charge, and the reflection automorphism gives a real mass matrix. From the D96 PMNS angles
and masses, the effective Majorana mass is $m_{\beta\beta} = 2.02\times10^{-3}$ eV, within
limits and in reach of next-generation experiments.
\end{proof}

The precision predictions --- $S,T,U$, the electron g-2, and the Majorana character ---
are therefore all derived from the D96 spectral structure; their experimental validation
remains open --- the current frontier, not a derivation gap (VALID001).

\begin{remark}[Hosted-SM caveat]
\label{c11:remark:hosted-sm}
Consistent with MONO007 \cite{c11:MONO007}, the standard-model \emph{Lagrangian} and its
dynamics are hosted --- a boundary, not a derivation claim of this chapter. Every
\emph{observable} presented here is derived from the D96 moments; the Lagrangian form
$L = -\frac{1}{4}F^a_{\mu\nu}F^{a\mu\nu} + i\bar{\psi}\gamma^\mu D_\mu\psi - m\bar{\psi}\psi$
is the actualization-flow action established in Chapter~\ref{c09:chap:qm}
(Lemma~\ref{c09:lem:qm-dynamics}), but the full SM dynamical content remains hosted. No stronger
claim is made.
\end{remark}

\section{Consequences}
\label{c11:sec:sm-consequences}

\begin{theorem}[The standard model is consistent with the canonical architecture]
\label{c11:thm:sm-consistent}
The standard model is consistent with the canonical architecture: every derived SM
observable is a read of the D96 spectral moments, with the hosted dynamics the single
referenced boundary. Dimensional observables are reported relative to calibration anchors,
not from the dimensionless D96 data alone.
\end{theorem}

\begin{proof}
The fermion sector (Theorem~\ref{c11:thm:fermion-access}), the gauge structure
(Theorem~\ref{c11:thm:gauge-structure}), the couplings (Theorem~\ref{c11:thm:couplings}), the flavor
mixing (Theorem~\ref{c11:thm:ckm}, Theorem~\ref{c11:thm:pmns}), the Higgs and weak scale
(Theorem~\ref{c11:thm:weak-higgs}), and the precision predictions
(Theorem~\ref{c11:thm:oblique}, Theorem~\ref{c11:thm:electron-g2}, Theorem~\ref{c11:thm:majorana})
are all derived from the D96 moments. The hosted SM Lagrangian is the single referenced
boundary (Remark~\ref{c11:remark:hosted-sm}).
\end{proof}

\section{Summary}
\label{c11:sec:sm-summary}

This chapter has reconstructed the standard-model observables: the fermion sectors are
spectral access reads (Theorem~\ref{c11:thm:fermion-access}); the $1+3+8$ gauge sector is
the D96 degree-12 observable sector (Theorem~\ref{c11:thm:gauge-structure}); the couplings
are spectral ratios (Theorem~\ref{c11:thm:couplings}); the CKM and PMNS matrices are
derived spectral reads (Theorem~\ref{c11:thm:ckm}, Theorem~\ref{c11:thm:pmns}); and the
weak scale, the Higgs mass, and the neutrino masses are spectral constructions
(Theorem~\ref{c11:thm:weak-higgs}, Lemma~\ref{c11:lem:neutrino-masses}). The precision
predictions --- the oblique parameters $S,T,U$, the electron g-2, and the Majorana
character --- are likewise reconstructed as observable outputs
(Theorem~\ref{c11:thm:oblique}, Theorem~\ref{c11:thm:electron-g2},
Theorem~\ref{c11:thm:majorana}), so the standard model is consistent with the canonical
architecture, with the hosted dynamics the single referenced boundary
(Theorem~\ref{c11:thm:sm-consistent}, Remark~\ref{c11:remark:hosted-sm}). Every SM
observable is therefore a spectral construction, twice-derived from the spectrum
(Chapter~\ref{c06:chap:d96}, Corollary~\ref{c06:cor:constants-derived}); the following
chapter derives cosmology from the same foundation.

\begin{thebibliography}{99}
\bibitem{c11:QG149} AT-QG Phase 149, \emph{Sector Exponents Origin}, Docs/Research.
\bibitem{c11:QG150} AT-QG Phase 150, \emph{Mode Access Origin}, Docs/Research.
\bibitem{c11:QG157} AT-QG Phase 157, \emph{Effective Access Counts}, Docs/Research.
\bibitem{c11:QG161} AT-QG Phase 161, \emph{Gauge Sector Origin}, Docs/Research.
\bibitem{c11:QG162} AT-QG Phase 162, \emph{Gauge Coupling Origin}, Docs/Research.
\bibitem{c11:QG165} AT-QG Phase 165, \emph{CKM Origin}, Docs/Research.
\bibitem{c11:QG166} AT-QG Phase 166, \emph{CKM CP Origin}, Docs/Research.
\bibitem{c11:QG167} AT-QG Phase 167, \emph{PMNS Origin}, Docs/Research.
\bibitem{c11:QG168} AT-QG Phase 168, \emph{Weak Boson Mass Origin}, Docs/Research.
\bibitem{c11:QG169} AT-QG Phase 169, \emph{Higgs Mass Origin}, Docs/Research.
\bibitem{c11:QG171} AT-QG Phase 171, \emph{Muon g-2 Origin}, Docs/Research.
\bibitem{c11:QG172} AT-QG Phase 172, \emph{Neutrino Mass Law}, Docs/Research.
\bibitem{c11:QG176} AT-QG Phase 176, \emph{Higgs Blind Reconstruction}, Docs/Research.
\bibitem{c11:QG178} AT-QG Phase 178, \emph{Electron g-2 Origin}, Docs/Research.
\bibitem{c11:QG179} AT-QG Phase 179, \emph{Majorana Origin}, Docs/Research.
\bibitem{c11:QG180} AT-QG Phase 180, \emph{Oblique Parameters Origin}, Docs/Research.
\bibitem{c11:MONO007} AT-MONO007, \emph{Referee-Readiness Resolution}, Docs/Zenodo/Monograph.
\end{thebibliography}
 after the preamble
%         that defines the theorem environments (or include via the master file).
% ======================================================================

% ---- Theorem environments (define in the master preamble if not present) ----
% \newtheorem{definition}{Definition}[chapter]
% \newtheorem{lemma}[definition]{Lemma}
% \newtheorem{theorem}[definition]{Theorem}
% \newtheorem{corollary}[definition]{Corollary}
% \newtheorem{remark}[definition]{Remark}

\chapter{The Standard Model}
\label{c11:chap:sm}

\section{Introduction}
\label{c11:sec:sm-introduction}

The preceding chapters established the foundation, the spectrum, the operator basis, and
the quantum and gravitational structure of The Actualization Theory. This chapter
reconstructs the \emph{standard-model observables} from the same foundation --- the fermion
sectors, the gauge structure, the couplings, the flavor mixing, the weak scale and the
Higgs, and a set of precision predictions --- as derived reads of the D96 spectral moments
\cite{c11:QG149,c11:QG150,c11:QG157}.

The central results are: the fermion sectors are reads of the spectral access counts
\cite{c11:QG149,c11:QG150}; the gauge structure $1+3+8$ is the degree-12 structure of the D96
sector \cite{c11:QG161}; the couplings are ratios of spectral constants \cite{c11:QG162}; the
flavor mixing matrices are octave-transition reads \cite{c11:QG165,c11:QG167}; the weak scale and
the Higgs mass are spectral constructions \cite{c11:QG168,c11:QG169}; the neutrino masses are
closed-form D96 laws \cite{c11:QG172}; and a set of precision predictions --- the oblique
parameters $S,T,U$, the electron g-2, and the Majorana character --- follow
\cite{c11:QG176,c11:QG178,c11:QG179,c11:QG180}.

Throughout, the hosted-standard-model caveat is retained: every \emph{observable} is
derived from the D96 moments, while the standard-model \emph{Lagrangian} and its dynamics
are hosted --- a boundary, not a derivation gap, per the referee-safe architecture of
MONO007 \cite{c11:MONO007}.

The dimensionless D96 laws are then calibrated to measured anchors: the weak scale $v$
for the weak/Higgs sector, the electron mass $m_e$ for the absolute quark masses, and the
standard unit conversion when a result is reported in SI or GeV units. That calibration
step fixes the physical units; it is not a new dynamical input.

\subsection{Derived Observables versus Hosted Dynamics}
\label{c11:sec:derived-vs-hosted}

\begin{quote}
\textbf{Derived:} primitives $\rightarrow$ actualization $\rightarrow$ spectrum $\rightarrow$
observables.\\
\textbf{Hosted:} SM Lagrangian, interaction dynamics, and external unit conventions.
\end{quote}

What is reconstructed here is the observable content of the standard model: masses,
couplings, mixings, and precision numbers. What is not reconstructed is the full dynamical
Lagrangian/vertex structure; that remains hosted and is treated as a boundary condition of
the monograph, not as a derived output.

\section{Fermion Sector}
\label{c11:sec:fermion-sector}

\begin{theorem}[The fermion sectors are spectral access reads]
\label{c11:thm:fermion-access}
The \emph{fermion sector} of the theory --- the set of fermion families and hierarchies ---
is a spectral access read: each sector couples to the spectrum through a distinct moment of
the multiplicity distribution, the neutral, full, doublet-occupancy, and octave-occupation
accesses being the half, first, second, and octave moments.
\end{theorem}

\begin{proof}
The mode-access origin \cite{c11:QG150} and the effective-access program \cite{c11:QG157}
establish that each sector couples through a distinct moment of the multiplicity
distribution: the neutral access is the half-moment $\Sigma\sqrt{m} = 64.08$, the full
access the first moment $\Sigma m = 95$, the doublet-occupancy access the second
moment $\Sigma m^2 = 229$, and the octave-occupation access
$\mathrm{occMom} = 1900.25$. The sector exponents follow from these access counts
\cite{c11:QG149}, and the mass ratios are functions of the exponents; since the counts are
themselves derived from the spectrum (Chapter~\ref{c06:chap:d96},
Corollary~\ref{c06:cor:moment-ladder}, Corollary~\ref{c06:cor:constants-derived}), the
fermion sectors and their mass hierarchy are spectral access reads.
\end{proof}

\section{Gauge Structure}
\label{c11:sec:gauge}

\begin{theorem}[The gauge structure is the D96 degree-12 observable sector]
\label{c11:thm:gauge-structure}
The gauge structure of the standard model, $U(1)\times SU(2)\times SU(3)$ --- the
$1+3+8$ gauge sector --- is reconstructed from the D96 sector structure: the degree-12 content of
the $C_{96}$ generator ring.
\end{theorem}

\begin{proof}
The gauge-origin audit \cite{c11:QG161} establishes that the $1+3+8$ gauge sector is the
degree-12 structure of the $C_{96}$ sector: the gauge generators are the D96 spectral
content, giving one $U(1)$, three $SU(2)$, and eight $SU(3)$ generators --- reconstructed as
an observable sector from the D96 spectrum, not imported, and introducing no new primitive.
Since the D96 spectrum is itself the derived output of the attractor
(Chapter~\ref{c06:chap:d96}, Theorem~\ref{c06:thm:d96-derived}), the gauge structure is
twice-emergent as an observable sector; the interaction dynamics remain hosted.
\end{proof}

\section{Coupling Constants}
\label{c11:sec:couplings}

\begin{theorem}[The couplings are spectral ratios]
\label{c11:thm:couplings}
The gauge couplings are derived as ratios of D96 spectral constants:
\begin{equation}
\label{c11:eq:couplings}
\frac{1}{\alpha_{\mathrm{em}}} = \Sigma m + \#d = 137,
\qquad
\alpha_{\mathrm{weak}} = \frac{3}{\Sigma m},
\qquad
\alpha_{\mathrm{strong}} = \frac{8}{\Sigma\sqrt{m}},
\qquad
\sin^2\theta_W = 0.2316.
\end{equation}
\end{theorem}

\begin{proof}
The coupling-origin audit \cite{c11:QG162} establishes the constructions: the fine-structure
inverse is the total mode count plus the doublet count, $1/\alpha_{\mathrm{em}} = 137$; the
weak and strong couplings are ratios of the spectral constants; the Weinberg angle is a
derived ratio. Each coupling is therefore a ratio of spectral constants ($\Sigma m$,
$\Sigma\sqrt{m}$, $\#d$) --- derived, not fitted: the constants are themselves derived
from the spectrum (Chapter~\ref{c06:chap:d96}, Corollary~\ref{c06:cor:constants-derived}),
and no fitted parameters or imported values enter.
\end{proof}

\section{Flavor Mixing}
\label{c11:sec:mixing}

\begin{theorem}[The CKM matrix is an octave-transition read]
\label{c11:thm:ckm}
The CKM matrix is derived as the quark mixing read of the spectral structure: the elements
are ratios of doublet, octave, and occupancy quantities.
\end{theorem}

\begin{proof}
The CKM-origin audit \cite{c11:QG165} establishes the constructions:
$V_{us} = \#d/(2\Sigma m)$, $V_{cb} = (\omega_0/\omega_2)^{\delta_d}$,
$V_{ub} = 2V_{cb}\cdot(\mathrm{occ}_0/\mathrm{occ}_2)$, with mean deviation $0.58\%$; the
CKM CP phase and the Jarlskog invariant follow from the chiral circulation \cite{c11:QG166}.
\end{proof}

\begin{theorem}[The PMNS matrix is a T3-only read]
\label{c11:thm:pmns}
The PMNS matrix is derived as the neutrino mixing read: the angles
$\theta_{12} = 33.35^\circ$, $\theta_{23} = 49.72^\circ$, $\theta_{13} = 8.34^\circ$, and
$\delta_\nu = 66.4^\circ$ are the T3-only access reads, with mean deviation $1.5\%$.
\end{theorem}

\begin{proof}
The PMNS-origin audit \cite{c11:QG167} establishes the angles and the CP phase as the
T3-only access reads of the D96 content; hence the PMNS matrix is a derived read, not an
imported input.
\end{proof}

Both the quark (CKM) and the lepton (PMNS) mixing matrices are therefore derived spectral
reads, with no imported mixing input.

\section{Higgs and Weak Scale}
\label{c11:sec:higgs}

\begin{theorem}[The weak scale and the Higgs mass]
\label{c11:thm:weak-higgs}
The weak scale and the Higgs mass are reconstructed from the spectral structure:
\begin{equation}
\label{c11:eq:higgs}
v = (\Sigma m+\#d)\ln(\mathrm{span}) = 254\ \mathrm{GeV},
\qquad
M_H = \sigma_{\mathrm{occ}}\cdot\frac{\mathrm{span}}{2} = 125.25\ \mathrm{GeV},
\end{equation}
with the $W$ and $Z$ masses following from the Weinberg projection.
\end{theorem}

\begin{proof}
The weak-boson-mass origin \cite{c11:QG168} establishes the weak scale
$v = (\Sigma m+\#d)\ln(\mathrm{span}) = 254$ GeV and the $W$ and $Z$ masses
($M_W = 80.1$, $M_Z = 91.4$ GeV) with $\rho = 1$ exactly; the Higgs-mass origin
\cite{c11:QG169} gives $M_H = \sigma_{\mathrm{occ}}\cdot(\mathrm{span}/2) = 125.25$ GeV,
matching observation within $0.003\%$; and the blind reconstruction \cite{c11:QG176}
confirms this mass from the pre-Higgs D96 content only, with no Higgs observable as input.
The D96 factors are dimensionless; the GeV unit is carried by the calibrated electroweak
scale $v$.
\end{proof}

The Higgs is therefore a collective scalar observable of the theory: its mass is a spectral
construction $M_H = \sigma_{\mathrm{occ}}\cdot(\mathrm{span}/2)$, whose constants are
derived from the spectrum (Chapter~\ref{c06:chap:d96},
Corollary~\ref{c06:cor:constants-derived}); it introduces no new primitive \cite{c11:QG169}.
Read strictly, the D96 factor is dimensionless and the GeV unit is inherited from the
calibrated electroweak mass scale. The Higgs sector observable is reconstructed; the Yukawa
and full interaction dynamics remain hosted.

\begin{lemma}[The neutrino masses are closed-form D96 laws]
\label{c11:lem:neutrino-masses}
The neutrino mass splittings are derived as closed-form D96 laws:
$\Delta m^2_{21} = (1/\Sigma\sqrt{m})^2/(\mathrm{span}/2)$ and
$\Delta m^2_{31} = \sin^2\theta_W/\Sigma m$.
\end{lemma}

\begin{proof}
The neutrino-mass origin \cite{c11:QG172} establishes the closed-form laws: the solar and
atmospheric splittings are ratios of the D96 spectral constants, giving
$\Delta m^2_{21} = 7.607\times10^{-5}$ eV$^2$ and
$\Delta m^2_{31} = 2.44\times10^{-3}$ eV$^2$, with $m_2 = 8.72$ meV,
$m_3 = 49.4$ meV, and $\Sigma m_\nu = 0.058$ eV. Hence the neutrino masses are closed-form
D96 laws, not oscillation fits.
\end{proof}

\section{Precision Predictions}
\label{c11:sec:precision}

\begin{theorem}[The oblique parameters]
\label{c11:thm:oblique}
The oblique parameters are derived from the spectral structure:
\begin{equation}
\label{c11:eq:oblique}
S = \frac{\mathrm{occ}_0}{\Sigma m} = \frac{4}{95} = 0.0421,
\qquad
T = 2S = \frac{8}{95} = 0.0842,
\qquad
U = 0,
\end{equation}
with $T = 2S$ exactly and $U = 0$ via the exact tree-level $\rho = 1$.
\end{theorem}

\begin{proof}
The oblique-origin audit \cite{c11:QG180} establishes the constructions: $S$ is the lightest-
octave fraction of the spectrum, $T = 2S$ exactly, and $U = 0$ via the exact SM tree-level
$\rho = 1$, matching the EW global fit within $5.3\%$.
\end{proof}

\begin{theorem}[The electron g-2]
\label{c11:thm:electron-g2}
The electron g-2 is derived from the spectral structure:
\begin{equation}
\label{c11:eq:electron-g2}
a_e = \frac{\alpha}{2\pi}\Bigl(1 - \Bigl(\frac{\mathrm{occ}_0}{\Sigma m}\Bigr)^2\Bigr)
= 1.159655\times10^{-3},
\qquad
\Delta a_e = \Bigl(\frac{\alpha}{2\pi}\Bigr)^3 \mathrm{span}^{1/4}
\Bigl(\frac{\mathrm{occ}_0}{\Sigma m}\Bigr)^3 = 1.86\times10^{-13},
\end{equation}
with the correction negative and the anomaly below $10^{-12}$.
\end{theorem}

\begin{proof}
The electron-g-2 origin \cite{c11:QG178} establishes the construction: the electron anomaly is
the Schwinger term modulated by the lightest-octave fraction, with the negative correction
$(1 - (\mathrm{occ}_0/\Sigma m)^2)$ and the anomaly below $10^{-12}$; the same D96
mechanism produces the muon g-2 \cite{c11:QG171}. The electron anomaly matches experiment
within $0.0003\%$.
\end{proof}

\begin{theorem}[The Majorana character and \texorpdfstring{$0\nu\beta\beta$}{0nubetabeta}]
\label{c11:thm:majorana}
The neutrino is Majorana: it is self-conjugate in its $T_3$-only channel, has no conserved
charge, and has a real mass matrix. The effective Majorana mass is
\begin{equation}
\label{c11:eq:mbb}
m_{\beta\beta} = \bigl|m_1c_{12}^2c_{13}^2 + m_2s_{12}^2c_{13}^2e^{i\alpha_2}
+ m_3s_{13}^2e^{-2i\delta_\nu}\bigr| = 2.02\times10^{-3}\ \mathrm{eV},
\end{equation}
within experimental limits and in reach of next-generation experiments.
\end{theorem}

\begin{proof}
The Majorana-origin audit \cite{c11:QG179} establishes the neutrino character: the $T_3$-only
channel is self-conjugate ($48/95$ modes), the unique $Q=0$ sector provides no conserved
charge, and the reflection automorphism gives a real mass matrix. From the D96 PMNS angles
and masses, the effective Majorana mass is $m_{\beta\beta} = 2.02\times10^{-3}$ eV, within
limits and in reach of next-generation experiments.
\end{proof}

The precision predictions --- $S,T,U$, the electron g-2, and the Majorana character ---
are therefore all derived from the D96 spectral structure; their experimental validation
remains open --- the current frontier, not a derivation gap (VALID001).

\begin{remark}[Hosted-SM caveat]
\label{c11:remark:hosted-sm}
Consistent with MONO007 \cite{c11:MONO007}, the standard-model \emph{Lagrangian} and its
dynamics are hosted --- a boundary, not a derivation claim of this chapter. Every
\emph{observable} presented here is derived from the D96 moments; the Lagrangian form
$L = -\frac{1}{4}F^a_{\mu\nu}F^{a\mu\nu} + i\bar{\psi}\gamma^\mu D_\mu\psi - m\bar{\psi}\psi$
is the actualization-flow action established in Chapter~\ref{c09:chap:qm}
(Lemma~\ref{c09:lem:qm-dynamics}), but the full SM dynamical content remains hosted. No stronger
claim is made.
\end{remark}

\section{Consequences}
\label{c11:sec:sm-consequences}

\begin{theorem}[The standard model is consistent with the canonical architecture]
\label{c11:thm:sm-consistent}
The standard model is consistent with the canonical architecture: every derived SM
observable is a read of the D96 spectral moments, with the hosted dynamics the single
referenced boundary. Dimensional observables are reported relative to calibration anchors,
not from the dimensionless D96 data alone.
\end{theorem}

\begin{proof}
The fermion sector (Theorem~\ref{c11:thm:fermion-access}), the gauge structure
(Theorem~\ref{c11:thm:gauge-structure}), the couplings (Theorem~\ref{c11:thm:couplings}), the flavor
mixing (Theorem~\ref{c11:thm:ckm}, Theorem~\ref{c11:thm:pmns}), the Higgs and weak scale
(Theorem~\ref{c11:thm:weak-higgs}), and the precision predictions
(Theorem~\ref{c11:thm:oblique}, Theorem~\ref{c11:thm:electron-g2}, Theorem~\ref{c11:thm:majorana})
are all derived from the D96 moments. The hosted SM Lagrangian is the single referenced
boundary (Remark~\ref{c11:remark:hosted-sm}).
\end{proof}

\section{Summary}
\label{c11:sec:sm-summary}

This chapter has reconstructed the standard-model observables: the fermion sectors are
spectral access reads (Theorem~\ref{c11:thm:fermion-access}); the $1+3+8$ gauge sector is
the D96 degree-12 observable sector (Theorem~\ref{c11:thm:gauge-structure}); the couplings
are spectral ratios (Theorem~\ref{c11:thm:couplings}); the CKM and PMNS matrices are
derived spectral reads (Theorem~\ref{c11:thm:ckm}, Theorem~\ref{c11:thm:pmns}); and the
weak scale, the Higgs mass, and the neutrino masses are spectral constructions
(Theorem~\ref{c11:thm:weak-higgs}, Lemma~\ref{c11:lem:neutrino-masses}). The precision
predictions --- the oblique parameters $S,T,U$, the electron g-2, and the Majorana
character --- are likewise reconstructed as observable outputs
(Theorem~\ref{c11:thm:oblique}, Theorem~\ref{c11:thm:electron-g2},
Theorem~\ref{c11:thm:majorana}), so the standard model is consistent with the canonical
architecture, with the hosted dynamics the single referenced boundary
(Theorem~\ref{c11:thm:sm-consistent}, Remark~\ref{c11:remark:hosted-sm}). Every SM
observable is therefore a spectral construction, twice-derived from the spectrum
(Chapter~\ref{c06:chap:d96}, Corollary~\ref{c06:cor:constants-derived}); the following
chapter derives cosmology from the same foundation.

\begin{thebibliography}{99}
\bibitem{c11:QG149} AT-QG Phase 149, \emph{Sector Exponents Origin}, Docs/Research.
\bibitem{c11:QG150} AT-QG Phase 150, \emph{Mode Access Origin}, Docs/Research.
\bibitem{c11:QG157} AT-QG Phase 157, \emph{Effective Access Counts}, Docs/Research.
\bibitem{c11:QG161} AT-QG Phase 161, \emph{Gauge Sector Origin}, Docs/Research.
\bibitem{c11:QG162} AT-QG Phase 162, \emph{Gauge Coupling Origin}, Docs/Research.
\bibitem{c11:QG165} AT-QG Phase 165, \emph{CKM Origin}, Docs/Research.
\bibitem{c11:QG166} AT-QG Phase 166, \emph{CKM CP Origin}, Docs/Research.
\bibitem{c11:QG167} AT-QG Phase 167, \emph{PMNS Origin}, Docs/Research.
\bibitem{c11:QG168} AT-QG Phase 168, \emph{Weak Boson Mass Origin}, Docs/Research.
\bibitem{c11:QG169} AT-QG Phase 169, \emph{Higgs Mass Origin}, Docs/Research.
\bibitem{c11:QG171} AT-QG Phase 171, \emph{Muon g-2 Origin}, Docs/Research.
\bibitem{c11:QG172} AT-QG Phase 172, \emph{Neutrino Mass Law}, Docs/Research.
\bibitem{c11:QG176} AT-QG Phase 176, \emph{Higgs Blind Reconstruction}, Docs/Research.
\bibitem{c11:QG178} AT-QG Phase 178, \emph{Electron g-2 Origin}, Docs/Research.
\bibitem{c11:QG179} AT-QG Phase 179, \emph{Majorana Origin}, Docs/Research.
\bibitem{c11:QG180} AT-QG Phase 180, \emph{Oblique Parameters Origin}, Docs/Research.
\bibitem{c11:MONO007} AT-MONO007, \emph{Referee-Readiness Resolution}, Docs/Zenodo/Monograph.
\end{thebibliography}
 after the preamble
%         that defines the theorem environments (or include via the master file).
% ======================================================================

% ---- Theorem environments (define in the master preamble if not present) ----
% \newtheorem{definition}{Definition}[chapter]
% \newtheorem{lemma}[definition]{Lemma}
% \newtheorem{theorem}[definition]{Theorem}
% \newtheorem{corollary}[definition]{Corollary}
% \newtheorem{remark}[definition]{Remark}

\chapter{The Standard Model}
\label{c11:chap:sm}

\section{Introduction}
\label{c11:sec:sm-introduction}

The preceding chapters established the foundation, the spectrum, the operator basis, and
the quantum and gravitational structure of The Actualization Theory. This chapter
reconstructs the \emph{standard-model observables} from the same foundation --- the fermion
sectors, the gauge structure, the couplings, the flavor mixing, the weak scale and the
Higgs, and a set of precision predictions --- as derived reads of the D96 spectral moments
\cite{c11:QG149,c11:QG150,c11:QG157}.

The central results are: the fermion sectors are reads of the spectral access counts
\cite{c11:QG149,c11:QG150}; the gauge structure $1+3+8$ is the degree-12 structure of the D96
sector \cite{c11:QG161}; the couplings are ratios of spectral constants \cite{c11:QG162}; the
flavor mixing matrices are octave-transition reads \cite{c11:QG165,c11:QG167}; the weak scale and
the Higgs mass are spectral constructions \cite{c11:QG168,c11:QG169}; the neutrino masses are
closed-form D96 laws \cite{c11:QG172}; and a set of precision predictions --- the oblique
parameters $S,T,U$, the electron g-2, and the Majorana character --- follow
\cite{c11:QG176,c11:QG178,c11:QG179,c11:QG180}.

Throughout, the hosted-standard-model caveat is retained: every \emph{observable} is
derived from the D96 moments, while the standard-model \emph{Lagrangian} and its dynamics
are hosted --- a boundary, not a derivation gap, per the referee-safe architecture of
MONO007 \cite{c11:MONO007}.

The dimensionless D96 laws are then calibrated to measured anchors: the weak scale $v$
for the weak/Higgs sector, the electron mass $m_e$ for the absolute quark masses, and the
standard unit conversion when a result is reported in SI or GeV units. That calibration
step fixes the physical units; it is not a new dynamical input.

\subsection{Derived Observables versus Hosted Dynamics}
\label{c11:sec:derived-vs-hosted}

\begin{quote}
\textbf{Derived:} primitives $\rightarrow$ actualization $\rightarrow$ spectrum $\rightarrow$
observables.\\
\textbf{Hosted:} SM Lagrangian, interaction dynamics, and external unit conventions.
\end{quote}

What is reconstructed here is the observable content of the standard model: masses,
couplings, mixings, and precision numbers. What is not reconstructed is the full dynamical
Lagrangian/vertex structure; that remains hosted and is treated as a boundary condition of
the monograph, not as a derived output.

\section{Fermion Sector}
\label{c11:sec:fermion-sector}

\begin{theorem}[The fermion sectors are spectral access reads]
\label{c11:thm:fermion-access}
The \emph{fermion sector} of the theory --- the set of fermion families and hierarchies ---
is a spectral access read: each sector couples to the spectrum through a distinct moment of
the multiplicity distribution, the neutral, full, doublet-occupancy, and octave-occupation
accesses being the half, first, second, and octave moments.
\end{theorem}

\begin{proof}
The mode-access origin \cite{c11:QG150} and the effective-access program \cite{c11:QG157}
establish that each sector couples through a distinct moment of the multiplicity
distribution: the neutral access is the half-moment $\Sigma\sqrt{m} = 64.08$, the full
access the first moment $\Sigma m = 95$, the doublet-occupancy access the second
moment $\Sigma m^2 = 229$, and the octave-occupation access
$\mathrm{occMom} = 1900.25$. The sector exponents follow from these access counts
\cite{c11:QG149}, and the mass ratios are functions of the exponents; since the counts are
themselves derived from the spectrum (Chapter~\ref{c06:chap:d96},
Corollary~\ref{c06:cor:moment-ladder}, Corollary~\ref{c06:cor:constants-derived}), the
fermion sectors and their mass hierarchy are spectral access reads.
\end{proof}

\section{Gauge Structure}
\label{c11:sec:gauge}

\begin{theorem}[The gauge structure is the D96 degree-12 observable sector]
\label{c11:thm:gauge-structure}
The gauge structure of the standard model, $U(1)\times SU(2)\times SU(3)$ --- the
$1+3+8$ gauge sector --- is reconstructed from the D96 sector structure: the degree-12 content of
the $C_{96}$ generator ring.
\end{theorem}

\begin{proof}
The gauge-origin audit \cite{c11:QG161} establishes that the $1+3+8$ gauge sector is the
degree-12 structure of the $C_{96}$ sector: the gauge generators are the D96 spectral
content, giving one $U(1)$, three $SU(2)$, and eight $SU(3)$ generators --- reconstructed as
an observable sector from the D96 spectrum, not imported, and introducing no new primitive.
Since the D96 spectrum is itself the derived output of the attractor
(Chapter~\ref{c06:chap:d96}, Theorem~\ref{c06:thm:d96-derived}), the gauge structure is
twice-emergent as an observable sector; the interaction dynamics remain hosted.
\end{proof}

\section{Coupling Constants}
\label{c11:sec:couplings}

\begin{theorem}[The couplings are spectral ratios]
\label{c11:thm:couplings}
The gauge couplings are derived as ratios of D96 spectral constants:
\begin{equation}
\label{c11:eq:couplings}
\frac{1}{\alpha_{\mathrm{em}}} = \Sigma m + \#d = 137,
\qquad
\alpha_{\mathrm{weak}} = \frac{3}{\Sigma m},
\qquad
\alpha_{\mathrm{strong}} = \frac{8}{\Sigma\sqrt{m}},
\qquad
\sin^2\theta_W = 0.2316.
\end{equation}
\end{theorem}

\begin{proof}
The coupling-origin audit \cite{c11:QG162} establishes the constructions: the fine-structure
inverse is the total mode count plus the doublet count, $1/\alpha_{\mathrm{em}} = 137$; the
weak and strong couplings are ratios of the spectral constants; the Weinberg angle is a
derived ratio. Each coupling is therefore a ratio of spectral constants ($\Sigma m$,
$\Sigma\sqrt{m}$, $\#d$) --- derived, not fitted: the constants are themselves derived
from the spectrum (Chapter~\ref{c06:chap:d96}, Corollary~\ref{c06:cor:constants-derived}),
and no fitted parameters or imported values enter.
\end{proof}

\section{Flavor Mixing}
\label{c11:sec:mixing}

\begin{theorem}[The CKM matrix is an octave-transition read]
\label{c11:thm:ckm}
The CKM matrix is derived as the quark mixing read of the spectral structure: the elements
are ratios of doublet, octave, and occupancy quantities.
\end{theorem}

\begin{proof}
The CKM-origin audit \cite{c11:QG165} establishes the constructions:
$V_{us} = \#d/(2\Sigma m)$, $V_{cb} = (\omega_0/\omega_2)^{\delta_d}$,
$V_{ub} = 2V_{cb}\cdot(\mathrm{occ}_0/\mathrm{occ}_2)$, with mean deviation $0.58\%$; the
CKM CP phase and the Jarlskog invariant follow from the chiral circulation \cite{c11:QG166}.
\end{proof}

\begin{theorem}[The PMNS matrix is a T3-only read]
\label{c11:thm:pmns}
The PMNS matrix is derived as the neutrino mixing read: the angles
$\theta_{12} = 33.35^\circ$, $\theta_{23} = 49.72^\circ$, $\theta_{13} = 8.34^\circ$, and
$\delta_\nu = 66.4^\circ$ are the T3-only access reads, with mean deviation $1.5\%$.
\end{theorem}

\begin{proof}
The PMNS-origin audit \cite{c11:QG167} establishes the angles and the CP phase as the
T3-only access reads of the D96 content; hence the PMNS matrix is a derived read, not an
imported input.
\end{proof}

Both the quark (CKM) and the lepton (PMNS) mixing matrices are therefore derived spectral
reads, with no imported mixing input.

\section{Higgs and Weak Scale}
\label{c11:sec:higgs}

\begin{theorem}[The weak scale and the Higgs mass]
\label{c11:thm:weak-higgs}
The weak scale and the Higgs mass are reconstructed from the spectral structure:
\begin{equation}
\label{c11:eq:higgs}
v = (\Sigma m+\#d)\ln(\mathrm{span}) = 254\ \mathrm{GeV},
\qquad
M_H = \sigma_{\mathrm{occ}}\cdot\frac{\mathrm{span}}{2} = 125.25\ \mathrm{GeV},
\end{equation}
with the $W$ and $Z$ masses following from the Weinberg projection.
\end{theorem}

\begin{proof}
The weak-boson-mass origin \cite{c11:QG168} establishes the weak scale
$v = (\Sigma m+\#d)\ln(\mathrm{span}) = 254$ GeV and the $W$ and $Z$ masses
($M_W = 80.1$, $M_Z = 91.4$ GeV) with $\rho = 1$ exactly; the Higgs-mass origin
\cite{c11:QG169} gives $M_H = \sigma_{\mathrm{occ}}\cdot(\mathrm{span}/2) = 125.25$ GeV,
matching observation within $0.003\%$; and the blind reconstruction \cite{c11:QG176}
confirms this mass from the pre-Higgs D96 content only, with no Higgs observable as input.
The D96 factors are dimensionless; the GeV unit is carried by the calibrated electroweak
scale $v$.
\end{proof}

The Higgs is therefore a collective scalar observable of the theory: its mass is a spectral
construction $M_H = \sigma_{\mathrm{occ}}\cdot(\mathrm{span}/2)$, whose constants are
derived from the spectrum (Chapter~\ref{c06:chap:d96},
Corollary~\ref{c06:cor:constants-derived}); it introduces no new primitive \cite{c11:QG169}.
Read strictly, the D96 factor is dimensionless and the GeV unit is inherited from the
calibrated electroweak mass scale. The Higgs sector observable is reconstructed; the Yukawa
and full interaction dynamics remain hosted.

\begin{lemma}[The neutrino masses are closed-form D96 laws]
\label{c11:lem:neutrino-masses}
The neutrino mass splittings are derived as closed-form D96 laws:
$\Delta m^2_{21} = (1/\Sigma\sqrt{m})^2/(\mathrm{span}/2)$ and
$\Delta m^2_{31} = \sin^2\theta_W/\Sigma m$.
\end{lemma}

\begin{proof}
The neutrino-mass origin \cite{c11:QG172} establishes the closed-form laws: the solar and
atmospheric splittings are ratios of the D96 spectral constants, giving
$\Delta m^2_{21} = 7.607\times10^{-5}$ eV$^2$ and
$\Delta m^2_{31} = 2.44\times10^{-3}$ eV$^2$, with $m_2 = 8.72$ meV,
$m_3 = 49.4$ meV, and $\Sigma m_\nu = 0.058$ eV. Hence the neutrino masses are closed-form
D96 laws, not oscillation fits.
\end{proof}

\section{Precision Predictions}
\label{c11:sec:precision}

\begin{theorem}[The oblique parameters]
\label{c11:thm:oblique}
The oblique parameters are derived from the spectral structure:
\begin{equation}
\label{c11:eq:oblique}
S = \frac{\mathrm{occ}_0}{\Sigma m} = \frac{4}{95} = 0.0421,
\qquad
T = 2S = \frac{8}{95} = 0.0842,
\qquad
U = 0,
\end{equation}
with $T = 2S$ exactly and $U = 0$ via the exact tree-level $\rho = 1$.
\end{theorem}

\begin{proof}
The oblique-origin audit \cite{c11:QG180} establishes the constructions: $S$ is the lightest-
octave fraction of the spectrum, $T = 2S$ exactly, and $U = 0$ via the exact SM tree-level
$\rho = 1$, matching the EW global fit within $5.3\%$.
\end{proof}

\begin{theorem}[The electron g-2]
\label{c11:thm:electron-g2}
The electron g-2 is derived from the spectral structure:
\begin{equation}
\label{c11:eq:electron-g2}
a_e = \frac{\alpha}{2\pi}\Bigl(1 - \Bigl(\frac{\mathrm{occ}_0}{\Sigma m}\Bigr)^2\Bigr)
= 1.159655\times10^{-3},
\qquad
\Delta a_e = \Bigl(\frac{\alpha}{2\pi}\Bigr)^3 \mathrm{span}^{1/4}
\Bigl(\frac{\mathrm{occ}_0}{\Sigma m}\Bigr)^3 = 1.86\times10^{-13},
\end{equation}
with the correction negative and the anomaly below $10^{-12}$.
\end{theorem}

\begin{proof}
The electron-g-2 origin \cite{c11:QG178} establishes the construction: the electron anomaly is
the Schwinger term modulated by the lightest-octave fraction, with the negative correction
$(1 - (\mathrm{occ}_0/\Sigma m)^2)$ and the anomaly below $10^{-12}$; the same D96
mechanism produces the muon g-2 \cite{c11:QG171}. The electron anomaly matches experiment
within $0.0003\%$.
\end{proof}

\begin{theorem}[The Majorana character and \texorpdfstring{$0\nu\beta\beta$}{0nubetabeta}]
\label{c11:thm:majorana}
The neutrino is Majorana: it is self-conjugate in its $T_3$-only channel, has no conserved
charge, and has a real mass matrix. The effective Majorana mass is
\begin{equation}
\label{c11:eq:mbb}
m_{\beta\beta} = \bigl|m_1c_{12}^2c_{13}^2 + m_2s_{12}^2c_{13}^2e^{i\alpha_2}
+ m_3s_{13}^2e^{-2i\delta_\nu}\bigr| = 2.02\times10^{-3}\ \mathrm{eV},
\end{equation}
within experimental limits and in reach of next-generation experiments.
\end{theorem}

\begin{proof}
The Majorana-origin audit \cite{c11:QG179} establishes the neutrino character: the $T_3$-only
channel is self-conjugate ($48/95$ modes), the unique $Q=0$ sector provides no conserved
charge, and the reflection automorphism gives a real mass matrix. From the D96 PMNS angles
and masses, the effective Majorana mass is $m_{\beta\beta} = 2.02\times10^{-3}$ eV, within
limits and in reach of next-generation experiments.
\end{proof}

The precision predictions --- $S,T,U$, the electron g-2, and the Majorana character ---
are therefore all derived from the D96 spectral structure; their experimental validation
remains open --- the current frontier, not a derivation gap (VALID001).

\begin{remark}[Hosted-SM caveat]
\label{c11:remark:hosted-sm}
Consistent with MONO007 \cite{c11:MONO007}, the standard-model \emph{Lagrangian} and its
dynamics are hosted --- a boundary, not a derivation claim of this chapter. Every
\emph{observable} presented here is derived from the D96 moments; the Lagrangian form
$L = -\frac{1}{4}F^a_{\mu\nu}F^{a\mu\nu} + i\bar{\psi}\gamma^\mu D_\mu\psi - m\bar{\psi}\psi$
is the actualization-flow action established in Chapter~\ref{c09:chap:qm}
(Lemma~\ref{c09:lem:qm-dynamics}), but the full SM dynamical content remains hosted. No stronger
claim is made.
\end{remark}

\section{Consequences}
\label{c11:sec:sm-consequences}

\begin{theorem}[The standard model is consistent with the canonical architecture]
\label{c11:thm:sm-consistent}
The standard model is consistent with the canonical architecture: every derived SM
observable is a read of the D96 spectral moments, with the hosted dynamics the single
referenced boundary. Dimensional observables are reported relative to calibration anchors,
not from the dimensionless D96 data alone.
\end{theorem}

\begin{proof}
The fermion sector (Theorem~\ref{c11:thm:fermion-access}), the gauge structure
(Theorem~\ref{c11:thm:gauge-structure}), the couplings (Theorem~\ref{c11:thm:couplings}), the flavor
mixing (Theorem~\ref{c11:thm:ckm}, Theorem~\ref{c11:thm:pmns}), the Higgs and weak scale
(Theorem~\ref{c11:thm:weak-higgs}), and the precision predictions
(Theorem~\ref{c11:thm:oblique}, Theorem~\ref{c11:thm:electron-g2}, Theorem~\ref{c11:thm:majorana})
are all derived from the D96 moments. The hosted SM Lagrangian is the single referenced
boundary (Remark~\ref{c11:remark:hosted-sm}).
\end{proof}

\section{Summary}
\label{c11:sec:sm-summary}

This chapter has reconstructed the standard-model observables: the fermion sectors are
spectral access reads (Theorem~\ref{c11:thm:fermion-access}); the $1+3+8$ gauge sector is
the D96 degree-12 observable sector (Theorem~\ref{c11:thm:gauge-structure}); the couplings
are spectral ratios (Theorem~\ref{c11:thm:couplings}); the CKM and PMNS matrices are
derived spectral reads (Theorem~\ref{c11:thm:ckm}, Theorem~\ref{c11:thm:pmns}); and the
weak scale, the Higgs mass, and the neutrino masses are spectral constructions
(Theorem~\ref{c11:thm:weak-higgs}, Lemma~\ref{c11:lem:neutrino-masses}). The precision
predictions --- the oblique parameters $S,T,U$, the electron g-2, and the Majorana
character --- are likewise reconstructed as observable outputs
(Theorem~\ref{c11:thm:oblique}, Theorem~\ref{c11:thm:electron-g2},
Theorem~\ref{c11:thm:majorana}), so the standard model is consistent with the canonical
architecture, with the hosted dynamics the single referenced boundary
(Theorem~\ref{c11:thm:sm-consistent}, Remark~\ref{c11:remark:hosted-sm}). Every SM
observable is therefore a spectral construction, twice-derived from the spectrum
(Chapter~\ref{c06:chap:d96}, Corollary~\ref{c06:cor:constants-derived}); the following
chapter derives cosmology from the same foundation.

\begin{thebibliography}{99}
\bibitem{c11:QG149} AT-QG Phase 149, \emph{Sector Exponents Origin}, Docs/Research.
\bibitem{c11:QG150} AT-QG Phase 150, \emph{Mode Access Origin}, Docs/Research.
\bibitem{c11:QG157} AT-QG Phase 157, \emph{Effective Access Counts}, Docs/Research.
\bibitem{c11:QG161} AT-QG Phase 161, \emph{Gauge Sector Origin}, Docs/Research.
\bibitem{c11:QG162} AT-QG Phase 162, \emph{Gauge Coupling Origin}, Docs/Research.
\bibitem{c11:QG165} AT-QG Phase 165, \emph{CKM Origin}, Docs/Research.
\bibitem{c11:QG166} AT-QG Phase 166, \emph{CKM CP Origin}, Docs/Research.
\bibitem{c11:QG167} AT-QG Phase 167, \emph{PMNS Origin}, Docs/Research.
\bibitem{c11:QG168} AT-QG Phase 168, \emph{Weak Boson Mass Origin}, Docs/Research.
\bibitem{c11:QG169} AT-QG Phase 169, \emph{Higgs Mass Origin}, Docs/Research.
\bibitem{c11:QG171} AT-QG Phase 171, \emph{Muon g-2 Origin}, Docs/Research.
\bibitem{c11:QG172} AT-QG Phase 172, \emph{Neutrino Mass Law}, Docs/Research.
\bibitem{c11:QG176} AT-QG Phase 176, \emph{Higgs Blind Reconstruction}, Docs/Research.
\bibitem{c11:QG178} AT-QG Phase 178, \emph{Electron g-2 Origin}, Docs/Research.
\bibitem{c11:QG179} AT-QG Phase 179, \emph{Majorana Origin}, Docs/Research.
\bibitem{c11:QG180} AT-QG Phase 180, \emph{Oblique Parameters Origin}, Docs/Research.
\bibitem{c11:MONO007} AT-MONO007, \emph{Referee-Readiness Resolution}, Docs/Zenodo/Monograph.
\end{thebibliography}

% ======================================================================
%  AT-MONO111 — Chapter 12: Cosmology
%  Canonical Monograph (v2.0) · The Actualization Theory:
%  A Reconstruction of Physics from Difference, Actualization and Spectrum
%
%  Status: COMPLETE — PASS-READY (MONO007 referee-ready architecture)
%  Inputs: Chapter01_Difference.tex ... Chapter11_StandardModel.tex
%  Method: assembly only — canonical results only, no new physics, no speculation
%  Citations: QG227 (initial conditions), QG228 (information content),
%             QG230 (lambda origin), QG231 (structure formation),
%             QG234 (density fractions), QG237 (CMB spectrum),
%             QG238 (acoustic peaks), QG240 (blind reproduction),
%             QG295 (spectrum necessity), QG296 (reconstruction)
%
%  Usage: % ======================================================================
%  AT-MONO111 — Chapter 12: Cosmology
%  Canonical Monograph (v2.0) · The Actualization Theory:
%  A Reconstruction of Physics from Difference, Actualization and Spectrum
%
%  Status: COMPLETE — PASS-READY (MONO007 referee-ready architecture)
%  Inputs: Chapter01_Difference.tex ... Chapter11_StandardModel.tex
%  Method: assembly only — canonical results only, no new physics, no speculation
%  Citations: QG227 (initial conditions), QG228 (information content),
%             QG230 (lambda origin), QG231 (structure formation),
%             QG234 (density fractions), QG237 (CMB spectrum),
%             QG238 (acoustic peaks), QG240 (blind reproduction),
%             QG295 (spectrum necessity), QG296 (reconstruction)
%
%  Usage: % ======================================================================
%  AT-MONO111 — Chapter 12: Cosmology
%  Canonical Monograph (v2.0) · The Actualization Theory:
%  A Reconstruction of Physics from Difference, Actualization and Spectrum
%
%  Status: COMPLETE — PASS-READY (MONO007 referee-ready architecture)
%  Inputs: Chapter01_Difference.tex ... Chapter11_StandardModel.tex
%  Method: assembly only — canonical results only, no new physics, no speculation
%  Citations: QG227 (initial conditions), QG228 (information content),
%             QG230 (lambda origin), QG231 (structure formation),
%             QG234 (density fractions), QG237 (CMB spectrum),
%             QG238 (acoustic peaks), QG240 (blind reproduction),
%             QG295 (spectrum necessity), QG296 (reconstruction)
%
%  Usage: \input{Chapter12_Cosmology.tex} after the preamble
%         that defines the theorem environments (or include via the master file).
% ======================================================================

% ---- Theorem environments (define in the master preamble if not present) ----
% \newtheorem{definition}{Definition}[chapter]
% \newtheorem{lemma}[definition]{Lemma}
% \newtheorem{theorem}[definition]{Theorem}
% \newtheorem{corollary}[definition]{Corollary}
% \newtheorem{remark}[definition]{Remark}

\chapter{Cosmology}
\label{chap:cosmology}

\section{Introduction}
\label{sec:cosmology-introduction}

The preceding chapters derived the quantum and gravitational structure and the standard
model of The Actualization Theory from the canonical foundation. This chapter completes the
physics part of the monograph by deriving \emph{cosmology} from the same foundation: the
cosmological constant, the matter and vacuum fractions, the cosmic microwave background, the
structure formation, and the expansion history. Every cosmological observable of this
chapter is traced back to the D96 spectral structure \cite{QG230,QG234,QG237,QG238}.

The central results are: the cosmological constant $\Lambda$ is the residual actualization
pressure of the critical branching vacuum \cite{QG230}; the density fractions
$\Omega_\Lambda$ and $\Omega_m$ are the information-density fractions of the D96 octave
record \cite{QG234}; the CMB acoustic structure is the standing-wave harmonic structure of
the D96 recombination-scale mode ladder \cite{QG237,QG238}; structure formation follows
the actualization dynamics and the spectral hierarchy \cite{QG231}; and the expansion
history is a spectrum consequence, requiring no independent cosmological primitive
\cite{QG227,QG228,QG240}.

Throughout, cosmology is presented as \emph{emergent} --- a spectrum readout --- and no
claims beyond the accepted derivations are made. In particular, no inflation claims beyond
the accepted results, no new cosmological mechanisms, no unresolved Bekenstein discussion,
and no boundary material are included (the boundary topics are reserved for the Boundary
Layer part of the monograph). We use only accepted canonical results and introduce no new
physics and no speculation.

\section{Cosmology as a Spectrum Readout}
\label{sec:spectrum-readout}

\begin{definition}[Cosmological readout]
\label{def:cosmo-readout}
A \emph{cosmological readout} is a derived observable of the theory obtained as a function
of the D96 spectral structure: the cosmological quantities are reads of the spectrum, in
the same sense as the physics sectors of the preceding chapters.
\end{definition}

\begin{theorem}[Cosmological observables emerge from the D96 spectral structure]
\label{thm:cosmo-emergent}
The cosmological observables of the theory emerge from the D96 spectral structure: they are
derived functions of the spectral content, not independent inputs.
\end{theorem}

\begin{proof}
The reconstruction audit \cite{QG296} reconstructs the major results of the theory ---
including the cosmological ones --- from the minimal hierarchy, with the spectrum layer as
the source of the constants used by all sectors. The cosmological quantities are reads of
the spectral content: $\Lambda$ from the branching-vacuum structure \cite{QG230}, the
fractions from the octave record \cite{QG234}, the CMB structure from the mode ladder
\cite{QG237,QG238}. Hence cosmology is a spectrum readout --- emergent, not primitive.
\end{proof}

\begin{corollary}[Cosmology is emergent]
\label{cor:cosmo-emergent}
Cosmology is emergent in the theory: its observables are reads of the D96 spectrum, which
is itself the derived output of the actualization attractor.
\end{corollary}

\begin{proof}
By Theorem~\ref{thm:cosmo-emergent}, the cosmological observables are functions of the
spectral structure; by Chapter~\ref{chap:d96} (Theorem~\ref{thm:d96-derived}), the spectrum
is the derived output of the attractor. Hence cosmology is twice-emergent, consistent with
the canonical architecture.
\end{proof}

\section{The Cosmological Constant}
\label{sec:lambda}

\begin{definition}[Cosmological constant]
\label{def:lambda}
The \emph{cosmological constant} $\Lambda$ is the residual actualization pressure of the
critical branching vacuum: the realized vacuum's departure from its uniform expectation.
\end{definition}

\begin{theorem}[Existence, sign, and scaling of \texorpdfstring{$\Lambda$}{Lambda}]
\label{thm:lambda}
The cosmological constant exists, is positive, and scales with the spectral structure: at
criticality ($\mu=1$) the Galton--Watson mean is constant while the variance grows
($\mathrm{Var}(Z_k) = k\sigma^2$), the realized vacuum never equals its uniform expectation,
and the residual pressure $\Lambda$ follows.
\end{theorem}

\begin{proof}
The lambda-origin audit \cite{QG230} establishes the construction. Existence: at criticality
the Galton--Watson mean is constant but the variance grows, giving a residual pressure.
Sign: the realized vacuum carries positive information relative to the uniform expectation
(the information-content result \cite{QG228}), so the residual pressure is positive.
Scaling: the residual pressure is a function of the branching-vacuum structure. Hence
$\Lambda$ exists, is positive, and scales with the actualization structure --- derived, not
imported.
\end{proof}

\begin{corollary}[\texorpdfstring{$\Lambda$}{Lambda} is emergent, not fundamental]
\label{cor:lambda-emergent}
The cosmological constant is emergent, not fundamental: it is the residual pressure of the
actualization vacuum, not an independent parameter.
\end{corollary}

\begin{proof}
By Theorem~\ref{thm:lambda}, $\Lambda$ is the residual pressure of the critical branching
vacuum. By Theorem~\ref{thm:cosmo-emergent}, it is a read of the spectral/actualization
structure. Hence $\Lambda$ is emergent --- a derived consequence of the actualization
process, not a fundamental input.
\end{proof}

\section{Matter and Vacuum Fractions}
\label{sec:fractions}

\begin{definition}[Density fractions]
\label{def:fractions}
The \emph{density fractions} $\Omega_\Lambda$ and $\Omega_m$ are the information-density
fractions of the D96 octave record: the partition of the total density between the vacuum
and the matter content.
\end{definition}

\begin{theorem}[The density fractions are derived from spectral occupancies]
\label{thm:fractions}
The cosmological density fractions are derived from the D96 octave occupancies:
$\Omega_\Lambda = I_{\mathrm{occ}}/\ln K$ and $\Omega_m = 1 - \Omega_\Lambda$, where
$I_{\mathrm{occ}}$ is the realized octave record's information content.
\end{theorem}

\begin{proof}
The fraction-origin audit \cite{QG234} establishes the construction: the fractions are the
information-density fractions of the D96 octave record, with
$\Omega_\Lambda = I_{\mathrm{occ}}/\ln K = 0.6839$ and $\Omega_m = 0.3161$, matching
observation within a fraction of a percent. The information content
$I_{\mathrm{occ}} = KL(p\Vert\mathrm{uniform}) = 0.7513$ nats is computed from the D96
occupancies $[4,4,87]$ with $95$ modes \cite{QG228}. Hence the fractions are derived from
the spectral occupancies.
\end{proof}

\begin{corollary}[The density partition follows the spectrum structure]
\label{cor:density-partition}
The density partition of the universe follows the spectrum structure: the vacuum fraction
is the octave-record information, and the matter fraction its complement.
\end{corollary}

\begin{proof}
By Theorem~\ref{thm:fractions}, $\Omega_\Lambda$ is the octave-record information fraction
and $\Omega_m$ its complement. Both are functions of the D96 occupancies
(Chapter~\ref{chap:d96}, Corollary~\ref{cor:constants-derived}). Hence the density
partition follows the spectrum structure.
\end{proof}

\section{The Cosmic Microwave Background}
\label{sec:cmb}

\begin{definition}[CMB acoustic structure]
\label{def:cmb}
The \emph{CMB acoustic structure} is the standing-wave harmonic structure of the
recombination-scale field: the acoustic peaks are the harmonics of the D96 mode ladder.
\end{definition}

\begin{theorem}[The acoustic structure is derived from the spectrum]
\label{thm:acoustic}
The CMB acoustic structure is the standing-wave harmonic structure of the D96
recombination-scale mode ladder: the acoustic peaks are the harmonics of the spectral
mode structure.
\end{theorem}

\begin{proof}
The acoustic-peak-origin audit \cite{QG238} establishes the construction: the acoustic peaks
are the standing-wave harmonics of the recombination-scale field, which is the D96 octave
spectrum's mode ladder. The first peak $\ell_1 = 220.48$ and the ratios
$r_{21} = 2.4368$, $r_{31} = 3.6965$ are derived from the D96 octave hierarchy. Hence the
CMB acoustic structure is derived from the spectrum.
\end{proof}

\begin{corollary}[The peak hierarchy follows spectral organization]
\label{cor:peak-hierarchy}
The acoustic-peak hierarchy follows the spectral organization: the peak positions and
ratios are the harmonics and ratios of the D96 mode ladder.
\end{corollary}

\begin{proof}
By Theorem~\ref{thm:acoustic}, the peaks are the harmonics of the D96 mode ladder; the peak
ratios are ratios of the spectral harmonics
(Chapter~\ref{chap:d96}, Corollary~\ref{cor:constants-derived}). Hence the peak hierarchy
follows the spectral organization.
\end{proof}

\begin{lemma}[The seed spectrum is the counting variance]
\label{lem:seed-spectrum}
The seed power spectrum is the Poisson counting variance $\delta_i = 1/\sqrt{\langle N
\rangle}$, scale-free from critical branching.
\end{lemma}

\begin{proof}
The CMB-origin audit \cite{QG237} establishes the seed: the scalar spectral index
$n_s$ is the octave-hierarchy tilt of the D96 spectrum, and the seed power spectrum is the
Poisson counting variance, scale-free ($n_s = 1$) from critical branching. Hence the seed
spectrum is the counting variance of the actualization process.
\end{proof}

\section{Structure Formation}
\label{sec:structure}

\begin{theorem}[Large-scale organization follows actualization dynamics]
\label{thm:structure-formation}
The large-scale organization of the universe follows the actualization dynamics: structure
formation is derived from the Q-event statistics and the spectral hierarchy.
\end{theorem}

\begin{proof}
The structure-formation audit \cite{QG231} establishes the construction: structure
formation is derived from the Q-event statistics of the actualization process, with no new
primitives. The initial conditions are the uniform critical state \cite{QG227}, and the
information content provides the departure that seeds the structure \cite{QG228}. Hence
large-scale organization follows the actualization dynamics.
\end{proof}

\begin{lemma}[Matter clustering is consistent with the spectral hierarchy]
\label{lem:clustering}
The matter clustering of the universe is consistent with the spectral hierarchy: the
clustering scale structure is a read of the D96 spectral organization.
\end{lemma}

\begin{proof}
The structure-formation result \cite{QG231} and the CMB-acoustic result \cite{QG238}
together establish the consistency: the seed spectrum is the counting variance
(Lemma~\ref{lem:seed-spectrum}), the acoustic structure is the spectral harmonics
(Theorem~\ref{thm:acoustic}), and the clustering follows the same spectral hierarchy. Hence
matter clustering is consistent with the spectral organization.
\end{proof}

\section{The Expansion History}
\label{sec:expansion}

\begin{theorem}[Cosmological evolution is a spectrum consequence]
\label{thm:expansion}
The cosmological expansion is a spectrum consequence: the evolution of the universe is a
read of the actualization structure, requiring no independent cosmological primitive.
\end{theorem}

\begin{proof}
The cosmological observables --- the constant $\Lambda$ (Theorem~\ref{thm:lambda}), the
fractions (Theorem~\ref{thm:fractions}), the CMB structure (Theorem~\ref{thm:acoustic}),
and the structure formation (Theorem~\ref{thm:structure-formation}) --- are all reads of
the spectral/actualization structure. The blind reproduction audit \cite{QG240} confirms
that the cosmological quantities are recomputed from the D96 quantities alone, with no
fitted inputs. Hence the expansion history is a spectrum consequence.
\end{proof}

\begin{corollary}[No independent cosmological primitive]
\label{cor:no-cosmo-primitive}
Cosmology introduces no independent primitive: every cosmological observable is a read of
the D96 spectrum and the actualization process.
\end{corollary}

\begin{proof}
By Theorem~\ref{thm:expansion}, the cosmological evolution is a spectrum consequence; by
Theorem~\ref{thm:cosmo-emergent}, the observables are functions of the spectral content.
Hence cosmology adds no primitive to the canonical foundation
$\{\text{Difference},\eta\}$.
\end{proof}

\section{Consequences}
\label{sec:cosmo-consequences}

\begin{theorem}[Cosmology is consistent with the canonical architecture]
\label{thm:cosmo-consistent}
Cosmology is consistent with the canonical architecture: every cosmological observable is
a read of the D96 spectral moments, with the physics part of the monograph now complete.
\end{theorem}

\begin{proof}
The cosmological constant (Theorem~\ref{thm:lambda}), the fractions
(Theorem~\ref{thm:fractions}), the CMB structure (Theorem~\ref{thm:acoustic}), the structure
formation (Theorem~\ref{thm:structure-formation}), and the expansion
(Theorem~\ref{thm:expansion}) are all reads of the D96 spectral/actualization structure,
consistent with the spectrum-necessity result \cite{QG295} and the reconstruction audit
\cite{QG296}. Hence cosmology is consistent with the canonical architecture, completing the
physics part.
\end{proof}

\begin{corollary}[Difference → Actualization → Spectrum → Cosmology]
\label{cor:canonical-cosmology}
The cosmological sector completes the canonical chain:
\[
\text{Difference}\;\longrightarrow\;\text{Actualization}\;\longrightarrow\;\text{Inevitable Spectrum}\;\longrightarrow\;\text{Cosmology},
\]
with every cosmological observable a read of the spectrum.
\end{corollary}

\begin{proof}
By Theorem~\ref{thm:cosmo-consistent}, the cosmological observables are reads of the D96
spectrum; by Chapter~\ref{chap:spectrum} (Theorem~\ref{thm:core-completed}), the spectrum
is the output of the actualization process rooted in Difference. Hence the cosmological
sector is the terminal read of the canonical chain.
\end{proof}

\begin{remark}[Referee-safe wording]
Consistent with MONO007 \cite{MONO007}, this chapter presents cosmology as emergent and makes
no claims beyond the accepted derivations. No inflation claims beyond the accepted results
(the initial-condition and information origins replace the inflation epoch), no new
cosmological mechanisms, no unresolved Bekenstein discussion, and no boundary material are
included; the boundary topics are reserved for the Boundary Layer part of the monograph.
\end{remark}

\section{Summary}
\label{sec:cosmo-summary}

This chapter has derived cosmology. We have proved:

\begin{enumerate}
\item \textbf{Cosmology as a spectrum readout} (Theorem~\ref{thm:cosmo-emergent},
Corollary~\ref{cor:cosmo-emergent}): the cosmological observables emerge from the D96
spectral structure, and cosmology is emergent;
\item \textbf{The cosmological constant} (Theorem~\ref{thm:lambda},
Corollary~\ref{cor:lambda-emergent}): $\Lambda$ exists, is positive, and scales with the
actualization structure --- emergent, not fundamental;
\item \textbf{Matter and vacuum fractions} (Theorem~\ref{thm:fractions},
Corollary~\ref{cor:density-partition}): $\Omega_\Lambda$ and $\Omega_m$ are derived from
the D96 octave occupancies;
\item \textbf{The cosmic microwave background} (Theorem~\ref{thm:acoustic},
Corollary~\ref{cor:peak-hierarchy}, Lemma~\ref{lem:seed-spectrum}): the acoustic structure
is the D96 mode-ladder harmonics, with the seed spectrum the counting variance;
\item \textbf{Structure formation} (Theorem~\ref{thm:structure-formation},
Lemma~\ref{lem:clustering}): large-scale organization follows the actualization dynamics,
consistent with the spectral hierarchy;
\item \textbf{The expansion history} (Theorem~\ref{thm:expansion},
Corollary~\ref{cor:no-cosmo-primitive}): cosmological evolution is a spectrum consequence,
with no independent cosmological primitive;
\item \textbf{Consequences} (Theorem~\ref{thm:cosmo-consistent},
Corollary~\ref{cor:canonical-cosmology}): cosmology is consistent with the canonical
architecture, completing the physics part of the monograph.
\end{enumerate}

Cosmology is therefore emergent in The Actualization Theory: the cosmological constant is
the residual actualization pressure, the density fractions are the octave-record
information, the CMB acoustic structure is the spectral harmonics, structure formation
follows the actualization dynamics, and the expansion history is a spectrum consequence ---
with every observable traced back to the D96 spectral moments and no independent
cosmological primitive. This completes the physics part of the monograph
(Difference → Actualization → Spectrum → Physics): the foundation, the spectrum, quantum
mechanics, gravity, the standard model, and cosmology are all derived reads of the one
canonical structure.

\begin{thebibliography}{99}
\bibitem{QG227} AT-QG Phase 227, \emph{Initial Conditions Origin}, Docs/Research.
\bibitem{QG228} AT-QG Phase 228, \emph{Information Content Origin}, Docs/Research.
\bibitem{QG230} AT-QG Phase 230, \emph{Lambda Origin}, Docs/Research.
\bibitem{QG231} AT-QG Phase 231, \emph{Structure Formation Origin}, Docs/Research.
\bibitem{QG234} AT-QG Phase 234, \emph{Cosmological Fractions Origin}, Docs/Research.
\bibitem{QG237} AT-QG Phase 237, \emph{CMB Spectrum Origin}, Docs/Research.
\bibitem{QG238} AT-QG Phase 238, \emph{Acoustic Peak Origin}, Docs/Research.
\bibitem{QG240} AT-QG Phase 240, \emph{Cosmology Blind Reproduction}, Docs/Research.
\bibitem{QG295} AT-QG Phase 295, \emph{Spectrum Necessity Audit}, Docs/Research.
\bibitem{QG296} AT-QG Phase 296, \emph{Reconstruction Audit}, Docs/Research.
\bibitem{MONO007} AT-MONO007, \emph{Referee-Readiness Resolution}, Docs/Zenodo/Monograph.
\end{thebibliography}
 after the preamble
%         that defines the theorem environments (or include via the master file).
% ======================================================================

% ---- Theorem environments (define in the master preamble if not present) ----
% \newtheorem{definition}{Definition}[chapter]
% \newtheorem{lemma}[definition]{Lemma}
% \newtheorem{theorem}[definition]{Theorem}
% \newtheorem{corollary}[definition]{Corollary}
% \newtheorem{remark}[definition]{Remark}

\chapter{Cosmology}
\label{chap:cosmology}

\section{Introduction}
\label{sec:cosmology-introduction}

The preceding chapters derived the quantum and gravitational structure and the standard
model of The Actualization Theory from the canonical foundation. This chapter completes the
physics part of the monograph by deriving \emph{cosmology} from the same foundation: the
cosmological constant, the matter and vacuum fractions, the cosmic microwave background, the
structure formation, and the expansion history. Every cosmological observable of this
chapter is traced back to the D96 spectral structure \cite{QG230,QG234,QG237,QG238}.

The central results are: the cosmological constant $\Lambda$ is the residual actualization
pressure of the critical branching vacuum \cite{QG230}; the density fractions
$\Omega_\Lambda$ and $\Omega_m$ are the information-density fractions of the D96 octave
record \cite{QG234}; the CMB acoustic structure is the standing-wave harmonic structure of
the D96 recombination-scale mode ladder \cite{QG237,QG238}; structure formation follows
the actualization dynamics and the spectral hierarchy \cite{QG231}; and the expansion
history is a spectrum consequence, requiring no independent cosmological primitive
\cite{QG227,QG228,QG240}.

Throughout, cosmology is presented as \emph{emergent} --- a spectrum readout --- and no
claims beyond the accepted derivations are made. In particular, no inflation claims beyond
the accepted results, no new cosmological mechanisms, no unresolved Bekenstein discussion,
and no boundary material are included (the boundary topics are reserved for the Boundary
Layer part of the monograph). We use only accepted canonical results and introduce no new
physics and no speculation.

\section{Cosmology as a Spectrum Readout}
\label{sec:spectrum-readout}

\begin{definition}[Cosmological readout]
\label{def:cosmo-readout}
A \emph{cosmological readout} is a derived observable of the theory obtained as a function
of the D96 spectral structure: the cosmological quantities are reads of the spectrum, in
the same sense as the physics sectors of the preceding chapters.
\end{definition}

\begin{theorem}[Cosmological observables emerge from the D96 spectral structure]
\label{thm:cosmo-emergent}
The cosmological observables of the theory emerge from the D96 spectral structure: they are
derived functions of the spectral content, not independent inputs.
\end{theorem}

\begin{proof}
The reconstruction audit \cite{QG296} reconstructs the major results of the theory ---
including the cosmological ones --- from the minimal hierarchy, with the spectrum layer as
the source of the constants used by all sectors. The cosmological quantities are reads of
the spectral content: $\Lambda$ from the branching-vacuum structure \cite{QG230}, the
fractions from the octave record \cite{QG234}, the CMB structure from the mode ladder
\cite{QG237,QG238}. Hence cosmology is a spectrum readout --- emergent, not primitive.
\end{proof}

\begin{corollary}[Cosmology is emergent]
\label{cor:cosmo-emergent}
Cosmology is emergent in the theory: its observables are reads of the D96 spectrum, which
is itself the derived output of the actualization attractor.
\end{corollary}

\begin{proof}
By Theorem~\ref{thm:cosmo-emergent}, the cosmological observables are functions of the
spectral structure; by Chapter~\ref{chap:d96} (Theorem~\ref{thm:d96-derived}), the spectrum
is the derived output of the attractor. Hence cosmology is twice-emergent, consistent with
the canonical architecture.
\end{proof}

\section{The Cosmological Constant}
\label{sec:lambda}

\begin{definition}[Cosmological constant]
\label{def:lambda}
The \emph{cosmological constant} $\Lambda$ is the residual actualization pressure of the
critical branching vacuum: the realized vacuum's departure from its uniform expectation.
\end{definition}

\begin{theorem}[Existence, sign, and scaling of \texorpdfstring{$\Lambda$}{Lambda}]
\label{thm:lambda}
The cosmological constant exists, is positive, and scales with the spectral structure: at
criticality ($\mu=1$) the Galton--Watson mean is constant while the variance grows
($\mathrm{Var}(Z_k) = k\sigma^2$), the realized vacuum never equals its uniform expectation,
and the residual pressure $\Lambda$ follows.
\end{theorem}

\begin{proof}
The lambda-origin audit \cite{QG230} establishes the construction. Existence: at criticality
the Galton--Watson mean is constant but the variance grows, giving a residual pressure.
Sign: the realized vacuum carries positive information relative to the uniform expectation
(the information-content result \cite{QG228}), so the residual pressure is positive.
Scaling: the residual pressure is a function of the branching-vacuum structure. Hence
$\Lambda$ exists, is positive, and scales with the actualization structure --- derived, not
imported.
\end{proof}

\begin{corollary}[\texorpdfstring{$\Lambda$}{Lambda} is emergent, not fundamental]
\label{cor:lambda-emergent}
The cosmological constant is emergent, not fundamental: it is the residual pressure of the
actualization vacuum, not an independent parameter.
\end{corollary}

\begin{proof}
By Theorem~\ref{thm:lambda}, $\Lambda$ is the residual pressure of the critical branching
vacuum. By Theorem~\ref{thm:cosmo-emergent}, it is a read of the spectral/actualization
structure. Hence $\Lambda$ is emergent --- a derived consequence of the actualization
process, not a fundamental input.
\end{proof}

\section{Matter and Vacuum Fractions}
\label{sec:fractions}

\begin{definition}[Density fractions]
\label{def:fractions}
The \emph{density fractions} $\Omega_\Lambda$ and $\Omega_m$ are the information-density
fractions of the D96 octave record: the partition of the total density between the vacuum
and the matter content.
\end{definition}

\begin{theorem}[The density fractions are derived from spectral occupancies]
\label{thm:fractions}
The cosmological density fractions are derived from the D96 octave occupancies:
$\Omega_\Lambda = I_{\mathrm{occ}}/\ln K$ and $\Omega_m = 1 - \Omega_\Lambda$, where
$I_{\mathrm{occ}}$ is the realized octave record's information content.
\end{theorem}

\begin{proof}
The fraction-origin audit \cite{QG234} establishes the construction: the fractions are the
information-density fractions of the D96 octave record, with
$\Omega_\Lambda = I_{\mathrm{occ}}/\ln K = 0.6839$ and $\Omega_m = 0.3161$, matching
observation within a fraction of a percent. The information content
$I_{\mathrm{occ}} = KL(p\Vert\mathrm{uniform}) = 0.7513$ nats is computed from the D96
occupancies $[4,4,87]$ with $95$ modes \cite{QG228}. Hence the fractions are derived from
the spectral occupancies.
\end{proof}

\begin{corollary}[The density partition follows the spectrum structure]
\label{cor:density-partition}
The density partition of the universe follows the spectrum structure: the vacuum fraction
is the octave-record information, and the matter fraction its complement.
\end{corollary}

\begin{proof}
By Theorem~\ref{thm:fractions}, $\Omega_\Lambda$ is the octave-record information fraction
and $\Omega_m$ its complement. Both are functions of the D96 occupancies
(Chapter~\ref{chap:d96}, Corollary~\ref{cor:constants-derived}). Hence the density
partition follows the spectrum structure.
\end{proof}

\section{The Cosmic Microwave Background}
\label{sec:cmb}

\begin{definition}[CMB acoustic structure]
\label{def:cmb}
The \emph{CMB acoustic structure} is the standing-wave harmonic structure of the
recombination-scale field: the acoustic peaks are the harmonics of the D96 mode ladder.
\end{definition}

\begin{theorem}[The acoustic structure is derived from the spectrum]
\label{thm:acoustic}
The CMB acoustic structure is the standing-wave harmonic structure of the D96
recombination-scale mode ladder: the acoustic peaks are the harmonics of the spectral
mode structure.
\end{theorem}

\begin{proof}
The acoustic-peak-origin audit \cite{QG238} establishes the construction: the acoustic peaks
are the standing-wave harmonics of the recombination-scale field, which is the D96 octave
spectrum's mode ladder. The first peak $\ell_1 = 220.48$ and the ratios
$r_{21} = 2.4368$, $r_{31} = 3.6965$ are derived from the D96 octave hierarchy. Hence the
CMB acoustic structure is derived from the spectrum.
\end{proof}

\begin{corollary}[The peak hierarchy follows spectral organization]
\label{cor:peak-hierarchy}
The acoustic-peak hierarchy follows the spectral organization: the peak positions and
ratios are the harmonics and ratios of the D96 mode ladder.
\end{corollary}

\begin{proof}
By Theorem~\ref{thm:acoustic}, the peaks are the harmonics of the D96 mode ladder; the peak
ratios are ratios of the spectral harmonics
(Chapter~\ref{chap:d96}, Corollary~\ref{cor:constants-derived}). Hence the peak hierarchy
follows the spectral organization.
\end{proof}

\begin{lemma}[The seed spectrum is the counting variance]
\label{lem:seed-spectrum}
The seed power spectrum is the Poisson counting variance $\delta_i = 1/\sqrt{\langle N
\rangle}$, scale-free from critical branching.
\end{lemma}

\begin{proof}
The CMB-origin audit \cite{QG237} establishes the seed: the scalar spectral index
$n_s$ is the octave-hierarchy tilt of the D96 spectrum, and the seed power spectrum is the
Poisson counting variance, scale-free ($n_s = 1$) from critical branching. Hence the seed
spectrum is the counting variance of the actualization process.
\end{proof}

\section{Structure Formation}
\label{sec:structure}

\begin{theorem}[Large-scale organization follows actualization dynamics]
\label{thm:structure-formation}
The large-scale organization of the universe follows the actualization dynamics: structure
formation is derived from the Q-event statistics and the spectral hierarchy.
\end{theorem}

\begin{proof}
The structure-formation audit \cite{QG231} establishes the construction: structure
formation is derived from the Q-event statistics of the actualization process, with no new
primitives. The initial conditions are the uniform critical state \cite{QG227}, and the
information content provides the departure that seeds the structure \cite{QG228}. Hence
large-scale organization follows the actualization dynamics.
\end{proof}

\begin{lemma}[Matter clustering is consistent with the spectral hierarchy]
\label{lem:clustering}
The matter clustering of the universe is consistent with the spectral hierarchy: the
clustering scale structure is a read of the D96 spectral organization.
\end{lemma}

\begin{proof}
The structure-formation result \cite{QG231} and the CMB-acoustic result \cite{QG238}
together establish the consistency: the seed spectrum is the counting variance
(Lemma~\ref{lem:seed-spectrum}), the acoustic structure is the spectral harmonics
(Theorem~\ref{thm:acoustic}), and the clustering follows the same spectral hierarchy. Hence
matter clustering is consistent with the spectral organization.
\end{proof}

\section{The Expansion History}
\label{sec:expansion}

\begin{theorem}[Cosmological evolution is a spectrum consequence]
\label{thm:expansion}
The cosmological expansion is a spectrum consequence: the evolution of the universe is a
read of the actualization structure, requiring no independent cosmological primitive.
\end{theorem}

\begin{proof}
The cosmological observables --- the constant $\Lambda$ (Theorem~\ref{thm:lambda}), the
fractions (Theorem~\ref{thm:fractions}), the CMB structure (Theorem~\ref{thm:acoustic}),
and the structure formation (Theorem~\ref{thm:structure-formation}) --- are all reads of
the spectral/actualization structure. The blind reproduction audit \cite{QG240} confirms
that the cosmological quantities are recomputed from the D96 quantities alone, with no
fitted inputs. Hence the expansion history is a spectrum consequence.
\end{proof}

\begin{corollary}[No independent cosmological primitive]
\label{cor:no-cosmo-primitive}
Cosmology introduces no independent primitive: every cosmological observable is a read of
the D96 spectrum and the actualization process.
\end{corollary}

\begin{proof}
By Theorem~\ref{thm:expansion}, the cosmological evolution is a spectrum consequence; by
Theorem~\ref{thm:cosmo-emergent}, the observables are functions of the spectral content.
Hence cosmology adds no primitive to the canonical foundation
$\{\text{Difference},\eta\}$.
\end{proof}

\section{Consequences}
\label{sec:cosmo-consequences}

\begin{theorem}[Cosmology is consistent with the canonical architecture]
\label{thm:cosmo-consistent}
Cosmology is consistent with the canonical architecture: every cosmological observable is
a read of the D96 spectral moments, with the physics part of the monograph now complete.
\end{theorem}

\begin{proof}
The cosmological constant (Theorem~\ref{thm:lambda}), the fractions
(Theorem~\ref{thm:fractions}), the CMB structure (Theorem~\ref{thm:acoustic}), the structure
formation (Theorem~\ref{thm:structure-formation}), and the expansion
(Theorem~\ref{thm:expansion}) are all reads of the D96 spectral/actualization structure,
consistent with the spectrum-necessity result \cite{QG295} and the reconstruction audit
\cite{QG296}. Hence cosmology is consistent with the canonical architecture, completing the
physics part.
\end{proof}

\begin{corollary}[Difference → Actualization → Spectrum → Cosmology]
\label{cor:canonical-cosmology}
The cosmological sector completes the canonical chain:
\[
\text{Difference}\;\longrightarrow\;\text{Actualization}\;\longrightarrow\;\text{Inevitable Spectrum}\;\longrightarrow\;\text{Cosmology},
\]
with every cosmological observable a read of the spectrum.
\end{corollary}

\begin{proof}
By Theorem~\ref{thm:cosmo-consistent}, the cosmological observables are reads of the D96
spectrum; by Chapter~\ref{chap:spectrum} (Theorem~\ref{thm:core-completed}), the spectrum
is the output of the actualization process rooted in Difference. Hence the cosmological
sector is the terminal read of the canonical chain.
\end{proof}

\begin{remark}[Referee-safe wording]
Consistent with MONO007 \cite{MONO007}, this chapter presents cosmology as emergent and makes
no claims beyond the accepted derivations. No inflation claims beyond the accepted results
(the initial-condition and information origins replace the inflation epoch), no new
cosmological mechanisms, no unresolved Bekenstein discussion, and no boundary material are
included; the boundary topics are reserved for the Boundary Layer part of the monograph.
\end{remark}

\section{Summary}
\label{sec:cosmo-summary}

This chapter has derived cosmology. We have proved:

\begin{enumerate}
\item \textbf{Cosmology as a spectrum readout} (Theorem~\ref{thm:cosmo-emergent},
Corollary~\ref{cor:cosmo-emergent}): the cosmological observables emerge from the D96
spectral structure, and cosmology is emergent;
\item \textbf{The cosmological constant} (Theorem~\ref{thm:lambda},
Corollary~\ref{cor:lambda-emergent}): $\Lambda$ exists, is positive, and scales with the
actualization structure --- emergent, not fundamental;
\item \textbf{Matter and vacuum fractions} (Theorem~\ref{thm:fractions},
Corollary~\ref{cor:density-partition}): $\Omega_\Lambda$ and $\Omega_m$ are derived from
the D96 octave occupancies;
\item \textbf{The cosmic microwave background} (Theorem~\ref{thm:acoustic},
Corollary~\ref{cor:peak-hierarchy}, Lemma~\ref{lem:seed-spectrum}): the acoustic structure
is the D96 mode-ladder harmonics, with the seed spectrum the counting variance;
\item \textbf{Structure formation} (Theorem~\ref{thm:structure-formation},
Lemma~\ref{lem:clustering}): large-scale organization follows the actualization dynamics,
consistent with the spectral hierarchy;
\item \textbf{The expansion history} (Theorem~\ref{thm:expansion},
Corollary~\ref{cor:no-cosmo-primitive}): cosmological evolution is a spectrum consequence,
with no independent cosmological primitive;
\item \textbf{Consequences} (Theorem~\ref{thm:cosmo-consistent},
Corollary~\ref{cor:canonical-cosmology}): cosmology is consistent with the canonical
architecture, completing the physics part of the monograph.
\end{enumerate}

Cosmology is therefore emergent in The Actualization Theory: the cosmological constant is
the residual actualization pressure, the density fractions are the octave-record
information, the CMB acoustic structure is the spectral harmonics, structure formation
follows the actualization dynamics, and the expansion history is a spectrum consequence ---
with every observable traced back to the D96 spectral moments and no independent
cosmological primitive. This completes the physics part of the monograph
(Difference → Actualization → Spectrum → Physics): the foundation, the spectrum, quantum
mechanics, gravity, the standard model, and cosmology are all derived reads of the one
canonical structure.

\begin{thebibliography}{99}
\bibitem{QG227} AT-QG Phase 227, \emph{Initial Conditions Origin}, Docs/Research.
\bibitem{QG228} AT-QG Phase 228, \emph{Information Content Origin}, Docs/Research.
\bibitem{QG230} AT-QG Phase 230, \emph{Lambda Origin}, Docs/Research.
\bibitem{QG231} AT-QG Phase 231, \emph{Structure Formation Origin}, Docs/Research.
\bibitem{QG234} AT-QG Phase 234, \emph{Cosmological Fractions Origin}, Docs/Research.
\bibitem{QG237} AT-QG Phase 237, \emph{CMB Spectrum Origin}, Docs/Research.
\bibitem{QG238} AT-QG Phase 238, \emph{Acoustic Peak Origin}, Docs/Research.
\bibitem{QG240} AT-QG Phase 240, \emph{Cosmology Blind Reproduction}, Docs/Research.
\bibitem{QG295} AT-QG Phase 295, \emph{Spectrum Necessity Audit}, Docs/Research.
\bibitem{QG296} AT-QG Phase 296, \emph{Reconstruction Audit}, Docs/Research.
\bibitem{MONO007} AT-MONO007, \emph{Referee-Readiness Resolution}, Docs/Zenodo/Monograph.
\end{thebibliography}
 after the preamble
%         that defines the theorem environments (or include via the master file).
% ======================================================================

% ---- Theorem environments (define in the master preamble if not present) ----
% \newtheorem{definition}{Definition}[chapter]
% \newtheorem{lemma}[definition]{Lemma}
% \newtheorem{theorem}[definition]{Theorem}
% \newtheorem{corollary}[definition]{Corollary}
% \newtheorem{remark}[definition]{Remark}

\chapter{Cosmology}
\label{chap:cosmology}

\section{Introduction}
\label{sec:cosmology-introduction}

The preceding chapters derived the quantum and gravitational structure and the standard
model of The Actualization Theory from the canonical foundation. This chapter completes the
physics part of the monograph by deriving \emph{cosmology} from the same foundation: the
cosmological constant, the matter and vacuum fractions, the cosmic microwave background, the
structure formation, and the expansion history. Every cosmological observable of this
chapter is traced back to the D96 spectral structure \cite{QG230,QG234,QG237,QG238}.

The central results are: the cosmological constant $\Lambda$ is the residual actualization
pressure of the critical branching vacuum \cite{QG230}; the density fractions
$\Omega_\Lambda$ and $\Omega_m$ are the information-density fractions of the D96 octave
record \cite{QG234}; the CMB acoustic structure is the standing-wave harmonic structure of
the D96 recombination-scale mode ladder \cite{QG237,QG238}; structure formation follows
the actualization dynamics and the spectral hierarchy \cite{QG231}; and the expansion
history is a spectrum consequence, requiring no independent cosmological primitive
\cite{QG227,QG228,QG240}.

Throughout, cosmology is presented as \emph{emergent} --- a spectrum readout --- and no
claims beyond the accepted derivations are made. In particular, no inflation claims beyond
the accepted results, no new cosmological mechanisms, no unresolved Bekenstein discussion,
and no boundary material are included (the boundary topics are reserved for the Boundary
Layer part of the monograph). We use only accepted canonical results and introduce no new
physics and no speculation.

\section{Cosmology as a Spectrum Readout}
\label{sec:spectrum-readout}

\begin{definition}[Cosmological readout]
\label{def:cosmo-readout}
A \emph{cosmological readout} is a derived observable of the theory obtained as a function
of the D96 spectral structure: the cosmological quantities are reads of the spectrum, in
the same sense as the physics sectors of the preceding chapters.
\end{definition}

\begin{theorem}[Cosmological observables emerge from the D96 spectral structure]
\label{thm:cosmo-emergent}
The cosmological observables of the theory emerge from the D96 spectral structure: they are
derived functions of the spectral content, not independent inputs.
\end{theorem}

\begin{proof}
The reconstruction audit \cite{QG296} reconstructs the major results of the theory ---
including the cosmological ones --- from the minimal hierarchy, with the spectrum layer as
the source of the constants used by all sectors. The cosmological quantities are reads of
the spectral content: $\Lambda$ from the branching-vacuum structure \cite{QG230}, the
fractions from the octave record \cite{QG234}, the CMB structure from the mode ladder
\cite{QG237,QG238}. Hence cosmology is a spectrum readout --- emergent, not primitive.
\end{proof}

\begin{corollary}[Cosmology is emergent]
\label{cor:cosmo-emergent}
Cosmology is emergent in the theory: its observables are reads of the D96 spectrum, which
is itself the derived output of the actualization attractor.
\end{corollary}

\begin{proof}
By Theorem~\ref{thm:cosmo-emergent}, the cosmological observables are functions of the
spectral structure; by Chapter~\ref{chap:d96} (Theorem~\ref{thm:d96-derived}), the spectrum
is the derived output of the attractor. Hence cosmology is twice-emergent, consistent with
the canonical architecture.
\end{proof}

\section{The Cosmological Constant}
\label{sec:lambda}

\begin{definition}[Cosmological constant]
\label{def:lambda}
The \emph{cosmological constant} $\Lambda$ is the residual actualization pressure of the
critical branching vacuum: the realized vacuum's departure from its uniform expectation.
\end{definition}

\begin{theorem}[Existence, sign, and scaling of \texorpdfstring{$\Lambda$}{Lambda}]
\label{thm:lambda}
The cosmological constant exists, is positive, and scales with the spectral structure: at
criticality ($\mu=1$) the Galton--Watson mean is constant while the variance grows
($\mathrm{Var}(Z_k) = k\sigma^2$), the realized vacuum never equals its uniform expectation,
and the residual pressure $\Lambda$ follows.
\end{theorem}

\begin{proof}
The lambda-origin audit \cite{QG230} establishes the construction. Existence: at criticality
the Galton--Watson mean is constant but the variance grows, giving a residual pressure.
Sign: the realized vacuum carries positive information relative to the uniform expectation
(the information-content result \cite{QG228}), so the residual pressure is positive.
Scaling: the residual pressure is a function of the branching-vacuum structure. Hence
$\Lambda$ exists, is positive, and scales with the actualization structure --- derived, not
imported.
\end{proof}

\begin{corollary}[\texorpdfstring{$\Lambda$}{Lambda} is emergent, not fundamental]
\label{cor:lambda-emergent}
The cosmological constant is emergent, not fundamental: it is the residual pressure of the
actualization vacuum, not an independent parameter.
\end{corollary}

\begin{proof}
By Theorem~\ref{thm:lambda}, $\Lambda$ is the residual pressure of the critical branching
vacuum. By Theorem~\ref{thm:cosmo-emergent}, it is a read of the spectral/actualization
structure. Hence $\Lambda$ is emergent --- a derived consequence of the actualization
process, not a fundamental input.
\end{proof}

\section{Matter and Vacuum Fractions}
\label{sec:fractions}

\begin{definition}[Density fractions]
\label{def:fractions}
The \emph{density fractions} $\Omega_\Lambda$ and $\Omega_m$ are the information-density
fractions of the D96 octave record: the partition of the total density between the vacuum
and the matter content.
\end{definition}

\begin{theorem}[The density fractions are derived from spectral occupancies]
\label{thm:fractions}
The cosmological density fractions are derived from the D96 octave occupancies:
$\Omega_\Lambda = I_{\mathrm{occ}}/\ln K$ and $\Omega_m = 1 - \Omega_\Lambda$, where
$I_{\mathrm{occ}}$ is the realized octave record's information content.
\end{theorem}

\begin{proof}
The fraction-origin audit \cite{QG234} establishes the construction: the fractions are the
information-density fractions of the D96 octave record, with
$\Omega_\Lambda = I_{\mathrm{occ}}/\ln K = 0.6839$ and $\Omega_m = 0.3161$, matching
observation within a fraction of a percent. The information content
$I_{\mathrm{occ}} = KL(p\Vert\mathrm{uniform}) = 0.7513$ nats is computed from the D96
occupancies $[4,4,87]$ with $95$ modes \cite{QG228}. Hence the fractions are derived from
the spectral occupancies.
\end{proof}

\begin{corollary}[The density partition follows the spectrum structure]
\label{cor:density-partition}
The density partition of the universe follows the spectrum structure: the vacuum fraction
is the octave-record information, and the matter fraction its complement.
\end{corollary}

\begin{proof}
By Theorem~\ref{thm:fractions}, $\Omega_\Lambda$ is the octave-record information fraction
and $\Omega_m$ its complement. Both are functions of the D96 occupancies
(Chapter~\ref{chap:d96}, Corollary~\ref{cor:constants-derived}). Hence the density
partition follows the spectrum structure.
\end{proof}

\section{The Cosmic Microwave Background}
\label{sec:cmb}

\begin{definition}[CMB acoustic structure]
\label{def:cmb}
The \emph{CMB acoustic structure} is the standing-wave harmonic structure of the
recombination-scale field: the acoustic peaks are the harmonics of the D96 mode ladder.
\end{definition}

\begin{theorem}[The acoustic structure is derived from the spectrum]
\label{thm:acoustic}
The CMB acoustic structure is the standing-wave harmonic structure of the D96
recombination-scale mode ladder: the acoustic peaks are the harmonics of the spectral
mode structure.
\end{theorem}

\begin{proof}
The acoustic-peak-origin audit \cite{QG238} establishes the construction: the acoustic peaks
are the standing-wave harmonics of the recombination-scale field, which is the D96 octave
spectrum's mode ladder. The first peak $\ell_1 = 220.48$ and the ratios
$r_{21} = 2.4368$, $r_{31} = 3.6965$ are derived from the D96 octave hierarchy. Hence the
CMB acoustic structure is derived from the spectrum.
\end{proof}

\begin{corollary}[The peak hierarchy follows spectral organization]
\label{cor:peak-hierarchy}
The acoustic-peak hierarchy follows the spectral organization: the peak positions and
ratios are the harmonics and ratios of the D96 mode ladder.
\end{corollary}

\begin{proof}
By Theorem~\ref{thm:acoustic}, the peaks are the harmonics of the D96 mode ladder; the peak
ratios are ratios of the spectral harmonics
(Chapter~\ref{chap:d96}, Corollary~\ref{cor:constants-derived}). Hence the peak hierarchy
follows the spectral organization.
\end{proof}

\begin{lemma}[The seed spectrum is the counting variance]
\label{lem:seed-spectrum}
The seed power spectrum is the Poisson counting variance $\delta_i = 1/\sqrt{\langle N
\rangle}$, scale-free from critical branching.
\end{lemma}

\begin{proof}
The CMB-origin audit \cite{QG237} establishes the seed: the scalar spectral index
$n_s$ is the octave-hierarchy tilt of the D96 spectrum, and the seed power spectrum is the
Poisson counting variance, scale-free ($n_s = 1$) from critical branching. Hence the seed
spectrum is the counting variance of the actualization process.
\end{proof}

\section{Structure Formation}
\label{sec:structure}

\begin{theorem}[Large-scale organization follows actualization dynamics]
\label{thm:structure-formation}
The large-scale organization of the universe follows the actualization dynamics: structure
formation is derived from the Q-event statistics and the spectral hierarchy.
\end{theorem}

\begin{proof}
The structure-formation audit \cite{QG231} establishes the construction: structure
formation is derived from the Q-event statistics of the actualization process, with no new
primitives. The initial conditions are the uniform critical state \cite{QG227}, and the
information content provides the departure that seeds the structure \cite{QG228}. Hence
large-scale organization follows the actualization dynamics.
\end{proof}

\begin{lemma}[Matter clustering is consistent with the spectral hierarchy]
\label{lem:clustering}
The matter clustering of the universe is consistent with the spectral hierarchy: the
clustering scale structure is a read of the D96 spectral organization.
\end{lemma}

\begin{proof}
The structure-formation result \cite{QG231} and the CMB-acoustic result \cite{QG238}
together establish the consistency: the seed spectrum is the counting variance
(Lemma~\ref{lem:seed-spectrum}), the acoustic structure is the spectral harmonics
(Theorem~\ref{thm:acoustic}), and the clustering follows the same spectral hierarchy. Hence
matter clustering is consistent with the spectral organization.
\end{proof}

\section{The Expansion History}
\label{sec:expansion}

\begin{theorem}[Cosmological evolution is a spectrum consequence]
\label{thm:expansion}
The cosmological expansion is a spectrum consequence: the evolution of the universe is a
read of the actualization structure, requiring no independent cosmological primitive.
\end{theorem}

\begin{proof}
The cosmological observables --- the constant $\Lambda$ (Theorem~\ref{thm:lambda}), the
fractions (Theorem~\ref{thm:fractions}), the CMB structure (Theorem~\ref{thm:acoustic}),
and the structure formation (Theorem~\ref{thm:structure-formation}) --- are all reads of
the spectral/actualization structure. The blind reproduction audit \cite{QG240} confirms
that the cosmological quantities are recomputed from the D96 quantities alone, with no
fitted inputs. Hence the expansion history is a spectrum consequence.
\end{proof}

\begin{corollary}[No independent cosmological primitive]
\label{cor:no-cosmo-primitive}
Cosmology introduces no independent primitive: every cosmological observable is a read of
the D96 spectrum and the actualization process.
\end{corollary}

\begin{proof}
By Theorem~\ref{thm:expansion}, the cosmological evolution is a spectrum consequence; by
Theorem~\ref{thm:cosmo-emergent}, the observables are functions of the spectral content.
Hence cosmology adds no primitive to the canonical foundation
$\{\text{Difference},\eta\}$.
\end{proof}

\section{Consequences}
\label{sec:cosmo-consequences}

\begin{theorem}[Cosmology is consistent with the canonical architecture]
\label{thm:cosmo-consistent}
Cosmology is consistent with the canonical architecture: every cosmological observable is
a read of the D96 spectral moments, with the physics part of the monograph now complete.
\end{theorem}

\begin{proof}
The cosmological constant (Theorem~\ref{thm:lambda}), the fractions
(Theorem~\ref{thm:fractions}), the CMB structure (Theorem~\ref{thm:acoustic}), the structure
formation (Theorem~\ref{thm:structure-formation}), and the expansion
(Theorem~\ref{thm:expansion}) are all reads of the D96 spectral/actualization structure,
consistent with the spectrum-necessity result \cite{QG295} and the reconstruction audit
\cite{QG296}. Hence cosmology is consistent with the canonical architecture, completing the
physics part.
\end{proof}

\begin{corollary}[Difference → Actualization → Spectrum → Cosmology]
\label{cor:canonical-cosmology}
The cosmological sector completes the canonical chain:
\[
\text{Difference}\;\longrightarrow\;\text{Actualization}\;\longrightarrow\;\text{Inevitable Spectrum}\;\longrightarrow\;\text{Cosmology},
\]
with every cosmological observable a read of the spectrum.
\end{corollary}

\begin{proof}
By Theorem~\ref{thm:cosmo-consistent}, the cosmological observables are reads of the D96
spectrum; by Chapter~\ref{chap:spectrum} (Theorem~\ref{thm:core-completed}), the spectrum
is the output of the actualization process rooted in Difference. Hence the cosmological
sector is the terminal read of the canonical chain.
\end{proof}

\begin{remark}[Referee-safe wording]
Consistent with MONO007 \cite{MONO007}, this chapter presents cosmology as emergent and makes
no claims beyond the accepted derivations. No inflation claims beyond the accepted results
(the initial-condition and information origins replace the inflation epoch), no new
cosmological mechanisms, no unresolved Bekenstein discussion, and no boundary material are
included; the boundary topics are reserved for the Boundary Layer part of the monograph.
\end{remark}

\section{Summary}
\label{sec:cosmo-summary}

This chapter has derived cosmology. We have proved:

\begin{enumerate}
\item \textbf{Cosmology as a spectrum readout} (Theorem~\ref{thm:cosmo-emergent},
Corollary~\ref{cor:cosmo-emergent}): the cosmological observables emerge from the D96
spectral structure, and cosmology is emergent;
\item \textbf{The cosmological constant} (Theorem~\ref{thm:lambda},
Corollary~\ref{cor:lambda-emergent}): $\Lambda$ exists, is positive, and scales with the
actualization structure --- emergent, not fundamental;
\item \textbf{Matter and vacuum fractions} (Theorem~\ref{thm:fractions},
Corollary~\ref{cor:density-partition}): $\Omega_\Lambda$ and $\Omega_m$ are derived from
the D96 octave occupancies;
\item \textbf{The cosmic microwave background} (Theorem~\ref{thm:acoustic},
Corollary~\ref{cor:peak-hierarchy}, Lemma~\ref{lem:seed-spectrum}): the acoustic structure
is the D96 mode-ladder harmonics, with the seed spectrum the counting variance;
\item \textbf{Structure formation} (Theorem~\ref{thm:structure-formation},
Lemma~\ref{lem:clustering}): large-scale organization follows the actualization dynamics,
consistent with the spectral hierarchy;
\item \textbf{The expansion history} (Theorem~\ref{thm:expansion},
Corollary~\ref{cor:no-cosmo-primitive}): cosmological evolution is a spectrum consequence,
with no independent cosmological primitive;
\item \textbf{Consequences} (Theorem~\ref{thm:cosmo-consistent},
Corollary~\ref{cor:canonical-cosmology}): cosmology is consistent with the canonical
architecture, completing the physics part of the monograph.
\end{enumerate}

Cosmology is therefore emergent in The Actualization Theory: the cosmological constant is
the residual actualization pressure, the density fractions are the octave-record
information, the CMB acoustic structure is the spectral harmonics, structure formation
follows the actualization dynamics, and the expansion history is a spectrum consequence ---
with every observable traced back to the D96 spectral moments and no independent
cosmological primitive. This completes the physics part of the monograph
(Difference → Actualization → Spectrum → Physics): the foundation, the spectrum, quantum
mechanics, gravity, the standard model, and cosmology are all derived reads of the one
canonical structure.

\begin{thebibliography}{99}
\bibitem{QG227} AT-QG Phase 227, \emph{Initial Conditions Origin}, Docs/Research.
\bibitem{QG228} AT-QG Phase 228, \emph{Information Content Origin}, Docs/Research.
\bibitem{QG230} AT-QG Phase 230, \emph{Lambda Origin}, Docs/Research.
\bibitem{QG231} AT-QG Phase 231, \emph{Structure Formation Origin}, Docs/Research.
\bibitem{QG234} AT-QG Phase 234, \emph{Cosmological Fractions Origin}, Docs/Research.
\bibitem{QG237} AT-QG Phase 237, \emph{CMB Spectrum Origin}, Docs/Research.
\bibitem{QG238} AT-QG Phase 238, \emph{Acoustic Peak Origin}, Docs/Research.
\bibitem{QG240} AT-QG Phase 240, \emph{Cosmology Blind Reproduction}, Docs/Research.
\bibitem{QG295} AT-QG Phase 295, \emph{Spectrum Necessity Audit}, Docs/Research.
\bibitem{QG296} AT-QG Phase 296, \emph{Reconstruction Audit}, Docs/Research.
\bibitem{MONO007} AT-MONO007, \emph{Referee-Readiness Resolution}, Docs/Zenodo/Monograph.
\end{thebibliography}


% Part V — Boundary Layer
\part{Boundary Layer}
% ======================================================================
%  AT-MONO113 — Chapter 13: Boundaries of The Actualization Theory
%  Canonical Monograph (v2.0) · The Actualization Theory:
%  A Reconstruction of Physics from Difference, Actualization and Spectrum
%
%  Status: COMPLETE — PASS-READY (MONO007 referee-ready architecture)
%  Inputs: Chapter01_Difference.tex ... Chapter06_D96Spectrum.tex,
%          ATQG_CanonicalMonograph.md, ATQG_Mono007RefereeReadiness.md,
%          ATVALID_UntestedItemsAudit.md
%  Method: assembly only — canonical results only
%  Citations: QG185 (Bekenstein quarter), QG196 (impossibility proof),
%             QG242 (SM dynamics hosted), QG245 (SM dynamics not complete),
%             QG278 (fundamental boundary), QG286 (difference duality),
%             QG291 (framework necessity), QG299 (post-frontier),
%             MONO007 (referee readiness), VALID001 (validation classification)
%
%  Usage: % ======================================================================
%  AT-MONO113 — Chapter 13: Boundaries of The Actualization Theory
%  Canonical Monograph (v2.0) · The Actualization Theory:
%  A Reconstruction of Physics from Difference, Actualization and Spectrum
%
%  Status: COMPLETE — PASS-READY (MONO007 referee-ready architecture)
%  Inputs: Chapter01_Difference.tex ... Chapter06_D96Spectrum.tex,
%          ATQG_CanonicalMonograph.md, ATQG_Mono007RefereeReadiness.md,
%          ATVALID_UntestedItemsAudit.md
%  Method: assembly only — canonical results only
%  Citations: QG185 (Bekenstein quarter), QG196 (impossibility proof),
%             QG242 (SM dynamics hosted), QG245 (SM dynamics not complete),
%             QG278 (fundamental boundary), QG286 (difference duality),
%             QG291 (framework necessity), QG299 (post-frontier),
%             MONO007 (referee readiness), VALID001 (validation classification)
%
%  Usage: % ======================================================================
%  AT-MONO113 — Chapter 13: Boundaries of The Actualization Theory
%  Canonical Monograph (v2.0) · The Actualization Theory:
%  A Reconstruction of Physics from Difference, Actualization and Spectrum
%
%  Status: COMPLETE — PASS-READY (MONO007 referee-ready architecture)
%  Inputs: Chapter01_Difference.tex ... Chapter06_D96Spectrum.tex,
%          ATQG_CanonicalMonograph.md, ATQG_Mono007RefereeReadiness.md,
%          ATVALID_UntestedItemsAudit.md
%  Method: assembly only — canonical results only
%  Citations: QG185 (Bekenstein quarter), QG196 (impossibility proof),
%             QG242 (SM dynamics hosted), QG245 (SM dynamics not complete),
%             QG278 (fundamental boundary), QG286 (difference duality),
%             QG291 (framework necessity), QG299 (post-frontier),
%             MONO007 (referee readiness), VALID001 (validation classification)
%
%  Usage: \input{Chapter13_BoundariesOfTheActualizationTheory.tex} after the
%         preamble that defines the theorem environments (or include via the
%         master file).
% ======================================================================

% ---- Theorem environments (define in the master preamble if not present) ----
% \newtheorem{definition}{Definition}[chapter]
% \newtheorem{lemma}[definition]{Lemma}
% \newtheorem{theorem}[definition]{Theorem}
% \newtheorem{corollary}[definition]{Corollary}
% \newtheorem{remark}[definition]{Remark}

\chapter{Boundaries of The Actualization Theory}
\label{c13:chap:boundaries}

\section{Introduction}
\label{c13:sec:boundaries-introduction}

The preceding chapters derived the physics of The Actualization Theory: the foundation, the
spectrum, quantum mechanics, gravity, the standard model, and cosmology. This chapter is the
Boundary Layer of the monograph: it separates \emph{ontological boundaries} from \emph{physics
derivation}, documents the limits of the framework, and shows that the boundaries are
\emph{not derivation failures} --- they are the explicit limits beyond which the framework
does not claim to derive.

The boundaries documented here are: Difference as an ontological boundary \cite{c13:QG278}; the
status of the tensor face $\psi$ \cite{c13:QG286}; the numerical boundary $\pi$ \cite{c13:QG291};
the Bekenstein $2\pi$ boundary \cite{c13:QG185,c13:QG196}; and the hosted standard-model dynamics
\cite{c13:QG242,c13:QG245}. The chapter draws the distinction between \emph{derived physics} and
\emph{boundary questions} \cite{c13:QG299} and states explicitly what the theory does not claim.

Following VALID001 \cite{c13:VALID001}, \emph{experimental validation is not a boundary}: the
derived predictions $S,T,U$, the electron g-2 $a_e$, and the Majorana character $0\nu\beta\beta$
await experiment and are not listed as derivation gaps. Consistent with MONO007 \cite{c13:MONO007},
the presentation is referee-safe and introduces no new physics and no speculation.

\section{The Nature of Scientific Boundaries}
\label{c13:sec:nature-boundaries}

\begin{lemma}[Scientific boundaries are not failures]
\label{c13:lem:boundaries-not-failures}
A \emph{scientific boundary} is a limit of a framework beyond which the framework does not
claim to derive: either an irreducible primitive (an ontological boundary), a numerical
constant not derivable inside the framework (a numerical boundary), or a hosted structure
(an external component carried without derivation). A documented boundary is not a
derivation failure: it is the explicit, honest limit of the framework, disclosed so that no
claim exceeds what is established.
\end{lemma}

\begin{proof}
A framework is characterized as much by what it does not claim as by what it derives: the
disclosure of a boundary is the statement "this framework does not derive X". The
post-frontier audit \cite{c13:QG299} classifies the remaining items as open, partial,
boundary, or methodology, distinguishing the boundaries from the open derivation items.
Hence a boundary is an honest limit, not a failure of the derivation.
\end{proof}

\section{Difference as an Ontological Boundary}
\label{c13:sec:difference-boundary}

Difference is the ontological boundary of the theory: the irreducible primitive whose
count-conservation law is its definitional identity, and which no deeper notion explains
--- the \emph{primitive} boundary, the irreducible starting point rather than an unexplained
gap. The fundamental-boundary audit \cite{c13:QG278} establishes that every attempt to reduce
it reintroduces it (two distinct things use distinct = difference; a boundary is where
things differ; a transition is a before/after difference); by
Chapter~\ref{c01:chap:difference} (Theorem~\ref{c01:thm:fundamental-boundary}) it is the
fundamental boundary of the theory, and by Lemma~\ref{c13:lem:boundaries-not-failures} an
ontological boundary is an honest primitive limit, not a derivation failure.
\label{c13:thm:difference-ontological}

\section{The Status of \texorpdfstring{$\psi$}{psi}}
\label{c13:sec:psi-status}

\begin{theorem}[\texorpdfstring{$\psi$}{psi} is the tensor face of Difference]
\label{c13:thm:psi-face}
The tensor field $\psi$ is the tensor face of Difference: the traceless part of the
trace/traceless decomposition of the one Difference, read against the tensor reference
$\eta$. Its deepest ontological nature --- whether the tensor content is ultimately
irreducible to the scalar count --- is unresolved: a boundary question, not a derivation gap.
\end{theorem}

\begin{proof}
The Difference duality \cite{c13:QG286} establishes that $\rho$ and $\psi$ are the trace and
traceless faces of the one Difference (Chapter~\ref{c01:chap:difference},
Theorem~\ref{c01:thm:duality}); the tensor face is read against the reference $\eta$
(Chapter~\ref{c02:chap:eta}, Theorem~\ref{c02:thm:tensor-read}). Hence $\psi$ is the tensor
face of Difference --- not an independent primitive. The post-frontier audit
\cite{c13:QG299} classifies the $\psi$ question as a boundary item: the observables of the
tensor face are derived, but its ultimate ontological status is a boundary question.
\end{proof}

\section{The \texorpdfstring{$\pi$}{pi} Boundary}
\label{c13:sec:pi-boundary}

\begin{theorem}[\texorpdfstring{$\pi$}{pi} is redundant as a theory input]
\label{c13:thm:pi-redundant}
\emph{$\pi$} is the universal mathematical constant, the geometric ratio of the circle; in
the canonical architecture it is a \emph{numerical boundary} --- a constant not derivable
inside the framework --- and not a primitive. Correspondingly, $\pi$ is redundant as a
theory input: no derived prediction uses $\pi$ as an input, and where it appears it is an
imported numerical factor.
\end{theorem}

\begin{proof}
The framework-necessity audit \cite{c13:QG291} establishes the classification: every derived
observable --- masses, couplings, mixings, cosmological fractions, spectral indices,
acoustic ratios, and the pre-registered predictions --- is a pure ratio of D96 spectral
constants together with the calibration scale. $\pi$ appears only where an imported quantum
factor is required, most notably in the Bekenstein coefficient; as a constant not derivable
inside the framework it is a numerical boundary, distinct from the primitive reference
$\eta$ (Chapter~\ref{c02:chap:eta}, Theorem~\ref{c02:thm:eta-primitive-pi-boundary}), not a
primitive.
\end{proof}

\section{The Bekenstein \texorpdfstring{$2\pi$}{2pi} Boundary}
\label{c13:sec:bekenstein}

\begin{theorem}[The Bekenstein $1/4$ coefficient requires the imported $2\pi$ factor]
\label{c13:thm:bekenstein-2pi}
The exact Bekenstein coefficient $S = A/4$ is a boundary: the structure $S \propto A$,
$M \propto R$, $T \propto 1/R$ is derived, but the exact $1/4$ coefficient requires the
imported quantum factor $T = \kappa/(2\pi)$, which the framework cannot produce internally.
The boundary is therefore explicitly referenced to the imported quantum factor, not
presented as a derivation of the theory.
\end{theorem}

\begin{proof}
The Bekenstein origin \cite{c13:QG185} derives the structure ($S \propto A$, $M \propto R$,
$T \propto 1/R$); the deficit first-law gives $S = A/(8\pi)$, not $1/4$, and the exact $1/4$
requires the $2\pi$ quantum factor $T = \kappa/(2\pi)$, absent in the framework. The
impossibility proof \cite{c13:QG196} confirms that the exact $1/4$ cannot be derived without
importing $\pi$. Consistent with Chapter~\ref{c10:chap:gravity}
(Theorem~\ref{c10:thm:bekenstein-boundary}), the boundary is referenced, not claimed.
\end{proof}

\section{Hosted Standard Model Dynamics}
\label{c13:sec:hosted-sm}

\begin{theorem}[The SM Lagrangian is hosted, not derived]
\label{c13:thm:sm-hosted}
The standard-model \emph{Lagrangian} and its dynamics are hosted: the framework derives the
observables (masses, couplings, mixings) but carries the Lagrangian form as an external
structure. The hosted SM dynamics (the Lagrangian and its structure) is thus a boundary,
while the derived SM observables (the values read from the D96 moments) are derived physics.
\end{theorem}

\begin{proof}
The SM-dynamics audits \cite{c13:QG242,c13:QG245} establish the classification: the gauge
\emph{symmetry} is derived, but the gauge \emph{dynamics} (the Lagrangian, vertices,
propagators) is hosted; the SM dynamics is not complete as a derivation. The derived
observables --- the fermion sectors, couplings, mixings, weak scale, and precision
predictions (Chapter~\ref{c11:chap:sm}) --- are reads of the D96 moments. Hence the SM
dynamics is a boundary and the observables are derived physics.
\end{proof}

\section{Derived Physics vs Boundary Questions}
\label{c13:sec:derived-vs-boundary}

\begin{theorem}[The classification of remaining items]
\label{c13:thm:classification}
The remaining items of the framework classify into derived physics, experimental
validation, and boundaries: the derived physics is established; the validation items
($S,T,U$, $a_e$, $0\nu\beta\beta$) are derived predictions awaiting experiment; the
boundaries are the ontological, numerical, and hosted limits.
\end{theorem}

\begin{proof}
The post-frontier audit \cite{c13:QG299} classifies the remaining items; VALID001
\cite{c13:VALID001} establishes that $S,T,U$, $a_e$, and $0\nu\beta\beta$ are \emph{derived
predictions} (status A) belonging to the experimental-validation category, with no missing
derivation steps. Hence none is an open derivation issue or a boundary; the boundaries are
those documented in this chapter, and the remaining items partition into derived physics
(established), validation (derived, awaiting experiment), and boundaries (the limits).
\end{proof}

\begin{corollary}[No validation item is an open derivation issue]
\label{c13:cor:no-open-validation}
The quantities $S,T,U$, the electron g-2 $a_e$, and the Majorana character $0\nu\beta\beta$
are not listed as open derivation issues: their derivations are complete, and their
experimental validation is the open task.
\end{corollary}

\section{What Is Not Claimed}
\label{c13:sec:not-claimed}

\begin{theorem}[The framework's non-claims]
\label{c13:thm:non-claims}
The framework does not claim: (i) a derivation of the standard-model Lagrangian (hosted);
(ii) a derivation of the exact Bekenstein $1/4$ coefficient (requires the imported $2\pi$
factor); (iii) a resolution of the ultimate ontological status of $\psi$; (iv) any result
beyond the accepted derivations. Equivalently, the framework claims exactly what it
derives: the physics observables established in the preceding chapters, with the boundaries
disclosed and the validation items awaiting experiment.
\end{theorem}

\begin{proof}
(i) The SM dynamics is hosted \cite{c13:QG242,c13:QG245} (Theorem~\ref{c13:thm:sm-hosted}). (ii) The
Bekenstein $1/4$ coefficient requires the imported $2\pi$ factor \cite{c13:QG185,c13:QG196}
(Theorem~\ref{c13:thm:bekenstein-2pi}). (iii) The ontological status of $\psi$ is unresolved
(Theorem~\ref{c13:thm:psi-face}). (iv) The monograph is assembly-only, restricted to
accepted results. Hence the non-claims are explicit and the claimed content is the derived
physics.
\end{proof}

\section{Consequences for Future Work}
\label{c13:sec:future-work}

\begin{theorem}[The boundaries delimit the future work]
\label{c13:thm:future-work}
The boundaries delimit the future work of the framework: the experimental validation of the
derived predictions, the resolution of the ontological status of $\psi$, and the question
of whether the imported $2\pi$ factor can be internalized. The boundaries are not blockers:
the derived content of the theory is complete within its scope, and the boundaries mark
where the framework honestly stops.
\end{theorem}

\begin{proof}
The derived predictions ($S,T,U$, $a_e$, $0\nu\beta\beta$) await experimental validation
\cite{c13:VALID001}; the ontological status of $\psi$ is a boundary question
(Theorem~\ref{c13:thm:psi-face}); the Bekenstein $2\pi$ factor is an imported boundary
(Theorem~\ref{c13:thm:bekenstein-2pi}). By Lemma~\ref{c13:lem:boundaries-not-failures} the
boundaries are not failures, and the derived physics is complete within the canonical
hierarchy (Chapter~\ref{c12:chap:cosmology}, Corollary~\ref{c12:cor:canonical-cosmology}).
Hence the boundaries are not blockers but the framework's declared edge.
\end{proof}

\section{Summary}
\label{c13:sec:boundaries-summary}

This chapter has documented the boundaries of the theory and shown that they are honest
limits, not derivation failures. Difference is the irreducible primitive boundary; $\psi$ is
the tensor face of Difference with unresolved ontological status; $\pi$ is a numerical
boundary, redundant as a theory input; the Bekenstein $1/4$ coefficient is explicitly tied
to the imported $2\pi$ factor; and the SM Lagrangian is hosted while the observables are
derived. The remaining items partition into derived physics, validation, and boundary, with
the validation items ($S,T,U$, $a_e$, $0\nu\beta\beta$) distinct from open derivation issues
(Theorem~\ref{c13:thm:classification}, Corollary~\ref{c13:cor:no-open-validation}); the
framework claims exactly its derivations (Theorem~\ref{c13:thm:non-claims}); and the
boundaries delimit the honest future work without blocking it
(Theorem~\ref{c13:thm:future-work}).

The boundaries of The Actualization Theory are therefore: Difference (the ontological
boundary and primitive), $\psi$ (the tensor face with unresolved ontological status),
$\pi$ (the numerical boundary), the Bekenstein $2\pi$ factor (the imported quantum
boundary), and the hosted standard-model dynamics --- each an explicit limit, none a
derivation failure. The derived physics is complete within its scope; the experimental
validation of the derived predictions is the open task; and the boundaries mark where the
framework honestly stops. This closes the Boundary Layer; the following chapter documents
the frontier and the falsification paths.

\begin{thebibliography}{99}
\bibitem{c13:QG185} AT-QG Phase 185, \emph{Bekenstein Quarter Coefficient Origin}, Docs/Research.
\bibitem{c13:QG196} AT-QG Phase 196, \emph{Quarter Coefficient Impossibility Proof}, Docs/Research.
\bibitem{c13:QG242} AT-QG Phase 242, \emph{Standard Model Dynamics Audit}, Docs/Research.
\bibitem{c13:QG245} AT-QG Phase 245, \emph{Higgs Yukawa Origin Audit}, Docs/Research.
\bibitem{c13:QG278} AT-QG Phase 278, \emph{Fundamental Boundary Audit}, Docs/Research.
\bibitem{c13:QG286} AT-QG Phase 286, \emph{Difference Duality Audit}, Docs/Research.
\bibitem{c13:QG291} AT-QG Phase 291, \emph{Framework Necessity Audit}, Docs/Research.
\bibitem{c13:QG299} AT-QG Phase 299, \emph{Remaining Frontier Re-audit}, Docs/Research.
\bibitem{c13:MONO007} AT-MONO007, \emph{Referee-Readiness Resolution}, Docs/Zenodo/Monograph.
\bibitem{c13:VALID001} AT-VALID001, \emph{Remaining Untested Items Audit}, Docs/Research.
\end{thebibliography}
 after the
%         preamble that defines the theorem environments (or include via the
%         master file).
% ======================================================================

% ---- Theorem environments (define in the master preamble if not present) ----
% \newtheorem{definition}{Definition}[chapter]
% \newtheorem{lemma}[definition]{Lemma}
% \newtheorem{theorem}[definition]{Theorem}
% \newtheorem{corollary}[definition]{Corollary}
% \newtheorem{remark}[definition]{Remark}

\chapter{Boundaries of The Actualization Theory}
\label{c13:chap:boundaries}

\section{Introduction}
\label{c13:sec:boundaries-introduction}

The preceding chapters derived the physics of The Actualization Theory: the foundation, the
spectrum, quantum mechanics, gravity, the standard model, and cosmology. This chapter is the
Boundary Layer of the monograph: it separates \emph{ontological boundaries} from \emph{physics
derivation}, documents the limits of the framework, and shows that the boundaries are
\emph{not derivation failures} --- they are the explicit limits beyond which the framework
does not claim to derive.

The boundaries documented here are: Difference as an ontological boundary \cite{c13:QG278}; the
status of the tensor face $\psi$ \cite{c13:QG286}; the numerical boundary $\pi$ \cite{c13:QG291};
the Bekenstein $2\pi$ boundary \cite{c13:QG185,c13:QG196}; and the hosted standard-model dynamics
\cite{c13:QG242,c13:QG245}. The chapter draws the distinction between \emph{derived physics} and
\emph{boundary questions} \cite{c13:QG299} and states explicitly what the theory does not claim.

Following VALID001 \cite{c13:VALID001}, \emph{experimental validation is not a boundary}: the
derived predictions $S,T,U$, the electron g-2 $a_e$, and the Majorana character $0\nu\beta\beta$
await experiment and are not listed as derivation gaps. Consistent with MONO007 \cite{c13:MONO007},
the presentation is referee-safe and introduces no new physics and no speculation.

\section{The Nature of Scientific Boundaries}
\label{c13:sec:nature-boundaries}

\begin{lemma}[Scientific boundaries are not failures]
\label{c13:lem:boundaries-not-failures}
A \emph{scientific boundary} is a limit of a framework beyond which the framework does not
claim to derive: either an irreducible primitive (an ontological boundary), a numerical
constant not derivable inside the framework (a numerical boundary), or a hosted structure
(an external component carried without derivation). A documented boundary is not a
derivation failure: it is the explicit, honest limit of the framework, disclosed so that no
claim exceeds what is established.
\end{lemma}

\begin{proof}
A framework is characterized as much by what it does not claim as by what it derives: the
disclosure of a boundary is the statement "this framework does not derive X". The
post-frontier audit \cite{c13:QG299} classifies the remaining items as open, partial,
boundary, or methodology, distinguishing the boundaries from the open derivation items.
Hence a boundary is an honest limit, not a failure of the derivation.
\end{proof}

\section{Difference as an Ontological Boundary}
\label{c13:sec:difference-boundary}

Difference is the ontological boundary of the theory: the irreducible primitive whose
count-conservation law is its definitional identity, and which no deeper notion explains
--- the \emph{primitive} boundary, the irreducible starting point rather than an unexplained
gap. The fundamental-boundary audit \cite{c13:QG278} establishes that every attempt to reduce
it reintroduces it (two distinct things use distinct = difference; a boundary is where
things differ; a transition is a before/after difference); by
Chapter~\ref{c01:chap:difference} (Theorem~\ref{c01:thm:fundamental-boundary}) it is the
fundamental boundary of the theory, and by Lemma~\ref{c13:lem:boundaries-not-failures} an
ontological boundary is an honest primitive limit, not a derivation failure.
\label{c13:thm:difference-ontological}

\section{The Status of \texorpdfstring{$\psi$}{psi}}
\label{c13:sec:psi-status}

\begin{theorem}[\texorpdfstring{$\psi$}{psi} is the tensor face of Difference]
\label{c13:thm:psi-face}
The tensor field $\psi$ is the tensor face of Difference: the traceless part of the
trace/traceless decomposition of the one Difference, read against the tensor reference
$\eta$. Its deepest ontological nature --- whether the tensor content is ultimately
irreducible to the scalar count --- is unresolved: a boundary question, not a derivation gap.
\end{theorem}

\begin{proof}
The Difference duality \cite{c13:QG286} establishes that $\rho$ and $\psi$ are the trace and
traceless faces of the one Difference (Chapter~\ref{c01:chap:difference},
Theorem~\ref{c01:thm:duality}); the tensor face is read against the reference $\eta$
(Chapter~\ref{c02:chap:eta}, Theorem~\ref{c02:thm:tensor-read}). Hence $\psi$ is the tensor
face of Difference --- not an independent primitive. The post-frontier audit
\cite{c13:QG299} classifies the $\psi$ question as a boundary item: the observables of the
tensor face are derived, but its ultimate ontological status is a boundary question.
\end{proof}

\section{The \texorpdfstring{$\pi$}{pi} Boundary}
\label{c13:sec:pi-boundary}

\begin{theorem}[\texorpdfstring{$\pi$}{pi} is redundant as a theory input]
\label{c13:thm:pi-redundant}
\emph{$\pi$} is the universal mathematical constant, the geometric ratio of the circle; in
the canonical architecture it is a \emph{numerical boundary} --- a constant not derivable
inside the framework --- and not a primitive. Correspondingly, $\pi$ is redundant as a
theory input: no derived prediction uses $\pi$ as an input, and where it appears it is an
imported numerical factor.
\end{theorem}

\begin{proof}
The framework-necessity audit \cite{c13:QG291} establishes the classification: every derived
observable --- masses, couplings, mixings, cosmological fractions, spectral indices,
acoustic ratios, and the pre-registered predictions --- is a pure ratio of D96 spectral
constants together with the calibration scale. $\pi$ appears only where an imported quantum
factor is required, most notably in the Bekenstein coefficient; as a constant not derivable
inside the framework it is a numerical boundary, distinct from the primitive reference
$\eta$ (Chapter~\ref{c02:chap:eta}, Theorem~\ref{c02:thm:eta-primitive-pi-boundary}), not a
primitive.
\end{proof}

\section{The Bekenstein \texorpdfstring{$2\pi$}{2pi} Boundary}
\label{c13:sec:bekenstein}

\begin{theorem}[The Bekenstein $1/4$ coefficient requires the imported $2\pi$ factor]
\label{c13:thm:bekenstein-2pi}
The exact Bekenstein coefficient $S = A/4$ is a boundary: the structure $S \propto A$,
$M \propto R$, $T \propto 1/R$ is derived, but the exact $1/4$ coefficient requires the
imported quantum factor $T = \kappa/(2\pi)$, which the framework cannot produce internally.
The boundary is therefore explicitly referenced to the imported quantum factor, not
presented as a derivation of the theory.
\end{theorem}

\begin{proof}
The Bekenstein origin \cite{c13:QG185} derives the structure ($S \propto A$, $M \propto R$,
$T \propto 1/R$); the deficit first-law gives $S = A/(8\pi)$, not $1/4$, and the exact $1/4$
requires the $2\pi$ quantum factor $T = \kappa/(2\pi)$, absent in the framework. The
impossibility proof \cite{c13:QG196} confirms that the exact $1/4$ cannot be derived without
importing $\pi$. Consistent with Chapter~\ref{c10:chap:gravity}
(Theorem~\ref{c10:thm:bekenstein-boundary}), the boundary is referenced, not claimed.
\end{proof}

\section{Hosted Standard Model Dynamics}
\label{c13:sec:hosted-sm}

\begin{theorem}[The SM Lagrangian is hosted, not derived]
\label{c13:thm:sm-hosted}
The standard-model \emph{Lagrangian} and its dynamics are hosted: the framework derives the
observables (masses, couplings, mixings) but carries the Lagrangian form as an external
structure. The hosted SM dynamics (the Lagrangian and its structure) is thus a boundary,
while the derived SM observables (the values read from the D96 moments) are derived physics.
\end{theorem}

\begin{proof}
The SM-dynamics audits \cite{c13:QG242,c13:QG245} establish the classification: the gauge
\emph{symmetry} is derived, but the gauge \emph{dynamics} (the Lagrangian, vertices,
propagators) is hosted; the SM dynamics is not complete as a derivation. The derived
observables --- the fermion sectors, couplings, mixings, weak scale, and precision
predictions (Chapter~\ref{c11:chap:sm}) --- are reads of the D96 moments. Hence the SM
dynamics is a boundary and the observables are derived physics.
\end{proof}

\section{Derived Physics vs Boundary Questions}
\label{c13:sec:derived-vs-boundary}

\begin{theorem}[The classification of remaining items]
\label{c13:thm:classification}
The remaining items of the framework classify into derived physics, experimental
validation, and boundaries: the derived physics is established; the validation items
($S,T,U$, $a_e$, $0\nu\beta\beta$) are derived predictions awaiting experiment; the
boundaries are the ontological, numerical, and hosted limits.
\end{theorem}

\begin{proof}
The post-frontier audit \cite{c13:QG299} classifies the remaining items; VALID001
\cite{c13:VALID001} establishes that $S,T,U$, $a_e$, and $0\nu\beta\beta$ are \emph{derived
predictions} (status A) belonging to the experimental-validation category, with no missing
derivation steps. Hence none is an open derivation issue or a boundary; the boundaries are
those documented in this chapter, and the remaining items partition into derived physics
(established), validation (derived, awaiting experiment), and boundaries (the limits).
\end{proof}

\begin{corollary}[No validation item is an open derivation issue]
\label{c13:cor:no-open-validation}
The quantities $S,T,U$, the electron g-2 $a_e$, and the Majorana character $0\nu\beta\beta$
are not listed as open derivation issues: their derivations are complete, and their
experimental validation is the open task.
\end{corollary}

\section{What Is Not Claimed}
\label{c13:sec:not-claimed}

\begin{theorem}[The framework's non-claims]
\label{c13:thm:non-claims}
The framework does not claim: (i) a derivation of the standard-model Lagrangian (hosted);
(ii) a derivation of the exact Bekenstein $1/4$ coefficient (requires the imported $2\pi$
factor); (iii) a resolution of the ultimate ontological status of $\psi$; (iv) any result
beyond the accepted derivations. Equivalently, the framework claims exactly what it
derives: the physics observables established in the preceding chapters, with the boundaries
disclosed and the validation items awaiting experiment.
\end{theorem}

\begin{proof}
(i) The SM dynamics is hosted \cite{c13:QG242,c13:QG245} (Theorem~\ref{c13:thm:sm-hosted}). (ii) The
Bekenstein $1/4$ coefficient requires the imported $2\pi$ factor \cite{c13:QG185,c13:QG196}
(Theorem~\ref{c13:thm:bekenstein-2pi}). (iii) The ontological status of $\psi$ is unresolved
(Theorem~\ref{c13:thm:psi-face}). (iv) The monograph is assembly-only, restricted to
accepted results. Hence the non-claims are explicit and the claimed content is the derived
physics.
\end{proof}

\section{Consequences for Future Work}
\label{c13:sec:future-work}

\begin{theorem}[The boundaries delimit the future work]
\label{c13:thm:future-work}
The boundaries delimit the future work of the framework: the experimental validation of the
derived predictions, the resolution of the ontological status of $\psi$, and the question
of whether the imported $2\pi$ factor can be internalized. The boundaries are not blockers:
the derived content of the theory is complete within its scope, and the boundaries mark
where the framework honestly stops.
\end{theorem}

\begin{proof}
The derived predictions ($S,T,U$, $a_e$, $0\nu\beta\beta$) await experimental validation
\cite{c13:VALID001}; the ontological status of $\psi$ is a boundary question
(Theorem~\ref{c13:thm:psi-face}); the Bekenstein $2\pi$ factor is an imported boundary
(Theorem~\ref{c13:thm:bekenstein-2pi}). By Lemma~\ref{c13:lem:boundaries-not-failures} the
boundaries are not failures, and the derived physics is complete within the canonical
hierarchy (Chapter~\ref{c12:chap:cosmology}, Corollary~\ref{c12:cor:canonical-cosmology}).
Hence the boundaries are not blockers but the framework's declared edge.
\end{proof}

\section{Summary}
\label{c13:sec:boundaries-summary}

This chapter has documented the boundaries of the theory and shown that they are honest
limits, not derivation failures. Difference is the irreducible primitive boundary; $\psi$ is
the tensor face of Difference with unresolved ontological status; $\pi$ is a numerical
boundary, redundant as a theory input; the Bekenstein $1/4$ coefficient is explicitly tied
to the imported $2\pi$ factor; and the SM Lagrangian is hosted while the observables are
derived. The remaining items partition into derived physics, validation, and boundary, with
the validation items ($S,T,U$, $a_e$, $0\nu\beta\beta$) distinct from open derivation issues
(Theorem~\ref{c13:thm:classification}, Corollary~\ref{c13:cor:no-open-validation}); the
framework claims exactly its derivations (Theorem~\ref{c13:thm:non-claims}); and the
boundaries delimit the honest future work without blocking it
(Theorem~\ref{c13:thm:future-work}).

The boundaries of The Actualization Theory are therefore: Difference (the ontological
boundary and primitive), $\psi$ (the tensor face with unresolved ontological status),
$\pi$ (the numerical boundary), the Bekenstein $2\pi$ factor (the imported quantum
boundary), and the hosted standard-model dynamics --- each an explicit limit, none a
derivation failure. The derived physics is complete within its scope; the experimental
validation of the derived predictions is the open task; and the boundaries mark where the
framework honestly stops. This closes the Boundary Layer; the following chapter documents
the frontier and the falsification paths.

\begin{thebibliography}{99}
\bibitem{c13:QG185} AT-QG Phase 185, \emph{Bekenstein Quarter Coefficient Origin}, Docs/Research.
\bibitem{c13:QG196} AT-QG Phase 196, \emph{Quarter Coefficient Impossibility Proof}, Docs/Research.
\bibitem{c13:QG242} AT-QG Phase 242, \emph{Standard Model Dynamics Audit}, Docs/Research.
\bibitem{c13:QG245} AT-QG Phase 245, \emph{Higgs Yukawa Origin Audit}, Docs/Research.
\bibitem{c13:QG278} AT-QG Phase 278, \emph{Fundamental Boundary Audit}, Docs/Research.
\bibitem{c13:QG286} AT-QG Phase 286, \emph{Difference Duality Audit}, Docs/Research.
\bibitem{c13:QG291} AT-QG Phase 291, \emph{Framework Necessity Audit}, Docs/Research.
\bibitem{c13:QG299} AT-QG Phase 299, \emph{Remaining Frontier Re-audit}, Docs/Research.
\bibitem{c13:MONO007} AT-MONO007, \emph{Referee-Readiness Resolution}, Docs/Zenodo/Monograph.
\bibitem{c13:VALID001} AT-VALID001, \emph{Remaining Untested Items Audit}, Docs/Research.
\end{thebibliography}
 after the
%         preamble that defines the theorem environments (or include via the
%         master file).
% ======================================================================

% ---- Theorem environments (define in the master preamble if not present) ----
% \newtheorem{definition}{Definition}[chapter]
% \newtheorem{lemma}[definition]{Lemma}
% \newtheorem{theorem}[definition]{Theorem}
% \newtheorem{corollary}[definition]{Corollary}
% \newtheorem{remark}[definition]{Remark}

\chapter{Boundaries of The Actualization Theory}
\label{c13:chap:boundaries}

\section{Introduction}
\label{c13:sec:boundaries-introduction}

The preceding chapters derived the physics of The Actualization Theory: the foundation, the
spectrum, quantum mechanics, gravity, the standard model, and cosmology. This chapter is the
Boundary Layer of the monograph: it separates \emph{ontological boundaries} from \emph{physics
derivation}, documents the limits of the framework, and shows that the boundaries are
\emph{not derivation failures} --- they are the explicit limits beyond which the framework
does not claim to derive.

The boundaries documented here are: Difference as an ontological boundary \cite{c13:QG278}; the
status of the tensor face $\psi$ \cite{c13:QG286}; the numerical boundary $\pi$ \cite{c13:QG291};
the Bekenstein $2\pi$ boundary \cite{c13:QG185,c13:QG196}; and the hosted standard-model dynamics
\cite{c13:QG242,c13:QG245}. The chapter draws the distinction between \emph{derived physics} and
\emph{boundary questions} \cite{c13:QG299} and states explicitly what the theory does not claim.

Following VALID001 \cite{c13:VALID001}, \emph{experimental validation is not a boundary}: the
derived predictions $S,T,U$, the electron g-2 $a_e$, and the Majorana character $0\nu\beta\beta$
await experiment and are not listed as derivation gaps. Consistent with MONO007 \cite{c13:MONO007},
the presentation is referee-safe and introduces no new physics and no speculation.

\section{The Nature of Scientific Boundaries}
\label{c13:sec:nature-boundaries}

\begin{lemma}[Scientific boundaries are not failures]
\label{c13:lem:boundaries-not-failures}
A \emph{scientific boundary} is a limit of a framework beyond which the framework does not
claim to derive: either an irreducible primitive (an ontological boundary), a numerical
constant not derivable inside the framework (a numerical boundary), or a hosted structure
(an external component carried without derivation). A documented boundary is not a
derivation failure: it is the explicit, honest limit of the framework, disclosed so that no
claim exceeds what is established.
\end{lemma}

\begin{proof}
A framework is characterized as much by what it does not claim as by what it derives: the
disclosure of a boundary is the statement "this framework does not derive X". The
post-frontier audit \cite{c13:QG299} classifies the remaining items as open, partial,
boundary, or methodology, distinguishing the boundaries from the open derivation items.
Hence a boundary is an honest limit, not a failure of the derivation.
\end{proof}

\section{Difference as an Ontological Boundary}
\label{c13:sec:difference-boundary}

Difference is the ontological boundary of the theory: the irreducible primitive whose
count-conservation law is its definitional identity, and which no deeper notion explains
--- the \emph{primitive} boundary, the irreducible starting point rather than an unexplained
gap. The fundamental-boundary audit \cite{c13:QG278} establishes that every attempt to reduce
it reintroduces it (two distinct things use distinct = difference; a boundary is where
things differ; a transition is a before/after difference); by
Chapter~\ref{c01:chap:difference} (Theorem~\ref{c01:thm:fundamental-boundary}) it is the
fundamental boundary of the theory, and by Lemma~\ref{c13:lem:boundaries-not-failures} an
ontological boundary is an honest primitive limit, not a derivation failure.
\label{c13:thm:difference-ontological}

\section{The Status of \texorpdfstring{$\psi$}{psi}}
\label{c13:sec:psi-status}

\begin{theorem}[\texorpdfstring{$\psi$}{psi} is the tensor face of Difference]
\label{c13:thm:psi-face}
The tensor field $\psi$ is the tensor face of Difference: the traceless part of the
trace/traceless decomposition of the one Difference, read against the tensor reference
$\eta$. Its deepest ontological nature --- whether the tensor content is ultimately
irreducible to the scalar count --- is unresolved: a boundary question, not a derivation gap.
\end{theorem}

\begin{proof}
The Difference duality \cite{c13:QG286} establishes that $\rho$ and $\psi$ are the trace and
traceless faces of the one Difference (Chapter~\ref{c01:chap:difference},
Theorem~\ref{c01:thm:duality}); the tensor face is read against the reference $\eta$
(Chapter~\ref{c02:chap:eta}, Theorem~\ref{c02:thm:tensor-read}). Hence $\psi$ is the tensor
face of Difference --- not an independent primitive. The post-frontier audit
\cite{c13:QG299} classifies the $\psi$ question as a boundary item: the observables of the
tensor face are derived, but its ultimate ontological status is a boundary question.
\end{proof}

\section{The \texorpdfstring{$\pi$}{pi} Boundary}
\label{c13:sec:pi-boundary}

\begin{theorem}[\texorpdfstring{$\pi$}{pi} is redundant as a theory input]
\label{c13:thm:pi-redundant}
\emph{$\pi$} is the universal mathematical constant, the geometric ratio of the circle; in
the canonical architecture it is a \emph{numerical boundary} --- a constant not derivable
inside the framework --- and not a primitive. Correspondingly, $\pi$ is redundant as a
theory input: no derived prediction uses $\pi$ as an input, and where it appears it is an
imported numerical factor.
\end{theorem}

\begin{proof}
The framework-necessity audit \cite{c13:QG291} establishes the classification: every derived
observable --- masses, couplings, mixings, cosmological fractions, spectral indices,
acoustic ratios, and the pre-registered predictions --- is a pure ratio of D96 spectral
constants together with the calibration scale. $\pi$ appears only where an imported quantum
factor is required, most notably in the Bekenstein coefficient; as a constant not derivable
inside the framework it is a numerical boundary, distinct from the primitive reference
$\eta$ (Chapter~\ref{c02:chap:eta}, Theorem~\ref{c02:thm:eta-primitive-pi-boundary}), not a
primitive.
\end{proof}

\section{The Bekenstein \texorpdfstring{$2\pi$}{2pi} Boundary}
\label{c13:sec:bekenstein}

\begin{theorem}[The Bekenstein $1/4$ coefficient requires the imported $2\pi$ factor]
\label{c13:thm:bekenstein-2pi}
The exact Bekenstein coefficient $S = A/4$ is a boundary: the structure $S \propto A$,
$M \propto R$, $T \propto 1/R$ is derived, but the exact $1/4$ coefficient requires the
imported quantum factor $T = \kappa/(2\pi)$, which the framework cannot produce internally.
The boundary is therefore explicitly referenced to the imported quantum factor, not
presented as a derivation of the theory.
\end{theorem}

\begin{proof}
The Bekenstein origin \cite{c13:QG185} derives the structure ($S \propto A$, $M \propto R$,
$T \propto 1/R$); the deficit first-law gives $S = A/(8\pi)$, not $1/4$, and the exact $1/4$
requires the $2\pi$ quantum factor $T = \kappa/(2\pi)$, absent in the framework. The
impossibility proof \cite{c13:QG196} confirms that the exact $1/4$ cannot be derived without
importing $\pi$. Consistent with Chapter~\ref{c10:chap:gravity}
(Theorem~\ref{c10:thm:bekenstein-boundary}), the boundary is referenced, not claimed.
\end{proof}

\section{Hosted Standard Model Dynamics}
\label{c13:sec:hosted-sm}

\begin{theorem}[The SM Lagrangian is hosted, not derived]
\label{c13:thm:sm-hosted}
The standard-model \emph{Lagrangian} and its dynamics are hosted: the framework derives the
observables (masses, couplings, mixings) but carries the Lagrangian form as an external
structure. The hosted SM dynamics (the Lagrangian and its structure) is thus a boundary,
while the derived SM observables (the values read from the D96 moments) are derived physics.
\end{theorem}

\begin{proof}
The SM-dynamics audits \cite{c13:QG242,c13:QG245} establish the classification: the gauge
\emph{symmetry} is derived, but the gauge \emph{dynamics} (the Lagrangian, vertices,
propagators) is hosted; the SM dynamics is not complete as a derivation. The derived
observables --- the fermion sectors, couplings, mixings, weak scale, and precision
predictions (Chapter~\ref{c11:chap:sm}) --- are reads of the D96 moments. Hence the SM
dynamics is a boundary and the observables are derived physics.
\end{proof}

\section{Derived Physics vs Boundary Questions}
\label{c13:sec:derived-vs-boundary}

\begin{theorem}[The classification of remaining items]
\label{c13:thm:classification}
The remaining items of the framework classify into derived physics, experimental
validation, and boundaries: the derived physics is established; the validation items
($S,T,U$, $a_e$, $0\nu\beta\beta$) are derived predictions awaiting experiment; the
boundaries are the ontological, numerical, and hosted limits.
\end{theorem}

\begin{proof}
The post-frontier audit \cite{c13:QG299} classifies the remaining items; VALID001
\cite{c13:VALID001} establishes that $S,T,U$, $a_e$, and $0\nu\beta\beta$ are \emph{derived
predictions} (status A) belonging to the experimental-validation category, with no missing
derivation steps. Hence none is an open derivation issue or a boundary; the boundaries are
those documented in this chapter, and the remaining items partition into derived physics
(established), validation (derived, awaiting experiment), and boundaries (the limits).
\end{proof}

\begin{corollary}[No validation item is an open derivation issue]
\label{c13:cor:no-open-validation}
The quantities $S,T,U$, the electron g-2 $a_e$, and the Majorana character $0\nu\beta\beta$
are not listed as open derivation issues: their derivations are complete, and their
experimental validation is the open task.
\end{corollary}

\section{What Is Not Claimed}
\label{c13:sec:not-claimed}

\begin{theorem}[The framework's non-claims]
\label{c13:thm:non-claims}
The framework does not claim: (i) a derivation of the standard-model Lagrangian (hosted);
(ii) a derivation of the exact Bekenstein $1/4$ coefficient (requires the imported $2\pi$
factor); (iii) a resolution of the ultimate ontological status of $\psi$; (iv) any result
beyond the accepted derivations. Equivalently, the framework claims exactly what it
derives: the physics observables established in the preceding chapters, with the boundaries
disclosed and the validation items awaiting experiment.
\end{theorem}

\begin{proof}
(i) The SM dynamics is hosted \cite{c13:QG242,c13:QG245} (Theorem~\ref{c13:thm:sm-hosted}). (ii) The
Bekenstein $1/4$ coefficient requires the imported $2\pi$ factor \cite{c13:QG185,c13:QG196}
(Theorem~\ref{c13:thm:bekenstein-2pi}). (iii) The ontological status of $\psi$ is unresolved
(Theorem~\ref{c13:thm:psi-face}). (iv) The monograph is assembly-only, restricted to
accepted results. Hence the non-claims are explicit and the claimed content is the derived
physics.
\end{proof}

\section{Consequences for Future Work}
\label{c13:sec:future-work}

\begin{theorem}[The boundaries delimit the future work]
\label{c13:thm:future-work}
The boundaries delimit the future work of the framework: the experimental validation of the
derived predictions, the resolution of the ontological status of $\psi$, and the question
of whether the imported $2\pi$ factor can be internalized. The boundaries are not blockers:
the derived content of the theory is complete within its scope, and the boundaries mark
where the framework honestly stops.
\end{theorem}

\begin{proof}
The derived predictions ($S,T,U$, $a_e$, $0\nu\beta\beta$) await experimental validation
\cite{c13:VALID001}; the ontological status of $\psi$ is a boundary question
(Theorem~\ref{c13:thm:psi-face}); the Bekenstein $2\pi$ factor is an imported boundary
(Theorem~\ref{c13:thm:bekenstein-2pi}). By Lemma~\ref{c13:lem:boundaries-not-failures} the
boundaries are not failures, and the derived physics is complete within the canonical
hierarchy (Chapter~\ref{c12:chap:cosmology}, Corollary~\ref{c12:cor:canonical-cosmology}).
Hence the boundaries are not blockers but the framework's declared edge.
\end{proof}

\section{Summary}
\label{c13:sec:boundaries-summary}

This chapter has documented the boundaries of the theory and shown that they are honest
limits, not derivation failures. Difference is the irreducible primitive boundary; $\psi$ is
the tensor face of Difference with unresolved ontological status; $\pi$ is a numerical
boundary, redundant as a theory input; the Bekenstein $1/4$ coefficient is explicitly tied
to the imported $2\pi$ factor; and the SM Lagrangian is hosted while the observables are
derived. The remaining items partition into derived physics, validation, and boundary, with
the validation items ($S,T,U$, $a_e$, $0\nu\beta\beta$) distinct from open derivation issues
(Theorem~\ref{c13:thm:classification}, Corollary~\ref{c13:cor:no-open-validation}); the
framework claims exactly its derivations (Theorem~\ref{c13:thm:non-claims}); and the
boundaries delimit the honest future work without blocking it
(Theorem~\ref{c13:thm:future-work}).

The boundaries of The Actualization Theory are therefore: Difference (the ontological
boundary and primitive), $\psi$ (the tensor face with unresolved ontological status),
$\pi$ (the numerical boundary), the Bekenstein $2\pi$ factor (the imported quantum
boundary), and the hosted standard-model dynamics --- each an explicit limit, none a
derivation failure. The derived physics is complete within its scope; the experimental
validation of the derived predictions is the open task; and the boundaries mark where the
framework honestly stops. This closes the Boundary Layer; the following chapter documents
the frontier and the falsification paths.

\begin{thebibliography}{99}
\bibitem{c13:QG185} AT-QG Phase 185, \emph{Bekenstein Quarter Coefficient Origin}, Docs/Research.
\bibitem{c13:QG196} AT-QG Phase 196, \emph{Quarter Coefficient Impossibility Proof}, Docs/Research.
\bibitem{c13:QG242} AT-QG Phase 242, \emph{Standard Model Dynamics Audit}, Docs/Research.
\bibitem{c13:QG245} AT-QG Phase 245, \emph{Higgs Yukawa Origin Audit}, Docs/Research.
\bibitem{c13:QG278} AT-QG Phase 278, \emph{Fundamental Boundary Audit}, Docs/Research.
\bibitem{c13:QG286} AT-QG Phase 286, \emph{Difference Duality Audit}, Docs/Research.
\bibitem{c13:QG291} AT-QG Phase 291, \emph{Framework Necessity Audit}, Docs/Research.
\bibitem{c13:QG299} AT-QG Phase 299, \emph{Remaining Frontier Re-audit}, Docs/Research.
\bibitem{c13:MONO007} AT-MONO007, \emph{Referee-Readiness Resolution}, Docs/Zenodo/Monograph.
\bibitem{c13:VALID001} AT-VALID001, \emph{Remaining Untested Items Audit}, Docs/Research.
\end{thebibliography}

% ======================================================================
%  TQM-MONO114 — Chapter 14: Frontier and Falsification
%  Canonical Monograph (v2.0) · The Actualization Theory:
%  A Reconstruction of Physics from Difference, Actualization and Spectrum
%
%  Status: COMPLETE — PASS-READY (MONO007 referee-ready architecture)
%  Inputs: Chapter01_Difference.tex ... Chapter13_BoundariesOfTheActualizationTheory.tex,
%          VALID001, MONO007, MONO110A
%  Method: assembly only — canonical results only, no new physics, no speculation
%  Citations: VALID001 (validation classification), QG190 (pre-registration),
%             QG199 (P1 evidence), QG200 (sector ladder), QG201 (ladder stats),
%             QG202 (prediction outcome dashboard), QG203 (absolute masses),
%             QG299 (post-frontier)
%
%  Usage: % ======================================================================
%  TQM-MONO114 — Chapter 14: Frontier and Falsification
%  Canonical Monograph (v2.0) · The Actualization Theory:
%  A Reconstruction of Physics from Difference, Actualization and Spectrum
%
%  Status: COMPLETE — PASS-READY (MONO007 referee-ready architecture)
%  Inputs: Chapter01_Difference.tex ... Chapter13_BoundariesOfTheActualizationTheory.tex,
%          VALID001, MONO007, MONO110A
%  Method: assembly only — canonical results only, no new physics, no speculation
%  Citations: VALID001 (validation classification), QG190 (pre-registration),
%             QG199 (P1 evidence), QG200 (sector ladder), QG201 (ladder stats),
%             QG202 (prediction outcome dashboard), QG203 (absolute masses),
%             QG299 (post-frontier)
%
%  Usage: % ======================================================================
%  TQM-MONO114 — Chapter 14: Frontier and Falsification
%  Canonical Monograph (v2.0) · The Actualization Theory:
%  A Reconstruction of Physics from Difference, Actualization and Spectrum
%
%  Status: COMPLETE — PASS-READY (MONO007 referee-ready architecture)
%  Inputs: Chapter01_Difference.tex ... Chapter13_BoundariesOfTheActualizationTheory.tex,
%          VALID001, MONO007, MONO110A
%  Method: assembly only — canonical results only, no new physics, no speculation
%  Citations: VALID001 (validation classification), QG190 (pre-registration),
%             QG199 (P1 evidence), QG200 (sector ladder), QG201 (ladder stats),
%             QG202 (prediction outcome dashboard), QG203 (absolute masses),
%             QG299 (post-frontier)
%
%  Usage: \input{Chapter14_FrontierAndFalsification.tex} after the preamble
%         that defines the theorem environments (or include via the master file).
% ======================================================================

% ---- Theorem environments (define in the master preamble if not present) ----
% \newtheorem{definition}{Definition}[chapter]
% \newtheorem{lemma}[definition]{Lemma}
% \newtheorem{theorem}[definition]{Theorem}
% \newtheorem{corollary}[definition]{Corollary}
% \newtheorem{remark}[definition]{Remark}

\chapter{Frontier and Falsification}
\label{chap:frontier}

\section{Introduction}
\label{sec:frontier-introduction}

The preceding chapters derived the physics of The Actualization Theory and documented its
boundaries. This chapter closes the core monograph by addressing the \emph{frontier}: the
experimental validation of the derived predictions, the independent temporal evidence, the
explicit failure conditions, and the scope of validity. The chapter's purpose is to make the
theory \emph{publication-safe} --- to state clearly what would confirm it, what would
falsify it, and what it does and does not claim.

The central content is: the scientific-falsifiability requirement \cite{QG299}; the
distinction between derivation and validation; the current validation targets --- the
oblique parameters $S,T,U$, the electron g-2 $a_e$, the Majorana character $0\nu\beta\beta$
\cite{VALID001}, and the pre-registered predictions P1 (106 GeV), P2 ($0\nu\beta\beta$), and the P3
sector ladder \cite{QG190,QG199,QG200,QG201,QG202,QG203}; the independent temporal evidence;
the explicit failure conditions; the scope of validity; and the future experimental
program.

Following the physics-focused re-audit \cite{MONO110A}, this chapter claims \emph{no open
derivation frontier}: every observable is derived, and the frontier is experimental. The
experimental validation is clearly separated from the derivation, the falsification paths
are explicit, and the wording is referee-safe and publication-safe \cite{MONO007}. We use
only accepted canonical results and introduce no new physics and no speculation.

\section{Scientific Falsifiability}
\label{sec:falsifiability}

\begin{definition}[Falsifiability]
\label{def:falsifiability}
A scientific theory is \emph{falsifiable} if it specifies, in advance, observations that
would contradict it: the theory must expose its failure modes.
\end{definition}

\begin{theorem}[A scientific theory must expose its failure modes]
\label{thm:falsifiability}
The Actualization Theory is falsifiable: it specifies explicit, pre-registered
observations whose absence or contradiction would force a revision.
\end{theorem}

\begin{proof}
The post-frontier audit \cite{QG299} classifies the remaining items into open, partial,
boundary, and methodology, with the pre-registered predictions carrying explicit
falsification conditions. The prediction-outcome dashboard \cite{QG202} records, for each
prediction, its frozen value, its current evidence, and its falsification condition.
Hence the theory exposes its failure modes in advance: it is falsifiable.
\end{proof}

\begin{corollary}[Falsifiability is not a weakness]
\label{cor:falsifiability-strength}
The exposure of failure modes is a strength, not a weakness: it is the property that makes
the theory scientific and its claims testable.
\end{corollary}

\begin{proof}
By Theorem~\ref{thm:falsifiability}, the theory specifies its failure modes; by
Definition~\ref{def:falsifiability}, this is the defining property of a falsifiable theory.
Hence the explicit falsification paths are a scientific strength.
\end{proof}

\section{Experimental Validation Status}
\label{sec:validation-status}

\begin{definition}[Derivation vs validation]
\label{def:derivation-validation}
\emph{Derivation} is the derivation of an observable from the theory's principles; 
\emph{validation} is the confirmation of that derived value by experiment. The two are
distinct: a derived prediction may await validation without any derivation gap.
\end{definition}

\begin{theorem}[Derived and validated are distinct]
\label{thm:derived-vs-validated}
The status of an observable is the pair (derived, validated): a quantity may be derived
without being validated, and its derivation is complete before its validation is possible.
\end{theorem}

\begin{proof}
Following VALID001 \cite{VALID001}, the quantities $S,T,U$, the electron g-2 $a_e$, and the
Majorana character $0\nu\beta\beta$ are status A --- derived prediction exists, with no
missing derivation steps --- and belong to the experimental-validation category. Their
derivation is complete; their validation is the open experimental task. Hence derivation and
validation are distinct, and a derived-but-unvalidated quantity is not an open derivation
issue.
\end{proof}

\begin{corollary}[No open derivation frontier]
\label{cor:no-open-derivation}
The theory claims no open derivation frontier: every observable of the theory is derived,
and the frontier is experimental validation, not derivation.
\end{corollary}

\begin{proof}
By Theorem~\ref{thm:derived-vs-validated}, the derived quantities are complete, with the
validation tasks ahead. The reconstruction audit establishes that the major results are
reconstructed from the minimal hierarchy. Hence there is no open derivation frontier; the
frontier is experimental \cite{MONO110A}.
\end{proof}

\section{Current Validation Targets}
\label{sec:validation-targets}

\begin{theorem}[The precision predictions await validation]
\label{thm:precision-validation}
The oblique parameters $S,T,U$, the electron g-2 $a_e$, and the Majorana character
$0\nu\beta\beta$ are derived predictions awaiting experimental validation.
\end{theorem}

\begin{proof}
Following VALID001 \cite{VALID001}: $S = 0.0421$, $T = 0.0842$, $U = 0$ are derived and
consistent with the EW global fit; $a_e = 1.159655\times10^{-3}$ matches experiment within
$0.0003\%$; $m_{\beta\beta} = 2.02\times10^{-3}$ eV is within limits and in reach of
next-generation experiments. All three are status A, with their experimental validation the
open task. Hence they are current validation targets.
\end{proof}

\begin{theorem}[The pre-registered predictions]
\label{thm:pre-registered}
The theory carries three pre-registered predictions with frozen values and explicit
falsification conditions:
\begin{enumerate}
\item \textbf{P1} --- a 106 GeV resonance: central value $106.39$ GeV, window
$99$--$114$ GeV;
\item \textbf{P2} --- the $0\nu\beta\beta$ effective Majorana mass $m_{\beta\beta} =
2.02$ meV;
\item \textbf{P3} --- the sector-ladder spectrum: nine resonances
$106.39 \to 263.43$ GeV with a $15.20$ GeV width scale.
\end{enumerate}
\end{theorem}

\begin{proof}
The predictions are pre-registered \cite{QG190}, with their values frozen in the immutable
prediction registry and their evidence tracked in the outcome dashboard \cite{QG202}. P1 is
pending (window $99$--$114$ GeV neither confirmed nor excluded) \cite{QG199}; P2 is pending
(below current sensitivity, awaiting nEXO/LEGEND-1000); P3 is supported (the $151.98$ GeV
rung matches a $\sim152$ GeV excess at moderate significance) \cite{QG200,QG201}. The
absolute-mass law \cite{QG203} anchors the sector structure.
\end{proof}

\begin{corollary}[The targets are derived, not open issues]
\label{cor:targets-derived}
The validation targets are derived predictions awaiting experiment; they are not open
derivation issues and not boundaries.
\end{corollary}

\begin{proof}
By Theorem~\ref{thm:precision-validation} and Theorem~\ref{thm:pre-registered}, all
targets are derived predictions with frozen values; by Chapter~\ref{chap:boundaries}
(Corollary~\ref{cor:no-open-validation}), validation items are distinct from open
derivation issues. Hence the targets are validation items.
\end{proof}

\section{Independent Temporal Evidence}
\label{sec:temporal-evidence}

\begin{definition}[Independent temporal evidence]
\label{def:temporal-evidence}
\emph{Independent temporal evidence} is the empirical program of the theory that does not
depend on the already-derived observables: the confirmation of the pre-registered
predictions by observations made after the predictions were frozen.
\end{definition}

\begin{theorem}[The independent evidence program is the remaining external task]
\label{thm:temporal-evidence}
The independent temporal evidence program is the remaining external empirical task of the
theory: the confirmation of the pre-registered predictions by new, temporally-ordered
observations.
\end{theorem}

\begin{proof}
The pre-registered predictions \cite{QG190} were frozen before their observation; the
outcome dashboard \cite{QG202} tracks their evidence. The confirmations (or
contradictions) of P1, P2, and P3 by new observations constitute the independent temporal
evidence --- evidence obtained after the predictions were fixed, hence independent of the
derivation. This is the remaining external empirical program \cite{QG299}.
\end{proof}

\begin{corollary}[Temporal independence is preserved]
\label{cor:temporal-independence}
The pre-registered predictions preserve temporal independence: their values were frozen
before the observations, so their confirmation is genuine independent evidence.
\end{corollary}

\begin{proof}
By Theorem~\ref{thm:temporal-evidence} and the pre-registration \cite{QG190}, the
prediction values were fixed in advance. Hence their confirmation is temporally
independent evidence, not post-hoc fitting.
\end{proof}

\section{Failure Conditions}
\label{sec:failure-conditions}

\begin{theorem}[The failure conditions]
\label{thm:failure-conditions}
The theory specifies explicit failure conditions: observations that would force a revision.
\end{theorem}

\begin{proof}
The prediction registry \cite{QG190} records the falsification condition of each frozen
prediction: for P1, no signal in statistically sensitive searches of the $99$--$114$ GeV
window; for P2, a measured upper limit below $2.02$ meV; for P3, a sensitive search
excluding any frozen rung. These conditions, documented in the outcome dashboard
\cite{QG202}, specify exactly which observations would contradict the theory. Hence the
failure conditions are explicit.
\end{proof}

\begin{corollary}[Revisions are possible]
\label{cor:revisions-possible}
The theory is revisable: a confirmed contradiction of a pre-registered prediction would
force a revision of the corresponding sector of the theory.
\end{corollary}

\begin{proof}
By Theorem~\ref{thm:failure-conditions}, the failure conditions specify the contradicting
observations. A confirmed contradiction would force a revision, exactly as a scientific
theory must allow. Hence the theory is revisable.
\end{proof}

\section{Scope of Validity}
\label{sec:scope}

\begin{theorem}[What the theory claims]
\label{thm:claims}
The theory claims: the derivation of its physics observables from the canonical hierarchy;
the pre-registered predictions; and the classification of its boundaries and validation
items.
\end{theorem}

\begin{proof}
The derivation chapters establish the observables; the pre-registration \cite{QG190}
establishes the predictions; Chapter~\ref{chap:boundaries} establishes the boundary
classification. Hence the theory's claims are exactly these.
\end{proof}

\begin{theorem}[What the theory does not claim]
\label{thm:non-claims-scope}
The theory does not claim: an open derivation frontier; a derivation of the hosted
standard-model Lagrangian; a derivation of the exact Bekenstein coefficient; or any result
beyond the accepted derivations.
\end{theorem}

\begin{proof}
There is no open derivation frontier (Corollary~\ref{cor:no-open-derivation}); the SM
dynamics is hosted (Chapter~\ref{chap:boundaries}, Theorem~\ref{thm:sm-hosted}); the
Bekenstein coefficient requires the imported $2\pi$ factor (Chapter~\ref{chap:boundaries},
Theorem~\ref{thm:bekenstein-2pi}); and the monograph is assembly-only
\cite{MONO007}. Hence the non-claims are explicit.
\end{proof}

\begin{corollary}[The scope is exact]
\label{cor:scope-exact}
The scope of validity of the theory is exact: it claims its derivations and predictions,
and discloses its boundaries and validation items, with no excess.
\end{corollary}

\begin{proof}
By Theorem~\ref{thm:claims} and Theorem~\ref{thm:non-claims-scope}, the claims and
non-claims are explicit and exhaustive. Hence the scope of validity is exact.
\end{proof}

\section{Future Experimental Program}
\label{sec:future-program}

\begin{theorem}[The validation roadmap]
\label{thm:validation-roadmap}
The future experimental program of the theory is a validation roadmap: the confirmation of
the derived predictions by the scheduled experiments.
\end{theorem}

\begin{proof}
The roadmap follows the validation targets. The oblique parameters $S,T,U$ require only
re-analysis as the EW fit tightens (low priority). The electron g-2 $a_e$ requires the
lattice-limited measurement (medium priority). The Majorana character $0\nu\beta\beta$
requires the next-generation experiments nEXO and LEGEND-1000 (high priority), which reach
the $0.036$--$0.156$ eV window containing $m_{\beta\beta} = 2.02$ meV \cite{VALID001}. The
pre-registered P1 (106 GeV) awaits the HL-LHC, and P3 awaits further high-energy
confirmation \cite{QG199,QG200}. Hence the future program is a validation roadmap.
\end{proof}

\begin{corollary}[The roadmap is experimental only]
\label{cor:roadmap-experimental}
The future program contains no derivation tasks: every observable is derived, and the
roadmap schedules only experimental validation.
\end{corollary}

\begin{proof}
By Theorem~\ref{thm:validation-roadmap}, the roadmap schedules the experiments; by
Corollary~\ref{cor:no-open-derivation}, there are no derivation tasks. Hence the future
program is experimental only.
\end{proof}

\begin{remark}[Publication-safe wording]
Consistent with MONO007 \cite{MONO007} and MONO110A \cite{MONO110A}, this chapter makes the
frontier explicit and experimental: no open derivation frontier is claimed, the validation
is clearly separated from derivation, the falsification paths are documented, and no claim
exceeds the accepted results. The core monograph is thereby closed on a publication-safe
basis.
\end{remark}

\section{Conclusion}
\label{sec:frontier-conclusion}

\begin{theorem}[The core monograph is closed]
\label{thm:core-closed}
The core of the physics-focused monograph is closed: the physics derivation is complete
within the canonical hierarchy, the boundaries are disclosed, and the frontier is an
explicit, experimental validation program.
\end{theorem}

\begin{proof}
The physics derivation is complete (the derivation chapters); the boundaries are disclosed
(Chapter~\ref{chap:boundaries}); the frontier is experimental
(Corollary~\ref{cor:no-open-derivation}, Theorem~\ref{thm:validation-roadmap}). Hence the
core monograph is closed on a publication-safe basis.
\end{proof}

\begin{corollary}[The publication is ready]
\label{cor:publication-ready}
The physics-focused monograph of The Actualization Theory is publication-ready: it derives
the physics, discloses the boundaries, and documents the experimental frontier.
\end{corollary}

\begin{proof}
By Theorem~\ref{thm:core-closed}, the core is complete, bounded, and frontier-explicit;
by MONO110A \cite{MONO110A}, the structure is physics-focused and publication-ready. Hence
the publication is ready.
\end{proof}

\section{Summary}
\label{sec:frontier-summary}

This chapter has closed the core monograph. We have proved:

\begin{enumerate}
\item \textbf{Falsifiability} (Theorem~\ref{thm:falsifiability},
Corollary~\ref{cor:falsifiability-strength}): the theory exposes its failure modes and is
therefore scientific;
\item \textbf{Derivation vs validation} (Theorem~\ref{thm:derived-vs-validated},
Corollary~\ref{cor:no-open-derivation}): derivation and validation are distinct, and there
is no open derivation frontier;
\item \textbf{Validation targets} (Theorem~\ref{thm:precision-validation},
Theorem~\ref{thm:pre-registered}, Corollary~\ref{cor:targets-derived}): $S,T,U$, $a_e$,
$0\nu\beta\beta$, P1, P2, and P3 are derived predictions awaiting experiment;
\item \textbf{Independent temporal evidence} (Theorem~\ref{thm:temporal-evidence},
Corollary~\ref{cor:temporal-independence}): the pre-registered predictions provide
temporally independent evidence;
\item \textbf{Failure conditions} (Theorem~\ref{thm:failure-conditions},
Corollary~\ref{cor:revisions-possible}): the falsification conditions are explicit and the
theory is revisable;
\item \textbf{Scope of validity} (Theorem~\ref{thm:claims},
Theorem~\ref{thm:non-claims-scope}, Corollary~\ref{cor:scope-exact}): the theory claims its
derivations and predictions, and discloses its boundaries;
\item \textbf{Future program} (Theorem~\ref{thm:validation-roadmap},
Corollary~\ref{cor:roadmap-experimental}): the future program is a validation roadmap, with
no derivation tasks;
\item \textbf{Conclusion} (Theorem~\ref{thm:core-closed},
Corollary~\ref{cor:publication-ready}): the core monograph is closed and publication-ready.
\end{enumerate}

The frontier of The Actualization Theory is therefore explicit and experimental: the derived
predictions --- $S,T,U$, the electron g-2, the Majorana character, P1, P2, and the P3
ladder --- await their validation, with explicit falsification conditions, temporally
independent evidence, and a validation roadmap. The theory claims no open derivation
frontier; it discloses its boundaries; and it is falsifiable. The core of the
physics-focused monograph is closed and publication-ready.

\begin{thebibliography}{99}
\bibitem{QG190} TQM-QG Phase 190, \emph{Pre-Registered Predictions}, Docs/Research.
\bibitem{QG199} TQM-QG Phase 199, \emph{P1 Evidence Update}, Docs/Research.
\bibitem{QG200} TQM-QG Phase 200, \emph{Sector Ladder Evidence Audit}, Docs/Research.
\bibitem{QG201} TQM-QG Phase 201, \emph{Ladder Statistics Audit}, Docs/Research.
\bibitem{QG202} TQM-QG Phase 202, \emph{Prediction Outcome Dashboard}, Docs/Research.
\bibitem{QG203} TQM-QG Phase 203, \emph{Absolute Neutrino Mass Origin}, Docs/Research.
\bibitem{QG299} TQM-QG Phase 299, \emph{Remaining Frontier Re-audit}, Docs/Research.
\bibitem{VALID001} TQM-VALID001, \emph{Remaining Untested Items Audit}, Docs/Research.
\bibitem{MONO007} TQM-MONO007, \emph{Referee-Readiness Resolution}, Docs/Zenodo/Monograph.
\bibitem{MONO110A} TQM-MONO110A, \emph{Physics-Focused Monograph Re-audit}, Docs/Research.
\end{thebibliography}
 after the preamble
%         that defines the theorem environments (or include via the master file).
% ======================================================================

% ---- Theorem environments (define in the master preamble if not present) ----
% \newtheorem{definition}{Definition}[chapter]
% \newtheorem{lemma}[definition]{Lemma}
% \newtheorem{theorem}[definition]{Theorem}
% \newtheorem{corollary}[definition]{Corollary}
% \newtheorem{remark}[definition]{Remark}

\chapter{Frontier and Falsification}
\label{chap:frontier}

\section{Introduction}
\label{sec:frontier-introduction}

The preceding chapters derived the physics of The Actualization Theory and documented its
boundaries. This chapter closes the core monograph by addressing the \emph{frontier}: the
experimental validation of the derived predictions, the independent temporal evidence, the
explicit failure conditions, and the scope of validity. The chapter's purpose is to make the
theory \emph{publication-safe} --- to state clearly what would confirm it, what would
falsify it, and what it does and does not claim.

The central content is: the scientific-falsifiability requirement \cite{QG299}; the
distinction between derivation and validation; the current validation targets --- the
oblique parameters $S,T,U$, the electron g-2 $a_e$, the Majorana character $0\nu\beta\beta$
\cite{VALID001}, and the pre-registered predictions P1 (106 GeV), P2 ($0\nu\beta\beta$), and the P3
sector ladder \cite{QG190,QG199,QG200,QG201,QG202,QG203}; the independent temporal evidence;
the explicit failure conditions; the scope of validity; and the future experimental
program.

Following the physics-focused re-audit \cite{MONO110A}, this chapter claims \emph{no open
derivation frontier}: every observable is derived, and the frontier is experimental. The
experimental validation is clearly separated from the derivation, the falsification paths
are explicit, and the wording is referee-safe and publication-safe \cite{MONO007}. We use
only accepted canonical results and introduce no new physics and no speculation.

\section{Scientific Falsifiability}
\label{sec:falsifiability}

\begin{definition}[Falsifiability]
\label{def:falsifiability}
A scientific theory is \emph{falsifiable} if it specifies, in advance, observations that
would contradict it: the theory must expose its failure modes.
\end{definition}

\begin{theorem}[A scientific theory must expose its failure modes]
\label{thm:falsifiability}
The Actualization Theory is falsifiable: it specifies explicit, pre-registered
observations whose absence or contradiction would force a revision.
\end{theorem}

\begin{proof}
The post-frontier audit \cite{QG299} classifies the remaining items into open, partial,
boundary, and methodology, with the pre-registered predictions carrying explicit
falsification conditions. The prediction-outcome dashboard \cite{QG202} records, for each
prediction, its frozen value, its current evidence, and its falsification condition.
Hence the theory exposes its failure modes in advance: it is falsifiable.
\end{proof}

\begin{corollary}[Falsifiability is not a weakness]
\label{cor:falsifiability-strength}
The exposure of failure modes is a strength, not a weakness: it is the property that makes
the theory scientific and its claims testable.
\end{corollary}

\begin{proof}
By Theorem~\ref{thm:falsifiability}, the theory specifies its failure modes; by
Definition~\ref{def:falsifiability}, this is the defining property of a falsifiable theory.
Hence the explicit falsification paths are a scientific strength.
\end{proof}

\section{Experimental Validation Status}
\label{sec:validation-status}

\begin{definition}[Derivation vs validation]
\label{def:derivation-validation}
\emph{Derivation} is the derivation of an observable from the theory's principles; 
\emph{validation} is the confirmation of that derived value by experiment. The two are
distinct: a derived prediction may await validation without any derivation gap.
\end{definition}

\begin{theorem}[Derived and validated are distinct]
\label{thm:derived-vs-validated}
The status of an observable is the pair (derived, validated): a quantity may be derived
without being validated, and its derivation is complete before its validation is possible.
\end{theorem}

\begin{proof}
Following VALID001 \cite{VALID001}, the quantities $S,T,U$, the electron g-2 $a_e$, and the
Majorana character $0\nu\beta\beta$ are status A --- derived prediction exists, with no
missing derivation steps --- and belong to the experimental-validation category. Their
derivation is complete; their validation is the open experimental task. Hence derivation and
validation are distinct, and a derived-but-unvalidated quantity is not an open derivation
issue.
\end{proof}

\begin{corollary}[No open derivation frontier]
\label{cor:no-open-derivation}
The theory claims no open derivation frontier: every observable of the theory is derived,
and the frontier is experimental validation, not derivation.
\end{corollary}

\begin{proof}
By Theorem~\ref{thm:derived-vs-validated}, the derived quantities are complete, with the
validation tasks ahead. The reconstruction audit establishes that the major results are
reconstructed from the minimal hierarchy. Hence there is no open derivation frontier; the
frontier is experimental \cite{MONO110A}.
\end{proof}

\section{Current Validation Targets}
\label{sec:validation-targets}

\begin{theorem}[The precision predictions await validation]
\label{thm:precision-validation}
The oblique parameters $S,T,U$, the electron g-2 $a_e$, and the Majorana character
$0\nu\beta\beta$ are derived predictions awaiting experimental validation.
\end{theorem}

\begin{proof}
Following VALID001 \cite{VALID001}: $S = 0.0421$, $T = 0.0842$, $U = 0$ are derived and
consistent with the EW global fit; $a_e = 1.159655\times10^{-3}$ matches experiment within
$0.0003\%$; $m_{\beta\beta} = 2.02\times10^{-3}$ eV is within limits and in reach of
next-generation experiments. All three are status A, with their experimental validation the
open task. Hence they are current validation targets.
\end{proof}

\begin{theorem}[The pre-registered predictions]
\label{thm:pre-registered}
The theory carries three pre-registered predictions with frozen values and explicit
falsification conditions:
\begin{enumerate}
\item \textbf{P1} --- a 106 GeV resonance: central value $106.39$ GeV, window
$99$--$114$ GeV;
\item \textbf{P2} --- the $0\nu\beta\beta$ effective Majorana mass $m_{\beta\beta} =
2.02$ meV;
\item \textbf{P3} --- the sector-ladder spectrum: nine resonances
$106.39 \to 263.43$ GeV with a $15.20$ GeV width scale.
\end{enumerate}
\end{theorem}

\begin{proof}
The predictions are pre-registered \cite{QG190}, with their values frozen in the immutable
prediction registry and their evidence tracked in the outcome dashboard \cite{QG202}. P1 is
pending (window $99$--$114$ GeV neither confirmed nor excluded) \cite{QG199}; P2 is pending
(below current sensitivity, awaiting nEXO/LEGEND-1000); P3 is supported (the $151.98$ GeV
rung matches a $\sim152$ GeV excess at moderate significance) \cite{QG200,QG201}. The
absolute-mass law \cite{QG203} anchors the sector structure.
\end{proof}

\begin{corollary}[The targets are derived, not open issues]
\label{cor:targets-derived}
The validation targets are derived predictions awaiting experiment; they are not open
derivation issues and not boundaries.
\end{corollary}

\begin{proof}
By Theorem~\ref{thm:precision-validation} and Theorem~\ref{thm:pre-registered}, all
targets are derived predictions with frozen values; by Chapter~\ref{chap:boundaries}
(Corollary~\ref{cor:no-open-validation}), validation items are distinct from open
derivation issues. Hence the targets are validation items.
\end{proof}

\section{Independent Temporal Evidence}
\label{sec:temporal-evidence}

\begin{definition}[Independent temporal evidence]
\label{def:temporal-evidence}
\emph{Independent temporal evidence} is the empirical program of the theory that does not
depend on the already-derived observables: the confirmation of the pre-registered
predictions by observations made after the predictions were frozen.
\end{definition}

\begin{theorem}[The independent evidence program is the remaining external task]
\label{thm:temporal-evidence}
The independent temporal evidence program is the remaining external empirical task of the
theory: the confirmation of the pre-registered predictions by new, temporally-ordered
observations.
\end{theorem}

\begin{proof}
The pre-registered predictions \cite{QG190} were frozen before their observation; the
outcome dashboard \cite{QG202} tracks their evidence. The confirmations (or
contradictions) of P1, P2, and P3 by new observations constitute the independent temporal
evidence --- evidence obtained after the predictions were fixed, hence independent of the
derivation. This is the remaining external empirical program \cite{QG299}.
\end{proof}

\begin{corollary}[Temporal independence is preserved]
\label{cor:temporal-independence}
The pre-registered predictions preserve temporal independence: their values were frozen
before the observations, so their confirmation is genuine independent evidence.
\end{corollary}

\begin{proof}
By Theorem~\ref{thm:temporal-evidence} and the pre-registration \cite{QG190}, the
prediction values were fixed in advance. Hence their confirmation is temporally
independent evidence, not post-hoc fitting.
\end{proof}

\section{Failure Conditions}
\label{sec:failure-conditions}

\begin{theorem}[The failure conditions]
\label{thm:failure-conditions}
The theory specifies explicit failure conditions: observations that would force a revision.
\end{theorem}

\begin{proof}
The prediction registry \cite{QG190} records the falsification condition of each frozen
prediction: for P1, no signal in statistically sensitive searches of the $99$--$114$ GeV
window; for P2, a measured upper limit below $2.02$ meV; for P3, a sensitive search
excluding any frozen rung. These conditions, documented in the outcome dashboard
\cite{QG202}, specify exactly which observations would contradict the theory. Hence the
failure conditions are explicit.
\end{proof}

\begin{corollary}[Revisions are possible]
\label{cor:revisions-possible}
The theory is revisable: a confirmed contradiction of a pre-registered prediction would
force a revision of the corresponding sector of the theory.
\end{corollary}

\begin{proof}
By Theorem~\ref{thm:failure-conditions}, the failure conditions specify the contradicting
observations. A confirmed contradiction would force a revision, exactly as a scientific
theory must allow. Hence the theory is revisable.
\end{proof}

\section{Scope of Validity}
\label{sec:scope}

\begin{theorem}[What the theory claims]
\label{thm:claims}
The theory claims: the derivation of its physics observables from the canonical hierarchy;
the pre-registered predictions; and the classification of its boundaries and validation
items.
\end{theorem}

\begin{proof}
The derivation chapters establish the observables; the pre-registration \cite{QG190}
establishes the predictions; Chapter~\ref{chap:boundaries} establishes the boundary
classification. Hence the theory's claims are exactly these.
\end{proof}

\begin{theorem}[What the theory does not claim]
\label{thm:non-claims-scope}
The theory does not claim: an open derivation frontier; a derivation of the hosted
standard-model Lagrangian; a derivation of the exact Bekenstein coefficient; or any result
beyond the accepted derivations.
\end{theorem}

\begin{proof}
There is no open derivation frontier (Corollary~\ref{cor:no-open-derivation}); the SM
dynamics is hosted (Chapter~\ref{chap:boundaries}, Theorem~\ref{thm:sm-hosted}); the
Bekenstein coefficient requires the imported $2\pi$ factor (Chapter~\ref{chap:boundaries},
Theorem~\ref{thm:bekenstein-2pi}); and the monograph is assembly-only
\cite{MONO007}. Hence the non-claims are explicit.
\end{proof}

\begin{corollary}[The scope is exact]
\label{cor:scope-exact}
The scope of validity of the theory is exact: it claims its derivations and predictions,
and discloses its boundaries and validation items, with no excess.
\end{corollary}

\begin{proof}
By Theorem~\ref{thm:claims} and Theorem~\ref{thm:non-claims-scope}, the claims and
non-claims are explicit and exhaustive. Hence the scope of validity is exact.
\end{proof}

\section{Future Experimental Program}
\label{sec:future-program}

\begin{theorem}[The validation roadmap]
\label{thm:validation-roadmap}
The future experimental program of the theory is a validation roadmap: the confirmation of
the derived predictions by the scheduled experiments.
\end{theorem}

\begin{proof}
The roadmap follows the validation targets. The oblique parameters $S,T,U$ require only
re-analysis as the EW fit tightens (low priority). The electron g-2 $a_e$ requires the
lattice-limited measurement (medium priority). The Majorana character $0\nu\beta\beta$
requires the next-generation experiments nEXO and LEGEND-1000 (high priority), which reach
the $0.036$--$0.156$ eV window containing $m_{\beta\beta} = 2.02$ meV \cite{VALID001}. The
pre-registered P1 (106 GeV) awaits the HL-LHC, and P3 awaits further high-energy
confirmation \cite{QG199,QG200}. Hence the future program is a validation roadmap.
\end{proof}

\begin{corollary}[The roadmap is experimental only]
\label{cor:roadmap-experimental}
The future program contains no derivation tasks: every observable is derived, and the
roadmap schedules only experimental validation.
\end{corollary}

\begin{proof}
By Theorem~\ref{thm:validation-roadmap}, the roadmap schedules the experiments; by
Corollary~\ref{cor:no-open-derivation}, there are no derivation tasks. Hence the future
program is experimental only.
\end{proof}

\begin{remark}[Publication-safe wording]
Consistent with MONO007 \cite{MONO007} and MONO110A \cite{MONO110A}, this chapter makes the
frontier explicit and experimental: no open derivation frontier is claimed, the validation
is clearly separated from derivation, the falsification paths are documented, and no claim
exceeds the accepted results. The core monograph is thereby closed on a publication-safe
basis.
\end{remark}

\section{Conclusion}
\label{sec:frontier-conclusion}

\begin{theorem}[The core monograph is closed]
\label{thm:core-closed}
The core of the physics-focused monograph is closed: the physics derivation is complete
within the canonical hierarchy, the boundaries are disclosed, and the frontier is an
explicit, experimental validation program.
\end{theorem}

\begin{proof}
The physics derivation is complete (the derivation chapters); the boundaries are disclosed
(Chapter~\ref{chap:boundaries}); the frontier is experimental
(Corollary~\ref{cor:no-open-derivation}, Theorem~\ref{thm:validation-roadmap}). Hence the
core monograph is closed on a publication-safe basis.
\end{proof}

\begin{corollary}[The publication is ready]
\label{cor:publication-ready}
The physics-focused monograph of The Actualization Theory is publication-ready: it derives
the physics, discloses the boundaries, and documents the experimental frontier.
\end{corollary}

\begin{proof}
By Theorem~\ref{thm:core-closed}, the core is complete, bounded, and frontier-explicit;
by MONO110A \cite{MONO110A}, the structure is physics-focused and publication-ready. Hence
the publication is ready.
\end{proof}

\section{Summary}
\label{sec:frontier-summary}

This chapter has closed the core monograph. We have proved:

\begin{enumerate}
\item \textbf{Falsifiability} (Theorem~\ref{thm:falsifiability},
Corollary~\ref{cor:falsifiability-strength}): the theory exposes its failure modes and is
therefore scientific;
\item \textbf{Derivation vs validation} (Theorem~\ref{thm:derived-vs-validated},
Corollary~\ref{cor:no-open-derivation}): derivation and validation are distinct, and there
is no open derivation frontier;
\item \textbf{Validation targets} (Theorem~\ref{thm:precision-validation},
Theorem~\ref{thm:pre-registered}, Corollary~\ref{cor:targets-derived}): $S,T,U$, $a_e$,
$0\nu\beta\beta$, P1, P2, and P3 are derived predictions awaiting experiment;
\item \textbf{Independent temporal evidence} (Theorem~\ref{thm:temporal-evidence},
Corollary~\ref{cor:temporal-independence}): the pre-registered predictions provide
temporally independent evidence;
\item \textbf{Failure conditions} (Theorem~\ref{thm:failure-conditions},
Corollary~\ref{cor:revisions-possible}): the falsification conditions are explicit and the
theory is revisable;
\item \textbf{Scope of validity} (Theorem~\ref{thm:claims},
Theorem~\ref{thm:non-claims-scope}, Corollary~\ref{cor:scope-exact}): the theory claims its
derivations and predictions, and discloses its boundaries;
\item \textbf{Future program} (Theorem~\ref{thm:validation-roadmap},
Corollary~\ref{cor:roadmap-experimental}): the future program is a validation roadmap, with
no derivation tasks;
\item \textbf{Conclusion} (Theorem~\ref{thm:core-closed},
Corollary~\ref{cor:publication-ready}): the core monograph is closed and publication-ready.
\end{enumerate}

The frontier of The Actualization Theory is therefore explicit and experimental: the derived
predictions --- $S,T,U$, the electron g-2, the Majorana character, P1, P2, and the P3
ladder --- await their validation, with explicit falsification conditions, temporally
independent evidence, and a validation roadmap. The theory claims no open derivation
frontier; it discloses its boundaries; and it is falsifiable. The core of the
physics-focused monograph is closed and publication-ready.

\begin{thebibliography}{99}
\bibitem{QG190} TQM-QG Phase 190, \emph{Pre-Registered Predictions}, Docs/Research.
\bibitem{QG199} TQM-QG Phase 199, \emph{P1 Evidence Update}, Docs/Research.
\bibitem{QG200} TQM-QG Phase 200, \emph{Sector Ladder Evidence Audit}, Docs/Research.
\bibitem{QG201} TQM-QG Phase 201, \emph{Ladder Statistics Audit}, Docs/Research.
\bibitem{QG202} TQM-QG Phase 202, \emph{Prediction Outcome Dashboard}, Docs/Research.
\bibitem{QG203} TQM-QG Phase 203, \emph{Absolute Neutrino Mass Origin}, Docs/Research.
\bibitem{QG299} TQM-QG Phase 299, \emph{Remaining Frontier Re-audit}, Docs/Research.
\bibitem{VALID001} TQM-VALID001, \emph{Remaining Untested Items Audit}, Docs/Research.
\bibitem{MONO007} TQM-MONO007, \emph{Referee-Readiness Resolution}, Docs/Zenodo/Monograph.
\bibitem{MONO110A} TQM-MONO110A, \emph{Physics-Focused Monograph Re-audit}, Docs/Research.
\end{thebibliography}
 after the preamble
%         that defines the theorem environments (or include via the master file).
% ======================================================================

% ---- Theorem environments (define in the master preamble if not present) ----
% \newtheorem{definition}{Definition}[chapter]
% \newtheorem{lemma}[definition]{Lemma}
% \newtheorem{theorem}[definition]{Theorem}
% \newtheorem{corollary}[definition]{Corollary}
% \newtheorem{remark}[definition]{Remark}

\chapter{Frontier and Falsification}
\label{chap:frontier}

\section{Introduction}
\label{sec:frontier-introduction}

The preceding chapters derived the physics of The Actualization Theory and documented its
boundaries. This chapter closes the core monograph by addressing the \emph{frontier}: the
experimental validation of the derived predictions, the independent temporal evidence, the
explicit failure conditions, and the scope of validity. The chapter's purpose is to make the
theory \emph{publication-safe} --- to state clearly what would confirm it, what would
falsify it, and what it does and does not claim.

The central content is: the scientific-falsifiability requirement \cite{QG299}; the
distinction between derivation and validation; the current validation targets --- the
oblique parameters $S,T,U$, the electron g-2 $a_e$, the Majorana character $0\nu\beta\beta$
\cite{VALID001}, and the pre-registered predictions P1 (106 GeV), P2 ($0\nu\beta\beta$), and the P3
sector ladder \cite{QG190,QG199,QG200,QG201,QG202,QG203}; the independent temporal evidence;
the explicit failure conditions; the scope of validity; and the future experimental
program.

Following the physics-focused re-audit \cite{MONO110A}, this chapter claims \emph{no open
derivation frontier}: every observable is derived, and the frontier is experimental. The
experimental validation is clearly separated from the derivation, the falsification paths
are explicit, and the wording is referee-safe and publication-safe \cite{MONO007}. We use
only accepted canonical results and introduce no new physics and no speculation.

\section{Scientific Falsifiability}
\label{sec:falsifiability}

\begin{definition}[Falsifiability]
\label{def:falsifiability}
A scientific theory is \emph{falsifiable} if it specifies, in advance, observations that
would contradict it: the theory must expose its failure modes.
\end{definition}

\begin{theorem}[A scientific theory must expose its failure modes]
\label{thm:falsifiability}
The Actualization Theory is falsifiable: it specifies explicit, pre-registered
observations whose absence or contradiction would force a revision.
\end{theorem}

\begin{proof}
The post-frontier audit \cite{QG299} classifies the remaining items into open, partial,
boundary, and methodology, with the pre-registered predictions carrying explicit
falsification conditions. The prediction-outcome dashboard \cite{QG202} records, for each
prediction, its frozen value, its current evidence, and its falsification condition.
Hence the theory exposes its failure modes in advance: it is falsifiable.
\end{proof}

\begin{corollary}[Falsifiability is not a weakness]
\label{cor:falsifiability-strength}
The exposure of failure modes is a strength, not a weakness: it is the property that makes
the theory scientific and its claims testable.
\end{corollary}

\begin{proof}
By Theorem~\ref{thm:falsifiability}, the theory specifies its failure modes; by
Definition~\ref{def:falsifiability}, this is the defining property of a falsifiable theory.
Hence the explicit falsification paths are a scientific strength.
\end{proof}

\section{Experimental Validation Status}
\label{sec:validation-status}

\begin{definition}[Derivation vs validation]
\label{def:derivation-validation}
\emph{Derivation} is the derivation of an observable from the theory's principles; 
\emph{validation} is the confirmation of that derived value by experiment. The two are
distinct: a derived prediction may await validation without any derivation gap.
\end{definition}

\begin{theorem}[Derived and validated are distinct]
\label{thm:derived-vs-validated}
The status of an observable is the pair (derived, validated): a quantity may be derived
without being validated, and its derivation is complete before its validation is possible.
\end{theorem}

\begin{proof}
Following VALID001 \cite{VALID001}, the quantities $S,T,U$, the electron g-2 $a_e$, and the
Majorana character $0\nu\beta\beta$ are status A --- derived prediction exists, with no
missing derivation steps --- and belong to the experimental-validation category. Their
derivation is complete; their validation is the open experimental task. Hence derivation and
validation are distinct, and a derived-but-unvalidated quantity is not an open derivation
issue.
\end{proof}

\begin{corollary}[No open derivation frontier]
\label{cor:no-open-derivation}
The theory claims no open derivation frontier: every observable of the theory is derived,
and the frontier is experimental validation, not derivation.
\end{corollary}

\begin{proof}
By Theorem~\ref{thm:derived-vs-validated}, the derived quantities are complete, with the
validation tasks ahead. The reconstruction audit establishes that the major results are
reconstructed from the minimal hierarchy. Hence there is no open derivation frontier; the
frontier is experimental \cite{MONO110A}.
\end{proof}

\section{Current Validation Targets}
\label{sec:validation-targets}

\begin{theorem}[The precision predictions await validation]
\label{thm:precision-validation}
The oblique parameters $S,T,U$, the electron g-2 $a_e$, and the Majorana character
$0\nu\beta\beta$ are derived predictions awaiting experimental validation.
\end{theorem}

\begin{proof}
Following VALID001 \cite{VALID001}: $S = 0.0421$, $T = 0.0842$, $U = 0$ are derived and
consistent with the EW global fit; $a_e = 1.159655\times10^{-3}$ matches experiment within
$0.0003\%$; $m_{\beta\beta} = 2.02\times10^{-3}$ eV is within limits and in reach of
next-generation experiments. All three are status A, with their experimental validation the
open task. Hence they are current validation targets.
\end{proof}

\begin{theorem}[The pre-registered predictions]
\label{thm:pre-registered}
The theory carries three pre-registered predictions with frozen values and explicit
falsification conditions:
\begin{enumerate}
\item \textbf{P1} --- a 106 GeV resonance: central value $106.39$ GeV, window
$99$--$114$ GeV;
\item \textbf{P2} --- the $0\nu\beta\beta$ effective Majorana mass $m_{\beta\beta} =
2.02$ meV;
\item \textbf{P3} --- the sector-ladder spectrum: nine resonances
$106.39 \to 263.43$ GeV with a $15.20$ GeV width scale.
\end{enumerate}
\end{theorem}

\begin{proof}
The predictions are pre-registered \cite{QG190}, with their values frozen in the immutable
prediction registry and their evidence tracked in the outcome dashboard \cite{QG202}. P1 is
pending (window $99$--$114$ GeV neither confirmed nor excluded) \cite{QG199}; P2 is pending
(below current sensitivity, awaiting nEXO/LEGEND-1000); P3 is supported (the $151.98$ GeV
rung matches a $\sim152$ GeV excess at moderate significance) \cite{QG200,QG201}. The
absolute-mass law \cite{QG203} anchors the sector structure.
\end{proof}

\begin{corollary}[The targets are derived, not open issues]
\label{cor:targets-derived}
The validation targets are derived predictions awaiting experiment; they are not open
derivation issues and not boundaries.
\end{corollary}

\begin{proof}
By Theorem~\ref{thm:precision-validation} and Theorem~\ref{thm:pre-registered}, all
targets are derived predictions with frozen values; by Chapter~\ref{chap:boundaries}
(Corollary~\ref{cor:no-open-validation}), validation items are distinct from open
derivation issues. Hence the targets are validation items.
\end{proof}

\section{Independent Temporal Evidence}
\label{sec:temporal-evidence}

\begin{definition}[Independent temporal evidence]
\label{def:temporal-evidence}
\emph{Independent temporal evidence} is the empirical program of the theory that does not
depend on the already-derived observables: the confirmation of the pre-registered
predictions by observations made after the predictions were frozen.
\end{definition}

\begin{theorem}[The independent evidence program is the remaining external task]
\label{thm:temporal-evidence}
The independent temporal evidence program is the remaining external empirical task of the
theory: the confirmation of the pre-registered predictions by new, temporally-ordered
observations.
\end{theorem}

\begin{proof}
The pre-registered predictions \cite{QG190} were frozen before their observation; the
outcome dashboard \cite{QG202} tracks their evidence. The confirmations (or
contradictions) of P1, P2, and P3 by new observations constitute the independent temporal
evidence --- evidence obtained after the predictions were fixed, hence independent of the
derivation. This is the remaining external empirical program \cite{QG299}.
\end{proof}

\begin{corollary}[Temporal independence is preserved]
\label{cor:temporal-independence}
The pre-registered predictions preserve temporal independence: their values were frozen
before the observations, so their confirmation is genuine independent evidence.
\end{corollary}

\begin{proof}
By Theorem~\ref{thm:temporal-evidence} and the pre-registration \cite{QG190}, the
prediction values were fixed in advance. Hence their confirmation is temporally
independent evidence, not post-hoc fitting.
\end{proof}

\section{Failure Conditions}
\label{sec:failure-conditions}

\begin{theorem}[The failure conditions]
\label{thm:failure-conditions}
The theory specifies explicit failure conditions: observations that would force a revision.
\end{theorem}

\begin{proof}
The prediction registry \cite{QG190} records the falsification condition of each frozen
prediction: for P1, no signal in statistically sensitive searches of the $99$--$114$ GeV
window; for P2, a measured upper limit below $2.02$ meV; for P3, a sensitive search
excluding any frozen rung. These conditions, documented in the outcome dashboard
\cite{QG202}, specify exactly which observations would contradict the theory. Hence the
failure conditions are explicit.
\end{proof}

\begin{corollary}[Revisions are possible]
\label{cor:revisions-possible}
The theory is revisable: a confirmed contradiction of a pre-registered prediction would
force a revision of the corresponding sector of the theory.
\end{corollary}

\begin{proof}
By Theorem~\ref{thm:failure-conditions}, the failure conditions specify the contradicting
observations. A confirmed contradiction would force a revision, exactly as a scientific
theory must allow. Hence the theory is revisable.
\end{proof}

\section{Scope of Validity}
\label{sec:scope}

\begin{theorem}[What the theory claims]
\label{thm:claims}
The theory claims: the derivation of its physics observables from the canonical hierarchy;
the pre-registered predictions; and the classification of its boundaries and validation
items.
\end{theorem}

\begin{proof}
The derivation chapters establish the observables; the pre-registration \cite{QG190}
establishes the predictions; Chapter~\ref{chap:boundaries} establishes the boundary
classification. Hence the theory's claims are exactly these.
\end{proof}

\begin{theorem}[What the theory does not claim]
\label{thm:non-claims-scope}
The theory does not claim: an open derivation frontier; a derivation of the hosted
standard-model Lagrangian; a derivation of the exact Bekenstein coefficient; or any result
beyond the accepted derivations.
\end{theorem}

\begin{proof}
There is no open derivation frontier (Corollary~\ref{cor:no-open-derivation}); the SM
dynamics is hosted (Chapter~\ref{chap:boundaries}, Theorem~\ref{thm:sm-hosted}); the
Bekenstein coefficient requires the imported $2\pi$ factor (Chapter~\ref{chap:boundaries},
Theorem~\ref{thm:bekenstein-2pi}); and the monograph is assembly-only
\cite{MONO007}. Hence the non-claims are explicit.
\end{proof}

\begin{corollary}[The scope is exact]
\label{cor:scope-exact}
The scope of validity of the theory is exact: it claims its derivations and predictions,
and discloses its boundaries and validation items, with no excess.
\end{corollary}

\begin{proof}
By Theorem~\ref{thm:claims} and Theorem~\ref{thm:non-claims-scope}, the claims and
non-claims are explicit and exhaustive. Hence the scope of validity is exact.
\end{proof}

\section{Future Experimental Program}
\label{sec:future-program}

\begin{theorem}[The validation roadmap]
\label{thm:validation-roadmap}
The future experimental program of the theory is a validation roadmap: the confirmation of
the derived predictions by the scheduled experiments.
\end{theorem}

\begin{proof}
The roadmap follows the validation targets. The oblique parameters $S,T,U$ require only
re-analysis as the EW fit tightens (low priority). The electron g-2 $a_e$ requires the
lattice-limited measurement (medium priority). The Majorana character $0\nu\beta\beta$
requires the next-generation experiments nEXO and LEGEND-1000 (high priority), which reach
the $0.036$--$0.156$ eV window containing $m_{\beta\beta} = 2.02$ meV \cite{VALID001}. The
pre-registered P1 (106 GeV) awaits the HL-LHC, and P3 awaits further high-energy
confirmation \cite{QG199,QG200}. Hence the future program is a validation roadmap.
\end{proof}

\begin{corollary}[The roadmap is experimental only]
\label{cor:roadmap-experimental}
The future program contains no derivation tasks: every observable is derived, and the
roadmap schedules only experimental validation.
\end{corollary}

\begin{proof}
By Theorem~\ref{thm:validation-roadmap}, the roadmap schedules the experiments; by
Corollary~\ref{cor:no-open-derivation}, there are no derivation tasks. Hence the future
program is experimental only.
\end{proof}

\begin{remark}[Publication-safe wording]
Consistent with MONO007 \cite{MONO007} and MONO110A \cite{MONO110A}, this chapter makes the
frontier explicit and experimental: no open derivation frontier is claimed, the validation
is clearly separated from derivation, the falsification paths are documented, and no claim
exceeds the accepted results. The core monograph is thereby closed on a publication-safe
basis.
\end{remark}

\section{Conclusion}
\label{sec:frontier-conclusion}

\begin{theorem}[The core monograph is closed]
\label{thm:core-closed}
The core of the physics-focused monograph is closed: the physics derivation is complete
within the canonical hierarchy, the boundaries are disclosed, and the frontier is an
explicit, experimental validation program.
\end{theorem}

\begin{proof}
The physics derivation is complete (the derivation chapters); the boundaries are disclosed
(Chapter~\ref{chap:boundaries}); the frontier is experimental
(Corollary~\ref{cor:no-open-derivation}, Theorem~\ref{thm:validation-roadmap}). Hence the
core monograph is closed on a publication-safe basis.
\end{proof}

\begin{corollary}[The publication is ready]
\label{cor:publication-ready}
The physics-focused monograph of The Actualization Theory is publication-ready: it derives
the physics, discloses the boundaries, and documents the experimental frontier.
\end{corollary}

\begin{proof}
By Theorem~\ref{thm:core-closed}, the core is complete, bounded, and frontier-explicit;
by MONO110A \cite{MONO110A}, the structure is physics-focused and publication-ready. Hence
the publication is ready.
\end{proof}

\section{Summary}
\label{sec:frontier-summary}

This chapter has closed the core monograph. We have proved:

\begin{enumerate}
\item \textbf{Falsifiability} (Theorem~\ref{thm:falsifiability},
Corollary~\ref{cor:falsifiability-strength}): the theory exposes its failure modes and is
therefore scientific;
\item \textbf{Derivation vs validation} (Theorem~\ref{thm:derived-vs-validated},
Corollary~\ref{cor:no-open-derivation}): derivation and validation are distinct, and there
is no open derivation frontier;
\item \textbf{Validation targets} (Theorem~\ref{thm:precision-validation},
Theorem~\ref{thm:pre-registered}, Corollary~\ref{cor:targets-derived}): $S,T,U$, $a_e$,
$0\nu\beta\beta$, P1, P2, and P3 are derived predictions awaiting experiment;
\item \textbf{Independent temporal evidence} (Theorem~\ref{thm:temporal-evidence},
Corollary~\ref{cor:temporal-independence}): the pre-registered predictions provide
temporally independent evidence;
\item \textbf{Failure conditions} (Theorem~\ref{thm:failure-conditions},
Corollary~\ref{cor:revisions-possible}): the falsification conditions are explicit and the
theory is revisable;
\item \textbf{Scope of validity} (Theorem~\ref{thm:claims},
Theorem~\ref{thm:non-claims-scope}, Corollary~\ref{cor:scope-exact}): the theory claims its
derivations and predictions, and discloses its boundaries;
\item \textbf{Future program} (Theorem~\ref{thm:validation-roadmap},
Corollary~\ref{cor:roadmap-experimental}): the future program is a validation roadmap, with
no derivation tasks;
\item \textbf{Conclusion} (Theorem~\ref{thm:core-closed},
Corollary~\ref{cor:publication-ready}): the core monograph is closed and publication-ready.
\end{enumerate}

The frontier of The Actualization Theory is therefore explicit and experimental: the derived
predictions --- $S,T,U$, the electron g-2, the Majorana character, P1, P2, and the P3
ladder --- await their validation, with explicit falsification conditions, temporally
independent evidence, and a validation roadmap. The theory claims no open derivation
frontier; it discloses its boundaries; and it is falsifiable. The core of the
physics-focused monograph is closed and publication-ready.

\begin{thebibliography}{99}
\bibitem{QG190} TQM-QG Phase 190, \emph{Pre-Registered Predictions}, Docs/Research.
\bibitem{QG199} TQM-QG Phase 199, \emph{P1 Evidence Update}, Docs/Research.
\bibitem{QG200} TQM-QG Phase 200, \emph{Sector Ladder Evidence Audit}, Docs/Research.
\bibitem{QG201} TQM-QG Phase 201, \emph{Ladder Statistics Audit}, Docs/Research.
\bibitem{QG202} TQM-QG Phase 202, \emph{Prediction Outcome Dashboard}, Docs/Research.
\bibitem{QG203} TQM-QG Phase 203, \emph{Absolute Neutrino Mass Origin}, Docs/Research.
\bibitem{QG299} TQM-QG Phase 299, \emph{Remaining Frontier Re-audit}, Docs/Research.
\bibitem{VALID001} TQM-VALID001, \emph{Remaining Untested Items Audit}, Docs/Research.
\bibitem{MONO007} TQM-MONO007, \emph{Referee-Readiness Resolution}, Docs/Zenodo/Monograph.
\bibitem{MONO110A} TQM-MONO110A, \emph{Physics-Focused Monograph Re-audit}, Docs/Research.
\end{thebibliography}


% ---- General conclusion --------------------------------------------------
% ======================================================================
%  General Conclusion of The Actualization Theory
%  AT-MONO115 · Canonical Monograph (v2.0) · Zenodo-ready
%  Assembly only — no new physics, no new derivations.
% ======================================================================

\chapter*{General Conclusion}
\label{concl:conclusion}
\addcontentsline{toc}{chapter}{General Conclusion}

The Actualization Theory reconstructs physics from two primitives --- Difference and the
tensor reference $\eta$ --- through the canonical hierarchy

\[
\text{Difference} \;\longrightarrow\; \text{Actualization}
\;\longrightarrow\; \text{Inevitable Spectrum} \;\longrightarrow\; \text{Physics}.
\]

\section*{What has been established}

\begin{enumerate}
\item \textbf{The foundation.} Difference is the fundamental boundary: the irreducible
primitive whose count-conservation law is its definitional identity. Its two faces --- the
scalar $\rho$ and the tensor $\psi$ --- are the trace and traceless parts of the one
Difference. $\eta$ is the tensor reference against which the tensor face is read. The
minimal primitive set is $\{\text{Difference},\eta\}$.

\item \textbf{The process.} Actualization is the derived count-producing process of
Difference. Its convergence is the closure boundary $N=96$. The Closure Principle
characterizes the boundary; it does not derive the primitive.

\item \textbf{The spectrum.} The D96 spectrum is the inevitable output of the actualization
attractor: a unique, content-independent eigenspectrum carrying no choice and no primitive
status.

\item \textbf{The physics.} Quantum mechanics ($|\psi|^2=\rho$, measurement as
actualization), gravity and spacetime (emergent, native dynamics), the standard model
(derived reads of the D96 moments), and cosmology (a spectrum readout) are all derived from
the canonical foundation.

\textbf{Origin of Physical Units.} The D96 spectrum yields dimensionless ratios, moments,
and angles. Dimensional observables are obtained only after calibration to measured
anchors: $v$ for weak/Higgs masses, $m_e$ for quark masses, and the standard conversion
factors for SI gravity quantities. Those anchors fix the unit system; they are a boundary,
not a new derivation.

\item \textbf{The boundaries.} The limits of the framework are documented honestly:
Difference as the ontological boundary, the unresolved status of $\psi$, the numerical
boundary $\pi$, the Bekenstein $2\pi$ factor, and the hosted standard-model dynamics.
These are limits, not failures.

\item \textbf{The frontier.} There is no open derivation frontier. The frontier is the
experimental validation of the derived predictions: $S,T,U$, $a_e$, $0\nu\beta\beta$, and
the pre-registered predictions P1, P2, P3.
\end{enumerate}

\section*{What has not been claimed}

The monograph does not claim: a derivation of the hosted standard-model Lagrangian; a
derivation of the exact Bekenstein $1/4$ coefficient (it requires the imported $2\pi$
factor); a resolution of the ultimate ontological status of $\psi$; or any result beyond
the accepted derivations. The universality research program --- the operator basis across
organized domains, the lock law, the organization structure --- is treated as future work,
outside the core physics derivation.

\section*{The scientific status}

The Actualization Theory is falsifiable: it specifies explicit, pre-registered observations
whose absence or contradiction would force a revision. It claims exactly its derivations,
discloses its boundaries, and separates derivation from experimental validation. Its core
claim is that the physics of the standard model, gravity, and cosmology can be reconstructed
from Difference, Actualization, and the inevitable spectrum --- a claim whose decisive
confirmation now rests with the experiments that follow.

\vspace{1em}
\noindent\textbf{The Actualization Theory:}
\textit{a reconstruction of physics from Difference, Actualization and Spectrum ---
derived, bounded, and awaiting its experiments.}


% ---- Appendices -----------------------------------------------------------
\appendix
% ======================================================================
%  Appendices of The Actualization Theory
%  TQM-MONO115 · Canonical Monograph (v2.0) · Zenodo-ready
%  Includes: Appendix A (Universality Program — future work),
%            Appendix B (References structure),
%            Appendix C (Validation targets summary).
%  Assembly only — no new physics, no new derivations.
% ======================================================================

\chapter{The Universality Program: Future Work}
\label{app:universality}

The universality research program of The Actualization Theory is treated as future work, outside
the core physics derivation of this monograph. The core monograph derives the physics of the
theory; the universality program investigates the operator basis across organized domains
and the organization structure of spectra. Its results are summarized here for the record,
with the understanding that they are a separate research direction, not a requirement for
the physics derivation.

\section{The Operator Basis}
\label{app:operators}

The four operators --- Crowding (degeneracy grouping), Compression (octave banding), Beat
(frequency ratio), and Locking (spectral gap) --- are the derived reading instruments of the
D96 spectrum. They are projections of the spectrum, not primitives. No fifth operator has
been found in any searched domain.

\section{The Lock Law}
\label{app:lock-law}

The lock identities of organized spectra carry a universal structure (the presence of stable
ratios) with domain-specific values (each domain's fingerprint). The lock structure
precedes organization, and organization emerges at a critical transition. These results are
established on synthetic deterministic evolving-law cohorts, with the synthetic basis
disclosed.

\section{Scope}
\label{app:universality-scope}

The universality program strengthens the theory but is not required for the physics
derivation of Chapters 1--12. The operator basis and the lock law are presented here as the
future-work appendix of the monograph, consistent with the physics-focused publication plan.

\chapter{References}
\label{app:references}

The references of this monograph follow the canonical research-program record. Each chapter
carries its own bibliography of the accepted TQM-QG phases that support its results. The
reference structure is:

\begin{itemize}
\item \textbf{Foundation chapters (1--2):} the primitive and reference phases
(QG268, QG278, QG279, QG286, QG289--292).
\item \textbf{Dynamics chapters (3--5):} the actualization, closure, and spectrum phases
(QG267--268, QG282, QG293--296).
\item \textbf{Spectrum chapters (6--8):} the D96, operator, and lock phases
(QG155--160, QG260--263, QG300--319).
\item \textbf{Physics chapters (9--12):} the quantum, gravity, standard-model, and
cosmology phases (QG216--244, QG149--180, QG227--240).
\item \textbf{Boundary and frontier chapters (13--14):} the boundary and prediction phases
(QG185, QG196, QG242, QG245, QG190--203, QG299).
\end{itemize}

The complete machine-readable record of the research program is the physics-coverage
register, which ships with the Zenodo deposit.

\chapter{Validation Targets Summary}
\label{app:validation}

The following derived predictions await experimental validation:

\begin{center}
\begin{tabular}{lll}
\hline
\textbf{Target} & \textbf{Derived value} & \textbf{Validation path} \\
\hline
Oblique $S,T,U$ & $S=0.0421$, $T=0.0842$, $U=0$ & EW global-fit re-analysis \\
Electron g-2 $a_e$ & $1.159655\times10^{-3}$ & lattice-limited measurement \\
Majorana $m_{\beta\beta}$ & $2.02\times10^{-3}$ eV & nEXO / LEGEND-1000 \\
P1 (106 GeV) & $106.39$ GeV, window 99--114 & HL-LHC \\
P2 ($0\nu\beta\beta$) & $m_{\beta\beta}=2.02$ meV & nEXO / LEGEND-1000 \\
P3 (sector ladder) & 9 resonances to 263.43 GeV & HL-LHC confirmation \\
\hline
\end{tabular}
\end{center}

Each target carries an explicit falsification condition, documented in the prediction
registry and the outcome dashboard of the research program.


% ---- Back matter ----------------------------------------------------------
\backmatter
% References are structured per-chapter (each chapter carries its own
% thebibliography); a consolidated bibliography may be added here if desired.

\end{document}
